% !TEX program = lualatex
\documentclass[10pt,letterpaper]{article}
\usepackage{iftex}
\ifPDFTeX
  \usepackage[utf8]{inputenc}
  \usepackage[T1]{fontenc}
  \usepackage{mathptmx}
  \usepackage[scaled=0.92]{helvet}
  \usepackage{courier}
  \newcommand{\STAUnicodeAppendix}{\ttfamily}
\else
  \usepackage{fontspec}
  \setmainfont{TeX Gyre Termes}
  \setsansfont{TeX Gyre Heros}[Scale=MatchLowercase]
  \setmonofont{Latin Modern Mono}[Scale=MatchLowercase]
  \newfontfamily{\STAUnicodeAppendix}{Latin Modern Mono}[Scale=MatchLowercase]
\fi
\usepackage[
  letterpaper,
  top=0.75in,
  bottom=1in,
  left=0.625in,
  right=0.625in
]{geometry}
\setlength{\columnsep}{0.24in}
\usepackage[tracking=true]{microtype}
\usepackage{cite}
\usepackage{amsmath,amssymb,mathtools,bm,booktabs,xcolor,hyperref,fancyvrb,fvextra,array,longtable,caption,titlesec,tabularx,multicol}
\DeclareMathSizes{7.75}{7.75}{5}{5}
\allowdisplaybreaks[1]
\interdisplaylinepenalty=2500
\frenchspacing
\setlength{\emergencystretch}{1.5em}
\urlstyle{same}
\hypersetup{hidelinks}
\let\STAamsmathtext\text
\newcommand{\eqtext}[1]{\STAamsmathtext{\normalfont\rmfamily #1}}
\newcommand{\vectortext}{\STAamsmathtext{\normalfont\rmfamily\bfseries vector}}
\renewcommand{\text}[1]{\STAamsmathtext{\normalfont\rmfamily #1}}

\newcommand{\R}{\mathbb{R}}
\newcommand{\C}{\mathbb{C}}
% Semantic appearance hooks.  Paravectors remain visually neutral until an
% accent-free distinguishing style is chosen.  Vector and matrix wrappers
% reproduce the existing bold forms while exposing object roles in source.
\newcommand{\pv}[1]{#1}
\newcommand{\vct}[1]{\mathbf{#1}}
\newcommand{\mat}[1]{\mathbf{#1}}
\newcommand{\sigv}{\vct{\pv{\bm{\sigma}}}}
\newcommand{\bPsi}{\mat{\bm{\Psi}}}
\newcommand{\Wmat}{\mat{W}}
\newcommand{\Dmat}{\mat{D}}
\newcommand{\Tmat}{\mat{\bm{\mathsf T}}}
\newcommand{\Qmat}{\mat{Q}}
\newcommand{\Nmat}{\mat{N}}
\newcommand{\delop}{\partial}
\newcommand{\delr}{\delop_{\vct{r}}}
\newcommand{\delv}{\delop_{\vct{v}}}
\newcommand{\gradv}{\delr}
\newcommand{\herm}[1]{#1^{\mathsf H}}
\newcommand{\conj}[1]{#1^{*}}
\newcommand{\cherm}[1]{#1^{*\mathsf H}}
\newcommand{\mconj}[1]{#1^{c}}
\newcommand{\mherm}[1]{#1^{cT}}
\newcommand{\com}[2]{\operatorname{com}\!\left(#1,#2\right)}
\newcommand{\acom}[2]{\operatorname{acom}\!\left(#1,#2\right)}
\newcommand{\detf}[1]{\operatorname{det}\!\left(#1\right)}
\newcommand{\adjf}[1]{\operatorname{adj}\!\left(#1\right)}
\newcommand{\verifiedmark}{\textcolor{green!50!black}{verified}}
\newcommand{\flag}[1]{\textcolor{red!65!black}{#1}}
\newcommand{\dd}{\mathrm d}
\newcommand{\I}{\mat{I}}
\newcolumntype{C}[1]{>{\centering\arraybackslash}p{#1}}
\newcolumntype{L}[1]{>{\raggedright\arraybackslash}p{#1}}
\captionsetup{font=normalsize,labelfont=bf,labelsep=period,skip=4pt}
\newenvironment{tableblock}{%
  \par\addvspace{0.5\baselineskip}%
  \noindent\begin{minipage}{\columnwidth}%
  \captionsetup{type=table}%
  \centering
}{%
  \end{minipage}%
  \par\addvspace{0.5\baselineskip}%
}
\newenvironment{componenttable}[2][]{%
  \begin{tableblock}
  \captionof{table}{Components of $#2$.}%
  \if\relax\detokenize{#1}\relax\else\label{#1}\fi
  \renewcommand{\arraystretch}{1.10}%
  \setlength{\tabcolsep}{2pt}%
  \normalsize
  \begin{tabular}{@{}r@{:\quad}>{\raggedright\arraybackslash}p{0.48\columnwidth}>{\raggedright\arraybackslash}p{0.16\columnwidth}@{}}
  \toprule
}{%
  \bottomrule
  \end{tabular}
  \end{tableblock}
}
\newcommand{\componentrow}[3]{\eqtext{#1} & \(\displaystyle #2\) & #3 \\}
\newenvironment{sigmanumbertable}[2][]{%
  \begin{tableblock}
  \captionof{table}{#2}%
  \if\relax\detokenize{#1}\relax\else\label{#1}\fi
  \renewcommand{\arraystretch}{1.10}%
  \setlength{\tabcolsep}{3pt}%
  \normalsize
  \begin{tabular}{@{}r@{:\quad}>{\raggedright\arraybackslash}p{0.58\columnwidth}@{}}
  \toprule
}{%
  \bottomrule
  \end{tabular}
  \end{tableblock}
}
\newcommand{\sigmanumberrow}[2]{\eqtext{#1} & \(\displaystyle #2\) \\}
\renewcommand{\thesection}{\arabic{section}}
\renewcommand{\thesubsection}{\thesection.\arabic{subsection}}
\renewcommand{\thesubsubsection}{\thesubsection.\arabic{subsubsection}}
\setcounter{tocdepth}{3}
\titleformat{\section}
  {\normalfont\Large\bfseries}
  {\thesection}{0.65em}{}
\titleformat{name=\section,numberless}
  {\normalfont\Large\bfseries}
  {}{0pt}{}
\titleformat{\subsection}
  {\normalfont\large\bfseries}
  {\thesubsection}{0.65em}{}
\titleformat{name=\subsection,numberless}
  {\normalfont\large\bfseries}
  {}{0pt}{}
\titleformat{\subsubsection}[hang]
  {\normalfont\normalsize\bfseries}
  {\thesubsubsection}{0.65em}{}
\titleformat{name=\subsubsection,numberless}[hang]
  {\normalfont\normalsize\bfseries}
  {}{0pt}{}
\titlespacing*{\section}{0pt}{2.1ex plus .7ex minus .2ex}{0.9ex plus .2ex}
\titlespacing*{\subsection}{0pt}{1.7ex plus .5ex minus .2ex}{0.65ex plus .15ex}
\titlespacing*{\subsubsection}{0pt}{1.3ex plus .4ex minus .1ex}{0.45ex}
\pagestyle{plain}
\makeatletter
\newcommand{\STACompactTOC}{%
  \begingroup
  \footnotesize
  \setlength{\parskip}{0pt}%
  \setlength{\columnsep}{0.18in}%
  \begin{multicols}{2}%
  \@starttoc{toc}%
  \end{multicols}%
  \endgroup
}
\makeatother

\newcommand{\STATitle}{A Compact Paravector Form of Space-Time Algebra}
\newcommand{\STAAuthor}{Luke A. Wendt}
\newcommand{\STAAbstractText}{A notation-light paravector presentation is developed for the eight-real-dimensional algebra commonly called the real Pauli algebra or the three-dimensional real Clifford algebra. The complex structure is supplied by a central pseudoscalar, so every algebra element is written as one complex scalar plus one complex three-vector. The determinant quadratic form, spatial rotations, Lorentz boosts, relativistic kinematics, Maxwell electrodynamics, Klein--Gordon and Dirac factorizations, and the Pauli limit are derived within this single product calculus. The construction is presented as a self-contained reformulation of established flat-space physics, not as a proposal for new dynamics. Particular care is given to the distinct transformation laws for events, fields, and wave functions and to the projected one-column carrier required for one Dirac field.}
\newcommand{\STAIndexTerms}{space-time algebra, geometric algebra, Lorentz transformations, Maxwell electrodynamics, Dirac equation, Pauli equation}
\newcommand{\STAMakeTitleBlock}{%
  \begin{center}
  {\LARGE\bfseries \STATitle \par}
  \vspace{0.55em}
  {\normalsize \STAAuthor \par}
  \end{center}
  \vspace{0.55em}
  \noindent\textbf{Abstract}---\STAAbstractText
  \par\vspace{0.35em}
  \noindent\textit{Index Terms}---\STAIndexTerms.
  \par\vspace{0.7em}
}
\hypersetup{
  pdftitle={\STATitle},
  pdfauthor={\STAAuthor},
  pdfsubject={A compact paravector formulation of flat-space relativistic physics},
  pdfkeywords={space-time algebra, geometric algebra, Maxwell, Dirac, Klein-Gordon, Pauli}
}

\begin{document}
\pagenumbering{arabic}
\STAMakeTitleBlock
\thispagestyle{plain}

\section*{Introduction}\label{sec:orientation}
\addcontentsline{toc}{section}{Introduction}
A compact paravector language is developed for several standard parts of flat-space physics. Three anticommuting unit elements \(\pv{\sigma}_n\) are taken as generators. An eight-real-dimensional associative algebra is spanned by their products. In established terminology, this algebra is called the real Pauli algebra or the three-dimensional real Clifford algebra. A faithful representation is supplied by two-by-two complex matrices, and the same information is carried by the even products in the usual flat space-time algebra. These names are provided only as context; every symbol and operation used in the calculations is defined below.

By use of the central pseudoscalar \(i=\pv{\sigma}_1\pv{\sigma}_2\pv{\sigma}_3\), with \(i^2=-1\), every algebra element is written uniquely as
\[
\pv{Z}=S+\vct{V}\cdot\sigv,
\qquad S\in\C,\quad \vct{V}\in\C^3.
\]
Thus one complex scalar part and one complex three-vector part are assigned to each element. Multiplication is always taken to be the product defined from the sigma basis, not componentwise multiplication of those parts. The calculations are kept short without suppression of their physical content by retaining visible scalar/vector and real/imaginary pieces.

The presentation is organized by three involutions: sigma conjugation
\((\cdot)^*\), the Hermitian operation \((\cdot)^{\mathsf H}\), and their
composition \((\cdot)^{*\mathsf H}\). Intrinsically, sigma-product order is
reversed by \(H\); in the two-by-two matrix realization it is conjugate
transpose. The operation \(H\) is used in the positive norm, whereas
\(\pv{Z}\cherm{\pv{Z}}=S^2-\vct{V}^2\). This product is defined below as
the determinant; on real events it is the Minkowski interval. Rotations, boosts, relativistic
kinematics, Maxwell theory, and relativistic wave equations are then obtained
from the same calculus.

The physical results derived here are established results expressed in deliberately economical notation. No new interaction or change to the underlying physics is proposed; the common algebraic carrier of familiar constructions is instead exhibited explicitly. The presentation is informed by the STA and paravector literature of Hestenes, Baylis, and Doran--Lasenby \cite{hestenes1966,baylis1999,doran2003}, with historical roots in Hamilton's quaternions and Clifford's geometric algebra \cite{hamilton1844,clifford1878}.

Ordinary three-vectors are denoted by bold symbols. In natural units, the
reduced Planck constant and the speed of light are set equal to \(1\). Rationalized
Heaviside--Lorentz electromagnetic units and metric signature
\((+,-,-,-)\) are used throughout. The product formula for scalar
coefficients is extended directly to complex-valued functions; when differential
operators occur, their displayed order and the Leibniz rule are part of the
definition. In an approximation, terms of order \(x^n\) or smaller in the
stated limiting parameter are denoted by \(O(x^n)\).

\subsection*{Reading Guide}
\begingroup
\footnotesize
\setlength{\columnsep}{0.22in}
\begin{multicols}{2}
\begin{enumerate}
\setlength{\itemsep}{0.08em}
\setlength{\parsep}{0pt}
\setlength{\topsep}{0.15em}
\setlength{\partopsep}{0pt}
\item The principal definitions and results are collected on the Equation Reference Sheet on the next page in the order of their construction.
\item The identities, sigma basis, three involutions, complex paravectors, extraction functions, vector and paravector products, positive norm, adjugate, determinant, inverse, projectors, elementary functions, transform laws, quaternion correspondence, and matrix realizations are defined in Section \ref{sec:algebra}.
\item Real paravectors are specialized to events, and intervals, velocity, acceleration, energy, momentum, action, and the nonrelativistic limit are developed in Section \ref{sec:spacetime}.
\item The event law is applied to Lorentz boosts, length contraction, momentum, and boost--rotor compositions in Section \ref{sec:transforms}, where it is contrasted with the field similarity law.
\item Spatial and paravector derivatives, complex fields, the wave operator, and differentiation rules for noncommuting products are defined in Section \ref{sec:differential}.
\item Fields and potentials, all four Maxwell equations, boost laws, gauge freedom, field and potential wave equations, the Lorentz force, energy and flux, and the full linear stress-energy map are developed in Section \ref{sec:em}.
\item Scalar waves, the positive- and negative-frequency split, minimal electromagnetic coupling, the ordered Klein--Gordon factorization, the projected two-block Dirac field, its conserved pairing and gauge law, the squared Dirac system, and the controlled Pauli and helicity reductions are developed in Section \ref{sec:qm}.
\item The established results are separated from explicitly labeled open extensions and points requiring sharper foundations in Sections \ref{sec:conclusions} and \ref{sec:future-work}.
\item Independent symbolic checks and standard correspondences are recorded in Appendix \ref{sec:verification}; historical context is provided by the references without any change of notation.
\end{enumerate}
\end{multicols}
\endgroup

\clearpage
\thispagestyle{plain}
\begin{center}
\textbf{Equation Reference Sheet (Summary)}
\end{center}
\vspace{-0.65em}
\begingroup
\fontsize{7.75}{7.80}\selectfont
\setlength{\jot}{0.12em}
\setlength{\tabcolsep}{0pt}
\renewcommand{\arraystretch}{1.00}
\newcommand{\refhead}[1]{\multicolumn{3}{@{}l}{\textbf{#1}}\\[0.10em]}
\newcommand{\refrow}[3]{#1 & \(\textstyle #2\) & #3\\[0.04em]}
\newcommand{\refsub}[3]{\hspace*{0.9em}#1 & \hspace*{0.5em}\(\textstyle #2\) & #3\\[0.03em]}
\noindent Natural units: reduced Planck constant \(=1\), speed of light \(=1\).
Signature: \((+,-,-,-)\).
All space at fixed time is covered by unmarked spatial integrals with measure
\(\dd r_1\,\dd r_2\,\dd r_3\).
\par\vspace{0.12em}
\noindent
\begin{tabular}{@{}p{0.245\textwidth}@{\hspace{0.012\textwidth}}p{0.647\textwidth}@{\hspace{0.010\textwidth}}>{\raggedleft\arraybackslash}p{0.076\textwidth}@{}}
\textbf{Name} & \textbf{Formula} & \textbf{Ref.}\\[0.10em]
\refhead{Custom functions, maps, and identities}
\refrow{Identity elements}{\pv{Z}+0=\pv{Z},\quad \pv{Z}1=1\pv{Z}=\pv{Z}}{\eqref{eq:algebra-identities}}
\refrow{Complex paravector}{\pv{Z}=(a+ib)+(\vct{A}+i\vct{B})\cdot\sigv=S+\vct{V}\cdot\sigv=\pv{X}+i\pv{Y}}{\eqref{eq:paravector-cartesian}--\eqref{eq:paravector}}
\refsub{Real and img functions}{\eqtext{real}(\pv{Z})=\pv{X},\quad \eqtext{img}(\pv{Z})=\pv{Y}}{\eqref{eq:real-function}, \eqref{eq:img-function}}
\refsub{Scalar and vector functions}{\eqtext{scalar}(\pv{Z})=S,\quad \eqtext{vector}(\pv{Z})=\vct{V}}{\eqref{eq:scalar-function}, \eqref{eq:vector-function}}
\refsub{Axis functions}{\eqtext{par}(\vct{A})=(\vct{A}\cdot\vct{u})\vct{u},\quad \eqtext{prp}(\vct{A})=\vct{A}-\eqtext{par}(\vct{A}),\quad \vct{u}^2=1}{\eqref{eq:par-function}, \eqref{eq:prp-function}}
\refsub{Order functions}{\com{\pv{Z}}{\pv{Z}'}=\pv{Z}\pv{Z}'-\pv{Z}'\pv{Z},\quad \acom{\pv{Z}}{\pv{Z}'}=\pv{Z}\pv{Z}'+\pv{Z}'\pv{Z}}{\eqref{eq:com-function}, \eqref{eq:acom-function}}
\refsub{Adjugate function}{\adjf{\pv{Z}}=\cherm{\pv{Z}}=S-\vct{V}\cdot\sigv,\quad \herm{\pv{Z}}=\conj{\adjf{\pv{Z}}}=\adjf{\conj{\pv{Z}}}}{\eqref{eq:adjugate}, \eqref{eq:hermitian-adjugate-identity}}
\refsub{Determinant function}{\detf{\pv{Z}}=\adjf{\pv{Z}}\pv{Z}=\pv{Z}\adjf{\pv{Z}}=S^2-\vct{V}^2}{\eqref{eq:determinant-function}, \eqref{eq:minkowski-form}}
\refsub{Inverse}{\pv{Z}^{-1}=\adjf{\pv{Z}}/\detf{\pv{Z}},\quad \detf{\pv{Z}}\ne0}{\eqref{eq:paravector-inverse}}
\refsub{Matrix functions}{\begin{gathered}\I=\{1,0;0,1\},\quad \eqtext{tr}(\{a,b;c,d\})=a+d,\quad \eqtext{diag}(A,B)=\{A,0;0,B\},\\[-0.1em]\{a,b;c,d\}^{T}=\{a,c;b,d\},\quad \mconj{\mat{M}}=\eqtext{entrywise conjugate},\quad \mherm{\mat{M}}=(\mconj{\mat{M}})^T\end{gathered}}{\eqref{eq:matrix-identity}--\eqref{eq:matrix-hermitian}}
\refsub{Auxiliary block function}{\Wmat(\pv{Z})=\{0,\pv{Z};\adjf{\pv{Z}},0\}}{\eqref{eq:W-map}}
\refhead{Algebraic progression}
\refrow{Pseudoscalar}{i=\pv{\sigma}_1\pv{\sigma}_2\pv{\sigma}_3,\quad i^2=-1,\quad i\pv{\sigma}_n=\pv{\sigma}_n i}{\eqref{eq:basis}}
\refsub{Sigma product}{\pv{\sigma}_n\pv{\sigma}_m=\delta_{nm}+\sum_{k=1}^{3}\epsilon_{nmk}i\pv{\sigma}_k}{\eqref{eq:sigma-product}}
\refsub{Sigma/H operations}{\conj{\pv{X}\pv{Y}}=\conj{\pv{X}}\conj{\pv{Y}},\quad \herm{\pv{X}\pv{Y}}=\herm{\pv{Y}}\herm{\pv{X}},\quad \cherm{\pv{X}\pv{Y}}=\cherm{\pv{Y}}\cherm{\pv{X}}}{\eqref{eq:operator-rules}}
\refsub{Involution actions}{\conj{\pv{Z}}=(a-ib)+(-\vct{A}+i\vct{B})\cdot\sigv,\quad \herm{\pv{Z}}=(a-ib)+(\vct{A}-i\vct{B})\cdot\sigv,\quad \cherm{\pv{Z}}=S-\vct{V}\cdot\sigv}{\eqref{eq:paravector-conjugate}--\eqref{eq:paravector-reverse-conjugate}}
\refsub{Matrix bridge: \(*\ne c\)}{\Tmat=\{0,1;-1,0\},\quad \conj{\mat{M}}=\Tmat\mconj{\mat{M}}\Tmat^{-1},\quad \herm{\mat{M}}=\mherm{\mat{M}},\quad \cherm{\mat{M}}=\adjf{\mat{M}}=\Tmat\mat{M}^{T}\Tmat^{-1}}{\eqref{eq:matrix-T}, \eqref{eq:matrix-operation-dictionary}}
\refsub{Vector product}{(\vct{A}\cdot\sigv)(\vct{B}\cdot\sigv)=\vct{A}\cdot\vct{B}+i(\vct{A}\times\vct{B})\cdot\sigv}{\eqref{eq:dot-cross}}
\refsub{Paravector product}{\pv{Z}\pv{Z}'=SS'+\vct{V}\cdot\vct{V}'+(S\vct{V}'+S'\vct{V}+i\vct{V}\times\vct{V}')\cdot\sigv}{\eqref{eq:paravector-product}}
\refsub{Positive norms}{|S|^2=SS^*,\quad\|\vct{V}\|^2=\vct{V}\cdot\vct{V}^*,\quad\|\pv{Z}\|^2=\eqtext{scalar}(\pv{Z}\herm{\pv{Z}})=|S|^2+\|\vct{V}\|^2}{\eqref{eq:scalar-norm}--\eqref{eq:euclidean-form}}
\refsub{Determinant products}{\detf{\pv{Z}\pv{Z}'}=\detf{\pv{Z}}\detf{\pv{Z}'},\quad \detf{\adjf{\pv{Z}}}=\detf{\pv{Z}}}{\eqref{eq:metric-product}, \eqref{eq:adjugate-identities}}
\refrow{Axis projectors}{\pv{\Pi}(\pm)=\tfrac12(1\pm\vct{u}\cdot\sigv),\quad \pv{Z}(\pm)=\pv{\Pi}(\pm)\pv{Z},\quad \pv{Z}=\pv{Z}(+)+\pv{Z}(-)}{\eqref{eq:projectors}, \eqref{eq:projected-part-function}}
\refsub{Ladder function}{\pv{a}(\pm)=\tfrac12(\pv{\sigma}_1\pm i\pv{\sigma}_2),\quad \pv{a}(\pm)^2=0,\quad \com{\pv{\sigma}_3}{\pv{a}(\pm)}=\pm2\pv{a}(\pm)}{\eqref{eq:ladder}}
\refsub{Functional calculus}{f(\vct{u}\cdot\sigv)=f(+)\pv{\Pi}(+)+f(-)\pv{\Pi}(-)}{\eqref{eq:functional-calculus}}
\refsub{Vector exponential}{\exp(\theta\vct{u}\cdot\sigv)=\cosh(\theta)+\sinh(\theta)\vct{u}\cdot\sigv}{\eqref{eq:exp-vector}}
\refsub{Bivector exponential}{\exp(\theta\vct{u}\cdot i\sigv)=\cos(\theta)+\sin(\theta)\vct{u}\cdot i\sigv}{\eqref{eq:exp-bivector}}
\refsub{Paravector logarithm}{\log(\pv{Y})=\log(\lambda(\pv{Y},+))\pv{\Pi}(+)+\log(\lambda(\pv{Y},-))\pv{\Pi}(-)}{\eqref{eq:log-paravector}}
\refhead{Transforms and kinematics}
\refrow{Similarity transport}{\pv{T}\pv{Z}'=\pv{Z}\pv{T}\quad\eqtext{implies}\quad\pv{Z}'=\pv{T}^{-1}\pv{Z}\pv{T}}{\eqref{eq:transform-intertwining}, \eqref{eq:general-transform}}
\refsub{Event congruence}{\pv{X}'=\pv{T}^{*\mathsf H}\pv{X}\pv{T}^*,\quad \detf{\pv{T}}=1}{\eqref{eq:general-event-congruence}}
\refsub{Active rotor}{\pv{R}=\exp(\theta\vct{u}\cdot i\sigv/2),\quad \pv{X}'=\herm{\pv{R}}\pv{X}\pv{R}}{\eqref{eq:rotor-element}, \eqref{eq:rotor}}
\refsub{Passive event boost}{\pv{L}=\exp(\theta\vct{u}\cdot\sigv/2),\quad \pv{X}'=\conj{\pv{L}}\pv{X}\conj{\pv{L}}}{\eqref{eq:boost-element}, \eqref{eq:boost}}
\refsub{Passive field boost}{\pv{F}'=\conj{\pv{L}}\pv{F}\pv{L}}{\eqref{eq:field-transform}}
\refrow{Space-time interval}{\dd s^2=\detf{\dd\pv{X}}=\dd t^2-\dd\vct{r}^{\,2}}{\eqref{eq:interval}}
\refsub{Velocity and momentum}{\pv{U}=\gamma(1+\vct{v}\cdot\sigv),\quad m\pv{U}=E+\vct{P}\cdot\sigv}{\eqref{eq:proper-velocity}, \eqref{eq:proper-momentum}}
\refsub{Acceleration}{\pv{A}=\dd\pv{U}/\dd s,\quad \eqtext{scalar}(\pv{A}\conj{\pv{U}})=0}{\eqref{eq:proper-acceleration}, \eqref{eq:proper-acceleration-orthogonality}}
\refhead{Differential and electromagnetic fields}
\refrow{Spatial derivatives}{\gradv=\{\partial_1,\partial_2,\partial_3\},\quad \partial_n=\partial/\partial r_n,\quad \gradv^2=\partial_1^2+\partial_2^2+\partial_3^2}{\eqref{eq:gradient-notation}, \eqref{eq:laplacian}}
\refrow{Paravector derivatives}{\pv{\partial}=\partial_t+\sigv\cdot\gradv,\quad \conj{\pv{\partial}}=\partial_t-\sigv\cdot\gradv}{\eqref{eq:partial-operator}, \eqref{eq:partial-conjugate}}
\refsub{Wave operator}{\Box=\pv{\partial}\conj{\pv{\partial}}=\conj{\pv{\partial}}\pv{\partial}=\partial_t^2-\gradv^2}{\eqref{eq:dalembertian}}
\refsub{Ordered product rule}{\pv{\partial}(\pv{Z}\pv{Z}')=(\pv{\partial}\pv{Z})\pv{Z}'+\pv{Z}(\pv{\partial}\pv{Z}')+\sum_n\com{\pv{\sigma}_n}{\pv{Z}}\,\partial_n\pv{Z}'}{\eqref{eq:paravector-leibniz}}
\refrow{Field and potential}{\pv{F}=(\vct{E}+i\vct{B})\cdot\sigv,\quad \pv{\Phi}=V+\vct{A}\cdot\sigv}{\eqref{eq:F-Phi}, \eqref{eq:potential-paravector}}
\refsub{Potential gradient}{\pv{\partial}\pv{\Phi}=S+\conj{\pv{F}},\quad S=\partial_tV+\gradv\cdot\vct{A}}{\eqref{eq:potential-scalar}, \eqref{eq:potential-gradient}}
\refsub{Maxwell equation}{\pv{\partial}\pv{F}=\rho-\vct{J}\cdot\sigv}{\eqref{eq:maxwell}}
\refsub{Maxwell components}{\gradv\cdot\vct{E}=\rho,\ \gradv\cdot\vct{B}=0,\ \gradv\times\vct{B}=\vct{J}+\partial_t\vct{E},\ \gradv\times\vct{E}=-\partial_t\vct{B}}{\eqref{eq:maxwell-scalars}, \eqref{eq:maxwell-vectors}}
\refsub{Gauge functions}{\pv{\Phi}(\lambda)=\pv{\Phi}+\conj{\pv{\partial}}\lambda,\quad \pv{F}(\lambda)=\pv{F}}{\eqref{eq:gauge}, \eqref{eq:gauge-field}}
\refsub{Lorentz force}{m\,\dd\pv{U}/\dd s=q\,\eqtext{real}(\pv{F}\pv{U})}{\eqref{eq:lorentz-force}}
\refsub{Energy and flux}{\pv{T}=\tfrac12\pv{F}\herm{\pv{F}}=\tfrac12(\vct{E}^2+\vct{B}^2)+(\vct{E}\times\vct{B})\cdot\sigv}{\eqref{eq:energy-momentum-paravector}}
\refsub{Full stress-energy function}{\pv{T}(\pv{a})=\tfrac12\pv{F}\pv{a}\herm{\pv{F}},\quad \pv{a}\in\R+\R^3,\quad \pv{T}=\pv{T}(1)}{\eqref{eq:stress-energy-map}}
\refhead{Relativistic wave equations}
\refrow{Frequency split}{\Omega=\sqrt{m^2-\gradv^2},\quad(\partial_t\pm i\Omega)\psi(\pm)=0}{\eqref{eq:frequency-operator}, \eqref{eq:frequency-split}}
\refsub{Covariant energy/momentum}{E=i\partial_t-qV,\quad \vct{P}=-i\gradv-q\vct{A}}{\eqref{eq:energy-operator}, \eqref{eq:momentum-operator}}
\refsub{Klein--Gordon}{(-E^2+\vct{P}^2+m^2)\psi=0}{\eqref{eq:kg-coupled}}
\refsub{Ordered factors}{(E+\sigv\cdot\vct{P})(E-\sigv\cdot\vct{P})=E^2-\vct{P}^2-iq\pv{F}}{\eqref{eq:EP-factorization}}
\refsub{Dirac system}{(\Wmat(i\pv{\partial}-q\conj{\pv{\Phi}})-m\I)\bPsi=0}{\eqref{eq:dirac-coupled}}
\refsub{Fixed-column carrier}{\pv{\Pi}(\eqtext{column})=\tfrac12(1+\pv{\sigma}_3),\quad \bPsi_n=\bPsi_n\pv{\Pi}(\eqtext{column}),\quad n\in\{1,2\}}{\eqref{eq:fixed-column-carrier}, \eqref{eq:fixed-column-constraint}}
\refsub{Dirac carrier transform}{\bPsi_1'(\pv{X}')=\pv{T}^{\mathsf H}\bPsi_1(\pv{X}),\quad \bPsi_2'(\pv{X}')=\pv{T}^{-1}\bPsi_2(\pv{X})}{\eqref{eq:dirac-transform}, \eqref{eq:dirac-transform-lower}}
\refsub{Squared blocks}{(-E^2+\vct{P}^2+m^2+iq\conj{\pv{F}})\bPsi_1=0,\quad(-E^2+\vct{P}^2+m^2+iq\pv{F})\bPsi_2=0}{\eqref{eq:dirac-square}}
\refsub{Dirac current}{\pv{J}(\bPsi)=2\{\bPsi_2\herm{\bPsi_2}+\conj{(\bPsi_1\herm{\bPsi_1})}\}=\rho(\bPsi)+\vct{j}(\bPsi)\cdot\sigv,\quad \partial_t\rho(\bPsi)+\gradv\cdot\vct{j}(\bPsi)=0}{\eqref{eq:dirac-current-paravector}}
\refsub{Klein--Gordon pairing}{\rho(\phi,\psi)=\{\phi^*E\psi+(E\phi)^*\psi\}/(2m),\quad \langle\phi\mid\psi\rangle=\iiint\rho(\phi,\psi)\,\dd r_1\,\dd r_2\,\dd r_3}{\eqref{eq:kg-mixed-current}, \eqref{eq:kg-pairing}}
\refsub{Projected-wave pairings}{\begin{gathered}\langle\pv{\phi}\mid\pv{\psi}\rangle=2\iiint\eqtext{scalar}(\herm{\pv{\phi}}\pv{\psi})\,\dd r_1\,\dd r_2\,\dd r_3,\\[-0.1em]\langle\bPsi(\eqtext{left})\mid\bPsi(\eqtext{right})\rangle=2\iiint\sum_{n=1}^{2}\eqtext{scalar}\!\left(\herm{\bPsi_n(\eqtext{left})}\bPsi_n(\eqtext{right})\right)\dd r_1\,\dd r_2\,\dd r_3\end{gathered}}{\eqref{eq:adjoint-paravector-braket}, \eqref{eq:dirac-pairing}}
\refsub{Leading Pauli limit}{E\pv{\phi}(\eqtext{large})=(\vct{P}^2-q\vct{B}\cdot\sigv)\pv{\phi}(\eqtext{large})/(2m)+O(m^{-2})}{\eqref{eq:pauli-operator}}
\refsub{Pauli spin split}{\pv{\phi}(\eqtext{large},\pm)=\pv{\Pi}(\pm)\pv{\phi}(\eqtext{large}),\quad i\partial_t\pv{\phi}(\eqtext{large},\pm)=\{\vct{P}^2\mp q\|\vct{B}\|\}\pv{\phi}(\eqtext{large},\pm)/(2m)+qV\pv{\phi}(\eqtext{large},\pm)}{\eqref{eq:projectors}, \eqref{eq:pauli-spin-decomposition}, \eqref{eq:pauli-spin-split}}
\refsub{Free energy/axis split}{\begin{gathered}\mat{\mathcal H}\{\bPsi_1;\bPsi_2\}=\{-(\vct{p}\cdot\sigv)\bPsi_1+m\bPsi_2;\ m\bPsi_1+(\vct{p}\cdot\sigv)\bPsi_2\},\quad E(\vct{p})=\sqrt{m^2+\vct{p}^2},\\[-0.1em]\bPsi(e)=\tfrac12\{\bPsi+e\mat{\mathcal H}\bPsi/E(\vct{p})\},\quad \bPsi(e,s)=\{\pv{\Pi}(s)\bPsi_1(e);\pv{\Pi}(s)\bPsi_2(e)\}\end{gathered}}{\eqref{eq:free-energy-function}--\eqref{eq:energy-spin-projectors}}
\end{tabular}
\endgroup
\clearpage
\thispagestyle{plain}
\phantomsection
\pdfbookmark[1]{Contents}{contents}
\begin{center}
\textbf{Contents}
\end{center}
\vspace{-0.55em}
\STACompactTOC
\clearpage

\twocolumn
\section{Algebraic Foundations}\label{sec:algebra}\label{sec:foundations}\label{sec:functions}
The algebra is fixed before physical meanings are assigned to particular elements. The product is associative and unital. All coefficients in the master product formula are taken to be central complex scalars unless an ordered differential-operator calculation is stated explicitly.

\subsection{Basis and pseudoscalar}
The element \(0\) is the additive identity and \(1\) is the multiplicative identity:
\begin{equation}
\pv{Z}+0=0+\pv{Z}=\pv{Z},
\qquad
\pv{Z}1=1\pv{Z}=\pv{Z}.
\label{eq:algebra-identities}
\end{equation}
The two identities must not be confused: an inverse is defined by a product equal to \(1\), whereas the absence of a multiplicative inverse is indicated by a vanishing determinant.

For \(n\in\{1,2,3\}\), the algebra is generated over \(\R\) by three basis elements \(\pv{\sigma}_n\). Scalars, vectors, two-generator elements, and a three-generator pseudoscalar are generated by their products. The resulting eight-real-dimensional algebra is commonly called the real Pauli algebra or the three-dimensional real Clifford algebra.
The basis relations are
\begin{equation}
\begin{aligned}
\pv{\sigma}_n^2&=1,\qquad \pv{\sigma}_n\pv{\sigma}_m=-\pv{\sigma}_m\pv{\sigma}_n\quad (n\neq m),\\
i&:=\pv{\sigma}_1\pv{\sigma}_2\pv{\sigma}_3,\qquad i^2=-1.
\end{aligned}
\label{eq:basis}
\end{equation}
The pseudoscalar classification is established by \(i\pv{\sigma}_n=\pv{\sigma}_n i\), and the scalar \(-1\) is obtained by squaring \(i\):
\begin{align*}
(\pv{\sigma}_1\pv{\sigma}_2\pv{\sigma}_3)^2
&= \pv{\sigma}_1\pv{\sigma}_2(\pv{\sigma}_3\pv{\sigma}_1)\pv{\sigma}_2\pv{\sigma}_3,\notag\\
&= -\pv{\sigma}_1(\pv{\sigma}_2\pv{\sigma}_1)(\pv{\sigma}_3\pv{\sigma}_2)\pv{\sigma}_3,\notag\\
&= -\pv{\sigma}_1^2\pv{\sigma}_2^2\pv{\sigma}_3^2,\notag\\
&= -1.
\end{align*}
With the Kronecker delta and Levi-Civita symbol,
\begin{equation}
\pv{\sigma}_n\pv{\sigma}_m=\delta_{nm}+\sum_{k=1}^3\epsilon_{nmk}\,i\pv{\sigma}_k,
\label{eq:sigma-product}
\end{equation}
where
\begin{align*}
\delta_{nm} &:= \begin{cases}
1,& n=m,\\
0,& n\neq m,
\end{cases}\\[0.25em]
\epsilon_{nmk} &:= \begin{cases}
+1,& (n,m,k)\text{ even permutation of }(1,2,3),\\
-1,& (n,m,k)\text{ odd permutation of }(1,2,3),\\
0,& \text{otherwise}.
\end{cases}
\end{align*}

\subsection{The Conjugate and Hermitian Operators}
The sigma-conjugate operator is defined by
\begin{equation}
(\cdot)^* := \eqtext{flip the signs of } \pv{\sigma}_n.
\label{eq:star-definition}
\end{equation}
The Hermitian operator is defined by
\begin{equation}
(\cdot)^{\mathsf H} := \eqtext{reverse the order of multiplied } \pv{\sigma}_n.
\label{eq:reversion-definition}
\end{equation}
The name Hermitian is justified by the matrix realization developed in
Section \ref{sec:matrix-realization}: there \(H\) is represented by ordinary
entrywise complex conjugation followed by transpose. Within the intrinsic
algebra it is defined without matrices, solely by reversing sigma-product
order. The star is a real-linear automorphism and \(H\) is a real-linear
anti-automorphism. Because the central pseudoscalar \(i\) is sent to
\(-i\) by each operation, both are antilinear with respect to the internal complex structure.
Their composition is a complex-linear anti-automorphism. In operator form,
\begin{equation}
\begin{aligned}
{(\pv{X}\pv{Y})}^{*}&=\pv{X}^{*}\pv{Y}^{*},\qquad
{(\pv{X}\pv{Y})}^{\mathsf H}=\pv{Y}^{\mathsf H}\pv{X}^{\mathsf H},\\
\pv{X}^{*\mathsf H}&=(\pv{X}^{\mathsf H})^{*}=(\pv{X}^{*})^{\mathsf H},\qquad
(\pv{X}^*)^*=(\pv{X}^{\mathsf H})^{\mathsf H}=\pv{X}.
\end{aligned}
\label{eq:operator-rules}
\end{equation}
When applied to the generators and the pseudoscalar, these rules are reduced to
\begin{alignat}{3}
\conj{\pv{\sigma}_n} &{}= -\pv{\sigma}_n, \qquad&
\herm{\pv{\sigma}_n} &{}= \pv{\sigma}_n, \qquad&
\cherm{\pv{\sigma}_n} &{}= -\pv{\sigma}_n,
\label{eq:sigma-operator-signs}\\
\conj{i} &{}= -i, \qquad&
\herm{i} &{}= -i, \qquad&
\cherm{i} &{}= i.
\label{eq:i-operator-signs}
\end{alignat}
For a two-generator bivector pair with \(n\neq m\),
\begin{align}
(\pv{\sigma}_n\pv{\sigma}_m)^*
&= \pv{\sigma}_n^*\pv{\sigma}_m^*
= (-\pv{\sigma}_n)(-\pv{\sigma}_m)
= \pv{\sigma}_n\pv{\sigma}_m,\notag\\
(\pv{\sigma}_n\pv{\sigma}_m)^{\mathsf H}
&= \pv{\sigma}_m^{\mathsf H}\pv{\sigma}_n^{\mathsf H}
= \pv{\sigma}_m\pv{\sigma}_n
= -\pv{\sigma}_n\pv{\sigma}_m,\notag\\
(\pv{\sigma}_n\pv{\sigma}_m)^{*\mathsf H}
&= (\pv{\sigma}_m\pv{\sigma}_n)^*
= \pv{\sigma}_m^*\pv{\sigma}_n^* \notag\\
&= (-\pv{\sigma}_m)(-\pv{\sigma}_n)
= \pv{\sigma}_m\pv{\sigma}_n
= -\pv{\sigma}_n\pv{\sigma}_m.
\label{eq:bivector-operator-signs}
\end{align}
The same sign pattern on the sigma-number-two elements is recorded in
Eq.\ \eqref{eq:bivector-operator-signs}. These are the standard intrinsic involutions
of the real Pauli algebra. No assertion that every \(\pv{Z}^{\mathsf H}\) is
itself a Hermitian object is implied by the symbol \(H\); rather, a represented
matrix is Hermitian precisely when \(\pv{Z}^{\mathsf H}=\pv{Z}\). In the
conventional even space-time interpretation, the role of reversing
space-time products is assigned to the combined operation
\((\cdot)^{*\mathsf H}\), rather than to \((\cdot)^{\mathsf H}\) alone
\cite{hestenes1966,doran2003}.
No additional notation or rule is introduced by this contextual statement.

The sigma number of a reduced basis product is defined as the number of sigma generators that remain after simplification. Sigma numbers \(0\), \(1\), \(2\), and \(3\) are assigned respectively to \(1\), \(\pv{\sigma}_n\), \(\pv{\sigma}_n\pv{\sigma}_m\) with \(n\neq m\), and \(i\). Any unreduced product with sigma number larger than \(3\) is simplified to one of these four cases. The basis elements and operator signs are summarized by sigma number in Tables \ref{tab:sigma-number-components} and \ref{tab:sigma-number-signs}.
\noindent These tables list the reduced basis elements and involution signs by sigma number.
\begin{tableblock}
\captionsetup{type=table,font=normalsize,labelfont=bf,skip=4pt}
\captionof{table}{Sigma basis elements by sigma number.}\label{tab:sigma-number-components}
\renewcommand{\arraystretch}{1.12}
\setlength{\tabcolsep}{4pt}
\normalsize
\begin{tabularx}{\columnwidth}{@{}l >{\raggedright\arraybackslash}X@{}}
\toprule
sigma number & representative basis elements \\
\midrule
\(0\) (scalar) & \(1\) \\
\(1\) (vector) & \(\{\sigma_1,\ \sigma_2,\ \sigma_3\}\) \\
\(2\) (bivector) & \(\{\sigma_2\sigma_3,\ \sigma_3\sigma_1,\ \sigma_1\sigma_2\}\) \\
\(3\) (pseudoscalar) & \(\sigma_1\sigma_2\sigma_3\) \\
\bottomrule
\end{tabularx}
\end{tableblock}

\begin{tableblock}
\captionsetup{type=table,font=normalsize,labelfont=bf,skip=4pt}
\captionof{table}{Involution signs by sigma number.}\label{tab:sigma-number-signs}
\renewcommand{\arraystretch}{1.12}
\setlength{\tabcolsep}{4pt}
\normalsize
\begin{tabularx}{\columnwidth}{@{}l >{\raggedright\arraybackslash}X >{\raggedright\arraybackslash}X >{\raggedright\arraybackslash}X@{}}
\toprule
sigma number & \((\cdot)^*\) & \((\cdot)^{\mathrm R}\) & \((\cdot)^{*\mathrm R}\) \\
\midrule
\(0\) (scalar) & \(+\) & \(+\) & \(+\) \\
\(1\) (vector) & \(-\) & \(+\) & \(-\) \\
\(2\) (bivector) & \(+\) & \(-\) & \(-\) \\
\(3\) (pseudoscalar) & \(-\) & \(-\) & \(+\) \\
\bottomrule
\end{tabularx}
\end{tableblock}


\subsection{Generic Complex Paravector}
The three generators are collected into the sigma-vector
\begin{equation}
\sigv := \{ \pv{\sigma}_1, \pv{\sigma}_2, \pv{\sigma}_3 \},
\label{eq:sigma-vector}
\end{equation}
which is a vector of algebraic objects rather than a vector of scalar components. The associated bivector triple is obtained by multiplication by the pseudoscalar:
\begin{equation}
i\sigv := \{ i\pv{\sigma}_1,\ i\pv{\sigma}_2,\ i\pv{\sigma}_3 \}
= \{ \pv{\sigma}_2\pv{\sigma}_3,\ \pv{\sigma}_3\pv{\sigma}_1,\ \pv{\sigma}_1\pv{\sigma}_2 \},
\label{eq:isigma-vector}
\end{equation}
because \(i=\pv{\sigma}_1\pv{\sigma}_2\pv{\sigma}_3\) and each product \(i\pv{\sigma}_n\) is reduced to the complementary ordered two-generator product. The canonical sigma-number-two basis built from the generators is therefore recorded in Eq.\ \eqref{eq:isigma-vector}.
The descriptor \(\R+\R^3\) is used for one real scalar plus one real
three-vector, and \(\C+\C^3\) is used for one complex scalar plus one complex
three-vector. This decomposition is recorded by the plus sign in a type
descriptor; no second multiplication rule is introduced.
With this basis fixed, for real scalars \(a,b\in\R\) and real vectors \(\vct{A},\vct{B}\in\R^3\), with \(\vct{A}\cdot\sigv=A_1\pv{\sigma}_1+A_2\pv{\sigma}_2+A_3\pv{\sigma}_3\), the real paravectors \(\pv{X},\pv{Y}\in\R+\R^3\) are defined by
\begin{equation}
\pv{X}:=a+\vct{A}\cdot\sigv,
\label{eq:real-paravector-x}
\end{equation}
\begin{equation}
\pv{Y}:=b+\vct{B}\cdot\sigv.
\label{eq:real-paravector-y}
\end{equation}
The complex scalar and vector components are defined first:
\begin{equation}
S:=a+ib\in\C.
\label{eq:paravector-scalar}
\end{equation}
\begin{equation}
\vct{V}:=\vct{A}+i\vct{B}\in\C^3.
\label{eq:paravector-vector}
\end{equation}
Every algebra element can then be defined as a complex paravector
\(\pv{Z}\in\C+\C^3\), written as a real paravector plus an imaginary
paravector:
\begin{equation}
\pv{Z}:=\pv{X}+i\pv{Y}.
\label{eq:paravector-xy}
\end{equation}
The equivalent Cartesian and scalar--vector forms are
\begin{equation}
\pv{Z}=(a+ib)+(\vct{A}+i\vct{B})\cdot\sigv.
\label{eq:paravector-cartesian}
\end{equation}
\begin{equation}
\pv{Z}=S+\vct{V}\cdot\sigv.
\label{eq:paravector}
\end{equation}
On coefficient objects, the induced actions of either basic involution are
\(
S^*=S^{\mathsf H}=a-ib
\)
and
\(
\vct{V}^*=\vct{V}^{\mathsf H}=\vct{A}-i\vct{B}
\).
When the full algebra term \(\vct{V}\cdot\sigv\) is acted upon by an
involution, its action on the sigma basis is included as well,
according to Eqs.\ \eqref{eq:operator-rules} and
\eqref{eq:sigma-operator-signs}.
The real function is defined by
\begin{equation}
\eqtext{real}(\pv{Z}):=\pv{X}=\frac{\pv{Z}+\herm{\pv{Z}}}{2}.
\label{eq:real-function}
\end{equation}
The imaginary function is defined by
\begin{equation}
\eqtext{img}(\pv{Z}):=\pv{Y}=\frac{\pv{Z}-\herm{\pv{Z}}}{2i}.
\label{eq:img-function}
\end{equation}
The scalar function is defined by
\begin{equation}
\eqtext{scalar}(\pv{Z}):=S=\frac{\pv{Z}+\cherm{\pv{Z}}}{2}.
\label{eq:scalar-function}
\end{equation}
The vector function is defined by
\begin{equation}
\eqtext{vector}(\pv{Z}):=\vct{V}.
\label{eq:vector-function}
\end{equation}
The algebra element carried by that vector is recovered as
\begin{equation}
\eqtext{vector}(\pv{Z})\cdot\sigv
=\frac{\pv{Z}-\cherm{\pv{Z}}}{2}.
\label{eq:vector-recovery}
\end{equation}
The actions of the conjugate and Hermitian operators on the generic complex paravector are
\begin{align}
\conj{\pv{Z}} &= (a-ib)+(-\vct{A}+i\vct{B})\cdot\sigv,
\label{eq:paravector-conjugate}\\
\herm{\pv{Z}} &= \pv{X}-i\pv{Y}=(a-ib)+(\vct{A}-i\vct{B})\cdot\sigv,
\label{eq:paravector-reverse}\\
\cherm{\pv{Z}} &= (a+ib)-(\vct{A}+i\vct{B})\cdot\sigv.
\label{eq:paravector-reverse-conjugate}
\end{align}
For the real paravector \(\pv{X}\) from Eq.\ \eqref{eq:real-paravector-x}, these are reduced to
\begin{align}
\conj{\pv{X}} &= a-\vct{A}\cdot\sigv,
\label{eq:real-paravector-conjugate}\\
\herm{\pv{X}} &= \pv{X}=a+\vct{A}\cdot\sigv,
\label{eq:real-paravector-reverse}\\
\cherm{\pv{X}} &= a-\vct{A}\cdot\sigv=\conj{\pv{X}}.
\label{eq:real-paravector-combined}
\end{align}
Equations \eqref{eq:paravector-conjugate}--\eqref{eq:paravector-reverse-conjugate}
are the generic formulas used below. In particular, conjugation
of the complex coefficients is already encoded by \(H\); no overbar
notation is needed. The matrix-only superscript \(c\), introduced later, is
therefore not a fourth intrinsic operation. Equations
\eqref{eq:real-paravector-conjugate}--\eqref{eq:real-paravector-combined}
are the real specialization.
The real and complex paravector pieces and extractors used throughout the later algebra are fixed by Eqs.\ \eqref{eq:real-paravector-x}--\eqref{eq:vector-function}. The type assignments are recorded in Table \ref{tab:paravector-types}, the recovery formulas for the extractor definitions are recorded in Table \ref{tab:components-paravector-recovery}, and the same split is recorded in components in Table \ref{tab:components-paravector-real-imag}.

\begin{tableblock}
\captionof{table}{Type assignments for the generic complex paravector notation.}
\label{tab:paravector-types}
\normalsize
\begin{tabular}{@{}lll@{}}
\toprule
Type name & Underlying set & Example \\
\midrule
real scalar & \(\R\) & \(a,b\) \\
real vector & \(\R^3\) & \(\vct{A},\vct{B}\) \\
complex scalar & \(\C\) & \(S\) \\
complex vector & \(\C^3\) & \(\vct{V}\) \\
real paravector & \(\R+\R^3\) & \(\pv{X},\pv{Y}\) \\
complex paravector & \(\C+\C^3\) & \(\pv{Z}\) \\
\bottomrule
\end{tabular}
\end{tableblock}

\begin{componenttable}[tab:components-paravector-recovery]{\pv{Z}}
\componentrow{scalar}{S=(\pv{Z}+\pv{Z}^{*\mathsf H})/2}{}
\componentrow{vector}{\vct{V}\cdot\sigv=(\pv{Z}-\pv{Z}^{*\mathsf H})/2}{}
\componentrow{real}{\pv{X}=(\pv{Z}+\pv{Z}^{\mathsf H})/2}{}
\componentrow{img}{\pv{Y}=(\pv{Z}-\pv{Z}^{\mathsf H})/(2i)}{}
\end{componenttable}

\begin{tableblock}
\captionof{table}{Real and imaginary scalar/vector parts of $\pv{Z}$.}
\label{tab:components-paravector-real-imag}
\renewcommand{\arraystretch}{1.10}%
\setlength{\tabcolsep}{2pt}%
\normalsize
\begin{tabular}{@{}r@{:\quad}>{\raggedright\arraybackslash}p{0.48\columnwidth}>{\raggedright\arraybackslash}p{0.16\columnwidth}@{}}
\toprule
\componentrow{real scalar}{\begin{aligned}[t] a&=(S+S^{\mathsf H})/2\\ &=(\pv{X}+\pv{X}^*)/2 \end{aligned}}{}
\componentrow{img scalar}{\begin{aligned}[t] b&=(S-S^{\mathsf H})/(2i)\\ &=(\pv{Y}+\pv{Y}^*)/2 \end{aligned}}{}
\componentrow{real vector}{\begin{aligned}[t] \vct{A}\cdot\sigv&=(\vct{V}+\vct{V}^{\mathsf H})\cdot\sigv/2\\ &=(\pv{X}-\pv{X}^*)/2 \end{aligned}}{}
\componentrow{img vector}{\begin{aligned}[t] \vct{B}\cdot\sigv&=(\vct{V}-\vct{V}^{\mathsf H})\cdot\sigv/(2i)\\ &=(\pv{Y}-\pv{Y}^*)/2 \end{aligned}}{}
\bottomrule
\end{tabular}
\end{tableblock}

\subsection{Dot products and cross products}
For real vectors \(\vct{A},\vct{B}\in\R^3\), the dot product is defined by
\begin{align}
\vct{A}\cdot\vct{B}
&:= \sum_{n=1}^3\sum_{m=1}^3\delta_{nm}A_nB_m \notag\\
&= A_1B_1+A_2B_2+A_3B_3.
\label{eq:vector-dot-norm}
\end{align}
The square of a vector is defined by
\begin{equation}
\vct{A}^2:=\vct{A}\cdot\vct{A}.
\label{eq:vector-square}
\end{equation}
\noindent For the cross product, the Levi-Civita form is
\begin{equation}
(\vct{A}\times\vct{B})_k
\,:= \sum_{n=1}^3\sum_{m=1}^3\epsilon_{nmk}A_nB_m,
\label{eq:cross-product-definition}
\end{equation}
from which the vector is obtained:
\begin{equation}
\vct{A}\times\vct{B}=
\left\{
\begin{aligned}
&A_2B_3-A_3B_2,\\
&A_3B_1-A_1B_3,\\
&A_1B_2-A_2B_1
\end{aligned}
\right\}.
\label{eq:cross-product}
\end{equation}
Dot products are commutative, whereas cross products are antisymmetric:
\begin{align*}
\vct{A}\cdot\vct{B} &= +\vct{B}\cdot\vct{A},\\
\vct{A}\times\vct{B} &= -\vct{B}\times\vct{A}.
\end{align*}
For a fixed unit vector \(\vct{u}\), where \(\vct{u}^2=1\), the parallel
component function is defined by
\begin{equation}
\eqtext{par}(\vct{A}):=(\vct{A}\cdot\vct{u})\vct{u}.
\label{eq:par-function}
\end{equation}
The perpendicular component function is defined by
\begin{equation}
\eqtext{prp}(\vct{A}):=\vct{A}-\eqtext{par}(\vct{A}).
\label{eq:prp-function}
\end{equation}
With Eqs.\ \eqref{eq:basis} and \eqref{eq:sigma-product},
\begin{align*}
(\vct{A}\cdot\sigv)(\vct{B}\cdot\sigv)
&= \sum_{n=1}^3\sum_{m=1}^3 A_n B_m \pv{\sigma}_n \pv{\sigma}_m \notag\\
&= \sum_{n=1}^3 A_n B_n \notag\\
&\quad + \sum_{n<m}(A_n B_m-A_m B_n)\pv{\sigma}_n \pv{\sigma}_m \notag\\
&= A_1B_1+A_2B_2+A_3B_3 \notag\\
&\quad + (A_2B_3-A_3B_2)i\pv{\sigma}_1 \notag\\
&\quad - (A_1B_3-A_3B_1)i\pv{\sigma}_2 \notag\\
&\quad + (A_1B_2-A_2B_1)i\pv{\sigma}_3.
\end{align*}
This can be efficiently expressed in terms of the dot product and cross product as
\begin{equation}
(\vct{A}\cdot\sigv)(\vct{B}\cdot\sigv)=\vct{A}\cdot\vct{B}+i(\vct{A}\times\vct{B})\cdot\sigv.
\label{eq:dot-cross}
\end{equation}
The same product is recorded in component form in Table \ref{tab:components-vector-product}. The dot and cross products are extended complex-bilinearly to \(\C^3\). Formula \eqref{eq:dot-cross} may also be used for commuting scalar-valued functions. A separate ordered calculation is required for differential-operator coefficients because everything to their right is differentiated according to the Leibniz rule.

Only the antisymmetric term is flipped by the reversed order:
\begin{align*}
(\vct{B}\cdot\sigv)(\vct{A}\cdot\sigv)
&= \vct{A}\cdot\vct{B}-i(\vct{A}\times\vct{B})\cdot\sigv,\\
(\vct{A}\cdot\sigv)(\vct{B}\cdot\sigv)+(\vct{B}\cdot\sigv)(\vct{A}\cdot\sigv)
&= 2\,\vct{A}\cdot\vct{B},\\
(\vct{A}\cdot\sigv)(\vct{B}\cdot\sigv)-(\vct{B}\cdot\sigv)(\vct{A}\cdot\sigv)
&= 2i(\vct{A}\times\vct{B})\cdot\sigv.
\end{align*}
The following identities are obtained immediately from the same product:
\begin{align*}
(\vct{A}\cdot\sigv)^2 &= \vct{A}^2,\\
(\vct{A}\cdot\sigv)^{2n} &= (\vct{A}^2)^n,\\
(\vct{A}\cdot\sigv)^{2n+1} &= (\vct{A}^2)^n(\vct{A}\cdot\sigv),\\
(\vct{A}\cdot\sigv)\sigv &= \vct{A}+i\,\vct{A}\times\sigv.
\end{align*}
These are the basic sigma-product identities from which the later paravector formulas are built.

\subsection{Products}
With the scalar/vector split for real vectors fixed by Eq.\ \eqref{eq:dot-cross}, the fully complex product is obtained as the master formula from which the simpler cases are recovered by restriction.

For \(\pv{Z}\) from Eq.\ \eqref{eq:paravector} and \(\pv{Z}'=S'+\vct{V}'\cdot\sigv\) defined analogously to Eqs.\ \eqref{eq:paravector} and \eqref{eq:paravector-vector},
\begin{equation}
\pv{Z}\pv{Z}'=(SS'+\vct{V}\cdot\vct{V}')+\big(S\vct{V}'+S'\vct{V}+i(\vct{V}\times\vct{V}')\big)\cdot\sigv.
\label{eq:paravector-product}
\end{equation}
The generic complex paravector and its master multiplication rule are fixed by Eqs.\ \eqref{eq:paravector} and \eqref{eq:paravector-product}: pure scalar, real vector, real paravector, and mixed products are all obtained by setting the unused parts to zero.

\subsubsection*{Vector form and explicit sigma products}
Equation \eqref{eq:sigma-vector} turns a real vector into an algebra element by contraction against the sigma basis:
\begin{align*}
\mathbf A\cdot\sigv &= A_1\sigma_1+A_2\sigma_2+A_3\sigma_3,\\
\mathbf A\cdot\sigv+\mathbf B\cdot\sigv &= (\mathbf A+\mathbf B)\cdot\sigv.
\end{align*}
With Eqs.\ \eqref{eq:basis} and \eqref{eq:sigma-product},
\begin{align*}
(\mathbf A\cdot\sigv)(\mathbf B\cdot\sigv)
&= \sum_{n,m} A_n B_m \sigma_n \sigma_m \notag\\
&= \sum_n A_n B_n \sigma_n^2 \notag\\
&\quad + \sum_{n<m}(A_n B_m-A_m B_n)\sigma_n \sigma_m \notag\\
&= \mathbf A\cdot\mathbf B \notag\\
&\quad + (A_2B_3-A_3B_2)i\sigma_1 \notag\\
&\quad - (A_1B_3-A_3B_1)i\sigma_2 \notag\\
&\quad + (A_1B_2-A_2B_1)i\sigma_3 \notag\\
&= \mathbf A\cdot\mathbf B \notag\\
&\quad + i(\mathbf A\times\mathbf B)\cdot\sigv.
\end{align*}
Table \ref{tab:components-vector-product} records this component split.

The reversed order flips only the antisymmetric term:
\begin{align*}
(\mathbf B\cdot\sigv)(\mathbf A\cdot\sigv)
&= \mathbf A\cdot\mathbf B-i(\mathbf A\times\mathbf B)\cdot\sigv,\\
\frac{(\mathbf A\cdot\sigv)(\mathbf B\cdot\sigv)+(\mathbf B\cdot\sigv)(\mathbf A\cdot\sigv)}{2}
&= \mathbf A\cdot\mathbf B,\\
\frac{(\mathbf A\cdot\sigv)(\mathbf B\cdot\sigv)-(\mathbf B\cdot\sigv)(\mathbf A\cdot\sigv)}{2}
&= i(\mathbf A\times\mathbf B)\cdot\sigv.
\end{align*}
The same product immediately gives
\begin{align*}
(\mathbf A\cdot\sigv)^2 &= \mathbf A^2,\\
(\mathbf A\cdot\sigv)^{2n} &= (\mathbf A^2)^n,\\
(\mathbf A\cdot\sigv)^{2n+1} &= (\mathbf A^2)^n(\mathbf A\cdot\sigv),\\
(\mathbf A\cdot\sigv)\sigv &= \mathbf A+i\,\mathbf A\times\sigv.
\end{align*}
These are the basic sigma-product identities from which the later paravector formulas are built.

\subsubsection*{Products of real scalars, vectors, and paravectors}
Let
\begin{equation*}
X=a+\mathbf A\cdot\sigv,\qquad
Y=b+\mathbf B\cdot\sigv
\qquad
(a,b\in\R,\ \mathbf A,\mathbf B\in\R^3).
\end{equation*}
Using Eq.\ \eqref{eq:dot-cross},
\begin{align*}
XY
&= (a+\mathbf A\cdot\sigv)(b+\mathbf B\cdot\sigv)\notag\\
&= (ab+\mathbf A\cdot\mathbf B)
+ \bigl(a\mathbf B+b\mathbf A+i(\mathbf A\times\mathbf B)\bigr)\cdot\sigv,\\
YX
&= (ab+\mathbf A\cdot\mathbf B)
+ \bigl(a\mathbf B+b\mathbf A-i(\mathbf A\times\mathbf B)\bigr)\cdot\sigv \notag\\
&= (XY)^{\mathrm R}.
\end{align*}
Table \ref{tab:components-real-paravector-product} gives the same product in component form. The scalar part is symmetric, while the imaginary-vector part changes sign with the order.

Two immediate consequences are
\begin{align*}
[X,Y] &= 2i(\mathbf A\times\mathbf B)\cdot\sigv,\\
XY\in\R\oplus\R^3 &\iff \mathbf A\times\mathbf B=0.
\end{align*}
So two real paravectors commute exactly when their vector parts are parallel.

Setting \(Y=X\) removes the antisymmetric term automatically:
\begin{align*}
X^2
&= a^2+\mathbf A^2+2a\,\mathbf A\cdot\sigv.
\end{align*}
If \(X^2=b+\mathbf B\cdot\sigv\), then matching scalar and vector parts gives
\begin{align*}
b &= a^2+\mathbf A^2,\\
\mathbf B &= 2a\,\mathbf A,
\end{align*}
hence
\begin{align*}
a^2 &= \frac{1}{2}\Bigl(b\pm\sqrt{b^2-\mathbf B^2}\Bigr),\\
\mathbf A^2 &= b-a^2.
\end{align*}
This is the real-paravector square-root relation obtained directly from the product decomposition.

\subsubsection*{Products of complex paravectors}
The master product law \eqref{eq:paravector-product} is the complex extension of the real formulas above. Writing
\begin{equation*}
Z=X+iY=S+\mathbf V\cdot\sigv,\qquad
Z'=X'+iY'=S'+\mathbf V'\cdot\sigv,
\end{equation*}
one obtains every later specialization by setting the unused scalar or vector parts to zero. Table \ref{tab:components-paravector-product} gives the full component decomposition of \(ZZ'\), Table \ref{tab:components-paravector-square} gives \(Z^2\), and Tables \ref{tab:components-zzstar}--\ref{tab:components-zcrevz} record the order-sensitive products built from the conjugate and reverse operators.

The scalar part of a product is cyclic:
\begin{equation*}
\eqtext{scalar}(ZZ')=\eqtext{scalar}(Z'Z).
\end{equation*}
In particular, if \(Z'=MZ^{*\mathrm R}\), then
\begin{align*}
\eqtext{scalar}(Z M Z^{*\mathrm R})
&=\eqtext{scalar}(M Z^{*\mathrm R} Z)\notag\\
&=\eqtext{scalar}(M)\,(Z^{*\mathrm R} Z),
\end{align*}
because \(Z^{*\mathrm R} Z\) is scalar.

The order-sensitive variants fall into three algebraic types:
\begin{align*}
ZZ^* &\eqtext{ is a real scalar plus a complex vector},\\
ZZ^{\mathrm R} &\eqtext{ is a real scalar plus a real vector},\\
ZZ^{\mathrm{R}*} &\eqtext{ is a complex scalar}.
\end{align*}
Those three cases are precisely the ones needed later for quadratic forms, inverses, and transformation formulas.

\noindent The remaining algebra tables collect the recovery formulas, product decompositions, and order-dependent component expansions used in later calculations.
\begin{componenttable}[tab:components-paravector-recovery]{Z}
\componentrow{scalar}{a+ib=(Z+Z^{\mathrm{R}*})/2}{}
\componentrow{vector}{(\mathbf A+i\mathbf B)\cdot\sigv=(Z-Z^{\mathrm{R}*})/2}{}
\componentrow{real part}{a+\mathbf A\cdot\sigv=(Z+Z^{\mathrm R})/2}{}
\componentrow{imaginary part}{b+\mathbf B\cdot\sigv=(Z-Z^{\mathrm R})/(2i)}{}
\componentrow{real scalar}{a=(Z+Z^*+Z^{\mathrm R}+Z^{\mathrm{R}*})/4}{}
\componentrow{imaginary scalar}{b=(Z-Z^*-Z^{\mathrm R}+Z^{\mathrm{R}*})/(4i)}{}
\componentrow{real vector}{\mathbf A\cdot\sigv=(Z-Z^*+Z^{\mathrm R}-Z^{\mathrm{R}*})/4}{}
\componentrow{imaginary vector}{\mathbf B\cdot\sigv=(Z+Z^*-Z^{\mathrm R}-Z^{\mathrm{R}*})/(4i)}{}
\end{componenttable}

\noindent The following two tables collect the full component expansions for the generic product \(ZZ'\) and the special case \(Z^2\).
\begin{componenttable}[tab:components-paravector-product]{Z Z'}
\componentrow{real scalar}{aa'-bb'+\mathbf A\cdot\mathbf A'-\mathbf B\cdot\mathbf B'}{}
\componentrow{imaginary scalar}{ab'+ba'+\mathbf A\cdot\mathbf B'+\mathbf B\cdot\mathbf A'}{}
\componentrow{real vector}{a\mathbf A'-b\mathbf B'+a'\mathbf A-b'\mathbf B-\mathbf A\times\mathbf B'-\mathbf B\times\mathbf A'}{}
\componentrow{imaginary vector}{a\mathbf B'+b\mathbf A'+a'\mathbf B+b'\mathbf A+\mathbf A\times\mathbf A'-\mathbf B\times\mathbf B'}{}
\end{componenttable}

\begin{componenttable}[tab:components-paravector-square]{Z^{2}}
\componentrow{real scalar}{a^{2}-b^{2}+\mathbf A^{2}-\mathbf B^{2}}{}
\componentrow{imaginary scalar}{2ab+2\mathbf A\cdot\mathbf B}{}
\componentrow{real vector}{2(a\mathbf A-b\mathbf B)}{}
\componentrow{imaginary vector}{2(a\mathbf B+b\mathbf A)}{}
\end{componenttable}

\noindent The following two tables make the order dependence of \(ZZ^{*}\) and \(Z^{*}Z\) explicit.
\begin{componenttable}[tab:components-zzstar]{Z Z^{*}}
\componentrow{real scalar}{a^{2}+b^{2}-\mathbf A^{2}-\mathbf B^{2}}{}
\componentrow{imaginary scalar}{0}{}
\componentrow{real vector}{-2\,\mathbf A\times\mathbf B}{\(\leftarrow\) differs from \(Z^{*}Z\)}
\componentrow{imaginary vector}{2(a\mathbf B-b\mathbf A)}{}
\end{componenttable}

\begin{componenttable}[tab:components-zstarz]{Z^{*} Z}
\componentrow{real scalar}{a^{2}+b^{2}-\mathbf A^{2}-\mathbf B^{2}}{}
\componentrow{imaginary scalar}{0}{}
\componentrow{real vector}{2\,\mathbf A\times\mathbf B}{\(\leftarrow\) differs from \(ZZ^{*}\)}
\componentrow{imaginary vector}{2(a\mathbf B-b\mathbf A)}{}
\end{componenttable}

\noindent The following two tables play the same role for products involving the reverse operator.
\begin{componenttable}[tab:components-zzrev]{Z Z^{\mathrm R}}
\componentrow{real scalar}{a^{2}+b^{2}+\mathbf A^{2}+\mathbf B^{2}}{}
\componentrow{imaginary scalar}{0}{}
\componentrow{real vector}{2(a\mathbf A+b\mathbf B+\mathbf A\times\mathbf B)}{\(\leftarrow\) differs from \(Z^{\mathrm R}Z\)}
\componentrow{imaginary vector}{\mathbf 0}{}
\end{componenttable}

\begin{componenttable}[tab:components-zrevz]{Z^{\mathrm R} Z}
\componentrow{real scalar}{a^{2}+b^{2}+\mathbf A^{2}+\mathbf B^{2}}{}
\componentrow{imaginary scalar}{0}{}
\componentrow{real vector}{2(a\mathbf A+b\mathbf B-\mathbf A\times\mathbf B)}{\(\leftarrow\) differs from \(ZZ^{\mathrm R}\)}
\componentrow{imaginary vector}{\mathbf 0}{}
\end{componenttable}

\noindent The following two tables show that the conjugate-reverse products collapse to purely scalar and imaginary-scalar parts.
\begin{componenttable}[tab:components-zzcrev]{Z Z^{\mathrm{R}*}}
\componentrow{real scalar}{a^{2}-b^{2}-\mathbf A^{2}+\mathbf B^{2}}{}
\componentrow{imaginary scalar}{2ab-2\mathbf A\cdot\mathbf B}{}
\componentrow{real vector}{\mathbf 0}{}
\componentrow{imaginary vector}{\mathbf 0}{}
\end{componenttable}

\begin{componenttable}[tab:components-zcrevz]{Z^{\mathrm{R}*} Z}
\componentrow{real scalar}{a^{2}-b^{2}-\mathbf A^{2}+\mathbf B^{2}}{}
\componentrow{imaginary scalar}{2ab-2\mathbf A\cdot\mathbf B}{}
\componentrow{real vector}{\mathbf 0}{}
\componentrow{imaginary vector}{\mathbf 0}{}
\end{componenttable}


\subsubsection*{Order dependence and commutators}
The part of a product that depends on order is its commutator. For generic complex paravectors \(Z,Z'\in\C\oplus\C^3\),
\begin{equation}
[Z,Z']
:= ZZ'-Z'Z
= 2i(\mathbf V\times\mathbf V')\cdot\sigv
\label{eq:paravector-commutator}
\end{equation}
Equation \eqref{eq:paravector-commutator} is the compact commutator law. In real and imaginary vector components it expands to
\begin{align*}
[Z,Z']
&= -2(\mathbf A\times\mathbf B'+\mathbf B\times\mathbf A')\cdot\sigv \notag\\
&\quad +2i(\mathbf A\times\mathbf A'-\mathbf B\times\mathbf B')\cdot\sigv.
\end{align*}
The real-paravector commutator
\begin{equation*}
[X,Y]=2i(\mathbf A\times\mathbf B)\cdot\sigv
\end{equation*}
is the corresponding specialization. In every case the scalar terms cancel, so the commutator isolates the cross-product part of the multiplication law.

The same order dependence appears already in the sigma basis itself. The pseudoscalar commutes with the generators,
\begin{align*}
[i,\sigma_n] &= 0,\\
i\sigv &= \sigv i = \{\sigma_2\sigma_3,\ \sigma_3\sigma_1,\ \sigma_1\sigma_2\},
\end{align*}
while distinct sigma generators satisfy
\begin{equation*}
[\sigma_n,\sigma_m]=2\sigma_n\sigma_m=2\epsilon_{nmk}i\sigma_k
\qquad (n\neq m).
\end{equation*}
Thus the oriented sigma-number-two basis is exactly the commutator closure of the sigma-number-one generators.

\subsubsection*{Paravector metric and quadratic forms}
The quadratic forms \eqref{eq:euclidean-form} and \eqref{eq:minkowski-form} are product formulas rather than added structure. For a real paravector \(X=a+\mathbf A\cdot\sigv\),
\begin{align*}
\langle\!\langle X\rangle\!\rangle^2
&= a^2-\mathbf A^2\\
&= XX^*=X^*X.
\end{align*}
The indefinite form is therefore the product of a real paravector with its conjugate.

For a complex paravector \(Z=S+\mathbf V\cdot\sigv=X+iY\),
\begin{align*}
\|Z\|^2
&= |S|^2+\mathbf V\cdot\mathbf V^*\\
&= \eqtext{scalar}(Z^{\mathrm R}Z)
= \eqtext{scalar}(ZZ^{\mathrm R}),
\end{align*}
while the indefinite form is
\begin{align*}
\langle\!\langle Z\rangle\!\rangle^2
&= S^2-\mathbf V^2\\
&= ZZ^{\mathrm{R}*}
= Z^{\mathrm{R}*}Z \notag\\
&= (a^2-b^2-\mathbf A^2+\mathbf B^2)
+ 2i(ab-\mathbf A\cdot\mathbf B).
\end{align*}
Tables \ref{tab:components-zzrev}--\ref{tab:components-zcrevz} show these same product types in component form.

The same formulas also show how the two quadratic forms scale:
\begin{align*}
\|r e^{i\theta} Z\|^2 &= r^2\|Z\|^2,\\
\langle\!\langle r e^{i\theta} Z\rangle\!\rangle^2
&= r^2 e^{2i\theta}\langle\!\langle Z\rangle\!\rangle^2.
\end{align*}
The Euclidean form is insensitive to phase, while the indefinite form retains the phase square.

\subsubsection*{Inverse formulas}
The inverse formulas are the algebraic consequence of the quadratic forms. If \(\mathbf A\times\mathbf B=0\), then
\begin{equation*}
\langle\!\langle XY\rangle\!\rangle^2
= \langle\!\langle X\rangle\!\rangle^2\langle\!\langle Y\rangle\!\rangle^2.
\end{equation*}
In particular, when \(a^2-\mathbf A^2\neq 0\),
\begin{align*}
X^{-1}
&= \frac{X^*}{XX^*}
= \frac{a-\mathbf A\cdot\sigv}{a^2-\mathbf A^2}.
\end{align*}
When \(\langle\!\langle Z\rangle\!\rangle^2\neq 0\), the same argument gives
\begin{align}
Z^{-1}
&= \frac{\crev{Z}}{\langle\!\langle Z\rangle\!\rangle^2}
= \frac{\crev{Z}}{\crev{Z}Z}.
\label{eq:paravector-inverse}
\end{align}
If the corresponding quadratic form vanishes, the inverse does not exist.

Differentiating \(X^{-1}X=1\) gives the standard inverse-product rule
\begin{align*}
\dd(X^{-1}) &= -X^{-1}(\dd X)X^{-1},
\end{align*}
and when \(\mathbf A\times\mathbf B=0\) one also has
\begin{equation*}
X^{-1}Y=Y X^{-1}.
\end{equation*}


\subsection{Positive Norm, Adjugate, and Determinant}
Two distinct quadratic constructions are used: a positive coefficient norm and a generally complex determinant.
For a complex scalar \(S=a+ib\in\C\), the scalar norm is defined by
\begin{equation}
|S|:=\sqrt{SS^*}.
\label{eq:scalar-norm}
\end{equation}
\begin{equation*}
\begin{aligned}
|S|^2&=SS^*\\
&=a^2+b^2.
\end{aligned}
\end{equation*}
If \(S\in\R\), then the ordinary absolute value \(|S|=\sqrt{S^2}\) is obtained from Eq.\ \eqref{eq:scalar-norm}.
For a complex vector \(\vct{V}=\vct{A}+i\vct{B}\in\C^3\), the vector norm is defined by
\begin{equation}
\|\vct{V}\|:=\sqrt{\vct{V}\cdot\vct{V}^*}.
\label{eq:vector-norm}
\end{equation}
\begin{equation*}
\begin{aligned}
\|\vct{V}\|^2&=\vct{V}\cdot\vct{V}^*\\
&=\vct{A}^2+\vct{B}^2.
\end{aligned}
\end{equation*}
The quadratic product \(\vct{V}^2=\vct{V}\cdot\vct{V}\) is not the same object as the positive norm square \(\vct{V}\cdot\vct{V}^*\).
For the complex paravector \(\pv{Z}\) from Eqs.\ \eqref{eq:paravector}--\eqref{eq:paravector-cartesian}, the paravector norm is defined by
\begin{equation}
\|\pv{Z}\|:=\sqrt{|S|^2+\|\vct{V}\|^2}.
\label{eq:euclidean-form}
\end{equation}
Its square is reduced directly to the coefficient norm:
\begin{align}
\|\pv{Z}\|^2
&= |S|^2+\|\vct{V}\|^2\notag\\
&= a^2+b^2+\vct{A}^2+\vct{B}^2\notag\\
&= \eqtext{scalar}(\pv{Z}\herm{\pv{Z}})
= \eqtext{scalar}(\herm{\pv{Z}}\pv{Z}).
\label{eq:euclidean-norm-expansion}
\end{align}
This is a nonnegative real norm: it is the sum of the squared magnitudes of the scalar and vector coefficients.

The determinant is obtained through construction of the multiplicative inverse. For
\(\pv{Z}=S+\vct{V}\cdot\sigv\), the adjugate is obtained by reversing the sign of the vector part:
\begin{equation}
\adjf{\pv{Z}}
:=\cherm{\pv{Z}}
=S-\vct{V}\cdot\sigv.
\label{eq:adjugate}
\end{equation}
By commutativity of the two basic involutions and involutivity of each,
the following useful consistency chain is also obtained from Eq.\ \eqref{eq:adjugate}:
\begin{equation}
\herm{\pv{Z}}
=\conj{\bigl(\cherm{\pv{Z}}\bigr)}
=\cherm{\bigl(\conj{\pv{Z}}\bigr)}
=\conj{\adjf{\pv{Z}}}
=\adjf{\conj{\pv{Z}}}.
\label{eq:hermitian-adjugate-identity}
\end{equation}
In the compact superscript notation, this is
\(\pv{Z}^{\mathsf H}=(\pv{Z}^{*\mathsf H})^*
=(\pv{Z}^*)^{*\mathsf H}
=\adjf{\pv{Z}}^*=\adjf{\pv{Z}^*}\).
Because \(S\) is central and
\((\vct{V}\cdot\sigv)^2=\vct{V}^2\), the vector terms are canceled by multiplication by the adjugate:
\begin{align}
\pv{Z}\adjf{\pv{Z}}
&=(S+\vct{V}\cdot\sigv)(S-\vct{V}\cdot\sigv)\notag\\
&=S^2-\vct{V}^2\notag\\
&=\adjf{\pv{Z}}\pv{Z}.
\label{eq:adjugate-product}
\end{align}
The remaining central scalar is exactly the divisor required to normalize the product to the multiplicative identity \(1\). The determinant function is therefore defined by
\begin{equation}
\detf{\pv{Z}}:=\adjf{\pv{Z}}\pv{Z}.
\label{eq:determinant-function}
\end{equation}
Equation \eqref{eq:adjugate-product} shows that the order of the two factors
does not matter. The determinant is reduced to
\begin{equation}
\detf{\pv{Z}}=\pv{Z}\adjf{\pv{Z}}=S^2-\vct{V}^2.
\label{eq:minkowski-form}
\end{equation}
In real and imaginary components, the same result is
\begin{equation}
\detf{\pv{Z}}
=(a^2-b^2-\vct{A}^2+\vct{B}^2)
+2i(ab-\vct{A}\cdot\vct{B}).
\label{eq:determinant-components}
\end{equation}
In the two-by-two complex matrix realization, \(\adjf{\pv{Z}}\) is the ordinary matrix adjugate and Eq.\ \eqref{eq:adjugate-product} is represented by \(\detf{\pv{Z}}\I\). The designation of the scalar in Eq.\ \eqref{eq:minkowski-form} as the determinant is justified by this representation. No square root is included in its definition.

For the real paravector \(\pv{X}\) from Eq.\ \eqref{eq:real-paravector-x}, \(\cherm{\pv{X}}=\conj{\pv{X}}\), and the determinant is restricted to the Minkowski quadratic form
\begin{align}
\detf{\pv{X}}
&= a^2-\vct{A}^2\notag\\
&= \pv{X}\conj{\pv{X}}
=\conj{\pv{X}}\pv{X}.
\label{eq:real-paravector-metric}
\end{align}
Product order is reversed by the adjugate, so for complex paravectors \(\pv{Z}\) and \(\pv{Z}'\),
\begin{align}
\adjf{\adjf{\pv{Z}}}&=\pv{Z},\notag\\
\adjf{\pv{Z}\pv{Z}'}
&=\adjf{\pv{Z}'}\adjf{\pv{Z}},\notag\\
\detf{\adjf{\pv{Z}}}&=\detf{\pv{Z}}.
\label{eq:adjugate-identities}
\end{align}
Consequently, multiplicativity of the determinant is obtained without any commutativity assumption:
\begin{align}
\detf{\pv{Z}\pv{Z}'}
&=\pv{Z}\pv{Z}'\adjf{\pv{Z}'}\adjf{\pv{Z}}\notag\\
&=\pv{Z}\detf{\pv{Z}'}\adjf{\pv{Z}}\notag\\
&=\detf{\pv{Z}'}\detf{\pv{Z}}\notag\\
&=\detf{\pv{Z}}\detf{\pv{Z}'}
=\detf{\pv{Z}'\pv{Z}}.
\label{eq:metric-product}
\end{align}
Under pure phase multiplication, the positive norm and determinant are transformed as
\begin{align*}
\|\exp(i\theta)\pv{Z}\|^2 &= \|\pv{Z}\|^2,\\
\detf{\exp(i\theta)\pv{Z}}
&= \exp(2i\theta)\detf{\pv{Z}}.
\end{align*}
More generally, under complex scalar multiplication by \(r\exp(i\theta)\) with \(r\in\R\), the two quadratic constructions are scaled differently:
\begin{align}
\|r\exp(i\theta)\pv{Z}\|^2 &= r^2\|\pv{Z}\|^2,
\label{eq:norm-scaling}\\
\detf{r\exp(i\theta)\pv{Z}}
&= r^2\exp(2i\theta)\detf{\pv{Z}}.
\label{eq:metric-scaling}
\end{align}
The positive norm is phase-insensitive, while the phase square is retained by the determinant. The term ``Minkowski metric'' below is reserved for the real-paravector restriction \eqref{eq:real-paravector-metric} and its polarization.

\subsection{Inverse formulas}
The inverse construction is completed by Eqs.\ \eqref{eq:adjugate-product} and \eqref{eq:minkowski-form}. For \(\pv{Z}=S+\vct{V}\cdot\sigv\), when \(\detf{\pv{Z}}\neq0\), the inverse is defined by
\begin{equation}
\pv{Z}^{-1}:=\frac{\adjf{\pv{Z}}}{\detf{\pv{Z}}}.
\label{eq:paravector-inverse}
\end{equation}
The equivalent explicit forms are
\begin{equation}
\pv{Z}^{-1}
=\frac{\cherm{\pv{Z}}}{\pv{Z}\cherm{\pv{Z}}}
=\frac{\cherm{\pv{Z}}}{\cherm{\pv{Z}}\pv{Z}}
=\frac{S-\vct{V}\cdot\sigv}{S^2-\vct{V}^2}.
\label{eq:paravector-inverse-explicit}
\end{equation}
If \(\detf{\pv{Z}}=0\), the inverse does not exist.

By differentiation of the identity \(\pv{Z}^{-1}\pv{Z}=1\), the following result is obtained:
\begin{align}
\dd(\pv{Z}^{-1}) &= -\pv{Z}^{-1}(\dd \pv{Z})\pv{Z}^{-1}.
\label{eq:paravector-inverse-diff}
\end{align}
For \(\pv{Z}'=S'+\vct{V}'\cdot\sigv\), if \(\vct{V}\times\vct{V}'=\vct{0}\), then
\begin{equation}
\pv{Z}^{-1}\pv{Z}'=\pv{Z}'\pv{Z}^{-1}.
\label{eq:inverse-commuting}
\end{equation}

\subsection{Supporting Algebraic Derivations}
The explicit sigma-basis and operator calculations, together with one generic map decomposition, are recorded in the supporting derivations without repetition of the product development already collected above.
\subsubsection{Basis-Element Classification}
Table~\ref{tab:sigma-number-components} lists the four sigma numbers and the corresponding basis-element types, while Table~\ref{tab:sigma-number-signs} lists the signs acquired under the conjugate operator, the reverse operator, and the reverse-conjugate operator.

\subsubsection{Complex Paravector-Valued Maps}
This notation records how real scalar and vector functions combine into a generic complex paravector-valued map without imposing any physical interpretation:
\begin{align*}
Z(t,\mathbf r) &= \bigl(a(t,\mathbf r)+i\,b(t,\mathbf r)\bigr)\notag\\
&\quad + \bigl(\mathbf A(t,\mathbf r)+i\,\mathbf B(t,\mathbf r)\bigr)\cdot\sigv,\\
a,b &: \R\times\R^3 \Rightarrow \R,\\
\mathbf A,\mathbf B &: \R\times\R^3 \Rightarrow \R^3,\\
Z &: \R\times\R^3 \Rightarrow \C\oplus\C^3.
\end{align*}


\subsection{Projector Algebra}\label{sec:projectors}
Projectors are introduced next so that the later function formulas and ladder constructions can be organized. Once Eqs.\ \eqref{eq:sigma-product} and \eqref{eq:dot-cross} are fixed, the cleanest route from simple unit-direction elements to more general paravector decompositions is supplied by the projector calculus.

If \(\vct{A}\in\R^3\) is written as \(\vct{A}=\theta\vct{u}\) with
\(\theta=\|\vct{A}\|\), \(\vct{u}\in\R^3\), and \(\vct{u}^2=1\), the
idempotent function is defined by
\begin{equation}
\pv{\Pi}(\pm):=\frac{1}{2}(1\pm \vct{u}\cdot\sigv).
\label{eq:projectors}
\end{equation}
The corresponding left-projected part of any algebra element \(\pv{Z}\) is
defined by
\begin{equation}
\pv{Z}(\pm):=\pv{\Pi}(\pm)\pv{Z}.
\label{eq:projected-part-function}
\end{equation}
The identities \(\pv{\Pi}(\pm)^2=\pv{\Pi}(\pm)\),
\(\pv{\Pi}(+)\pv{\Pi}(-)=0\), \(\pv{\Pi}(+)+\pv{\Pi}(-)=1\),
\(\pv{Z}=\pv{Z}(+)+\pv{Z}(-)\), and
\((\vct{u}\cdot\sigv)\pv{\Pi}(\pm)=\pm\pv{\Pi}(\pm)\) are satisfied.
Algebraically, the parts aligned and anti-aligned with the direction \(\vct{u}\) are isolated by \(\pv{\Pi}(+)\) and \(\pv{\Pi}(-)\). Their role is purely structural: an object is split into two parts adapted to a chosen axis.

For a scalar-valued function \(f:\{\pm 1\}\to\C\),
\begin{equation}
f(\vct{u}\cdot\sigv):=f(+)\pv{\Pi}(+)+f(-)\pv{\Pi}(-).
\label{eq:functional-calculus}
\end{equation}
A compact projector-based formula for roots, inverses, exponentials, logarithms, and related functions is thereby obtained. Once the minimal polynomial of an element of this Pauli algebra is known, polynomial and suitably analytic functions can be reduced within the algebraic framework, with derivative data required at a defective root. The detailed formulas are presented first as projector and eigenvalue identities and then as ladder operations associated with a chosen axis.

About the \(z\) axis, the ladder elements are defined by
\begin{equation}
\pv{a}(\pm):=\frac{1}{2}(\pv{\sigma}_1\pm i\pv{\sigma}_2).
\label{eq:ladder}
\end{equation}
The associated commutator, nilpotence, and product identities are
\begin{equation}
\begin{gathered}
\com{\pv{\sigma}_3}{\pv{a}(\pm)}=\pm 2\pv{a}(\pm),\qquad
\pv{a}(\pm)^2=0,\\
\pv{a}(+)\pv{a}(-)=\pv{\Pi}(+),\qquad
\pv{a}(-)\pv{a}(+)=\pv{\Pi}(-).
\end{gathered}
\label{eq:ladder-properties}
\end{equation}
The nilpotent combinations of \(\pv{\sigma}_1\) and \(\pv{\sigma}_2\) by which the two projector parts selected by Eq.\ \eqref{eq:projectors} are exchanged are recorded in Eq.\ \eqref{eq:ladder}. The projector eigenvalue formulas, mixed projector products, and ladder identities built from Eqs.\ \eqref{eq:projectors}, \eqref{eq:functional-calculus}, and \eqref{eq:ladder} are collected next.
\subsubsection{Functions of \texorpdfstring{\(\mathbf u\cdot\sigv\)}{u.sigma}}
The projector pair from Eq.\ \eqref{eq:projectors}, together with the functional-calculus rule in Eq.\ \eqref{eq:functional-calculus}, gives the two-root decomposition of the unit-direction element \(\mathbf u\cdot\sigv\). The corresponding eigenvalue data are
\begin{align*}
\Pi(\pm) &= \tfrac{1}{2}(1\pm \mathbf u\cdot\sigv), \\
\Pi(\pm)^{2} &= \Pi(\pm), \\
\Pi(+)\Pi(-) &= 0, \\
\Pi(+) + \Pi(-) &= 1, \\
(\mathbf u\cdot\sigv)\Pi(\pm) &= \pm \Pi(\pm),
\end{align*}
so \(\mathbf u\cdot\sigv\) has the two eigenvalues \(+1\) and \(-1\).
Specified scalar values at those two roots determine the corresponding
projector function:
\begin{align*}
f(\mathbf u\cdot\sigv)
&= f(+1)\Pi(+) + f(-1)\Pi(-)\notag \\
&= \tfrac{1}{2}\bigl(f(+1)+f(-1)\bigr)
+ \tfrac{1}{2}\bigl(f(+1)-f(-1)\bigr)\,\mathbf u\cdot\sigv.
\end{align*}
These are the basic projector formulas for a unit-direction element. The first examples below collect the square root, logarithm branch data, rational inverse, Cayley transform, and the two exponential families for later reuse:
\begin{align*}
\sqrt{\mathbf u\cdot\sigv}
&= \epsilon(+)\Pi(+) + \epsilon(-) i\,\Pi(-), \\
\log(\mathbf u\cdot\sigv)
&= 2 i \pi\bigl(n\,\Pi(+) + (\tfrac12+m)\,\Pi(-)\bigr), \\
(z+\mathbf u\cdot\sigv)^{-1}
&= \frac{1}{z+1}\Pi(+) + \frac{1}{z-1}\Pi(-)\notag \\
&= \frac{z-\mathbf u\cdot\sigv}{z^{2}-1}, \\
\frac{i-\mathbf u\cdot\sigv}{i+\mathbf u\cdot\sigv}
&= i\,\mathbf u\cdot\sigv, \\
\exp(\theta\,\mathbf u\cdot\sigv)
&= \cosh(\theta)+\sinh(\theta)\,\mathbf u\cdot\sigv, \\
\exp(i\theta\,\mathbf u\cdot\sigv)
&= \cos(\theta)+i\sin(\theta)\,\mathbf u\cdot\sigv,
\end{align*}
where \(m,n\in\mathbb Z\) index independent logarithm branches, and the
inverse formula requires \(z\neq\pm1\). The signs
\(\epsilon(+),\epsilon(-)\in\{+1,-1\}\) are likewise independent; the conventional scalar
branches \(\sqrt{+1}=1\) and \(\sqrt{-1}=i\) correspond to
\(\epsilon(+)=\epsilon(-)=+1\). These identities follow by
applying the chosen scalar value separately on the two projector ranges.
The later exponential and logarithm subsections revisit the same formulas
with their domains stated for real and complex paravectors.

\subsubsection{Functions of X}
Let
\begin{align*}
X &= t + \mathbf r\cdot\sigv, \\
t &= \eqtext{scalar}(X)\in\mathbb R, \\
\mathbf r &= \eqtext{vector}(X)\in\mathbb R^{3}, \\
\mathbf u &:= \mathbf r/\lVert\mathbf r\rVert \qquad \text{when }\mathbf r\neq 0.
\end{align*}
If \(\mathbf r\neq 0\), write \(\mathbf r=\lVert\mathbf r\rVert \mathbf u\) with \(\mathbf u^{2}=1\). Then the minimal polynomial of \(X\) is
\begin{align*}
X^{2} - 2 t X + ( t^{2} - \mathbf r^{2} ) &= 0,\notag \\
\lambda(\pm) &= t \pm \lVert\mathbf r\rVert,\notag \\
0 &= (X-\lambda(-))(X-\lambda(+))\notag \\
&= X^{2} - \lambda(-)X\notag \\
&\quad - X\lambda(+) + \lambda(-)\lambda(+).
\end{align*}
With
\begin{equation*}
a:=\lambda(+),
\qquad
b:=\lambda(-),
\end{equation*}
the corresponding projector decomposition is
\begin{align*}
\Pi(\pm) &= \tfrac12(1\pm \mathbf u\cdot\sigv), \\
X &= a\,\Pi(+) + b\,\Pi(-), \\
\Pi(\pm)X &= \lambda(\pm)\Pi(\pm).
\end{align*}
If \(f\) is defined at the two eigenvalues, then
\begin{align*}
f(X)
&= f(a)\Pi(+) + f(b)\Pi(-)\notag \\
&= \tfrac12\bigl(f(a)+f(b)\bigr)
+ \tfrac12\bigl(f(a)-f(b)\bigr)\,\mathbf u\cdot\sigv.
\end{align*}
Writing \(X':=f(X)\), introduce its scalar and signed axial coefficients by
\begin{align*}
t' &= \tfrac12\bigl(f(a)+f(b)\bigr), \\
\rho'&:=\tfrac12\bigl(f(a)-f(b)\bigr),\\
X'&=t'+\rho'\,\mathbf u\cdot\sigv.
\end{align*}
For complex-valued \(f\), both \(t'\) and \(\rho'\) may be complex. Even
when they are real, \(\rho'\) is signed; the magnitude of the vector part is
\(\lvert\rho'\rvert\), not \(\rho'\) itself.
The projector summary may therefore be recorded as
\begin{align*}
\lambda(\pm)&=t\pm\lVert\mathbf r\rVert,\notag\\
\mathbf u&=\mathbf r/\lVert\mathbf r\rVert,\notag\\
\Pi(\pm)&=\tfrac12(1\pm\mathbf u\cdot\sigv),\notag\\
\Pi(\pm)X&=\lambda(\pm)\Pi(\pm),\notag\\
X&=\lambda(+)\Pi(+)+\lambda(-)\Pi(-),\notag\\
f(X)&=f\bigl(\lambda(+)\bigr)\Pi(+)+f\bigl(\lambda(-)\bigr)\Pi(-).
\end{align*}
The standard examples then follow immediately:
\begin{align*}
\exp(X)
&= \exp(t)\Bigl(\cosh(\lVert\mathbf r\rVert)+\sinh(\lVert\mathbf r\rVert)\,\mathbf u\cdot\sigv\Bigr), \\
\sqrt{X}
&= \epsilon(+)\sqrt{a}\,\Pi(+)+\epsilon(-)\sqrt{b}\,\Pi(-),\notag\\
&\qquad \epsilon(+),\epsilon(-)\in\{+1,-1\}, \\
\log(X) &= \log(a)\,\Pi(+) + \log(b)\,\Pi(-), \\
X^{-1}
&= \frac{\Pi(+)}{a} + \frac{\Pi(-)}{b}\notag \\
&= \frac{X^{*}}{t^{2}-\mathbf r^{2}}, \\
X^{n}
&= a^{n}\Pi(+) + b^{n}\Pi(-)\notag \\
&= \tfrac12(a^{n}+b^{n}) + \tfrac12(a^{n}-b^{n})\,\mathbf u\cdot\sigv.
\end{align*}
Here the scalar square roots and logarithms must be taken on specified
branches. A square root that is itself a real paravector exists when
\(a,b\geq0\), equivalently \(t\geq\lVert\mathbf r\rVert\); otherwise the
displayed roots are complex. The displayed logarithm requires \(a,b\neq0\); it is the real
principal logarithm when \(a,b>0\). The inverse requires \(ab\neq0\).
The power formula holds for \(n\geq0\), and also for negative integers when
\(ab\neq0\).
For two paravectors diagonalized along possibly different axes,
\begin{align*}
f(X)g(X')
&= (f(a)\Pi(+) + f(b)\Pi(-))\notag \\
&\quad \times (g(a')\Pi'(+)+g(b')\Pi'(-)),
\end{align*}
so
\begin{align*}
f(X)g(X')
&= f(a)g(a')\,\Pi(+)\Pi'(+)\notag\\
&\quad+f(a)g(b')\,\Pi(+)\Pi'(-)\notag \\
&\quad+f(b)g(a')\,\Pi(-)\Pi'(+)\notag\\
&\quad+f(b)g(b')\,\Pi(-)\Pi'(-).
\end{align*}
The mixed projector products are
\begin{align*}
\Pi(+)\Pi'(+)
&= \tfrac14(1+\mathbf u\cdot\mathbf u')\notag\\
&\quad + \tfrac14\bigl(\mathbf u'+\mathbf u+i(\mathbf u\times\mathbf u')\bigr)\cdot\sigv, \\
\Pi(+)\Pi'(-)
&= \tfrac14(1-\mathbf u\cdot\mathbf u')\notag\\
&\quad + \tfrac14\bigl(-\mathbf u'+\mathbf u-i(\mathbf u\times\mathbf u')\bigr)\cdot\sigv, \\
\Pi(-)\Pi'(+)
&= \tfrac14(1-\mathbf u\cdot\mathbf u')\notag\\
&\quad + \tfrac14\bigl(\mathbf u'-\mathbf u-i(\mathbf u\times\mathbf u')\bigr)\cdot\sigv, \\
\Pi(-)\Pi'(-)
&= \tfrac14(1+\mathbf u\cdot\mathbf u')\notag\\
&\quad + \tfrac14\bigl(-\mathbf u'-\mathbf u+i(\mathbf u\times\mathbf u')\bigr)\cdot\sigv.
\end{align*}
As a useful example, the Cayley-type transform is itself a projector exponential:
\begin{align*}
\frac{1-\mathbf r\cdot\sigv}{1+\mathbf r\cdot\sigv}
&= \exp\!\bigl(-2\,\operatorname{arctanh}(\mathbf r\cdot\sigv)\bigr)\notag \\
&= \exp\!\bigl(-2\,\operatorname{arctanh}(\lVert\mathbf r\rVert)\bigr)\Pi(+)\notag \\
&\quad + \exp\!\bigl(+2\,\operatorname{arctanh}(\lVert\mathbf r\rVert)\bigr)\Pi(-).
\end{align*}
The rational expression requires \(\lVert\mathbf r\rVert\neq1\).
For a real-valued \(\operatorname{arctanh}\) and the displayed real
exponentials, one further assumes \(\lVert\mathbf r\rVert<1\); otherwise a
complex branch must be selected.

\subsubsection{Minimal Polynomial and Projector Decomposition}
The two-root data of a complex paravector are collected here because the later function formulas, projector identities, and exponential identities all depend on this algebraic factorization; see Eqs.\ \eqref{eq:paravector-product} and \eqref{eq:functional-calculus}:
\begin{align*}
Z &= S + \mathbf V \cdot \sigv \in \C\oplus\C^{3}, \\
Z^{2} - 2 S Z + ( S^{2} - \mathbf V^{2} ) &= 0.
\end{align*}
When \(\mathbf V^{2}\neq0\), choose \(z\) with \(z^2=\mathbf V^2\) and set
\begin{align*}
\lambda(\pm)&:=S\pm z,\\
\Pi_Z(\pm)&:=\frac12\left(1\pm
\frac{\mathbf V\cdot\sigv}{z}\right).
\end{align*}
Then
\begin{align*}
( Z - \lambda (-) ) ( Z - \lambda (+) ) &= 0,\\
Z&=\lambda(+)\Pi_Z(+)+\lambda(-)\Pi_Z(-),\\
f(Z)&=f\bigl(\lambda(+)\bigr)\Pi_Z(+)
 \notag\\
&\quad+f\bigl(\lambda(-)\bigr)\Pi_Z(-)
\end{align*}
for any scalar function defined on the two eigenvalues. Changing the sign of
\(z\) merely exchanges the two labels.

The case \(\mathbf V^2=0\) must be treated separately. If
\(\mathbf V=\mathbf0\), then \(Z=S\) is scalar. If
\(\mathbf V\neq\mathbf0\), let \(N:=\mathbf V\cdot\sigv\); then \(N^2=0\),
\(Z=S+N\) is not diagonalizable, and for a function analytic at
\(S\),
\begin{equation*}
f(Z)=f(S)+f'(S)N.
\end{equation*}
This nilpotent formula is the continuous replacement for the two-projector
formula when the two roots coalesce.


\subsection{Exponential and Logarithmic Algebra}
With the projector calculus fixed, exponential and logarithmic formulas can be developed in one place. Hyperbolic half-angle transform elements are produced by the sigma-number-one exponential, and the circular half-angle elements used later for rotors and for the quaternion correspondence are produced by the sigma-number-two exponential.
\subsubsection{Real Exponential and Half-Angle Elements}
The sigma-number-one exponential is evaluated by applying the scalar
exponential separately to the two projector parts defined in
Eq.\ \eqref{eq:functional-calculus}. For the real paravector
\[
\pv{X}=a+\vct{A}\cdot\sigv\in\RPspace,
\]
the vector part is decomposed as
\begin{equation*}
\begin{aligned}
\vct{A}&=\theta \vct{u},\qquad \theta=\lVert \vct{A}\rVert,\\
\vct{u}&=\frac{\vct{A}}{\lVert \vct{A}\rVert}\ \ (\vct{A}\neq 0),\qquad
\vct{u}=\vthree{u_1}{u_2}{u_3},\\
\vct{u}^2&=1.
\end{aligned}
\end{equation*}
When \(\vct{A}=\vct{0}\), the vector exponential is \(1\) and no direction
is needed. The directional formulas below apply to
\(\vct{A}\neq\vct{0}\).
The exponential is therefore given by the projector-based function rule of
Eq.\ \eqref{eq:functional-calculus}:
\begin{align}
\exp(\pm\vct{A}\cdot\sigv)
&= \exp(\pm\theta\,\vct{u}\cdot\sigv)\notag \\
&= \exp(\pm\theta)\pv{\Pi}(+) + \exp(\mp\theta)\pv{\Pi}(-)\notag \\
&= \cosh(\theta) \pm \sinh(\theta)\,\vct{u}\cdot\sigv.
\label{eq:exp-vector}
\end{align}
An equivalent derivation is obtained from the power series by separating
the even and odd powers:
\begin{align*}
\exp(\vct{A}\cdot\sigv)
&= \exp(\theta\,\vct{u}\cdot\sigv)\notag \\
&= \sum_{n\geq0}\frac{\theta^n(\vct{u}\cdot\sigv)^n}{n!}\notag \\
&= \sum_{n\geq0}\frac{\theta^{2n}}{(2n)!}
+ \sum_{n\geq0}\frac{\theta^{2n+1}}{(2n+1)!}\,\vct{u}\cdot\sigv\notag \\
&= \cosh(\theta)+\sinh(\theta)\,\vct{u}\cdot\sigv.
\end{align*}
With
\begin{align*}
\gamma &:= \cosh(\theta),\notag \\
\vct{v} &:= \tanh(\theta)\,\vct{u},
\end{align*}
Eq.\ \eqref{eq:exp-vector} gives the same formula in terms of
\(\cosh(\theta)\), \(\sinh(\theta)\), and their ratio
\(\tanh(\theta)\):
\begin{align*}
\exp(\vct{A}\cdot\sigv) &= \gamma(1+\vct{v}\cdot\sigv),\notag \\
\exp(-\vct{A}\cdot\sigv) &= \gamma(1-\vct{v}\cdot\sigv).
\end{align*}
For the real paravector \(\pv{X}\) from Eq.\ \eqref{eq:real-paravector-x}, the scalar part commutes with the vector part, so
\begin{align*}
\exp(\pv{X}) &= \exp(a)\exp(\vct{A}\cdot\sigv),\notag \\
\exp(\pv{X}^{*}) &= \exp(a)\exp(-\vct{A}\cdot\sigv)=\exp(\pv{X})^{*},\notag \\
\exp(-\pv{X}) &= \exp(-a)\exp(-\vct{A}\cdot\sigv)=\exp(\pv{X})^{-1}.
\end{align*}
The corresponding half-angle element is
\begin{equation}
\pv{L}(\theta,\vct{u})
:=\exp\!\left(\frac{\theta}{2}\,\vct{u}\cdot\sigv\right).
\label{eq:boost-element}
\end{equation}
Its conjugation properties are
\[
\pv{L}^{\Hop}=\pv{L},
\qquad
\conj{\pv{L}}=\pv{L}^{-1}.
\]
The action of this element on real events is derived from the two-sided
event law in Eq.\ \eqref{eq:general-event-congruence}; it is not obtained
from the inverse-pair product in Eq.\ \eqref{eq:general-transform}. The
exponential and its inverse are specified by Eq.\ \eqref{eq:boost-element}. The two
projector weights are exchanged by the conjugate operator:
\begin{align*}
\exp(\vct{A}\cdot\sigv)^{*}
&= \exp(-\theta)\pv{\Pi}(+) + \exp(\theta)\pv{\Pi}(-)\notag \\
&= \gamma(1-\vct{v}\cdot\sigv)\notag \\
&= \exp(-\vct{A}\cdot\sigv)\notag \\
&= \exp\bigl((\vct{A}\cdot\sigv)^{*}\bigr)\notag \\
&= \exp(\vct{A}\cdot\sigv)^{-1}.
\end{align*}

\subsubsection{Real Logarithm of a Paravector}
The logarithm is the inverse problem for the same projector formulas. The
real paravector \(\pv{Y}\) from Eq.\ \eqref{eq:real-paravector-y} is assumed
to satisfy \(b>\lVert\vct{B}\rVert\), so that the two real values in
Eq.\ \eqref{eq:real-paravector-eigenvalues} are positive. The unique ordinary
real logarithm of each positive value is used; this choice is called the
principal logarithm. For \(\vct{B}\neq\vct{0}\), the local direction is
defined by
\begin{equation}
\vct{u}:=\frac{\vct{B}}{\lVert\vct{B}\rVert},
\qquad
\vct{u}^2=1.
\label{eq:log-paravector-axis}
\end{equation}
Thus \(\vct{B}=\lVert\vct{B}\rVert \vct{u}\). The
projector function from Eq.\ \eqref{eq:projectors} is then specialized
immediately after its direction:
\begin{equation}
\pv{\Pi}(\pm):=\tfrac12(1\pm \vct{u}\cdot\sigv).
\label{eq:log-paravector-projectors}
\end{equation}
The eigenvalue construction from
Eq.\ \eqref{eq:real-paravector-eigenvalues} is then specialized to
\(\pv{Y}\). Within this subsection, the two eigenvalues of \(\pv{Y}\) are
denoted locally by \(\lambda(\pm)\):
\begin{align}
\lambda(\pm) &:= b \pm \lVert\vct{B}\rVert,
\label{eq:log-paravector-eigenvalues}\\
\pv{Y} &= \lambda(+)\pv{\Pi}(+) + \lambda(-)\pv{\Pi}(-).\notag
\end{align}
The logarithm is therefore applied separately to the two eigenvalues:
\begin{align}
\log(\pv{Y})
&= \log\!\bigl(\lambda(+)\bigr)\pv{\Pi}(+)\notag\\
&\quad+\log\!\bigl(\lambda(-)\bigr)\pv{\Pi}(-)\notag \\
&= \tfrac12 \log\!\bigl(\lambda(+)\lambda(-)\bigr)
\notag\\
&\quad+ \tfrac12 \log\!\bigl(\lambda(+)/\lambda(-)\bigr)
\,\vct{u}\cdot\sigv.
\label{eq:log-paravector}
\end{align}
When \(\vct{B}=\vct{0}\), the axis is unnecessary and this formula is
reduced directly to \(\log(\pv{Y})=\log(b)\).
For \(\vct{B}\neq0\), if \(\pv{Y}=\exp(\pv{X})\) with \(\pv{X}\) from
Eq.\ \eqref{eq:real-paravector-x} and \(\vct{A}=\theta\vct{u}\), then
\(\lambda(\pm)=\exp(a\pm\theta)\). The relation is expanded as
\begin{align*}
\pv{Y}
&=\exp(\pv{X})\notag \\
&=\exp(a+\vct{A}\cdot\sigv)\notag \\
&=\exp(a)\exp(\vct{A}\cdot\sigv)\notag \\
&=\exp(a)\bigl(\cosh(\theta)+\sinh(\theta)\,\vct{u}\cdot\sigv\bigr)\notag \\
&=b+\lVert\vct{B}\rVert\,\vct{u}\cdot\sigv,
\end{align*}
and the recovered parameters are
\begin{align*}
\theta &= \operatorname{arctanh}\!\bigl(\lVert\vct{B}\rVert/b\bigr), \\
\exp(a) &= \sqrt{b^{2}-\vct{B}^{2}}, \\
a &= \tfrac12 \log\!\bigl(b^{2}-\vct{B}^{2}\bigr), \\
\vct{A} &= \theta\,\vct{u}
= \operatorname{arctanh}\!\bigl(\lVert\vct{B}\rVert/b\bigr)\,
\frac{\vct{B}}{\lVert\vct{B}\rVert}.
\end{align*}
The condition \(b>\lVert\vct{B}\rVert\) gives a positive scalar part and
\(\detf{\pv{Y}}=b^2-\vct{B}^2>0\). In special-relativity language this is
the future-directed timelike region \cite{minkowski1908,einstein1905}.
Thus its principal logarithm is again a real paravector,
\begin{equation*}
\log(\pv{Y})=a+\vct{A}\cdot\sigv.
\end{equation*}

\subsubsection{Circular Exponentials and Half-Angle Elements}
For this subsection, a nonzero real vector \(\vct{A}\) is written with the
local values
\[
\theta:=\|\vct{A}\|,
\qquad
\vct{u}:=\frac{\vct{A}}{\|\vct{A}\|},
\qquad
\vct{A}=\theta\vct{u},
\qquad
\vct{u}^2=1.
\]
Replacing \(\theta\) with \(i\theta\) changes the hyperbolic functions into
the ordinary sine and cosine functions. The resulting sigma-number-two
exponential is therefore called circular:
\begin{align}
\exp(\pm i\vct{A}\cdot\sigv)
&= \exp(\pm i\theta\,\vct{u}\cdot\sigv)\notag \\
&= \exp(\pm i\theta)\pv{\Pi}(+) + \exp(\mp i\theta)\pv{\Pi}(-)\notag \\
&= \cos(\theta) \pm i\sin(\theta)\,\vct{u}\cdot\sigv.
\label{eq:exp-bivector}
\end{align}
Hamilton's quaternion units are obtained by relabeling the same
sigma-number-two basis, without introducing new generators or a new product
\cite{hamilton1844}:
\begin{equation}
\begin{aligned}
\pv{q}_1&:=-\pv{\sigma}_2\pv{\sigma}_3=-i\pv{\sigma}_1,\qquad
\pv{q}_2:=-\pv{\sigma}_3\pv{\sigma}_1=-i\pv{\sigma}_2,\\
\pv{q}_3&:=-\pv{\sigma}_1\pv{\sigma}_2=-i\pv{\sigma}_3,\qquad
\vct{q}:=\vthree{\pv{q}_1}{\pv{q}_2}{\pv{q}_3}=-i\sigv.
\end{aligned}
\label{eq:quaternion-basis}
\end{equation}
The multiplication and intrinsic-operation rules are inherited directly from
Eqs.\ \eqref{eq:sigma-product} and \eqref{eq:bivector-operator-signs}:
\begin{equation}
\begin{aligned}
\pv{q}_n\pv{q}_m
&=-\delta_{nm}+\sum_{k=1}^{3}\epsilon_{nmk}\pv{q}_k,\\
\conj{\pv{q}_n}&=\pv{q}_n,\qquad
\herm{\pv{q}_n}=-\pv{q}_n.
\end{aligned}
\label{eq:quaternion-basis-properties}
\end{equation}
In particular,
\(\pv{q}_1^2=\pv{q}_2^2=\pv{q}_3^2
=\pv{q}_1\pv{q}_2\pv{q}_3=-1\).
The negative-sign exponential may therefore be written as
\begin{equation*}
\exp(-i\vct{A}\cdot\sigv)
=\cos(\theta)+\sin(\theta)\,\vct{u}\cdot(-i\sigv).
\end{equation*}
In the quaternion basis from Eq.\ \eqref{eq:quaternion-basis}, the same
formula is written as
\begin{equation*}
\exp(-i\vct{A}\cdot\sigv)
=\cos(\theta)+\sin(\theta)\,\vct{u}\cdot\vct{q}.
\end{equation*}
For a real angle \(\theta\in\R\) and a unit real vector \(\vct{u}\in\R^3\), the corresponding circular half-angle exponential is
\begin{equation}
\pv{R}(\theta,\vct{u}):=\exp\!\left(\frac{\theta}{2}\,\vct{u}\cdot i\sigv\right).
\label{eq:rotor-element}
\end{equation}
A real scalar coefficient and a purely imaginary vector coefficient are
carried, so \(\pv{R}\in\CPspace\).
Its conjugation properties are
\[
\herm{\pv{R}}\pv{R}=1,
\qquad
\conj{\pv{R}}=\pv{R}.
\]
Equation \eqref{eq:rotor-element} is the circular half-angle element built
from Eq.\ \eqref{eq:exp-bivector}. This element is called a rotor because
its two-sided action produces a spatial rotation. That action is derived in
the Rotation subsection from Eqs.\ \eqref{eq:rotor} and
\eqref{eq:similarity-components}, where the result is reduced to
Rodrigues's rotation formula \cite{doran2003}. The
equivalent quaternion form is obtained from the same sigma-number-two
algebra.

\subsubsection{Exponentials of Complex Paravectors}
The real paravectors are written as
\begin{align*}
\pv{X} &= a + \vct{A} \cdot \sigv \in \RPspace,\notag \\
\pv{Y} &= b + \vct{B} \cdot \sigv \in \RPspace,\notag \\
\pv{Z} &= S + \vct{V} \cdot \sigv \in \CPspace,\notag \\
&= \pv{X} + i \pv{Y},
\end{align*}
where
\begin{equation*}
S=a+ib\in\C,
\qquad
\vct{V}=\vct{A}+i\vct{B}\in\C^3.
\end{equation*}
Because the scalar part \(S\) commutes with \(\vct{V}\cdot\sigv\), the
exponential is separated into scalar and directional factors:
\begin{align*}
\exp(\pv{Z}) &= \exp( S + \vct{V} \cdot \sigv ),\notag \\
&= \exp(S) \exp( \vct{V} \cdot \sigv ).
\end{align*}
The associated star and inverse identities are
\begin{align*}
\exp(\pv{Z}^{*}) &= \exp(S^{*}) \exp( -\vct{V}^{*} \cdot \sigv ),\notag \\
\exp(\pv{Z})^{-1} &= \exp(-S) \exp( -\vct{V} \cdot \sigv ),\notag \\
\exp(\pv{Z}^{*})^{-1} &= \exp(-S^{*}) \exp( \vct{V}^{*} \cdot \sigv ),\notag \\
(\vct{V} \cdot \sigv )^{2} &= \vct{V} \cdot \vct{V} = \vct{V}^{2},\notag \\
&= \vct{A}^{2} - \vct{B}^{2} + 2 i \vct{A} \cdot \vct{B}.
\end{align*}
The key simplification is that all higher powers are reduced to the two basis
elements \(1\) and \(\vct{V}\cdot\sigv\):
\begin{align*}
(\vct{V} \cdot \sigv )^{2n} &= (\vct{V}^{2})^{n},\notag \\
z &= \sqrt{\vct{V}^{2}}\qquad \eqtext{complex scalar},\notag \\
\exp(\vct{V} \cdot \sigv ) &= \sum_{n\geq0} \frac{(\vct{V} \cdot \sigv )^{n}}{n!},\notag \\
&= \sum_{n\geq0} \frac{(\vct{V}^{2})^{n}}{(2n)!}
+ \sum_{n\geq0} \frac{(\vct{V}^{2})^{n}}{(2n+1)!}(\vct{V} \cdot \sigv ),\notag \\
&= \cosh(z) + \sinh(z) \frac{\vct{V} \cdot \sigv}{z}.
\end{align*}
The final expression is independent of the sign chosen for \(z\). When
\(\vct{V}^2=0\), the quotient \(\sinh(z)/z\) is understood by its
continuous value \(1\), and the formula is reduced to
\begin{equation*}
\exp(\pv{Z})=\exp(S)\bigl(1+\vct{V}\cdot\sigv\bigr).
\end{equation*}


\subsection{Generalized Transform Algebra}\label{sec:transform-algebra}
With the projector and exponential formulas fixed, the remaining purely mathematical transform material can be collected in one place. A generic-to-special order is used: the similarity law is given first, followed by the basis-transport identities and then by the circular and hyperbolic transforms built from the earlier exponential elements:

\begin{align}
\pv{T}\pv{Z}'&=\pv{Z}\pv{T},
\label{eq:transform-intertwining}\\
\pv{Z}'&=\pv{T}^{-1}\pv{Z}\pv{T}.
\label{eq:general-transform}
\end{align}
Here \(\pv{Z}\) is a generic complex paravector and \(\pv{T}\) is a generic invertible algebra element. The equivalence of the primed and unprimed algebra actions is defined by the intertwining requirement in Eq.\ \eqref{eq:transform-intertwining}. Its unique solution is the similarity \eqref{eq:general-transform}, which is used for basis changes, rotations, and electromagnetic fields. This similarity is also separated in the derivation below from the congruence required for a real event paravector.
\subsubsection{Intertwining and the General Similarity Transform}
Compatibility between two descriptions of the same algebraic action is
expressed by the starting relation \eqref{eq:transform-intertwining}. A
primed carrier is carried into the unprimed description by \(\pv{T}\). If an algebraic
action is represented by \(\pv{Z}'\) before that identification and by \(\pv{Z}\)
afterward, then commutation of the action with the change of description is required:
\begin{equation*}
\pv{T}(\pv{Z}'\pv{\psi}')=\pv{Z}(\pv{T}\pv{\psi}')
\qquad\text{for every carrier }\pv{\psi}'.
\end{equation*}
The relation \(\pv{T}\pv{Z}'=\pv{Z}\pv{T}\) is obtained by removing the
arbitrary carrier. Because \(\pv{T}\) is invertible, its unique solution is
\begin{equation*}
\pv{Z}'=\pv{T}^{-1}\pv{Z}\pv{T}.
\end{equation*}
Thus Eq.\ \eqref{eq:general-transform} is not an additional dynamical law:
it is the change-of-description law forced by the intertwining requirement.
Under the opposite carrier direction, \(\pv{T}\) is replaced everywhere by
\(\pv{T}^{-1}\), and the equivalent convention
\(\pv{Z}'=\pv{T}\pv{Z}\pv{T}^{-1}\) is obtained. Under the inverse-first
convention, the displayed positive sign in Rodrigues's formula is produced
by the positive exponent in Eq.\ \eqref{eq:rotor-element}.

The fully general complex paravectors are written as
\begin{align*}
\pv{Z}&=S+\vct{V}\cdot\sigv,
&S&\in\C,
&\vct{V}&\in\C^3,\\
\pv{T}&=f+\vct{g}\cdot\sigv,
&f&\in\C,
&\vct{g}&\in\C^3,
\end{align*}
where the dot and cross products are extended complex-bilinearly. The central
determinant of \(\pv{T}\) is \(\detf{\pv{T}}=f^2-\vct{g}^2\). When
\(\detf{\pv{T}}\neq0\), the inverse is given by the general formula
\eqref{eq:paravector-inverse}:
\begin{equation*}
\pv{T}^{-1}=\frac{\cherm{\pv{T}}}{\detf{\pv{T}}}
=\frac{f-\vct{g}\cdot\sigv}{f^2-\vct{g}^2}.
\end{equation*}
Because the scalar part \(S\) is central,
\begin{align*}
\pv{Z}'
&=S+\frac{(f-\vct{g}\cdot\sigv)
 (\vct{V}\cdot\sigv)(f+\vct{g}\cdot\sigv)}{\detf{\pv{T}}}.
\end{align*}
The first multiplication is
\begin{equation*}
(f-\vct{g}\cdot\sigv)(\vct{V}\cdot\sigv)
=-\vct{g}\cdot\vct{V}
+\bigl(f\vct{V}-i\,\vct{g}\times\vct{V}\bigr)\cdot\sigv.
\end{equation*}
After multiplication by the final factor, the scalar terms are canceled and the
remaining vector terms are
\begin{align}
\pv{Z}'&=S+\vct{V}'\cdot\sigv,\notag\\
\vct{V}'
&=\frac{(f^2+\vct{g}^2)\vct{V}
-2(\vct{g}\cdot\vct{V})\vct{g}
-2if(\vct{g}\times\vct{V})}{f^2-\vct{g}^2}.
\label{eq:similarity-components}
\end{align}
This is the similarity law for every complex paravector \(\pv{Z}\) and every
invertible complex paravector \(\pv{T}\). The complex scalar part is fixed.
Neither Eq.\ \eqref{eq:general-transform} nor
Eq.\ \eqref{eq:similarity-components} is changed when \(\pv{T}\) is multiplied
by a nonzero central complex scalar.

Products, inverses, and functions are governed by the same intertwining. For a
second algebra element \(\pv{W}\),
\begin{align}
(\pv{Z}\pv{W})'&=\pv{T}^{-1}\pv{Z}\pv{W}\pv{T}
=\pv{T}^{-1}\pv{Z}\pv{T}\,\pv{T}^{-1}\pv{W}\pv{T}=\pv{Z}'\pv{W}',\notag\\
(\pv{Z}')^{-1}&=\pv{T}^{-1}\pv{Z}^{-1}\pv{T}
\qquad\text{when }\pv{Z}^{-1}\text{ exists},\notag\\
(\pv{Z}')^n&=\pv{T}^{-1}\pv{Z}^n\pv{T},\notag\\
h(\pv{Z}')&=\pv{T}^{-1}h(\pv{Z})\pv{T}.
\label{eq:similarity-covariance}
\end{align}
The last line is valid for every convergent power-series function \(h\), and
therefore in particular
\begin{equation*}
\exp(\pv{Z}')=\pv{T}^{-1}\exp(\pv{Z})\pv{T}.
\end{equation*}
The same relation is obeyed by logarithms, roots, and rational functions
wherever the same branch and domain are defined.

Commutation of the adjugate with this similarity is also obtained. Since
\(\cherm{\pv{T}}=\detf{\pv{T}}\pv{T}^{-1}\) and
\(\cherm{(\pv{T}^{-1})}=\pv{T}/\detf{\pv{T}}\),
\begin{align*}
\cherm{\pv{Z}'}
&=\cherm{\pv{T}}\,\cherm{\pv{Z}}\,\cherm{(\pv{T}^{-1})}
=\pv{T}^{-1}\cherm{\pv{Z}}\pv{T},\\
\pv{Z}'\cherm{\pv{Z}'}
&=\pv{T}^{-1}(\pv{Z}\cherm{\pv{Z}})\pv{T}=\pv{Z}\cherm{\pv{Z}}.
\end{align*}
Thus the determinant form \(S^2-\vct{V}^2\), the spectrum, and the inverse
framework are preserved by similarity for every invertible \(\pv{T}\). The
positive coefficient norm is a different construction and is not generally
preserved by a nonunitary similarity.

Algebra elements whose action is represented by multiplication are transported
by the similarity law; this law is not automatically the Lorentz law for every
physical representation. A real event paravector is characterized by
\(\pv{X}^{\mathsf H}=\pv{X}\). Such paravectors are
spanned by Hermitian carrier bilinears
\(\pv{\psi}\pv{\psi}^{\mathsf H}\): concretely, in the Pauli realization every
Hermitian two-by-two matrix is a real linear
combination of bilinears of this form. The carrier identification is
\(\pv{\psi}'=\pv{T}^{-1}\pv{\psi}\), and therefore
\begin{equation*}
\pv{\psi}'\pv{\psi}'^{\mathsf H}
=\pv{T}^{-1}(\pv{\psi}\pv{\psi}^{\mathsf H})(\pv{T}^{-1})^{\mathsf H}.
\end{equation*}
By real linearity, the event law is consequently forced to be the
inverse-carried Hermitian congruence. For a normalized Lorentz element
\(\pv{T}\cherm{\pv{T}}=1\), it is
\begin{align}
\pv{X}'
&=\pv{T}^{-1}\pv{X}(\pv{T}^{-1})^{\mathsf H}\notag\\
&=\pv{T}^{*\mathsf H}\pv{X}\pv{T}^*.
\label{eq:general-event-congruence}
\end{align}
Distinct roles are assigned to the involutions. In the Pauli realization,
\({}^{\mathsf H}\) is Hermitian adjunction, whereas
\({}^{*\mathsf H}\) is the adjugate. The latter equals \(\pv{T}^{-1}\) only
after the determinant-one normalization \(\pv{T}\cherm{\pv{T}}=1\) has been imposed.
The real subspace is preserved because
\((\pv{X}')^{\mathsf H}=\pv{X}'\), and the event determinant is preserved because
\(\pv{T}^{-1}\) has unit determinant. When \(\pv{X}\) is invertible, the
transformed inverse is obtained without a new definition:
\begin{equation}
(\pv{X}')^{-1}=\pv{T}^{\mathsf H}\pv{X}^{-1}\pv{T},
\label{eq:event-inverse-transform}
\end{equation}
because the identity is obtained by direct multiplication with
\(\pv{X}'=\pv{T}^{-1}\pv{X}(\pv{T}^{-1})^{\mathsf H}\). The algebraic
pattern for the dual derivative paravector is also supplied by this
complementary law. The event intertwiner is obtained by multiplication of
Eq.\ \eqref{eq:general-event-congruence} on the left by \(\pv{T}\):
\begin{equation}
\pv{T}\pv{X}'=\pv{X}\pv{T}^*.
\label{eq:event-intertwining}
\end{equation}
The form \(\pv{T}\pv{X}'=\pv{X}\pv{T}\) is obtained for a circular rotor because
\(\pv{R}^*=\pv{R}\), but not for a boost. The similarity and event
intertwiners are therefore identical for circular rotors but remain distinct
for boosts.

In contrast, the action of a bivector such as the electromagnetic field is
represented by left multiplication on a carrier. The following similarity is
then required by the original intertwining argument:
\begin{equation}
\pv{F}'=\pv{T}^{-1}\pv{F}\pv{T}.
\label{eq:general-field-similarity}
\end{equation}
For a circular rotor, \(\pv{R}^*=\pv{R}\) and
\(\pv{R}^{-1}=\pv{R}^{\mathsf H}\), so Eq.\ \eqref{eq:general-event-congruence}
is reduced to the active spatial rotation
\begin{equation}
\pv{X}'=\herm{\pv{R}}\pv{X}\pv{R}.
\label{eq:rotor}
\end{equation}
For the hyperbolic element in Eq.\ \eqref{eq:boost-element},
\(\pv{L}^{\mathsf H}=\pv{L}\) and
\(\pv{L}^{-1}=\pv{L}^*=\pv{L}^{*\mathsf H}\), so the same general event law
is reduced to the passive boost
\begin{equation}
\pv{X}'=\conj{\pv{L}}\pv{X}\conj{\pv{L}}.
\label{eq:boost}
\end{equation}
By contrast,
Eq.\ \eqref{eq:general-field-similarity} is reduced to
\(\pv{F}'=\pv{L}^*\pv{F}\pv{L}\).
The unequal right factors are therefore consequences of the representations
being transported, not of an inconsistent choice of inverse.
Under the definition of \(\pv{L}(\theta,\vct{u})\), the sandwich
\(\pv{L}^{\mathsf H}\pv{X}\pv{L}=\pv{L}\pv{X}\pv{L}\) is the inverse-direction event boost, obtained by
replacing \(\theta\) with \(-\theta\). It is physically legitimate, but it
is not the passive \(+\theta\) convention used in Eq.\ \eqref{eq:boost}.

The hyperbolic reduction of the general complex formula is obtained
directly with
\[
f=\cosh(\theta/2),
\qquad
\vct{g}=\sinh(\theta/2)\vct{u},
\qquad
\vct{u}^2=1,
\]
Eq.\ \eqref{eq:similarity-components} is reduced to
\begin{equation}
\vct{V}'
=(\vct{u}\cdot\vct{V})\vct{u}
+\cosh(\theta)\bigl(\vct{V}-(\vct{u}\cdot\vct{V})\vct{u}\bigr)
-i\sinh(\theta)\,\vct{u}\times\vct{V}.
\label{eq:hyperbolic-similarity}
\end{equation}
For \(\vct{V}=\vct{E}+i\vct{B}\), its real and imaginary parts are
exactly the electric--magnetic boost formulas. For a real \(\vct{V}\),
however, the cross term is generally imaginary; hyperbolic similarity is
therefore a bivector or basis transformation rather than an event boost.

The distinction is exposed by the infinitesimal forms. If
\(\pv{T}=1+\pv{K}+O(\pv{K}^2)\), the order functions defined in
Eqs.\ \eqref{eq:com-function} and \eqref{eq:acom-function} give
\begin{align*}
\pv{Z}'&=\pv{Z}+\com{\pv{Z}}{\pv{K}}+O(\pv{K}^2),\\
\pv{X}'&=\pv{X}-\pv{K}\pv{X}-\pv{X}\pv{K}^{\mathsf H}+O(\pv{K}^2).
\end{align*}
For a circular generator \(\pv{K}^{\mathsf H}=-\pv{K}\), the event law is reduced
to the same commutator by which the similarity rotation is generated. For a
hyperbolic generator \(\pv{K}^{\mathsf H}=\pv{K}\), an anticommutator is obtained instead:
\(\pv{X}'=\pv{X}-\acom{\pv{K}}{\pv{X}}+O(\pv{K}^2)\). Spatial vectors are
rotated by the former, while the scalar time component is mixed with the
vector component parallel to the boost by the latter. Thus similarity
rotations are produced by circular generators, whereas event congruences
are produced by hyperbolic generators.

A transform element \(\pv{T}(\lambda)\) is taken to depend on one real
parameter \(\lambda\). All displayed derivatives are ordinary derivatives
with respect to that parameter. The similarity generator is defined by
\begin{equation}
\pv{\Omega}:=\pv{T}^{-1}\frac{\dd\pv{T}}{\dd\lambda}.
\label{eq:similarity-generator}
\end{equation}
Differentiation of the inverse and of the transported paravector then gives
\begin{align*}
\frac{\dd(\pv{T}^{-1})}{\dd\lambda}
&=-\pv{T}^{-1}\frac{\dd\pv{T}}{\dd\lambda}\pv{T}^{-1},\\
\frac{\dd\pv{Z}'}{\dd\lambda} &=
\pv{T}^{-1}\frac{\dd\pv{Z}}{\dd\lambda}\pv{T}
+\com{\pv{Z}'}{\pv{\Omega}},\notag \\
\frac{\dd^2\pv{Z}'}{\dd\lambda^2} &=
\pv{T}^{-1}\frac{\dd^2\pv{Z}}{\dd\lambda^2}\pv{T}\notag \\
&\quad +2\com{\dfrac{\dd\pv{Z}'}{\dd\lambda}}{\pv{\Omega}},\notag \\
&\quad +\com{\pv{Z}'}{\dfrac{\dd\pv{\Omega}}{\dd\lambda}},\notag \\
&\quad +\com{\pv{\Omega}}{\com{\pv{Z}'}{\pv{\Omega}}}.
\end{align*}
They are not space-time derivatives.

\subsubsection{Transformed Basis}
The parameter-dependent element is written as
\begin{align*}
\pv{T}(\lambda)&=f(\lambda)+\vct{g}(\lambda)\cdot\sigv,
\qquad f\in\C,\quad \vct{g}\in\C^3,\\
\pv{T}^{-1}&=\frac{f-\vct{g}\cdot\sigv}{f^2-\vct{g}^2},
\qquad f^2-\vct{g}^2\neq0.
\end{align*}
The transformed basis function is defined by
\begin{equation}
\pv{\sigma}_n':=\pv{T}^{-1}\pv{\sigma}_n\pv{T}.
\label{eq:transformed-basis-definition}
\end{equation}
Its inverse relation is
\[
\pv{\sigma}_n=\pv{T}\pv{\sigma}_n'\pv{T}^{-1}.
\]
Both real hyperbolic elements and imaginary-vector rotors are included by
allowing complex \(\vct{g}\). The following relations are obtained
immediately by similarity:
\begin{align*}
(\pv{\sigma}_n')^2&=1,\\
\pv{\sigma}_n'\pv{\sigma}_m'
&=\delta_{nm}+\sum_{k=1}^3\epsilon_{nmk}\,i\pv{\sigma}_k',\\
\pv{\sigma}_1'\pv{\sigma}_2'\pv{\sigma}_3'&=i.
\end{align*}
Thus the same Clifford algebra is satisfied by the transformed generators
for every invertible \(\pv{T}\).

The separate actions of the two basic involutions on the transformed
generators are represented by
\begin{align*}
\conj{\pv{\sigma}_n'}
&=\conj{(\pv{T}^{-1})}\conj{\pv{\sigma}_n}\conj{\pv{T}}
=-(\conj{\pv{T}})^{-1}\pv{\sigma}_n\conj{\pv{T}},\\
\herm{\pv{\sigma}_n'}
&=\herm{\pv{T}}\herm{\pv{\sigma}_n}\herm{(\pv{T}^{-1})}
=\herm{\pv{T}}\pv{\sigma}_n(\herm{\pv{T}})^{-1}.
\end{align*}
The scalar part is also preserved by similarity, so
\(\eqtext{scalar}(\pv{\sigma}_n')=\eqtext{scalar}(\pv{\sigma}_n)=0\).
The individual star and \(H\) actions are recorded separately by these
formulas and are combined by the adjugate relation.

For every invertible \(\pv{T}\), commutation of the adjugate with the
similarity is obtained from
\(\cherm{\pv{T}}=(f^2-\vct{g}^2)\pv{T}^{-1}\) and
\(\cherm{(\pv{T}^{-1})}=\pv{T}/(f^2-\vct{g}^2)\):
\begin{align*}
(\pv{\sigma}_n')^{*\mathsf H}
&=(f^2-\vct{g}^2)\pv{T}^{-1}(-\pv{\sigma}_n)
  \frac{\pv{T}}{f^2-\vct{g}^2}\\
&=-\pv{T}^{-1}\pv{\sigma}_n\pv{T}=-\pv{\sigma}_n'.
\end{align*}
Consequently a paravector written in the transformed basis,
\[
\pv{Z}=Z_s+\vct{Z}_{\vct{v}}\cdot\sigv',
\]
is assigned the same determinant quadratic form,
\[
\pv{Z}\pv{Z}^{*\mathsf H}=Z_s^2-\vct{Z}_{\vct{v}}^2.
\]

For differentiation with respect to \(\lambda\), the generator from
Eq.\ \eqref{eq:similarity-generator} gives
\begin{align*}
\frac{\dd\pv{\sigma}_n'}{\dd\lambda}
&=\com{\pv{\sigma}_n'}{\pv{\Omega}},&
\frac{\dd\sigv'}{\dd\lambda}
&=\com{\sigv'}{\pv{\Omega}}.
\end{align*}
The coefficient form is obtained directly from
\begin{align*}
\pv{\Omega}
&=\frac{f\,\dfrac{\dd f}{\dd\lambda}
        -\vct{g}\cdot\dfrac{\dd\vct{g}}{\dd\lambda}}
        {f^2-\vct{g}^2}\\
&\quad+
\frac{f\,\dfrac{\dd\vct{g}}{\dd\lambda}
      -\dfrac{\dd f}{\dd\lambda}\vct{g}
      -i\,\vct{g}\times\dfrac{\dd\vct{g}}{\dd\lambda}}
     {f^2-\vct{g}^2}\cdot\sigv.
\end{align*}
Relative to the primed basis, its vector coefficient is written uniquely as
\[
\pv{\Omega}=\Omega_s+\tfrac12(\vct{k}+i\vct{\boldsymbol\omega})\cdot\sigv',
\qquad \vct{k},\vct{\boldsymbol\omega}\in\R^3.
\]
The same generator carried into the unprimed basis is defined by
\begin{equation}
\pv{\Omega}':=\pv{T}\pv{\Omega}\pv{T}^{-1}.
\label{eq:carried-similarity-generator}
\end{equation}
The definition in Eq.\ \eqref{eq:similarity-generator} then gives
\[
\pv{\Omega}'=\frac{\dd\pv{T}}{\dd\lambda}\pv{T}^{-1}.
\]
The coefficients are recovered through scalar extraction.
For \(n\in\{1,2,3\}\), the coefficients are given by cyclicity of the
algebraic scalar part:
\[
k_n+i\omega_n
=2\,\eqtext{scalar}(\pv{\Omega}\pv{\sigma}_n')
=2\,\eqtext{scalar}(\pv{\Omega}'\pv{\sigma}_n).
\]
The scalar part is central with respect to the basis, while
\begin{align*}
\frac{\dd\sigv'}{\dd\lambda}
&=(-\vct{\boldsymbol\omega}+i\vct{k})\times\sigv',\\
\frac{\dd\pv{Z}}{\dd\lambda}
&=\frac{\dd Z_s}{\dd\lambda}+
\left(\frac{\dd\vct{Z}_{\vct{v}}}{\dd\lambda}
      -\vct{Z}_{\vct{v}}\times\vct{\boldsymbol\omega}
      +i\,\vct{Z}_{\vct{v}}\times\vct{k}\right)\cdot\sigv'.
\end{align*}
These are algebraic moving-basis identities. No space-time, energy, or
proper-time interpretation is implied by similarity alone; a specified
Lorentz congruence and frame geometry are required for such an interpretation.

\subsubsection{Rotation}
Equation \eqref{eq:rotor} is obtained as the rotation case of the general
transform law by specializing the transform element to \(\pv{T}=\pv{R}\).
Sigma number two is assigned to the generator, so the exponential is circular
rather than hyperbolic. The rotor function defined in
Eq.\ \eqref{eq:rotor-element} has the expanded form
\begin{align*}
\pv{R}(\theta ,\vct{u}) &= \exp( \tfrac{1}{2} \theta (\vct{u} \cdot i \sigv ) ),\notag \\
&= \cos( \tfrac{1}{2} \theta ) + \sin( \tfrac{1}{2} \theta ) (\vct{u} \cdot i \sigv )\qquad \text{rotor}.
\end{align*}
\begin{equation*}
\pv{R}^{-1} = \exp( -\tfrac{1}{2} \theta (\vct{u} \cdot i \sigv ) ).
\end{equation*}
\(\theta\in\R\), \(\vct{u}\in\R^3\) is a unit real vector, and \(\pv{R}\)
is generated by the sigma-number-two element \(\vct{u}\cdot i\sigv\). The
inverse relation is
\(\pv{R}^{-1}=\pv{R}^{\mathsf H}=\pv{R}^{*\mathsf H}\).
The orientation sign is fixed by the substitution
\begin{equation*}
f=\cos(\theta/2),
\qquad
\vct{g}=i\sin(\theta/2)\vct{u}
\end{equation*}
into Eq.\ \eqref{eq:similarity-components}. Then
\begin{align*}
f^2-\vct{g}^2&=1,\\
f^2+\vct{g}^2&=\cos(\theta),\\
-2(\vct{g}\cdot\vct{r})\vct{g}
&=(1-\cos(\theta))(\vct{u}\cdot\vct{r})\vct{u},\\
-2if(\vct{g}\times\vct{r})
&=\sin(\theta)(\vct{u}\times\vct{r}).
\end{align*}
Therefore the transformed paravector is
\begin{align*}
\pv{X}' &= \pv{R}^{\mathsf H} \pv{X} \pv{R},\notag \\
&= t' + \vct{r}' \cdot \sigv.
\end{align*}
\begin{align*}
t' &= t,\notag \\
\vct{r}' &= \vct{r} \cos(\theta ) + (\vct{u} \times \vct{r}) \sin(\theta ),\notag \\
&\quad + (\vct{u} \cdot \vct{r}) (1 - \cos(\theta )) \vct{u},\notag \\
\qquad \text{Rodrigues's formula}.
\end{align*}
For a parameter-dependent rotor, the specialization of the general generator
is \(\pv{R}^{\mathsf H}(\dd\pv{R}/\dd\lambda)\). Its form carried to the
other side is
\((\dd\pv{R}/\dd\lambda)\pv{R}^{\mathsf H}\).
The following identities are obtained by differentiating
\(\pv{R}^{\mathsf H}\pv{R}=1\):
\[
\left(\pv{R}^{\mathsf H}\frac{\dd\pv{R}}{\dd\lambda}\right)^{\mathsf H}
=-\pv{R}^{\mathsf H}\frac{\dd\pv{R}}{\dd\lambda},
\qquad
\left(\frac{\dd\pv{R}}{\dd\lambda}\pv{R}^{\mathsf H}\right)^{\mathsf H}
=-\frac{\dd\pv{R}}{\dd\lambda}\pv{R}^{\mathsf H}.
\]
The rotation transport law is the specialization of the general
one-parameter similarity formula:
\begin{align*}
\frac{\dd\pv{X}'}{\dd\lambda}
&=\pv{R}^{\mathsf H}\frac{\dd\pv{X}}{\dd\lambda}\pv{R}
+\com{\pv{X}'}{\pv{R}^{\mathsf H}\dfrac{\dd\pv{R}}{\dd\lambda}}\\
&=\pv{R}^{\mathsf H}
\left(\frac{\dd\pv{X}}{\dd\lambda}
+\com{\pv{X}}{\dfrac{\dd\pv{R}}{\dd\lambda}\pv{R}^{\mathsf H}}\right)\pv{R}.
\end{align*}

\subsubsection{Spin Rotors}
Each oriented complementary bivector is identified by the pseudoscalar
without an index-order ambiguity:
\[
i\pv{\sigma}_1=\pv{\sigma}_2\pv{\sigma}_3,\qquad
i\pv{\sigma}_2=\pv{\sigma}_3\pv{\sigma}_1,\qquad
i\pv{\sigma}_3=\pv{\sigma}_1\pv{\sigma}_2.
\]
For unit \(\vct{n},\vct{u}\in\R^3\), the rotor function defined in
Eq.\ \eqref{eq:rotor-element} is evaluated as
\begin{align*}
\pv{R}(\phi,\vct{n})
&=\exp\!\left(\frac{\phi}{2}\,\vct{n}\cdot i\sigv\right)\\
&=\cos(\phi/2)
 +\sin(\phi/2)\,\vct{n}\cdot i\sigv .
\end{align*}
Then
\[
\pv{R}^{-1}=\pv{R}^{\mathsf H}=\pv{R}^{*\mathsf H},\qquad
\pv{R}^{\mathsf H}\pv{R}=1,
\]
and
\[
\pv{R}^{\mathsf H}(\vct{u}\cdot\sigv)\pv{R}
=\vct{u}'\cdot\sigv
\]
with \(\vct{u}'\) defined by Rodrigues's formula. The rotor itself is
single-valued and \(4\pi\)-periodic:
\[
\pv{R}(\phi+2\pi,\vct{n})=-\pv{R}(\phi,\vct{n}),\qquad
\pv{R}(\phi+4\pi,\vct{n})=\pv{R}(\phi,\vct{n}).
\]
The induced spatial rotation is two-to-one because the same conjugation is
produced by \(\pv{R}\) and \(-\pv{R}\).

In the Pauli matrix realization, the two-by-two representation is transposed
and complex-conjugated by \(H\). A general
\(\pv{\psi}\in\C+\C^3\) is therefore a full two-by-two matrix, not one
selected two-entry column: its two columns are rotated independently by
left multiplication. One column is selected by a fixed projector on the
right when a single Pauli wave function is intended.
Under the active convention associated with this vector action,
\[
\pv{\psi}'=\pv{R}^{\mathsf H}\pv{\psi},\qquad
\pv{\psi}'\pv{\psi}'^{\mathsf H}
=\pv{R}^{\mathsf H}(\pv{\psi}\pv{\psi}^{\mathsf H})\pv{R}.
\]
A scalar pairing, normalization, expectation map, projected carrier, and
specified operator action are additionally required for a physical spin state. The
algebra action is thereby separated from the choice of a single physical
spin state.

\subsubsection{Spin-Axis Projectors}
The left action of a chosen unit axis \(\vct{u}\in\R^3\) is resolved by the
projector family from Eq.\ \eqref{eq:projectors}. For clarity, that
definition is restated here as an ordinary equality:
\[
\pv{\Pi}(\pm)=\tfrac12(1\pm\vct{u}\cdot\sigv),
\qquad
\vct{u}^{\,2}=1,
\qquad
\text{from Eq.\ \eqref{eq:projectors}.}
\]
In addition to the
idempotence, orthogonality, completeness, and axis-eigenvalue identities
listed there, the following involution actions are obtained:
\begin{align*}
\pv{\Pi}(\pm)^{\mathsf H}&=\pv{\Pi}(\pm),\\
\pv{\Pi}(\pm)^*&=\pv{\Pi}(\mp).
\end{align*}
Accordingly,
\begin{equation*}
\vct{u}\cdot\sigv=(+1)\pv{\Pi}(+)+(-1)\pv{\Pi}(-),
\end{equation*}
so the eigenvalues \(+1\) and \(-1\) of the axis operator are carried by the
two projector parts.

In the Pauli matrix realization, each \(\pv{\Pi}(\pm)\) is of rank one. A general
element of \(\C+\C^3\), however, is a full \(2\times2\) complex matrix,
not one selected two-entry column. The projection formulas are therefore
valid column by column. A single Pauli wave-function column is described after a
fixed projector is placed on the right; without that restriction, two
independent columns are contained in the full paravector.

If the axis is rotated by the rotor \(\pv{R}\) from Eq.\ \eqref{eq:rotor},
the transformed axis and its second projector pair satisfy
\begin{equation}
\vct{u}'\cdot\sigv=\herm{\pv{R}}(\vct{u}\cdot\sigv)\pv{R},
\qquad
\pv{\Pi}'(\pm)=\tfrac12(1\pm \vct{u}'\cdot\sigv)
=\herm{\pv{R}}\pv{\Pi}(\pm)\pv{R}.
\label{eq:rotated-axis-projectors}
\end{equation}
Then
\begin{equation*}
\vct{u}'\cdot\sigv=(+1)\pv{\Pi}'(+)+(-1)\pv{\Pi}'(-).
\end{equation*}

For an algebra element \(\pv{\psi}\in\C+\C^3\), specialize the left-projected
parts from Eq.\ \eqref{eq:projectors}:
\begin{equation*}
\pv{\psi}(\pm)=\pv{\Pi}(\pm)\pv{\psi}.
\end{equation*}
The following identities are obtained:
\begin{align*}
(\vct{u}\cdot\sigv)\pv{\psi}(\pm)&=\pm\pv{\psi}(\pm),\\
\pv{\psi}&=\pv{\psi}(+)+\pv{\psi}(-).
\end{align*}
Thus each spinor column is split into the two eigenspaces of the chosen axis.
The Hermitian operation is represented by
\begin{align*}
\herm{\pv{\psi}(\pm)}
&=\herm{\pv{\psi}}\,\pv{\Pi}(\pm)^{\mathsf H}\notag\\
&=\herm{\pv{\psi}}\,\pv{\Pi}(\pm),
\end{align*}
and therefore
\begin{align*}
\herm{\pv{\psi}(\pm)}\pv{\psi}(\pm)&=\herm{\pv{\psi}}\,\pv{\Pi}(\pm)\pv{\psi},\\
\herm{\pv{\psi}(\pm)}\pv{\psi}(\mp)&=\herm{\pv{\psi}}\,\pv{\Pi}(\pm)\pv{\Pi}(\mp)\pv{\psi}=0,\\
\herm{\pv{\psi}(\mp)}\pv{\psi}(\pm)&=\herm{\pv{\psi}}\,\pv{\Pi}(\mp)\pv{\Pi}(\pm)\pv{\psi}=0.
\end{align*}
For two algebra elements \(\pv{\phi},\pv{\psi}\), the pointwise algebraic
pairing associated with \(\pv{M}\) is
\(\eqtext{scalar}(\herm{\pv{\phi}}\pv{M}\pv{\psi})\). For a nonzero
\(\pv{\psi}\), the corresponding normalized diagonal matrix element is
\begin{equation*}
\frac{\eqtext{scalar}(\herm{\pv{\psi}}\pv{M}\pv{\psi})}
     {\eqtext{scalar}(\herm{\pv{\psi}}\pv{\psi})}.
\end{equation*}
When \(\pv{M}\) is Hermitian for the relevant pairing, this
diagonal element is real and is an expectation value. For a general
\(\pv{M}\), it is only a normalized matrix element and may be complex.
Scalar extraction is one half of the matrix trace. Consequently, for a full
matrix \(\pv{\psi}\) the numerator and denominator are each one half of the
sum of their two ordinary column pairings. After one column is selected by a
fixed right projector, they are each one half of the corresponding ordinary
selected-column quantity. The common factor is canceled in the normalized
matrix element. The spatially integrated quantum bracket is defined
separately in Eq.\ \eqref{eq:adjoint-paravector-braket}.

Rotational covariance is made explicit by the same convention. Under the
active rotation
\begin{equation*}
\pv{\psi}'=\herm{\pv{R}}\pv{\psi},
\qquad
\pv{M}'=\herm{\pv{R}}\pv{M}\pv{R},
\qquad
\pv{R}\herm{\pv{R}}=1,
\end{equation*}
both the numerator and denominator of the normalized matrix element are
unchanged:
\begin{equation}
\frac{\eqtext{scalar}(\herm{\pv{\psi}'}\pv{M}'\pv{\psi}')}
     {\eqtext{scalar}(\herm{\pv{\psi}'}\pv{\psi}')}
=
\frac{\eqtext{scalar}(\herm{\pv{\psi}}\pv{M}\pv{\psi})}
     {\eqtext{scalar}(\herm{\pv{\psi}}\pv{\psi})}.
\label{eq:rotor-expectation-covariance}
\end{equation}
Equivalently, if the observable \(\pv{M}\) is held fixed while only the
carrier is actively rotated, its value in the original carrier is computed
with the rotated-back observable \(\pv{R}\pv{M}\herm{\pv{R}}\). This is the
active-versus-passive equivalence expressed in
Eq.\ \eqref{eq:rotated-axis-projectors}.

For an axis observable, the following identity is obtained by cyclic scalar
extraction:
\begin{equation}
\eqtext{scalar}\!\left(
  \herm{\pv{\psi}}(\vct{u}\cdot\sigv)\pv{\psi}
\right)
=
\vct{u}\cdot
\eqtext{vector}\!\left(\pv{\psi}\herm{\pv{\psi}}\right).
\label{eq:axis-polarization-identity}
\end{equation}
Consequently its normalized pointwise value is
\begin{equation*}
\frac{\vct{u}\cdot
  \eqtext{vector}(\pv{\psi}\herm{\pv{\psi}})}
 {\eqtext{scalar}(\herm{\pv{\psi}}\pv{\psi})}.
\end{equation*}
In the selected-column matrix realization, twice the numerator is the
ordinary column matrix element and twice the denominator is the ordinary
column norm, so the common factor is canceled. For a spatial wave field the
same statement is valid with both quantities integrated as in
Eqs.\ \eqref{eq:adjoint-paravector-braket} and \eqref{eq:expectation-map}.

The unnormalized pointwise numerator for the rotated axis is expanded as
\begin{align*}
\eqtext{scalar}\!\left(\herm{\pv{\psi}}(\vct{u}'\cdot\sigv)\pv{\psi}\right)
&=\eqtext{scalar}\!\left(\herm{\pv{\psi}(+)}(\vct{u}'\cdot\sigv)\pv{\psi}(+)\right)\notag\\
&\quad+\eqtext{scalar}\!\left(\herm{\pv{\psi}(+)}(\vct{u}'\cdot\sigv)\pv{\psi}(-)\right)\notag\\
&\quad+\eqtext{scalar}\!\left(\herm{\pv{\psi}(-)}(\vct{u}'\cdot\sigv)\pv{\psi}(+)\right)\notag\\
&\quad+\eqtext{scalar}\!\left(\herm{\pv{\psi}(-)}(\vct{u}'\cdot\sigv)\pv{\psi}(-)\right).
\end{align*}

The algebra governing these same-side and cross-side terms is
\begin{align*}
(\vct{u}'\cdot\sigv)(\vct{u}\cdot\sigv)
&=\vct{u}'\cdot\vct{u}+i(\vct{u}'\times\vct{u})\cdot\sigv,\\
(\vct{u}\cdot\sigv)(\vct{u}'\cdot\sigv)
&=\vct{u}'\cdot\vct{u}-i(\vct{u}'\times\vct{u})\cdot\sigv,\\
(\vct{u}\cdot\sigv)(\vct{u}'\cdot\sigv)(\vct{u}\cdot\sigv)
&=\bigl(2(\vct{u}'\cdot\vct{u})\vct{u}-\vct{u}'\bigr)\cdot\sigv.
\end{align*}
The second axis is decomposed into its \(\eqtext{par}\) and
\(\eqtext{prp}\) parts by the axis functions from
Eqs.\ \eqref{eq:par-function} and \eqref{eq:prp-function}:
\begin{equation*}
\vct{u}'=\eqtext{par}(\vct{u}')+\eqtext{prp}(\vct{u}'),
\qquad
\eqtext{par}(\vct{u}')=(\vct{u}\cdot\vct{u}')\vct{u}.
\end{equation*}
Then the projector products are
\begin{equation*}
\pv{\Pi}(\pm)(\vct{u}'\cdot\sigv)\pv{\Pi}(\pm)=\pm(\vct{u}\cdot\vct{u}')\pv{\Pi}(\pm),
\end{equation*}
\begin{equation*}
\pv{\Pi}(\pm)(\vct{u}'\cdot\sigv)\pv{\Pi}(\mp)
=\frac12\Bigl(\eqtext{prp}(\vct{u}')\pm i\,\vct{u}\times\vct{u}'\Bigr)\cdot\sigv\,\pv{\Pi}(\mp).
\end{equation*}
The diagonal terms are therefore weighted by the axis overlap \(\vct{u}\cdot\vct{u}'\), while the off-diagonal terms are controlled by the \(\eqtext{prp}\) and commutator parts of the rotated axis.

If \(\theta\) is the angle between \(\vct{u}\) and \(\vct{u}'\), then
\begin{equation*}
\vct{u}'\cdot\vct{u}=\cos(\theta),
\end{equation*}
Because scalar extraction is one half of the matrix trace, the
trace-normalized rank-one projector overlap is
\begin{equation*}
2\,\eqtext{scalar}\!\bigl(\pv{\Pi}(s)\pv{\Pi}'(s')\bigr)
=\frac12\bigl(1+ss'\,\vct{u}\cdot\vct{u}'\bigr),
\qquad s,s'\in\{+1,-1\}.
\end{equation*}
For a normalized irreducible axis eigenstate this overlap is the transition
probability. In particular, the same-side and opposite-side probabilities
are
\begin{equation*}
\frac12(1+\vct{u}'\cdot\vct{u})=\cos^2(\theta/2),
\qquad
\frac12(1-\vct{u}'\cdot\vct{u})=\sin^2(\theta/2).
\end{equation*}
If \(\vct{u}'\) is obtained by rotating \(\vct{u}\) through angle \(\phi\)
about a unit axis \(\vct{n}\), the overlap is given by Rodrigues's formula:
\begin{align*}
\vct{u}'\cdot\vct{u}
&=(\vct{u}\cdot\vct{u})\cos(\phi)
+\bigl(\vct{u}\cdot(\vct{n}\times\vct{u})\bigr)\sin(\phi)\notag\\
&\quad+(\vct{n}\cdot\vct{u})\bigl(1-\cos(\phi)\bigr)(\vct{n}\cdot\vct{u})\notag\\
&=(\vct{n}\cdot\vct{u})^2+\bigl(1-(\vct{n}\cdot\vct{u})^2\bigr)\cos(\phi).
\end{align*}
In the perpendicular case \(\vct{n}\cdot\vct{u}=0\), this is reduced to
\begin{equation*}
\vct{u}'\cdot\vct{u}=\cos(\phi),
\end{equation*}
hence
\begin{equation*}
\frac12(1+\vct{u}'\cdot\vct{u})=\cos^2(\phi/2),
\qquad
\frac12(1-\vct{u}'\cdot\vct{u})=\sin^2(\phi/2).
\end{equation*}
These are the familiar spin-\(\tfrac12\) transition probabilities. Without
normalization and the irreducible-state qualification, only projector
overlaps are represented.

\subsubsection{Ladder Operations About the z Axis}
For raising and lowering about the \(z\) axis, the following axis and
projector pair are chosen:
\begin{equation*}
\vct{u}=\vct{e}_3=(0,0,1),
\qquad
\pv{\Pi}(\pm)=\tfrac12(1\pm \pv{\sigma}_3).
\end{equation*}
The ladder elements already defined in Eq.\ \eqref{eq:ladder} are used:
\begin{equation*}
\pv{a}(\pm)=\tfrac12(\pv{\sigma}_1\pm i\pv{\sigma}_2).
\end{equation*}
The commutator relation is satisfied:
\begin{equation*}
\com{\pv{\sigma}_3}{\pv{a}(\pm)}=\pm 2\pv{a}(\pm),
\end{equation*}
so the \(\pv{\sigma}_3\) eigenvalue is raised by \(\pv{a}(+)\) and lowered
by \(\pv{a}(-)\).

Their projector products are
\begin{equation*}
\pv{a}(+)\pv{a}(-)=\pv{\Pi}(+),
\qquad
\pv{a}(-)\pv{a}(+)=\pv{\Pi}(-),
\end{equation*}
and each ladder operator is nilpotent:
\begin{equation*}
\pv{a}(\pm)^2=0.
\end{equation*}
The one-sided projector actions are
\begin{equation*}
\begin{aligned}
\pv{a}(+)\pv{\Pi}(+) &= 0, &
\pv{a}(+)\pv{\Pi}(-) &= \pv{a}(+), \\
\pv{\Pi}(+)\pv{a}(+) &= \pv{a}(+), &
\pv{\Pi}(-)\pv{a}(+) &= 0, \\
\pv{a}(-)\pv{\Pi}(-) &= 0, &
\pv{a}(-)\pv{\Pi}(+) &= \pv{a}(-), \\
\pv{\Pi}(-)\pv{a}(-) &= \pv{a}(-), &
\pv{\Pi}(+)\pv{a}(-) &= 0,
\end{aligned}
\end{equation*}
by which the raising/lowering interpretation is made explicit. The original
sigma generators are recovered from the same ladder/projector data:
\begin{equation*}
\begin{aligned}
\pv{\sigma}_1&=\pv{a}(+)+\pv{a}(-),\qquad
\pv{\sigma}_2=\frac{\pv{a}(+)-\pv{a}(-)}{i},\\
\pv{\sigma}_3&=\pv{\Pi}(+)-\pv{\Pi}(-).
\end{aligned}
\end{equation*}

\subsubsection{Spin-Axis Noncommutativity}
For unit directions \(\vct{u},\vct{u}'\in\R^3\),
\begin{align*}
\acom{\vct{u}\cdot\sigv}{\vct{u}'\cdot\sigv}
&=2\,\vct{u}\cdot\vct{u}',\\
\com{\vct{u}\cdot\sigv}{\vct{u}'\cdot\sigv}
&=2i(\vct{u}\times\vct{u}')\cdot\sigv.
\end{align*}
The axis overlap is recorded by the anticommutator, while the oriented area
spanned by the two axes is recorded by the commutator. Thus exact commutation
of the axis operators is obtained when the directions are parallel or
antiparallel. These identities are purely algebraic. A normalized state,
expectation functional, and state-dependent variances are additionally
required for a Robertson or Robertson--Schr\"odinger uncertainty relation;
no such inequality is inferred from the
commutator alone \cite{robertson1929}.


\subsection{Quaternion Correspondence}
After development of the circular exponential and rotor law, the same sigma-number-two algebra can be restated in Hamilton's quaternion language without introduction of a new multiplication rule. The correspondence is collected here after the rotation discussion.
\subsubsection{Quaternion Basis Under Star and H}
Hamilton's quaternion units are also reproduced by the same sigma-number-two
basis used in Eq.\ \eqref{eq:exp-bivector}, so the rotation algebra can be
written in quaternion form without changing any multiplication rule
\cite{hamilton1844}:
\begin{equation}
\begin{aligned}
\pv{q}_1&=-\pv{\sigma}_2\pv{\sigma}_3,\qquad
\pv{q}_2=-\pv{\sigma}_3\pv{\sigma}_1,\\
\pv{q}_3&=-\pv{\sigma}_1\pv{\sigma}_2,\qquad
\vct{q}:=\vcol{\pv{q}_1,\pv{q}_2,\pv{q}_3}=-i\sigv.
\end{aligned}
\label{eq:quaternion-basis}
\end{equation}
The standard quaternion identities are satisfied by these basis elements:
\begin{equation*}
\pv{q}_1^2=\pv{q}_2^2=\pv{q}_3^2=\pv{q}_1\pv{q}_2\pv{q}_3=-1.
\end{equation*}
Thus the quaternion units are not new generators; the existing
sigma-number-two basis is only relabeled.

If \(\vct{B}=\vcol{B_1,B_2,B_3}\in\R^3\), then
\(\vct{B}\cdot\vct{q}=-i\,\vct{B}\cdot\sigv\). Since
\(\pv{q}_k=-i\pv{\sigma}_k\), the following relations are obtained
immediately from the previously defined star and \(H\) rules:
\begin{equation*}
\vct{q}^*=\vct{q},
\qquad
\vct{q}^{\mathsf H}=-\vct{q}.
\end{equation*}
Thus the quaternion basis is fixed by the conjugate operator, while its
sign is changed by the Hermitian operation.

\subsubsection{Quaternion Products}
For \(a\in\R\) and \(\vct{B}\in\R^3\), the quaternion is defined by
\begin{equation}
\pv{Q}:=a+\vct{B}\cdot\vct{q}
=a-i\,\vct{B}\cdot\sigv.
\label{eq:quaternion-paravector}
\end{equation}
A real scalar coefficient and a purely imaginary vector coefficient are
carried, so \(\pv{Q}\in\Pspace(\C)\).
In complex-paravector notation, the same quaternion may be written as
\begin{equation*}
\pv{Q}=\pv{X}-i\pv{Y},
\qquad
\pv{X}=a+\vct{A}\cdot\sigv,
\qquad
\pv{Y}=b+\vct{B}\cdot\sigv,
\end{equation*}
with the quaternion specialization
\begin{equation*}
\vct{A}=0,
\qquad
b=0.
\end{equation*}
Under this specialization, \(\pv{X}=a\) is purely scalar and
\(\pv{Y}=\vct{B}\cdot\sigv\) is purely vector. Because
\(\vct{q}^*=\vct{q}\),
\begin{equation*}
\pv{Q}^*=\pv{Q}.
\end{equation*}
Therefore
\begin{align}
\pv{Q}\pv{Q}^*&=\pv{Q}^2\notag\\
&=(a-i\,\vct{B}\cdot\sigv)^2\notag\\
&=a^2-2ia\,\vct{B}\cdot\sigv+(-i)^2(\vct{B}\cdot\sigv)^2\notag\\
&=a^2-\vct{B}^2-2ia\,\vct{B}\cdot\sigv.
\label{eq:quaternion-square}
\end{align}
The factor \(a\) must be carried by the linear term. This is just the
sigma-number-two multiplication law written in quaternion language.

Since \(\vct{q}^{\mathsf H}=-\vct{q}\), the Hermitian counterpart of \(\pv{Q}\) is
\begin{equation*}
\pv{Q}^{\mathsf H}=a-\vct{B}\cdot\vct{q}
=a+i\,\vct{B}\cdot\sigv.
\end{equation*}
The result of multiplication by the Hermitian counterpart is
\begin{equation}
\pv{Q}\pv{Q}^{\mathsf H}=a^2+\vct{B}^2.
\label{eq:quaternion-norm}
\end{equation}
Thus the standard positive quaternion norm square is recovered by the
Hermitian product. Here ``norm square'' means the nonnegative real number
that measures the coefficient magnitude. The ordinary algebra square is
instead recorded by
Eq.\ \eqref{eq:quaternion-square}.

\subsubsection{Quaternion Form of the Rotor Law}
A quaternion satisfying \(\pv{Q}\pv{Q}^{\mathsf H}=1\) is called a unit
quaternion. For such a quaternion and a real vector
\(\vct{r}\in\R^3\), the rotation law is
\begin{equation}
\vct{r}'\cdot\vct{q}
=\pv{Q}(\vct{r}\cdot\vct{q})\pv{Q}^{\mathsf H}.
\label{eq:quaternion-rotation}
\end{equation}
This is the quaternion version of the rotor action
\(\vct{r}'\cdot\sigv=\herm{\pv{R}}(\vct{r}\cdot\sigv)\pv{R}\). With the
choice
\begin{equation*}
\pv{Q}=\cos\!\left(\frac{\theta}{2}\right)
 +\sin\!\left(\frac{\theta}{2}\right)\vct{u}\cdot\vct{q}
\end{equation*}
the relation \(\pv{Q}=\pv{R}^{\mathsf H}\) is obtained because
\(\vct{q}=-i\sigv\). Consequently
\[
\pv{Q}(\vct{r}\cdot\vct{q})\pv{Q}^{\mathsf H}
=\pv{R}^{\mathsf H}(\vct{r}\cdot\vct{q})\pv{R},
\]
which has exactly the same orientation as the sigma-vector rotor law. The
equivalent two-sided quaternion product
\(\pv{Q}^{\mathsf H}(\vct{r}\cdot\vct{q})\pv{Q}\) is required under the
alternative choice \(\pv{Q}=\pv{R}\). Thus quaternion rotations are
the same circular half-angle algebra, with the order fixed by the selected
choice of positive rotation direction.


\subsection{Matrix and Auxiliary Realizations}
The algebra above is complete on its own, but two realizations are useful. In the Pauli representation, one paravector is identified with a two-by-two complex matrix and \(\pv{Z}\adjf{\pv{Z}}\) is identified with its determinant \cite{pauli1927}. An auxiliary two-by-two block map whose entries are paravectors is defined after the standard realization; after the Pauli realization it is represented by a four-by-four complex matrix. It is used as a packaging map, not as a second representation or an algebra homomorphism.
\subsubsection{Standard Two-by-Two Matrix Realization}
\label{sec:matrix-realization}
The algebra fixed by the three sigma generators and their multiplication
rule is commonly called the real Pauli algebra or the three-dimensional
real Clifford algebra. These names refer here only to the algebra already
defined. Each of its elements can be translated one-to-one into a
two-by-two complex matrix so that both addition and multiplication are
preserved. A one-to-one translation with these properties is called a
faithful matrix representation. The same matrix system also represents the
part of the usual flat space-time algebra built from products containing an
even number of space-time basis directions
\cite{clifford1878,pauli1927,hestenes1966}. Its eight real coefficients are
grouped by the internal pseudoscalar into four complex coefficients: one
scalar and one three-vector. Only the previously defined operations are
used below.

Square brackets and the physical row and column positions are used for
matrices. Thus the generic two-by-two matrix is
\(\smat{a&b\\c&d}\), where every entry is a complex scalar.
Each matrix operation used below is defined separately:
\begin{equation}
\I:=
\begin{bmatrix}
1&0\\
0&1
\end{bmatrix}.
\label{eq:matrix-identity}
\end{equation}
Multiplication by \(\I\) leaves any matrix of the matching size unchanged,
so \(\I\) is the matrix identity.
\begin{equation}
\eqtext{tr}\!\left(
\begin{bmatrix}
a&b\\
c&d
\end{bmatrix}
\right):=a+d.
\label{eq:matrix-trace}
\end{equation}
The sum of the diagonal entries is called the matrix trace.
\begin{equation}
\eqtext{diag}(A,B):=
\begin{bmatrix}
A&0\\
0&B
\end{bmatrix}.
\label{eq:matrix-diag}
\end{equation}
The function \(\eqtext{diag}(\cdot)\) places its inputs on the diagonal and
sets all other entries to zero.
The dimension of \(\I\) is fixed by its surrounding matrix product;
Eq.\ \eqref{eq:matrix-identity} displays the two-by-two case. When more
than two inputs are supplied to \(\eqtext{diag}(\cdot)\), the same rule is
used: the inputs are placed in order on the diagonal and every other entry
is set to zero.
\begin{equation}
\mtrans{\begin{bmatrix}a&b\\c&d\end{bmatrix}}
:=\begin{bmatrix}a&c\\b&d\end{bmatrix}.
\label{eq:matrix-transpose}
\end{equation}
The superscript \(\mathsf T\) exchanges rows and columns and is called the
transpose.
For a scalar entry \(z=x+iy\), the matrix-only conjugate is
\begin{equation}
\mconj{z}:=x-iy.
\label{eq:matrix-entry-conjugate}
\end{equation}
\begin{equation}
\mconj{\begin{bmatrix}a&b\\c&d\end{bmatrix}}
:=\begin{bmatrix}
\mconj{a}&\mconj{b}\\
\mconj{c}&\mconj{d}
\end{bmatrix}.
\label{eq:matrix-conjugate}
\end{equation}
The superscript \(\mathsf C\) changes \(i\) to \(-i\) in every entry and is
called entrywise complex conjugation.
\begin{equation}
\mherm{\begin{bmatrix}a&b\\c&d\end{bmatrix}}
:=\begin{bmatrix}
\mconj{a}&\mconj{c}\\
\mconj{b}&\mconj{d}
\end{bmatrix}.
\label{eq:matrix-hermitian}
\end{equation}
The superscript \(\mathsf{CT}\) applies entrywise conjugation and then
transpose; it is called the Hermitian operation.
\begin{equation}
\adjf{\begin{bmatrix}a&b\\c&d\end{bmatrix}}
:=\begin{bmatrix}d&-b\\-c&a\end{bmatrix}.
\label{eq:matrix-adjugate}
\end{equation}
The displayed exchange and sign change of the entries define the matrix
adjugate.
\begin{equation}
\detf{\begin{bmatrix}a&b\\c&d\end{bmatrix}}:=ad-bc.
\label{eq:matrix-determinant}
\end{equation}
The scalar \(ad-bc\) is the matrix determinant.
Entrywise conjugation and transpose commute, so either order gives the
matrix-only operation \(\mathsf{CT}\). The superscript \(\mathsf C\) is
reserved for matrices and their scalar entries. It is not an additional
operation on the sigma algebra itself. An operation defined by the sigma
rules without choosing a matrix translation is called intrinsic; star and
\(H\) are intrinsic in this sense.

The standard matrix translation is
\begin{equation}
\mat{\sigma}_1=\begin{bmatrix}0&1\\1&0\end{bmatrix},
\quad
\mat{\sigma}_2=\begin{bmatrix}0&-i\\i&0\end{bmatrix},
\quad
\mat{\sigma}_3=\begin{bmatrix}1&0\\0&-1\end{bmatrix}.
\label{eq:pauli-matrices}
\end{equation}
The entry \(i\) is the same central unit defined in
Eq.\ \eqref{eq:basis}; as a scalar matrix entry it obeys \(i^2=-1\).
No second imaginary unit is introduced. The distinction in the matrix
model lies in the operation being applied. Every displayed \(i\), including
the one built into \(\mat{\sigma}_2\), is conjugated by entrywise \(\mathsf C\).
Every sigma generator is additionally flipped by intrinsic star.

The pseudoscalar is represented by the matrix product
\(\mat{\sigma}_1\mat{\sigma}_2\mat{\sigma}_3=i\I\). In this realization,
\begin{equation*}
\mat{\sigma}_n^2=\I,
\qquad
\mat{\sigma}_m\mat{\sigma}_n=-\mat{\sigma}_n\mat{\sigma}_m\quad(m\neq n),
\qquad
\mat{\sigma}_1\mat{\sigma}_2\mat{\sigma}_3=i\I.
\end{equation*}
Thus Eq.\ \eqref{eq:pauli-matrices} is a concrete model of the sigma
generators and their pseudoscalar.

For the generic complex paravector
\(\pv{Z}=S+\vct{V}\cdot\sigv\), its represented matrix is denoted by
\(\mat{M}\):
\begin{equation}
\mat{M}:=
\begin{bmatrix}
S+V_3&V_1-iV_2\\
V_1+iV_2&S-V_3
\end{bmatrix}.
\label{eq:matrix-paravector-map}
\end{equation}
Thus \(\pv{Z}\) is represented by \(\mat{M}\).
Conversely, if \(\mat{M}=\smat{a&b\\c&d}\), then
\begin{equation}
\begin{aligned}
S&=\frac{a+d}{2},&
V_1&=\frac{b+c}{2},\\
V_2&=\frac{c-b}{2i},&
V_3&=\frac{a-d}{2}.
\end{aligned}
\label{eq:matrix-paravector-inverse-map}
\end{equation}
The realization in both directions is therefore given by
Eqs.\ \eqref{eq:matrix-paravector-map} and
\eqref{eq:matrix-paravector-inverse-map}.

For a real paravector
\(\pv{X}=t+\vct{r}\cdot\sigv\in\Pspace(\R)\), with
\(\vct{r}=\vcol{r_1,r_2,r_3}\in\R^3\), the matrix form is
\begin{equation}
\begin{gathered}
\pv{X}=t+\vct{r}\cdot\sigv\\
\eqtext{is represented by}\\
\begin{bmatrix}
t+r_3&r_1-ir_2\\
r_1+ir_2&t-r_3
\end{bmatrix},\\
\detf{\pv{X}}=t^2-\vct{r}^2.
\end{gathered}
\label{eq:matrix-X}
\end{equation}
The real space-time expression \(t^2-\vct{r}^2\), conventionally called the
Minkowski quadratic form and defined in
Eq.\ \eqref{eq:real-paravector-metric}, is reproduced by
Eq.\ \eqref{eq:matrix-X}, together with
\begin{align*}
\eqtext{tr}(\mat{M})&=2t,\\
x^2-\eqtext{tr}(\mat{M})x+\detf{\mat{M}}&=0,\\
x&=t\pm\lVert\vct{r}\rVert.
\end{align*}
Here \(\mat{M}\) denotes the displayed matrix representing \(\pv{X}\), and
\(x\) is only the polynomial input. The last two lines show that
\(t+\lVert\vct{r}\rVert\) and \(t-\lVert\vct{r}\rVert\) are the two roots
of the displayed matrix equation. For each root there is a nonzero column that is
merely multiplied by that root when acted on by the matrix. Scalars with
this property are called eigenvalues. They give the matrix form of the
split into scalar part \(t\) and vector magnitude
\(\lVert\vct{r}\rVert\).

The sigma-basis action not captured by entrywise \(\mathsf C\) is accounted for by
the fixed matrix
\begin{equation}
\Tmat:=
\begin{bmatrix}
0&1\\
-1&0
\end{bmatrix}.
\label{eq:matrix-T}
\end{equation}
It satisfies \(\Tmat^{-1}=-\Tmat\).
The complete intrinsic-to-matrix dictionary is
\begin{equation}
\begin{aligned}
\pv{Z}
&\quad\eqtext{is represented by}\quad \mat{M},\\
\conj{\pv{Z}}
&\quad\eqtext{is represented by}\quad
\conj{\mat{M}},\\
\conj{\mat{M}}
&=\Tmat\,\mconj{\mat{M}}\,\Tmat^{-1}
\\
&=\begin{bmatrix}
\mconj{d}&-\mconj{c}\\
-\mconj{b}&\mconj{a}
\end{bmatrix},\\
\herm{\pv{Z}}
&\quad\eqtext{is represented by}\quad
\herm{\mat{M}},\\
\herm{\mat{M}}
&=\mherm{\mat{M}}
\\
&=\begin{bmatrix}
\mconj{a}&\mconj{c}\\
\mconj{b}&\mconj{d}
\end{bmatrix},\\
\adjf{\pv{Z}}
&=\cherm{\pv{Z}},\\
\adjf{\pv{Z}}
&\quad\eqtext{is represented by}\quad
\adjf{\mat{M}},\\
\adjf{\mat{M}}
&=\Tmat\,\mtrans{\mat{M}}\,\Tmat^{-1}
\\
&=\begin{bmatrix}d&-b\\-c&a\end{bmatrix},\\
\detf{\pv{Z}}
&\quad\eqtext{is represented by}\quad
\detf{\mat{M}}=ad-bc,\\
0&\quad\eqtext{is represented by}\quad
\begin{bmatrix}0&0\\0&0\end{bmatrix},\\
1&\quad\eqtext{is represented by}\quad\I,\\
\pv{Z}^{-1}
&\quad\eqtext{is represented by}\quad
\mat{M}^{-1},\\
\mat{M}^{-1}
&=\frac{\adjf{\mat{M}}}{\detf{\mat{M}}},
\qquad \detf{\mat{M}}\neq0.
\end{aligned}
\label{eq:matrix-operation-dictionary}
\end{equation}
Thus \(\herm{\mat{M}}=\mherm{\mat{M}}\) is the ordinary Hermitian
adjoint. A matrix is Hermitian precisely when
\(\herm{\mat{M}}=\mat{M}\), while the matrix adjugate is always denoted by
\(\adjf{\mat{M}}\). Crucially, intrinsic star is not matrix \(\mathsf C\):
\(\conj{\mat{M}}\neq\mconj{\mat{M}}\) in general.

The star on the left is the intrinsic sigma conjugate transported into
the matrix realization; the \(\mathsf C\) on the right is only entrywise complex
conjugation. They agree on central complex scalars, but not on a generic
paravector. When \(\pv{Z}\) is represented by \(\mat{M}\), the direct bridge
is that \(\herm{\pv{Z}}\) is represented by \(\mherm{\mat{M}}\). Intrinsic star
instead requires the basis correction
\(\conj{\mat{M}}=\Tmat\mconj{\mat{M}}\Tmat^{-1}\).

The extraction functions have the equally direct matrix form
\begin{equation}
\begin{aligned}
\eqtext{scalar}(\pv{Z})
&\quad\eqtext{is represented by}\quad\tfrac12\eqtext{tr}(\mat{M}),\\
\eqtext{vector}(\pv{Z})\cdot\sigv
&\quad\eqtext{is represented by}\quad
\mat{M}-\tfrac12\eqtext{tr}(\mat{M})\I,\\
\eqtext{real}(\pv{Z})
&\quad\eqtext{is represented by}\quad
\tfrac12\bigl(\mat{M}+\herm{\mat{M}}\bigr),\\
\eqtext{img}(\pv{Z})
&\quad\eqtext{is represented by}\quad
\frac{\mat{M}-\herm{\mat{M}}}{2i}.
\end{aligned}
\label{eq:matrix-extraction-dictionary}
\end{equation}
Consequently, the algebraic determinant is exactly the
ordinary two-by-two matrix determinant:
\begin{equation}
\mat{M}\adjf{\mat{M}}
=\adjf{\mat{M}}\mat{M}
=\detf{\mat{M}}\I.
\label{eq:matrix-adjugate-product}
\end{equation}

The operations are distinguished by whether they conjugate complex scalar
coefficients and whether they reverse the order of a product. Matrix
\(\mathsf C\) conjugates the coefficients and preserves product order.
Matrix \(\mathsf T\) leaves the coefficients unchanged and reverses product
order. Their composition \(\mathsf{CT}\), represented in the sigma algebra
by \(H\), both conjugates coefficients and reverses product order. Intrinsic
star also conjugates coefficients but preserves product order. The combined
intrinsic operation \(*H\), represented by the matrix adjugate, leaves
coefficients unchanged and reverses product order. Consequently,
Eq.\ \eqref{eq:hermitian-adjugate-identity} is represented by
\begin{equation*}
\herm{\mat{M}}
=\conj{\adjf{\mat{M}}}
=\adjf{\conj{\mat{M}}},
\end{equation*}
where the star in \(\conj{\mat{M}}\) remains the induced sigma conjugate,
not entrywise \(\mathsf C\).

The distinction is made concrete by the basis actions:
\begin{equation}
\begin{aligned}
\mconj{\mat{\sigma}_1}&=\mat{\sigma}_1,
&
\mconj{\mat{\sigma}_2}&=-\mat{\sigma}_2,
&
\mconj{\mat{\sigma}_3}&=\mat{\sigma}_3,\\
\mtrans{\mat{\sigma}_1}&=\mat{\sigma}_1,
&
\mtrans{\mat{\sigma}_2}&=-\mat{\sigma}_2,
&
\mtrans{\mat{\sigma}_3}&=\mat{\sigma}_3,\\
\mherm{\mat{\sigma}_n}&=\mat{\sigma}_n,
&
\conj{\mat{\sigma}_n}&=-\mat{\sigma}_n,
&
\cherm{\mat{\sigma}_n}&=-\mat{\sigma}_n.
\end{aligned}
\label{eq:matrix-basis-actions}
\end{equation}
On these three individual generator matrices, the same sign table is obtained
from \(\mathsf C\) and \(\mathsf T\). Distinct operations are nevertheless represented:
\(\mconj{(\mat{A}\mat{B})}=\mconj{\mat{A}}\mconj{\mat{B}}\), whereas
\(\mtrans{(\mat{A}\mat{B})}=\mtrans{\mat{B}}\mtrans{\mat{A}}\). Only the represented second
generator is changed by entrywise \(\mathsf C\), while all three are changed by
intrinsic star. Every \(\mat{\sigma}_n\) is fixed by their composition
\(\mathsf{CT}\), exactly matching intrinsic \(H\). Thus \(\mathsf C\) depends
on the chosen matrix entries
and is distinct from both sigma-algebra operations. In particular,
it is not intrinsic star. For \(n\neq m\),
\[
\herm{\mat{\sigma}_n\mat{\sigma}_m}
=\mat{\sigma}_m\mat{\sigma}_n
=-\mat{\sigma}_n\mat{\sigma}_m,
\]
exactly as required by the intrinsic order-reversal rule.

\subsubsection{Auxiliary Block Map}
A scalar-coefficient paravector \(\pv{Z}\in\Pspace\) is packaged
into the \(2\times2\) block map
\begin{equation}
\label{eq:W-map}
\Wmat(\pv{Z}):=
\begin{bmatrix}
0&\pv{Z}\\
\adjf{\pv{Z}}&0
\end{bmatrix}.
\end{equation}
\begin{equation}
\Wmat(\pv{Z})^2
=\eqtext{diag}(\detf{\pv{Z}},\detf{\pv{Z}})
=\detf{\pv{Z}}\I.
\label{eq:W-square}
\end{equation}
The block map \(\Wmat(\pv{Z})\) is a \(2\times2\) matrix over the paravector
algebra. After each block is replaced by its Pauli matrix, it is a
\(4\times4\) complex matrix. Addition and multiplication by ordinary
scalars are preserved, so the packaging map is linear. General products are
not preserved, so it is not an additional representation of the algebra.
Equation \eqref{eq:W-square} is obtained because its
diagonal block products are reduced to the determinant in
Eq.\ \eqref{eq:minkowski-form}; the off-diagonal blocks are zero
\cite{hestenes1966,doran2003}.

If \(\detf{\pv{Z}}\neq0\), Eq.\ \eqref{eq:W-square} shows that the powers of
the block matrix are governed by the scalar equation
\(x^2-\detf{\pv{Z}}=0\). Its two roots are distinct; selecting one square
root fixes their labels. An invertible change of basis
can then put the block matrix into diagonal form; this property is called
diagonalizability. If
\(\detf{\pv{Z}}=0\) and \(\pv{Z}\neq0\), then
\(\Wmat(\pv{Z})^2=0\) while \(\Wmat(\pv{Z})\neq0\). A nonzero matrix whose
positive power is zero is called nilpotent; such a matrix cannot be put
into diagonal form. For blocks containing differential operations or other
quantities whose order cannot be exchanged, the unreduced statement is
\(\Wmat(\pv{Z})^2
=\eqtext{diag}(\pv{Z}\cherm{\pv{Z}},\cherm{\pv{Z}}\pv{Z})\).
A separate operator calculation is required before either ordered product can
be called a determinant or both can be reduced to one scalar multiple of the
identity.


\section{Space-Time Kinematics}\label{sec:spacetime}
On real event paravectors, the Minkowski quadratic form is obtained from the determinant form and four-velocities are represented by normalized timelike paravectors. Standard special-relativistic kinematics is thereby expressed in paravector form \cite{lorentz1904,einstein1905,minkowski1908,hestenes1966,baylis1999}.

\subsection{Differential space-time and four-velocity}
The differential event is obtained by specializing the real paravector
\(\pv{X}\) from Eq.\ \eqref{eq:real-paravector-x}:
\begin{equation}
\dd\pv{X}:=\dd t+\dd\vct{r}\cdot\sigv
\label{eq:dX}
\end{equation}
The interval is the differential version of Eq.\ \eqref{eq:minkowski-form}:
\begin{equation}
\dd s^2=\detf{\dd\pv{X}}
=\dd\pv{X}\,\conj{\dd\pv{X}}
=\dd t^2-\dd\vct{r}^2.
\label{eq:interval}
\end{equation}
Here \(t,s,\gamma\in\R\), \(\vct{r},\vct{v},\vct{P}\in\R^3\), and \(\pv{X},\pv{U}\) are real paravectors. Equation \eqref{eq:interval} is the real-event restriction of the determinant in Eq.\ \eqref{eq:minkowski-form}. For a future-directed timelike worldline, \(\dd t>0\) and \(\dd s=\sqrt{\dd t^2-\dd\vct{r}^2}>0\).

The coordinate velocity and Lorentz factor are defined by
\(\vct{v}:=\dd\vct{r}/\dd t\in\R^3\) and
\(\gamma:=(1-\vct{v}^2)^{-1/2}\in\R\). The four-velocity is defined by
\begin{equation}
\pv{U}:=\frac{\dd\pv{X}}{\dd s}
=\gamma(1+\vct{v}\cdot\sigv).
\label{eq:proper-velocity}
\end{equation}
Its normalization is
\begin{equation}
\detf{\pv{U}}=1.
\label{eq:proper-velocity-normalization}
\end{equation}
The energy and spatial momentum of a particle with real rest mass \(m>0\)
are defined by
\begin{equation}
E:=m\gamma.
\label{eq:particle-energy}
\end{equation}
\begin{equation}
\vct{P}:=m\gamma\vct{v}.
\label{eq:particle-momentum}
\end{equation}
The four-momentum is therefore
\begin{equation}
m\pv{U}=E+\vct{P}\cdot\sigv.
\label{eq:proper-momentum}
\end{equation}
from which $E^2-\vct{P}^2=m^2$ is obtained immediately.
The normalized tangent \(\pv{U}\) is conventionally called the four-velocity; only its spatial part \(\gamma\vct{v}\) is often denoted by ``proper velocity.'' The four-momentum of a massive particle is obtained by multiplication by \(m>0\). For a massless particle \(\dd s=0\), so \(\pv{U}\) is undefined and a null momentum paravector must instead be introduced directly.

\subsection{Four-acceleration and free motion}
The four-acceleration is the real paravector
\begin{equation}
\pv{A}:=\frac{\dd \pv{U}}{\dd s}
=\frac{\dd^2 \pv{X}}{\dd s^2}.
\label{eq:proper-acceleration}
\end{equation}
It obeys
\begin{equation}
\eqtext{scalar}(\conj{\pv{U}}\pv{A})=0.
\label{eq:proper-acceleration-orthogonality}
\end{equation}
The orthogonality condition in Eq.\ \eqref{eq:proper-acceleration} is the standard four-acceleration orthogonality relation. The invariant proper acceleration is the magnitude \(\sqrt{-\detf{\pv{A}}}\), not the paravector \(\pv{A}\) itself.

A free-particle Lagrangian is defined by
\begin{equation}
\mathcal L:=-\frac{m}{\gamma},
\label{eq:free-lagrangian}
\end{equation}
from whose Euler--Lagrange equations inertial motion and the usual nonrelativistic expansion are reproduced:
\begin{equation}
\mathcal L=-m+\frac{1}{2}m\vct{v}^2+O(\vct{v}^4).
\label{eq:free-lagrangian-low}
\end{equation}
Equation \eqref{eq:free-lagrangian-low} is displayed explicitly to exhibit the reduction of the relativistic action to the classical kinetic expression plus an additive rest-energy constant in the low-speed limit.

\subsection{Supporting Kinematic Derivations}
Equations \eqref{eq:interval}, \eqref{eq:proper-velocity}, and \eqref{eq:proper-momentum} are expanded in the supporting derivations, followed by the generalized-coordinate and metric remarks associated with the geometric part of the development.
\subsubsection{Space-Time Numbers}
A spacetime increment is the real paravector
\[
\dd\pv{X}=\dd t+\dd r_{1}\pv{\sigma}_{1}+\dd r_{2}\pv{\sigma}_{2}+\dd r_{3}\pv{\sigma}_{3}
=\dd t+\dd\vct{r}\cdot\sigv,
\]
with conjugate
\[
\conj{\dd\pv{X}}=\dd t-\dd r_{1}\pv{\sigma}_{1}-\dd r_{2}\pv{\sigma}_{2}-\dd r_{3}\pv{\sigma}_{3}
=\dd t-\dd\vct{r}\cdot\sigv.
\]
Here \(\dd t\in\R\) is the scalar part and \(\dd\vct{r}\in\R^{3}\) is the spatial vector part. Their recovery formulas are
\begin{equation*}
\frac{\dd\pv{X}+\conj{\dd\pv{X}}}{2}=\dd t,
\qquad
\frac{\dd\pv{X}-\conj{\dd\pv{X}}}{2}=\dd\vct{r}\cdot\sigv,
\end{equation*}
so
\begin{equation*}
\eqtext{scalar}(\dd\pv{X})=\dd t,
\qquad
\eqtext{vector}(\dd\pv{X})=\dd\vct{r}.
\end{equation*}
The determinant of the increment is
\begin{equation*}
\detf{\dd\pv{X}}
=\dd\pv{X}\conj{\dd\pv{X}}
=\conj{\dd\pv{X}}\dd\pv{X}
=\dd t^{2}-\dd\vct{r}^{2}
=\dd s^{2}.
\end{equation*}
The familiar Lorentzian identities are obeyed:
\begin{align*}
\dd t^2>\dd\vct{r}^2,\ \dd\pv{X}\neq 0
&\eqtext{ implies }\detf{\dd\pv{X}}>0,\notag\\
\detf{\dd\pv{X}+\dd\pv{X}'}
&-\detf{\dd\pv{X}}\notag\\
&\qquad -\detf{\dd\pv{X}'}
=2(\dd t\,\dd t'-\dd\vct{r}\cdot\dd\vct{r}'),\notag\\
\detf{a\,\dd\pv{X}}
&=a^2\detf{\dd\pv{X}},
\qquad a\in\R.
\end{align*}
The commutator of two spacetime increments is purely spatial:
\begin{align*}
\com{\dd\pv{X}}{\dd\pv{X}'}
&=\com{\dd\vct{r}\cdot\sigv}{\dd\vct{r}'\cdot\sigv}\notag\\
&=2i(\dd\vct{r}\times\dd\vct{r}')\cdot\sigv.
\end{align*}
Along a future-directed time-like path, so that \(\dd t>0\), the following relation is satisfied by the positive arc length:
\begin{align*}
\dd s
&=\sqrt{\dd t^{2}-\dd\vct{r}^{2}}\notag\\
&=\sqrt{1-\dd\vct{r}^{2}/\dd t^{2}}\,\dd t\notag\\
&=\sqrt{1-\vct{v}^{2}}\,\dd t,
\end{align*}
where
\begin{equation*}
\vct{v}:=\frac{\dd\vct{r}}{\dd t}=(v_{1},v_{2},v_{3})\in\R^{3},
\qquad
\gamma:=\frac{\dd t}{\dd s}=\frac{1}{\sqrt{1-\vct{v}^{2}}}.
\end{equation*}
For a real non-null increment, \(\herm{\dd\pv{X}}=\dd\pv{X}\), so the inverse exists and
\begin{align*}
(\dd\pv{X})^{-1}
&=\frac{\adjf{\dd\pv{X}}}{\detf{\dd\pv{X}}}\notag\\
&=\frac{\conj{\dd\pv{X}}}{\detf{\dd\pv{X}}}\notag\\
&=\frac{\conj{\dd\pv{X}}}{\dd\pv{X}\conj{\dd\pv{X}}}\notag\\
&=\frac{\conj{\dd\pv{X}}}{\dd s^{2}}.
\end{align*}
Null displacements are singled out algebraically because this inverse breaks down exactly when \(\dd s^{2}=0\).

\subsubsection{Interval Classes and Proper Time}
For a spacetime increment \(\dd\pv{X}=\dd t+\dd \vct{r}\cdot\sigv\), the interval is
\begin{equation*}
\detf{\dd\pv{X}}=\dd\pv{X}\,\conj{\dd\pv{X}}=\dd t^2-\dd \vct{r}^2=\dd s^2.
\end{equation*}
The three standard interval classes are obtained from this identity without any additional geometric machinery:
\begin{align*}
\dd s^2>0 &\quad\eqtext{implies a time-like interval},\notag \\
\dd s^2=0,\ \dd\pv{X}\neq0
&\quad\eqtext{implies a null interval},\notag \\
\dd s^2<0 &\quad\eqtext{implies a space-like interval}.
\end{align*}
If \(\detf{\dd\pv{X}}>0\), the interval is time-like. Then the arc length \(s\) equals the proper time \(\tau\) measured along the path:
\begin{align*}
\tau=s=\int \dd s=\int \dd\tau
&=\int \sqrt{1-\vct{v}^2}\,\dd t\notag\\
&=\int \frac{1}{\gamma}\,\dd t.
\end{align*}
This is the time-dilation relation in natural units. In SI units the same statement is
\begin{equation*}
\dd s=c\,\dd\tau.
\end{equation*}

If \(\detf{\dd\pv{X}}=0\) with \(\dd\pv{X}\neq0\), the interval is null or light-like. Then
\begin{equation*}
\dd t^2=\dd\vct{r}^2
\qquad\eqtext{implies}\qquad
\dd t=\pm \|\dd\vct{r}\|,
\end{equation*}
and \((\dd\pv{X})^{-1}\) does not exist.

If \(\detf{\dd\pv{X}}<0\), the interval is space-like.

In the paravector notation, the interval classification, the proper-time formula, and the null singularity of \((\dd\pv{X})^{-1}\) are all obtained from the single algebraic identity \(\dd\pv{X}\,\conj{\dd\pv{X}}=\dd t^2-\dd\vct{r}^2\).

\subsubsection{Four-Velocity}
The four-velocity paravector is the normalized tangent to a time-like worldline:
\begin{align*}
\frac{\dd\pv{X}}{\dd s}
&=\frac{\dd t}{\dd s}+\frac{\dd\vct{r}}{\dd s}\cdot\sigv\notag\\
&=\frac{\dd t+\sigv\cdot\dd\vct{r}}{\sqrt{\dd t^2-\dd\vct{r}^2}}\notag\\
&=\frac{\dd t/\dd t+\sigv\cdot(\dd\vct{r}/\dd t)}{\sqrt{(\dd t/\dd t)^2-(\dd\vct{r}/\dd t)^2}}\notag\\
&=\frac{1+\vct{v}\cdot\sigv}{\sqrt{1-\vct{v}^2}}\notag\\
&=\gamma(1+\vct{v}\cdot\sigv).
\end{align*}
Its scalar and vector parts are therefore
\begin{equation*}
\eqtext{scalar}\!\left(\frac{\dd\pv{X}}{\dd s}\right)=\gamma,
\qquad
\eqtext{vector}\!\left(\frac{\dd\pv{X}}{\dd s}\right)=\gamma\vct{v}.
\end{equation*}
Its determinant is unity:
\begin{align*}
\detf{\frac{\dd\pv{X}}{\dd s}}
&=\gamma^2(1+\vct{v}\cdot\sigv)(1-\vct{v}\cdot\sigv)\notag\\
&=\gamma^2(1-\vct{v}^2)\notag\\
&=1.
\end{align*}
Also,
\begin{equation*}
\frac{\dd\pv{X}}{\dd t}
=\frac{\dd\pv{X}}{\dd s}\frac{\dd s}{\dd t}
=\frac{1}{\gamma}\frac{\dd\pv{X}}{\dd s}.
\end{equation*}
The four-velocity paravector is defined by
\begin{align*}
\pv{U}&=\frac{\dd\pv{X}}{\dd s}\notag\\
&=\gamma+\gamma\vct{v}\cdot\sigv\notag\\
&=U_s+\vct{U}_{\vct{v}}\cdot\sigv,\\
U_s&=\gamma,\qquad
\vct{U}_{\vct{v}}=\gamma\vct{v}.
\end{align*}
Since \(\dd s^2=\dd\pv{X}\conj{\dd\pv{X}}\) with \(\dd s>0\),
\begin{align*}
\frac{\conj{\dd\pv{X}}}{\dd s}
&=(\dd\pv{X})^{-1}\dd s\notag\\
&=\left(\frac{\dd\pv{X}}{\dd s}\right)^{-1},
\end{align*}
so \(\pv{U}^{-1}=\pv{U}^*\) and
\begin{equation*}
\detf{\pv{U}}=\pv{U}\pv{U}^*=1.
\end{equation*}
The following relation is obtained by differentiation of this normalization:
\begin{align*}
\frac{\dd}{\dd s}\detf{\pv{U}}
&=\frac{\dd\pv{U}}{\dd s}\pv{U}^*+\pv{U}\left(\frac{\dd\pv{U}}{\dd s}\right)^*\notag\\
&=\left(\frac{\dd\pv{U}}{\dd s}\pv{U}^*\right)+\left(\frac{\dd\pv{U}}{\dd s}\pv{U}^*\right)^{*\mathsf H},
\end{align*}
hence
\begin{equation*}
\eqtext{scalar}\!\left(\pv{U}^*\frac{\dd\pv{U}}{\dd s}\right)=0.
\end{equation*}
The determinant identity is recovered as
\begin{align*}
\detf{\dd\pv{X}}
&=\dd\pv{X}\conj{\dd\pv{X}}\notag\\
&=\left(\frac{\dd\pv{X}}{\dd s}\dd s\right)\left(\left(\frac{\dd\pv{X}}{\dd s}\right)^*\dd s\right)\notag\\
&=\pv{U}\pv{U}^*\,\dd s^2\notag\\
&=\detf{\pv{U}}\dd s^2\notag\\
&=\dd s^2.
\end{align*}
The proper-acceleration formulas are based on this normalization and orthogonality structure.
\noindent The following table resolves the proper-velocity paravector into its scalar and spatial components.
\begin{componenttable}[tab:components-proper-velocity]{dX/ds}
\componentrow{scalar}{\gamma}{}
\componentrow{vector}{\gamma\mathbf v}{}
\end{componenttable}


\subsubsection{Four-Acceleration}
The coordinate acceleration is defined by
\begin{equation*}
\vct{a}:=\frac{\dd\vct{v}}{\dd t}=\{a_1,a_2,a_3\}\in\R^3.
\end{equation*}
With \(\gamma=(1-\vct{v}^2)^{-1/2}\), the following relations hold:
\begin{align*}
\frac{\dd\gamma}{\dd t}&=\gamma^3(\vct{v}\cdot\vct{a}),\\
\frac{\dd\gamma}{\dd s}&=\gamma^4(\vct{v}\cdot\vct{a}).
\end{align*}
The notation \(\delv:=\partial/\partial\vct{v}\) is used for the gradient
with respect to velocity. The scalar velocity derivative needed below is
\[
\delv{}\!\left(-\frac{1}{\gamma}\right)=\gamma\vct{v}.
\]
The following derivatives are obtained by direct differentiation of the vector product \(\gamma\vct{v}\):
\begin{align*}
\frac{\dd(\gamma\vct{v})}{\dd t}
&=\gamma\vct{a}+\frac{\dd\gamma}{\dd t}\vct{v}\notag\\
&=\gamma\vct{a}+\gamma^3(\vct{v}\cdot\vct{a})\vct{v},\\
\frac{\dd(\gamma\vct{v})}{\dd s}
&=\frac{\dd(\gamma\vct{v})}{\dd t}\frac{\dd t}{\dd s}\notag\\
&=\gamma^2\vct{a}+\gamma^4(\vct{v}\cdot\vct{a})\vct{v}.
\end{align*}

The four-velocity paravector is decomposed as
\begin{align*}
\pv{U}&=U_s+\vct{U}_{\vct{v}}\cdot\sigv,\\
U_s&=\gamma,\qquad
\vct{U}_{\vct{v}}=\gamma\vct{v}.
\end{align*}
Therefore the four-acceleration is
\begin{align*}
\pv{A}
&=\frac{\dd\pv{U}}{\dd s}\notag\\
&=\frac{\dd\gamma}{\dd s}+\frac{\dd(\gamma\vct{v})}{\dd s}\cdot\sigv\notag\\
&=\gamma^4(\vct{v}\cdot\vct{a})
+\Bigl(\gamma^4(\vct{v}\cdot\vct{a})\vct{v}+\gamma^2\vct{a}\Bigr)\cdot\sigv.
\end{align*}
Equivalently,
\begin{equation*}
\pv{A}=A_s+\vct{A}_{\vct{v}}\cdot\sigv,
\qquad
A_s=\gamma^4(\vct{v}\cdot\vct{a}),
\qquad
\vct{A}_{\vct{v}}=\gamma^4(\vct{v}\cdot\vct{a})\vct{v}+\gamma^2\vct{a}.
\end{equation*}

For \(\vct{v}\neq 0\), the unit direction is defined by
\(\vct{u}=\vct{v}/\|\vct{v}\|\), with \(\vct{u}^2=1\). The following
components are then obtained from the axis functions in
Eqs.\ \eqref{eq:par-function} and \eqref{eq:prp-function}:
\begin{align*}
\eqtext{par}(\vct{a})&=(\vct{a}\cdot\vct{u})\vct{u},\\
\eqtext{prp}(\vct{a})&=\vct{a}-\eqtext{par}(\vct{a}).
\end{align*}
Then
\begin{align*}
\vct{a}^2&=\eqtext{par}(\vct{a})^2+\eqtext{prp}(\vct{a})^2,\\
(\vct{v}\cdot\vct{a})^2&=\vct{v}^2\,\eqtext{par}(\vct{a})^2,
\end{align*}
so the proper-acceleration determinant is
\begin{align*}
\detf{\pv{A}}
&=-\gamma^4\Bigl(\vct{a}^2+\gamma^2(\vct{v}\cdot\vct{a})^2\Bigr)\notag\\
&=-\gamma^4\Bigl(\eqtext{prp}(\vct{a})^2+\gamma^2\eqtext{par}(\vct{a})^2\Bigr).
\end{align*}
This is strictly space-like unless the acceleration vanishes. Its invariant magnitude \(\sqrt{-\detf{\pv{A}}}\) is the proper acceleration measured in the instantaneous rest frame.

The products with \(\pv{U}\) are
\begin{align*}
\pv{U}\pv{A}
&=(U_sA_s+\vct{U}_{\vct{v}}\cdot\vct{A}_{\vct{v}})\notag\\
&\quad+\bigl(U_s\vct{A}_{\vct{v}}+A_s\vct{U}_{\vct{v}}+i(\vct{U}_{\vct{v}}\times\vct{A}_{\vct{v}})\bigr)\cdot\sigv,\\
\pv{U}^*\pv{A}
&=(U_sA_s-\vct{U}_{\vct{v}}\cdot\vct{A}_{\vct{v}})\notag\\
&\quad+\bigl(U_s\vct{A}_{\vct{v}}-A_s\vct{U}_{\vct{v}}-i(\vct{U}_{\vct{v}}\times\vct{A}_{\vct{v}})\bigr)\cdot\sigv.
\end{align*}
Their component forms are
\begin{align*}
\pv{U}\pv{A}
&=2\gamma^5(\vct{v}\cdot\vct{a})\notag\\
&\quad+\bigl(\gamma^3\vct{a}+2\gamma^5(\vct{v}\cdot\vct{a})\vct{v}\bigr)\cdot\sigv\notag\\
&\quad+i\,\gamma^3(\vct{v}\times\vct{a})\cdot\sigv,\\
\pv{U}^*\pv{A}
&=\gamma^3\vct{a}\cdot\sigv-i\,\gamma^3(\vct{v}\times\vct{a})\cdot\sigv.
\end{align*}
In particular,
\begin{equation*}
\eqtext{scalar}(\pv{U}^*\pv{A})=0,
\end{equation*}
which is the orthogonality relation between four-velocity and four-acceleration. The scalar and vector parts of \(\pv{U}\) are recorded in Table \ref{tab:components-proper-velocity}, while the corresponding scalar and vector parts of \(\pv{U}\pv{A}\) and \(\pv{U}^{*}\pv{A}\) are compared in Tables \ref{tab:components-ua} and \ref{tab:components-ustara}.
\noindent The products \(\pv{U}\pv{A}\) and \(\pv{U}^{*}\pv{A}\) are compared in the following two tables, and the orthogonality structure of proper acceleration is thereby made explicit.
\begin{componenttable}[tab:components-ua]{\pv{U}\pv{A}}
\componentrow{real scalar}{2\gamma^{5}(\vct{v}\cdot\vct{a})}{}
\componentrow{img scalar}{0}{}
\componentrow{real vector}{\gamma^{3}\vct{a}+2\gamma^{5}(\vct{v}\cdot\vct{a})\vct{v}}{}
\componentrow{img vector}{\gamma^{3}(\vct{v}\times\vct{a})}{}
\end{componenttable}

\begin{componenttable}[tab:components-ustara]{\pv{U}^{*}\pv{A}}
\componentrow{real scalar}{0}{}
\componentrow{img scalar}{0}{}
\componentrow{real vector}{\gamma^{3}\vct{a}}{}
\componentrow{img vector}{-\gamma^{3}(\vct{v}\times\vct{a})}{}
\end{componenttable}


\subsubsection{Four-Momentum}
For a massive particle with rest mass \(m>0\), the energy--momentum paravector is produced by multiplication of four-velocity by \(m\):
\begin{equation*}
m\pv{U}=E+\vct{P}\cdot\sigv.
\end{equation*}
Its scalar part is the relativistic energy,
\begin{align*}
E
&:=\eqtext{scalar}(m\pv{U})\notag\\
&=m\gamma.
\end{align*}
If \(\vct{v}=0\), then
\begin{equation*}
E=m
\qquad \text{rest energy}
\qquad (=mc^{2}\text{ in SI units}).
\end{equation*}
The kinetic energy is therefore \(E-m\).

Its vector part is the relativistic momentum,
\begin{align*}
\vct{P}
&:=\eqtext{vector}(m\pv{U})\notag\\
&=m\gamma\vct{v}.
\end{align*}

Because \(\detf{\pv{U}}=1\), the mass shell is obtained immediately:
\begin{align*}
\detf{m\pv{U}}
&=m^{2}\notag\\
&=(E+\vct{P}\cdot\sigv)(E-\vct{P}\cdot\sigv)\notag\\
&=E^{2}-\vct{P}^{2}.
\end{align*}
Hence
\begin{equation*}
E^{2}=\vct{P}^{2}+m^{2}.
\end{equation*}
The null branch \(E^{2}=\vct{P}^{2}\), appropriate to massless motion, is also admitted by the mass-shell equation. That branch must be introduced directly as a null energy--momentum paravector: Eq.\ \(m\pv{U}=E+\vct{P}\cdot\sigv\) is applicable only to massive time-like motion because \(\dd s=0\) and \(\pv{U}=\dd\pv{X}/\dd s\) is undefined on a null worldline.

\subsubsection{Flat Paravector Coordinate Map}
The scalar and vector maps are taken to have the domains
\begin{align*}
f&:\R\times\R^3\to\R,\\
\vct{g}&:\R\times\R^3\to\R^3.
\end{align*}
A real paravector-valued map into the fixed flat space-time algebra is then
defined by
\begin{equation}
\pv{X}:=f(t,\vct{r})+\vct{g}(t,\vct{r})\cdot\sigv.
\label{eq:flat-paravector-coordinate-map}
\end{equation}
The original coordinates are mapped into flat Minkowski paravector components; invariance of the Cartesian metric coefficients under an arbitrary coordinate change is not asserted. Its differential and conjugate are
\begin{equation*}
\dd\pv{X}=\dd f+\dd\vct{g}\cdot\sigv,
\qquad
\conj{\dd\pv{X}}=\dd f-\dd\vct{g}\cdot\sigv.
\end{equation*}
Their product is
\begin{align*}
\detf{\dd\pv{X}}
&=\dd\pv{X}\,\conj{\dd\pv{X}}\notag\\
&=(\dd f+\dd\vct{g}\cdot\sigv)(\dd f-\dd\vct{g}\cdot\sigv)\notag\\
&=\dd f^2-\dd\vct{g}^2.
\end{align*}
Along a future-directed time-like path, the conditions \(\dd f>0\) and
\(\dd f^2>\dd\vct{g}^2\) are assumed. Then
\begin{align*}
\dd s
&=\sqrt{\dd f^2-\dd\vct{g}^2}\notag\\
&=\sqrt{1-\left(\frac{\dd\vct{g}}{\dd f}\right)^2}\,\dd f.
\end{align*}
The corresponding four-velocity is
\begin{align*}
\pv{U}
&=\frac{\dd\pv{X}}{\dd s}\notag\\
&=\frac{\dd f+\dd\vct{g}\cdot\sigv}{\sqrt{\dd f^2-\dd\vct{g}^2}}\notag\\
&=\frac{1+\left(\frac{\dd\vct{g}}{\dd f}\right)\cdot\sigv}{\sqrt{1-\left(\frac{\dd\vct{g}}{\dd f}\right)^2}}\notag\\
&=\frac{1+\vct{v}\cdot\sigv}{\sqrt{1-\vct{v}^2}}\notag\\
&=\gamma(1+\vct{v}\cdot\sigv),
\end{align*}
where the velocity relative to the scalar component \(f\) and its Lorentz factor are
\begin{align*}
\vct{v}
&:=\frac{\dd\vct{g}}{\dd f}\notag\\
&=\frac{\dd\vct{g}/\dd t}{\dd f/\dd t}\notag\\
&=\frac{\dd\vct{g}/\dd s}{\dd f/\dd s},\\
\gamma
&:=\frac{\dd f}{\dd s}\notag\\
&=\frac{\dd f}{\sqrt{\dd f^2-\dd\vct{g}^2}}\notag\\
&=\frac{\dd f/\dd t}{\dd s/\dd t}\notag\\
&=\frac{1}{\sqrt{1-\vct{v}^2}}.
\end{align*}
Thus \(\vct{v}^2<1\). For a massive particle, the energy--momentum relation remains unchanged:
\begin{equation*}
m\pv{U}=E+\vct{P}\cdot\sigv.
\end{equation*}
The coordinate acceleration relative to \(f\) is defined, in direct analogy
with \(\vct{a}=\dd\vct{v}/\dd t\) above, by
\begin{align*}
\vct{a}
&:=\frac{\dd\vct{v}}{\dd f}\notag\\
&=\frac{\dd^2\vct{g}}{\dd f^2}.
\end{align*}
Because \(\dd/\dd s=\gamma\,\dd/\dd f\), the four-acceleration has the same form as before:
\begin{equation*}
\frac{\dd\pv{U}}{\dd s}
=\gamma^4(\vct{v}\cdot\vct{a})
+\bigl(\gamma^4(\vct{v}\cdot\vct{a})\vct{v}+\gamma^2\vct{a}\bigr)\cdot\sigv.
\end{equation*}
The commutator of four-velocity and four-acceleration is
\begin{align*}
\com{\pv{U}}{\dfrac{\dd\pv{U}}{\dd s}}
&=\com{\pv{U}}{\dfrac{\dd\pv{U}}{\dd f}}\frac{\dd f}{\dd s}\notag\\
&=\gamma\com{\pv{U}}{\dfrac{\dd\pv{U}}{\dd f}}\notag\\
&=2\gamma^3(\vct{v}\times\vct{a})\cdot(i\sigv).
\end{align*}
For a massive free particle in these flat target components,
\begin{equation*}
m\frac{\dd\pv{U}}{\dd s}=0,
\end{equation*}
which is equivalently the straight-worldline equation
\begin{equation*}
\frac{\dd^2\pv{X}}{\dd s^2}=0.
\end{equation*}
After division by the nonzero factor \(\gamma\), its scalar and vector parts are
\begin{align*}
\gamma^3(\vct{v}\cdot\vct{a})&=0,\\
\gamma^3(\vct{v}\cdot\vct{a})\vct{v}+\gamma\vct{a}&=0.
\end{align*}
For \(\vct{v}^2<1\), one has \(\gamma>0\). The relation
\(\vct{v}\cdot\vct{a}=0\) is therefore obtained from the scalar part, and
\(\vct{a}=0\) is then obtained from the vector part. Hence
\begin{equation*}
\frac{\dd^2\vct{g}}{\dd f^2}=0,
\end{equation*}
and along an inertial worldline
\begin{equation*}
\vct{g}=\vct{c}f+\vct{c}',
\qquad
\vct{v}=\frac{\dd\vct{g}}{\dd f}=\vct{c},
\end{equation*}
so inertial motion is linear in the flat target components \((f,\vct{g})\) with constant velocity. If the same straight worldline is described in arbitrary curvilinear coordinates, the corresponding connection terms are generally contained in its coordinate equation; no curved-space geodesic equation is asserted.

\subsubsection{Lagrangian Mechanics and Conserved Quantities}
For a particle trajectory written in terms of time, position, and velocity,
the action is taken to be
\begin{equation*}
S=\int \mathcal L(t,\vct{r},\vct{v})\,\dd t.
\end{equation*}
The Euler--Lagrange equation is then
\begin{align*}
\left(\gradv-\frac{\dd}{\dd t}\delv{}\right)\mathcal L&=0,\\
\frac{\dd}{\dd t}\delv{}\mathcal L&=\gradv\mathcal L.
\end{align*}
The standard conserved quantities are then obtained immediately. The canonical momentum is
\begin{equation*}
\delv{}\mathcal L,
\end{equation*}
so if \(\gradv\mathcal L=0\), then
\begin{equation*}
\frac{\dd}{\dd t}\delv{}\mathcal L=0,
\end{equation*}
and the canonical momentum is constant.

The Hamiltonian is defined by
\begin{align*}
\mathcal H
&=\vct{v}\cdot(\delv{}\mathcal L)-\mathcal L\notag\\
&=(\vct{v}\cdot\delv{}-1)\mathcal L.
\end{align*}
Its time derivative is
\begin{align*}
\frac{\dd\mathcal H}{\dd t}
&=\frac{\dd}{\dd t}(\delv{}\mathcal L)\cdot\vct{v}
+(\delv{}\mathcal L)\cdot\frac{\dd\vct{v}}{\dd t}
-\frac{\dd\mathcal L}{\dd t}\notag\\
&=(\gradv\mathcal L)\cdot\vct{v}
+(\delv{}\mathcal L)\cdot\frac{\dd\vct{v}}{\dd t}
-\frac{\dd\mathcal L}{\dd t}\notag\\
&=(\gradv\mathcal L)\cdot\vct{v}
+(\delv{}\mathcal L)\cdot\frac{\dd\vct{v}}{\dd t}\notag\\
&\qquad-\left(
\partial_t\mathcal L
+(\gradv\mathcal L)\cdot\vct{v}
+(\delv{}\mathcal L)\cdot\frac{\dd\vct{v}}{\dd t}
\right)\notag\\
&=-\partial_t\mathcal L.
\end{align*}
Therefore, if \(\partial_t\mathcal L=0\), then
\begin{equation*}
\frac{\dd\mathcal H}{\dd t}=0,
\end{equation*}
so the Hamiltonian is constant. These identities are applied directly to the free-particle and minimally coupled electromagnetic examples.

\subsubsection{Action of a Free Particle}
The free-particle action is the mass-weighted proper time along the world line. With \(\tau=\int \dd s\), \(\pv{U}=\dd\pv{X}/\dd s\), and \(\gamma=\dd t/\dd s\),
\begin{align*}
S
=\int \mathcal L\,\dd t
&=-m\tau\notag\\
&=-m\int \sqrt{\detf{\dd\pv{X}}}\notag\\
&=-m\int \dd s\notag\\
&=-m\int \sqrt{\detf{\frac{\dd\pv{X}}{\dd t}}}\,\dd t\notag\\
&=-m\int \frac{\dd s}{\dd t}\,\dd t\notag\\
&=-m\int \sqrt{\detf{\frac{\pv{U}}{\gamma}}}\,\dd t\notag\\
&=-m\int \frac{1}{\gamma}\,\dd t\notag\\
&=-m\int \sqrt{1-\vct{v}^2}\,\dd t.
\end{align*}
Therefore
\begin{equation*}
\mathcal L=-\frac{m}{\gamma}.
\end{equation*}
The canonical momentum is obtained by differentiation with respect to velocity:
\begin{align*}
\delv{} \mathcal L
&=m\gamma\vct{v}\notag\\
&=\vct{P},\\
(\delv{}\cdot\sigv)\mathcal L
&=\vct{P}\cdot\sigv.
\end{align*}
The Hamiltonian is
\begin{align*}
\mathcal H
&=(\delv{}\mathcal L)\cdot\vct{v}-\mathcal L\notag\\
&=m\gamma\vct{v}^2+\frac{m}{\gamma}\notag\\
&=m\gamma\left(\vct{v}^2+\frac{1}{\gamma^2}\right)\notag\\
&=m\gamma\notag\\
&=E.
\end{align*}
Because the free Lagrangian has no explicit time or position dependence,
\begin{equation*}
\partial_t\mathcal L=0,
\qquad
\gradv\mathcal L=0,
\end{equation*}
so the conserved quantities are the relativistic energy \(E=m\gamma\) and momentum \(\vct{P}=m\gamma\vct{v}\).

\subsubsection{Non-Relativistic Approximation}
For the non-relativistic approximation, the free Lagrangian is expanded
directly:
\begin{align*}
\mathcal L
&=-\frac{m}{\gamma}\notag\\
&=-m+\frac12 m\vct{v}^2+O(\vct{v}^4).
\end{align*}
The canonical momentum is then reduced to
\begin{align*}
\delv{}\mathcal L
&=\frac{\partial}{\partial \vct{v}}\left(\frac12 m\vct{v}^2+O(\vct{v}^4)\right)\notag\\
&=m\vct{v}+O(\vct{v}^3)\notag\\
&\approx m\vct{v}.
\end{align*}
The total energy is the kinetic term plus the rest-mass term:
\begin{align*}
\mathcal H
&=\bigl(m\vct{v}+O(\vct{v}^3)\bigr)\cdot\vct{v}
+m-\frac12m\vct{v}^2-O(\vct{v}^4)\notag\\
&=m+\frac12m\vct{v}^2+O(\vct{v}^4)\notag\\
&\approx m+\frac12m\vct{v}^2.
\end{align*}
Equivalently, the low-velocity series is
\begin{align*}
\mathcal L&=-m+\frac12 m\vct{v}^2+O(\vct{v}^4),\notag \\
\vct{P}&=m\vct{v}+O(\vct{v}^{3}),\notag \\
\mathcal H&=m+\frac12 m\vct{v}^{2}+O(\vct{v}^{4}).
\end{align*}
The Newtonian kinetic term together with the additive rest-energy constant is thereby shown to be contained in the relativistic action. Classical mechanics is therefore obtained as the approximation, not used as the starting point.

\subsubsection{Flat Pullback Metric and Diagonalization}
A flat metric in the original coordinate tuple is induced by the map \(\pv{X}=f+\vct{g}\cdot\sigv\):
\[
\vct{x}=(t,r_1,r_2,r_3)^T,
\qquad
\dd\vct{x}=(\dd t,\dd r_1,\dd r_2,\dd r_3)^T.
\]
The transpose operation was defined in Eq.\ \eqref{eq:matrix-transpose}.
The symmetric pullback matrix \(\Qmat(\vct{x})\) is the unique matrix defined
by expansion of \(\dd f\) and \(\dd\vct{g}\) in \(\dd\vct{x}\):
\begin{equation}
\dd f^2-\dd\vct{g}^2
=:\dd\vct{x}^T\Qmat(\vct{x})\,\dd\vct{x}.
\label{eq:pullback-matrix-function}
\end{equation}
Consequently,
\begin{equation*}
\dd s^2=\dd\vct{x}^T\Qmat(\vct{x})\,\dd\vct{x}.
\end{equation*}
If the map is locally invertible, the signature \((+,-,-,-)\) is preserved by Sylvester's law. Position-dependent entries of \(\Qmat\) can be produced by curvilinear coordinates even though the target geometry remains flat; curvature is not represented by these entries alone. If the map fails to be locally invertible, the pullback may instead be degenerate.

At every point where the map is locally invertible, the ordered eigenvalue
function is defined by
\begin{equation}
\lambda(n):=\eqtext{ordered eigenvalue }n\eqtext{ of }\Qmat,
\qquad n\in\{1,2,3,4\}.
\label{eq:metric-eigenvalue-function}
\end{equation}
An orthogonal decomposition is then obtained for the real symmetric matrix:
\[
\Qmat=\Nmat\eqtext{diag}(\lambda(1),\lambda(2),\lambda(3),\lambda(4))\Nmat^T,
\qquad
\Nmat^T\Nmat=\I.
\]
The eigenvalues are ordered so that
\[
\lambda(1)>0,
\qquad
\lambda(2)<0,
\qquad
\lambda(3)<0,
\qquad
\lambda(4)<0.
\]
The corresponding direction function is defined by
\begin{equation}
\vct{u}(n):=\eqtext{column }n\eqtext{ of }\Nmat.
\label{eq:metric-direction-function}
\end{equation}
The local differential directions are defined by
\begin{equation}
\omega(n):=\sqrt{|\lambda(n)|}\,\vct{u}(n)^T\dd\vct{x}.
\label{eq:local-differential-directions}
\end{equation}
The interval is consequently
\begin{equation*}
\dd s^2
=\omega(1)^2-\omega(2)^2-\omega(3)^2-\omega(4)^2.
\end{equation*}
The three minus signs cannot be dropped: a Lorentzian interval is not
turned into a Euclidean norm by diagonalization. When \(\Qmat\) varies with
position, the local differential directions \(\omega(n)\) need not be
derivatives of globally defined coordinates.


\section{Covariant Transform Applications}\label{sec:transforms}
The generic similarity law and the circular rotor law were established in Section \ref{sec:transform-algebra}. The Lorentz action on real event, velocity, and momentum paravectors is separated here from similarity and is applied explicitly.

For \(\pv{L}=\exp(\theta\vct{u}\cdot\sigv/2)\), the passive event boost is \(\pv{X}'=\conj{\pv{L}}\pv{X}\conj{\pv{L}}\), the \(\pv{T}=\pv{L}\) specialization of Eq.\ \eqref{eq:general-event-congruence}. The usual Lorentz formulas are supplied by its scalar and vector parts, and \(\detf{\pv{X}}\) is preserved. The distinct similarity law \eqref{eq:general-field-similarity} is later applied in Section \ref{sec:em} to the electromagnetic field paravector \(\pv{F}\). The spatial rotor convention used earlier is active; this boost convention is passive.
\subsubsection*{General Transformation Identities}
Equation \eqref{eq:general-transform} is expanded here into differential and commutator form.
\begin{align*}
X' &= T^{-1} X T,\notag \\
&= T^{-1} ( t + \mathbf r \cdot \sigv ) T,\notag \\
&= t + \mathbf r \cdot ( T^{-1}\sigv T ),\notag \\
&= t + \mathbf r \cdot \sigv '.
\end{align*}
\begin{equation*}
T^{-1} = T^{\mathrm{R}*} / (T^{\mathrm{R}*} T).
\end{equation*}
For boosts and rotors one has \(T^{\mathrm{R}*}T=1\), so the inverse simplifies to
\begin{align*}
T^{\mathrm{R}*} T &= 1,\qquad T^{-1} = T^{\mathrm{R}*},\notag \\
T^{-1} T &= 1,\notag \\
d(T^{-1}) T &= -T^{-1} dT,\notag \\
dX' &= d(T^{-1}) X T + T^{-1} dX T + T^{-1} X dT,\notag \\
&= T^{-1} dX T.
\end{align*}
Inserting \(TT^{-1}\) into the differentiated expression isolates the commutator contribution:
\begin{align*}
&= T^{-1} dX T + (-T^{-1} dT) X' + X' (T^{-1} dT),\notag \\
&= T^{-1} dX T + [X', T^{-1} dT].
\end{align*}
\begin{align*}
d\Theta &:= T^{-1} dT,\notag \\
dX' &= T^{-1} dX T + [ X', d\Theta ],\notag \\
ddX' &= d(T^{-1}) dX T + T^{-1} ddX T + T^{-1} dX dT,\notag \\
&\quad + [dX', d\Theta ] + [X', dd\Theta ].
\end{align*}
\begin{align*}
d(T^{-1}) dX T &= (-T^{-1} dT) (T^{-1} dX T),\notag \\
&= -d\Theta (dX' - [ X', d\Theta ]).
\end{align*}
\begin{align*}
T^{-1} dX dT &= (T^{-1} dX T) (T^{-1} dT),\notag \\
&= (dX' - [ X', d\Theta ]) d\Theta.
\end{align*}
\begin{align*}
ddX' &= T^{-1} ddX T,\notag \\
&\quad +2 [ dX', d\Theta ],\notag \\
&\quad + [ X', dd\Theta ],\notag \\
&\quad + [ d\Theta , [ X', d\Theta ] ].
\end{align*}
\begin{align*}
d\Theta /dt &= \Omega,\notag \\
dd\Theta /dt^{2} &= d\Omega /dt,\notag \\
dX'/dt &= T^{-1} (dX/dt) T + [ X', \Omega ],\notag \\
ddX'/dt^{2} &= T^{-1} (ddX/dt^{2}) T,\notag \\
&\quad +2 [ dX'/dt, \Omega ],\notag \\
&\quad + [ X', d\Omega /dt ],\notag \\
&\quad + [ \Omega , [ X', \Omega ] ].
\end{align*}

\subsubsection*{Transformed Basis Elements}
The transformed basis relations below use the involution rules of Eqs.\ \eqref{eq:sigma-involutions} and \eqref{eq:i-involutions}.
\begin{align*}
T(x) &= f(x) + \mathbf g(x) \cdot \sigv, \\
T^{-1} &= ( f - \mathbf g \cdot \sigv ) / ( f^{2} - \mathbf g^{2} ), \\
dT &= df + d\mathbf g \cdot \sigv, \\
\sigma _{n}' &:= T^{-1} \sigma _{n} T, \\
\sigma _{n} &= T \sigma _{n}' T^{-1}.
\end{align*}
Applying the sigma-conjugate and reverse operators gives
\begin{align*}
( \sigma _{n}' )^{*} &= (T^{-1})^{*} \sigma _{n}^{*} T^{*},\notag \\
&= -(T^{*})^{-1} \sigma _{n} T^{*}.
\end{align*}
\begin{align*}
( \sigma _{n}' )^{\mathrm{R}} &= T^{\mathrm{R}} \sigma _{n}^{\mathrm{R}} (T^{-1})^{\mathrm{R}},\notag \\
&= T^{\mathrm{R}} \sigma _{n} (T^{\mathrm{R}})^{-1}.
\end{align*}
\begin{align*}
( \sigma _{n}' )^{*\mathrm{R}} &= -T^{*\mathrm{R}} \sigma _{n} (T^{-1})^{*\mathrm{R}},\notag \\
&= -T^{-1}\sigma _{n} T \text{ if } T^{-1} = T^{*\mathrm{R}},\notag \\
&= -\sigma _{n}' \text{ if } T^{-1} = T^{*\mathrm{R}}.
\end{align*}
\begin{align*}
\sigv ' algebra &= \sigv algebra,\notag \\
(\sigma _{n}')^{2} &= T^{-1}\sigma _{n}TT^{-1}\sigma _{n}T,\notag \\
&= T^{-1}\sigma _{n}^{2} T = 1.
\end{align*}
For \(n\neq m\), the mixed products transform as follows.\par
\begin{align*}
T^{-1}\sigma _{n}\sigma _{m}T &= T^{-1}\sigma _{n}TT^{-1}\sigma _{m}T = \sigma _{n}'\sigma _{m}',\notag \\
&= T^{-1}\epsilon _{nmk} i \sigma _{k}T = \epsilon _{nmk} i \sigma _{k}'.
\end{align*}
\begin{align*}
\sigma _{1}'\sigma _{2}'\sigma _{3}' &= T^{-1}\sigma _{1}TT^{-1}\sigma _{2}TT^{-1}\sigma _{3}T,\notag \\
&= T^{-1}\sigma _{1}\sigma _{2}\sigma _{3} T = i T^{-1}T = i.
\end{align*}
\begin{align*}
d\sigma _{n}' &= [\sigma _{n}', T^{-1} dT],\notag \\
d\sigv ' &= [\sigv ', T^{-1} dT]\qquad \text{vector},\notag \\
Z &= Z_{s} + Z_{\mathbf v} \cdot \sigv '\qquad \text{paravector},\notag \\
\langle\!\langle Z \rangle\!\rangle^{2} &= ZZ^{*\mathrm{R}},\notag \\
&= (Z_{s} + Z_{\mathbf v} \cdot \sigv ' )(Z_{s} - Z_{\mathbf v} \cdot \sigv ' ) \text{ if } T^{-1} = T^{*\mathrm{R}},\notag \\
&= Z_{s}^{2} - Z_{\mathbf v}^{2} \text{ if } T^{-1} = T^{*\mathrm{R}}.
\end{align*}
\begin{align*}
dZ &= dZ_{s} + dZ_{\mathbf v} \cdot \sigv ' + Z_{\mathbf v} \cdot d\sigv ',\notag \\
&= dZ_{s} + dZ_{\mathbf v} \cdot \sigv ' + Z_{\mathbf v} \cdot [\sigv ', T^{-1} dT].
\end{align*}
\begin{align*}
T^{-1} dT &= (f df - g dg) / ( f^{2} - g^{2} ),\notag \\
&\quad + (f d\mathbf g - df \mathbf g - i \mathbf g \times d\mathbf g) \cdot \sigv / ( f^{2} - \mathbf g^{2} ).
\end{align*}
The connection-like quantity
\begin{align*}
\Omega &:= T^{-1} (dT/dt),\notag \\
\Omega ' &:= T \Omega T^{-1} = (dT/dt) T^{-1},\notag \\
\eqtext{scalar}(\sigma _{n}') &= \eqtext{scalar}(T^{-1} \sigma _{n} T),\notag \\
&= \eqtext{scalar}(\sigma _{n} TT^{-1})\qquad \text{cyclic property},\notag \\
&= \eqtext{scalar}(\sigma _{n}),\notag \\
&= 0.
\end{align*}
\begin{align*}
\eqtext{scalar}(\Omega \sigma _{n}'),\notag \\
&= \eqtext{scalar}(T^{-1} (dT/dt) T^{-1} \sigma _{n} T),\notag \\
&= \eqtext{scalar}((dT/dt) T^{-1} \sigma _{n} TT^{-1})\qquad \text{cyclic},\notag \\
&= \eqtext{scalar}((dT/dt) T^{-1} \sigma _{n} ),\notag \\
&= \eqtext{scalar}( \Omega ' \sigma _{n} ).
\end{align*}
\begin{align*}
k_{n} + i \omega _{n} &:= \eqtext{scalar}( 2 \Omega ' \sigma _{n} ),\notag \\
\Omega &= \Omega_s + \tfrac{1}{2} (k + i \omega ) \cdot \sigv ',\notag \\
[\sigv ', \omega \cdot ( i \sigv ' )] &= -2 \omega \times \sigv ',\notag \\
[\sigv ', \mathbf k \cdot \sigv '] &= 2 i \mathbf k \times \sigv ',\notag \\
d\sigv '/dt &= [\sigv ', \Omega ],\notag \\
&= i ( k + i \omega ) \times \sigv ',\notag \\
&= ( -\omega + i k ) \times \sigv '.
\end{align*}
\begin{align*}
Z_{\mathbf v} \cdot [\sigv ', \Omega ] &= ( -Z_{\mathbf v} \times \omega + i Z_{\mathbf v} \times k ) \cdot \sigv ',\notag \\
dZ/dt &= dZ_{s}/dt,\notag \\
&\quad + (dZ_{\mathbf v}/dt - Z_{\mathbf v} \times \omega + i Z_{\mathbf v} \times k ) \cdot \sigv '.
\end{align*}
The same decomposition is useful in non-inertial settings, where the transformed basis carries effective angular-velocity and boost-like terms:
\begin{align*}
X &= t + \mathbf r \cdot \sigv ',\notag \\
dX/dt &= dt/dt,\notag \\
&\quad + (d\mathbf r/dt - \mathbf r \times \omega + i \mathbf r \times \mathbf k ) \cdot \sigv '.
\end{align*}
The corresponding energy--momentum split is
\begin{align*}
m dX/dt &= E + \mathbf P' \cdot \sigv ', \\
E / m &= dt/dt, \\
\mathbf P' / m &= d\mathbf r/dt + \mathbf r \times (-\omega + i k).
\end{align*}
If \(T^{-1} = T^{*\mathrm{R}}\), then the transformed interval keeps the same Minkowski form:
\begin{align*}
\langle\!\langle dX \rangle\!\rangle^{2} &= dt^{2} - (d\mathbf r - \mathbf r \times \omega dt + i \mathbf r \times \mathbf k dt)^{2},\notag \\
&= (1 - (d\mathbf r/dt - \mathbf r \times \omega + i \mathbf r \times \mathbf k)^{2}) dt^{2}.
\end{align*}
\begin{equation*}
ds' := \sqrt{1 - (\mathbf v - \mathbf r \times \omega + i \mathbf r \times \mathbf k)^{2}} dt.
\end{equation*}

\subsubsection*{Lorentz Boosts}
The boost identities in this subsection unpack Eq.\ \eqref{eq:boost} and its component consequences.
\begin{align*}
\theta &= \operatorname{arctanh}(\beta ),\notag \\
\beta &= \tanh(\theta ),\notag \\
\gamma &= \cosh(\theta ),\notag \\
\gamma \beta &= \sinh(\theta ),\notag \\
\mathbf u &= (u_{1}, u_{2}, u_{3}), \mathbf u^{2} = 1,\notag \\
L(\theta ,\mathbf u) &:= \exp( \tfrac{1}{2} \theta (\mathbf u \cdot \sigv ) ),\notag \\
&= \cosh( \tfrac{1}{2} \theta ) + \sinh( \tfrac{1}{2} \theta ) (\mathbf u \cdot \sigv ).
\end{align*}
\begin{equation*}
L^{-1} = \exp( -\tfrac{1}{2} \theta (\mathbf u \cdot \sigv ) ).
\end{equation*}
The boost inverse satisfies \(L^{-1}=L^{*}=L^{*\mathrm{R}}\). Boosts along the same axis commute:
\begin{equation*}
[ L(\theta ,\mathbf u), L(\theta ',\mathbf u) ] = 0.
\end{equation*}
Consequently, collinear boosts add by rapidity:
\begin{align*}
L(\theta ', \mathbf u) &= L(\theta _{1}, \mathbf u) L(\theta _{2}, \mathbf u), \\
\beta ' &= (\beta _{1} + \beta _{2}) / (1 + \beta _{1} \beta _{2}), \\
\gamma ' &= \gamma _{1} \gamma _{2} (1 + \beta _{1} \beta _{2}).
\end{align*}
Acting on the spacetime paravector \(X=t+\mathbf r\cdot\sigv\), the boost gives
\begin{align*}
X &= t + \mathbf r \cdot \sigv,\notag \\
X' &= L^{*} X L,\notag \\
&= t' + \mathbf r' \cdot \sigv.
\end{align*}
\begin{align*}
t' &= t \cosh(\theta ) - \sinh(\theta ) (\mathbf u \cdot \mathbf r),\notag \\
&= \gamma ( t - \beta (\mathbf u \cdot \mathbf r) ).
\end{align*}
Splitting \(\mathbf r\) into parts parallel and perpendicular to the boost direction,
\begin{align*}
\eqtext{par}(\mathbf r) &:= (\mathbf r \cdot \mathbf u) \mathbf u\qquad \text{parallel to u},\notag \\
\eqtext{prp}(\mathbf r) &:= \mathbf r - \eqtext{par}(\mathbf r)\qquad \text{perpendicular to u},\notag \\
\mathbf r &= \eqtext{par}( \mathbf r ) + \eqtext{prp}( \mathbf r ),\notag \\
\eqtext{par}( \mathbf r' ) &= \gamma ( \eqtext{par}( \mathbf r ) - \beta t \mathbf u),\notag \\
\eqtext{prp}( \mathbf r' ) &= \eqtext{prp}( \mathbf r ),\notag \\
\mathbf r' &= \eqtext{par}( \mathbf r' ) + \eqtext{prp}( \mathbf r ' ),\notag \\
&= \eqtext{prp}( \mathbf r ) + \gamma ( \eqtext{par}( \mathbf r ) - \beta t \mathbf u),\notag \\
&= \mathbf r +( \gamma ( \mathbf r \cdot \mathbf u - \beta t ) - \mathbf r \cdot \mathbf u ) \mathbf u.
\end{align*}
In the small-velocity limit \(\beta\to 0\),
\begin{equation*}
\gamma \Rightarrow 1, t' \Rightarrow t, \mathbf r' \Rightarrow \mathbf r.
\end{equation*}
The boost preserves the paravector interval:
\begin{align*}
\langle\!\langle X' \rangle\!\rangle^{2} &= X' X'^{*},\notag \\
&= (L^{*} X L) (L^{*} X L)^{*},\notag \\
&= L^{*} X L L^{*} X^{*} L,\notag \\
&= L^{*} X X^{*} L,\notag \\
&= \langle\!\langle X \rangle\!\rangle^{2} L^{*} L,\notag \\
&= \langle\!\langle X \rangle\!\rangle^{2}.
\end{align*}

\subsubsection*{Spatial Rotations}
The rotor identities in this subsection unpack Eq.\ \eqref{eq:rotor} and its component consequences.
\begin{align*}
R(\theta ,\mathbf u) &:= \exp( \tfrac{1}{2} \theta (\mathbf u \cdot i \sigv ) ),\notag \\
&= \cos( \tfrac{1}{2} \theta ) + \sin( \tfrac{1}{2} \theta ) (\mathbf u \cdot i \sigv )\qquad \text{rotor}.
\end{align*}
\begin{equation*}
R^{-1} = \exp( -\tfrac{1}{2} \theta (\mathbf u \cdot i \sigv ) ).
\end{equation*}
The rotor inverse satisfies \(R^{-1}=R^{\mathrm{R}}=R^{*\mathrm{R}}\).
\begin{align*}
X' &= R^{\mathrm{R}} X R,\notag \\
&= t' + \mathbf r' \cdot \sigv.
\end{align*}
\begin{align*}
t' &= t,\notag \\
\mathbf r' &= \mathbf r \cos(\theta ) + (\mathbf u \times \mathbf r) \sin(\theta ),\notag \\
&\quad + (\mathbf u \cdot \mathbf r) (1 - \cos(\theta )) \mathbf u,\notag \\
\qquad \text{Rodrigues's formula}.
\end{align*}
Differentiating the transformed quantities gives
\begin{align*}
d(R^{\mathrm{R}}) &= (dR)^{\mathrm{R}} = dR^{\mathrm{R}}, \\
R^{\mathrm{R}} R &= 1, \\
dR^{\mathrm{R}} R + R^{\mathrm{R}} dR &= 0, \\
dR &= -R dR^{\mathrm{R}} R.
\end{align*}
The associated bivector-valued angular-velocity element is
\begin{align*}
\Omega &:= -R (dR^{\mathrm{R}}/ds) = (dR/ds) R^{\mathrm{R}},\notag \\
\Omega ^{\mathrm{R}} &= -\Omega,\notag \\
\Omega ' &= R^{\mathrm{R}} \Omega R = R (dR/ds),\notag \\
\Omega '^{\mathrm{R}} &= -\Omega ',\notag \\
dR/ds &= \Omega R,\notag \\
X' &= R^{\mathrm{R}} X R \Rightarrow R X' = X R,\notag \\
dR X' + R dX' &= dX R + X dR,\notag \\
R dX' &= dX R + X dR - dR X',\notag \\
&= dX R + X dR - dR X'.
\end{align*}
\begin{align*}
dX'/ds &= R^{\mathrm{R}} (dX/ds) R + [ X', \Omega ' ],\notag \\
&= R^{\mathrm{R}} ( dX/ds + [ X, \Omega ] ) R.
\end{align*}

\subsubsection*{Compositions of Boosts and Rotations}
\begin{align*}
T &= L R, \\
T^{-1} &= T^{*\mathrm{R}} = R^{*\mathrm{R}} L^{*\mathrm{R}} = R^{\mathrm{R}} L^{*}, \\
T &= R L, \\
T^{-1} &= T^{*\mathrm{R}} = L^{*\mathrm{R}} R^{*\mathrm{R}} = L^{*} R^{\mathrm{R}}, \\
T &= R_{1} R_{2}, \\
T^{-1} &= T^{*\mathrm{R}} = R_{2}^{\mathrm{R}} R_{1}^{\mathrm{R}}, \\
T &= L_{1} L_{2}, \\
T^{-1} &= T^{*\mathrm{R}} = L_{2}^{*}L_{1}^{*}.
\end{align*}


\section{Differential Operators}\label{sec:differential}
One fixed set of differential conventions is used in the electromagnetic and wave-equation sections. The conventions are collected here so that the same numbered equations can be referenced later without repeated operator definitions.

For spatial derivatives with respect to \(\vct{r}=(r_1,r_2,r_3)\), the
spatial derivative vector is defined by
\begin{equation}
\delr:=\left\{\frac{\partial}{\partial r_1},\frac{\partial}{\partial r_2},\frac{\partial}{\partial r_3}\right\}
=\{\partial_1,\partial_2,\partial_3\}.
\label{eq:gradient-notation}
\end{equation}
For shorthand, the coordinate symbol \(\vct{r}\) is dropped in the last form, so \(\partial_i:=\partial/\partial r_i\) for \(i\in\{1,2,3\}\). The corresponding spatial Laplacian is
\begin{equation}
\gradv^2:=\partial_1^2+\partial_2^2+\partial_3^2.
\label{eq:laplacian}
\end{equation}
The differential paravector and its conjugate are
\begin{equation}
\pv{\partial}:=\partial_t+\sigv\cdot\gradv
=\partial_t+\pv{\sigma}_1\partial_1+\pv{\sigma}_2\partial_2+\pv{\sigma}_3\partial_3.
\label{eq:partial-operator}
\end{equation}
\begin{equation}
\conj{\pv{\partial}}:=\partial_t-\sigv\cdot\gradv
=\partial_t-\pv{\sigma}_1\partial_1-\pv{\sigma}_2\partial_2-\pv{\sigma}_3\partial_3.
\label{eq:partial-conjugate}
\end{equation}
Because the coordinate derivatives commute on smooth fields, the scalar d'Alembertian is obtained from their ordered products:
\begin{equation}
\Box:=\pv{\partial}\conj{\pv{\partial}}.
\label{eq:dalembertian}
\end{equation}
The reversed order gives the same scalar operator:
\begin{equation}
\Box=\conj{\pv{\partial}}\pv{\partial}
=\partial_t^2-\gradv^2.
\label{eq:dalembertian-expansion}
\end{equation}
The ordinary two-term Leibniz rule must be corrected when the left field does not commute with the sigma basis. For smooth paravector fields \(\pv{Z}\) and \(\pv{Z}'\), the following result is obtained by direct expansion:
\begin{equation}
\pv{\partial}(\pv{Z}\pv{Z}')
=(\pv{\partial}\pv{Z})\pv{Z}'
+\pv{Z}(\pv{\partial}\pv{Z}')
+\sum_{n=1}^{3}\com{\pv{\sigma}_n}{\pv{Z}}\,\partial_n\pv{Z}'.
\label{eq:paravector-leibniz}
\end{equation}
The final term is generated when \(\pv{Z}\) is moved through the left-acting basis element \(\pv{\sigma}_n\). It is zero for a scalar \(\pv{Z}\), or whenever \(\com{\pv{\sigma}_n}{\pv{Z}}=0\) for every contributing direction, and the ordinary Leibniz rule is then recovered.
\subsection{Complex Fields}\label{subsec:differential-complex-fields}
The differential identities are first recorded for an arbitrary complex
paravector field. The letter \(\pv{Z}\) is used here so that \(\pv{F}\) can
be reserved for the electromagnetic field introduced later:
\begin{align*}
\pv{\partial} &= \partial_t+\gradv\cdot\sigv,\\
\pv{Z} &:= (a+ib)+(\vct{A}+i\vct{B})\cdot\sigv.
\end{align*}
At each time and position, \(a\) and \(b\) return real numbers, while
\(\vct{A}\) and \(\vct{B}\) return real three-component vectors:
\begin{align*}
a &: \R\times\R^3\to\R, &
\vct{A} &: \R\times\R^3\to\R^3,\\
b &: \R\times\R^3\to\R, &
\vct{B} &: \R\times\R^3\to\R^3,
\end{align*}
so that
\begin{equation*}
\pv{Z}:\R\times\R^3\to\CPspace.
\end{equation*}
The letters \(\vct{A}\) and \(\vct{B}\) in this subsection are generic
placeholders. Their electromagnetic meanings are assigned only in Section
\ref{sec:em}.
Direct multiplication gives the scalar and vector parts of
\(\pv{\partial}\pv{Z}\) and \(i\pv{\partial}\pv{Z}\):
\begin{componenttable}[tab:components-generic-derivative]{\pv{\partial} \pv{Z}}
\componentrow{real scalar}{\partial_t a+\gradv\cdot\vct{A}}{}
\componentrow{img scalar}{\partial_t b+\gradv\cdot\vct{B}}{}
\componentrow{real vector}{\gradv a+\partial_t\vct{A}-\gradv\times\vct{B}}{}
\componentrow{img vector}{\gradv b+\gradv\times\vct{A}+\partial_t\vct{B}}{}
\end{componenttable}

\begin{componenttable}[tab:components-generic-iderivative]{i\,\pv{\partial} \pv{Z}}
\componentrow{real scalar}{-\partial_t b-\gradv\cdot\vct{B}}{}
\componentrow{img scalar}{\partial_t a+\gradv\cdot\vct{A}}{}
\componentrow{real vector}{-\gradv b-\gradv\times\vct{A}-\partial_t\vct{B}}{}
\componentrow{img vector}{\gradv a+\partial_t\vct{A}-\gradv\times\vct{B}}{}
\end{componenttable}


\section{Electromagnetic Fields}\label{sec:em}
Electromagnetic fields, gauge structure, and the Lorentz force can be expressed in compact paravector forms. In physical content, Maxwell theory is recovered in rationalized Heaviside--Lorentz units \cite{maxwell1865}; the packaging is drawn from later STA presentations \cite{hestenes1966,baylis1999,doran2003}.

\subsection{Fields, potentials, and Maxwell equations}
The differential conventions are fixed once in Section \ref{sec:differential}; in particular, the spatial derivative vector, the shorthand \(\partial_i\), the paravector derivative, its conjugate, and the d'Alembertian are defined by Eqs.\ \eqref{eq:gradient-notation}--\eqref{eq:dalembertian}. With those operators fixed, and with \(V,\rho,\lambda,S,q\in\R\), \(\vct{E},\vct{B},\vct{A},\vct{J}\in\R^3\), the electromagnetic field vector is defined by
\begin{equation}
\vct{F}:=\vct{E}+i\vct{B}.
\label{eq:field-vector}
\end{equation}
The corresponding pure-vector paravector is defined by
\begin{equation}
\pv{F}:=\vct{F}\cdot\sigv.
\label{eq:F-Phi}
\end{equation}
The real potential paravector is defined separately by
\begin{equation}
\pv{\Phi}:=V+\vct{A}\cdot\sigv.
\label{eq:potential-paravector}
\end{equation}
Here \(\pv{\Phi}\) is a real paravector, \(\vct{F}\) is a complex three-vector, and \(\pv{F}\) is a pure-vector complex paravector. Electric and magnetic fields are packaged into \(\pv{F}\) by Eq.\ \eqref{eq:F-Phi}, and scalar and vector potentials are packaged into \(\pv{\Phi}\).

The scalar part generated by the potential is defined by
\begin{equation}
S:=\partial_t V+\gradv\cdot\vct{A}.
\label{eq:potential-scalar}
\end{equation}
The relations in Eq.\ \eqref{eq:E-B-from-potential} are then encoded by
\begin{equation}
\pv{\partial}\pv{\Phi}=S+\conj{\pv{F}}.
\label{eq:potential-gradient}
\end{equation}
with
\begin{equation}
\vct{E}=-\partial_t\vct{A}-\gradv V,
\qquad
\vct{B}=\gradv\times\vct{A}.
\label{eq:E-B-from-potential}
\end{equation}
The compact Maxwell equation is
\begin{equation}
\pv{\partial}\pv{F}=\rho-\vct{J}\cdot\sigv.
\label{eq:maxwell}
\end{equation}
By splitting the real/imaginary scalar/vector parts of \eqref{eq:maxwell}, the following equations are reproduced:
\begin{align}
\gradv\cdot\vct{E}&=\rho,
&
\gradv\cdot\vct{B}&=0,
\label{eq:maxwell-scalars}\\
\gradv\times\vct{B}&=\vct{J}+\partial_t\vct{E},
&
\gradv\times\vct{E}&=-\partial_t\vct{B}.
\label{eq:maxwell-vectors}
\end{align}
Agreement with standard space-time-algebra treatments is obtained in Eqs.\ \eqref{eq:maxwell-scalars} and \eqref{eq:maxwell-vectors}, subject to the sign conventions fixed by Eqs.\ \eqref{eq:partial-operator}, \eqref{eq:partial-conjugate}, \eqref{eq:dalembertian}, and \eqref{eq:F-Phi}. Maxwell's field equations are recovered component by component in the present signature and field convention \cite{maxwell1865,hestenes1966,baylis1999,jackson1998}.
Separate primitive divergence and curl operators are avoided at this stage in Eq.\ \eqref{eq:maxwell}. They are produced automatically when the scalar and vector parts of the single paravector equation are separated.

\subsection{Boost transformation of electromagnetic fields}
The field paravector is transformed by the same boost element \(\pv{L}\) from Eq.\ \eqref{eq:boost}:
\begin{equation}
\pv{F}'=\conj{\pv{L}}\pv{F}\pv{L}.
\label{eq:field-transform}
\end{equation}
Equation \eqref{eq:field-transform} is deliberately different from the event congruence \(\pv{X}'=\conj{\pv{L}}\pv{X}\conj{\pv{L}}\). The event and field laws were derived from different intertwining requirements in Section \ref{sec:transform-algebra}; the usual mixing of \(\eqtext{par}(\vct{E})\), \(\eqtext{prp}(\vct{E})\), \(\eqtext{par}(\vct{B})\), and \(\eqtext{prp}(\vct{B})\) is obtained from the resulting similarity.

\subsection{Wave equations, gauge, and Lorentz force}
The equivalent pair is obtained by applying \(\conj{\pv{\partial}}\) to
Eq.\ \eqref{eq:maxwell} and \(\pv{\partial}\) to its star-conjugate equation:
\begin{align}
\Box\pv{F}=\conj{\pv{\partial}}(\rho-\vct{J}\cdot\sigv),
\label{eq:wave-F}
\\
\Box\conj{\pv{F}}
=\pv{\partial}(\rho+\vct{J}\cdot\sigv).
\label{eq:wave-F-conjugate}
\end{align}
Charge conservation is obtained from the scalar part of the first relation:
\begin{equation}
\partial_t\rho+\gradv\cdot\vct{J}=0.
\label{eq:charge-conservation}
\end{equation}
Equation \eqref{eq:charge-conservation} is the continuity equation written in the present notation.
For a real scalar gauge function \(\lambda\in\R\), a function input is used so
that confusion with a Lorentz-frame prime is avoided:
\begin{equation}
\pv{\Phi}(\lambda):=\pv{\Phi}+\conj{\pv{\partial}}\lambda.
\label{eq:gauge}
\end{equation}
The scalar and vector potential functions contained in this paravector are
defined by
\begin{equation}
V(\lambda):=V+\partial_t\lambda.
\label{eq:gauge-scalar-potential}
\end{equation}
\begin{equation}
\vct{A}(\lambda):=\vct{A}-\gradv\lambda.
\label{eq:gauge-vector-potential}
\end{equation}
The corresponding field function is defined by
\begin{equation}
\pv{F}(\lambda):=\pv{F}.
\label{eq:gauge-field}
\end{equation}
The potential \(\pv{\Phi}\) is retained even when the field is physically primary because the redundancy is carried by \(\pv{\Phi}\), whereas the gauge-invariant derivative content is carried by \(\pv{F}\).

For a particle of real mass \(m>0\), real charge \(q\), and four-velocity \(\pv{U}\), the four-force formula is
\begin{equation}
m\frac{\dd \pv{U}}{\dd s}=q\,\eqtext{real}\bigl(\pv{F}\pv{U}\bigr)
=q\gamma\Big(\vct{v}\cdot\vct{E}+(\vct{E}+\vct{v}\times\vct{B})\cdot\sigv\Big),
\label{eq:lorentz-force}
\end{equation}
in which both relativistic power and the Lorentz three-force are packaged. Since \(\dd/\dd s=\gamma\,\dd/\dd t\), the equations \(\dd E/\dd t=q\vct{v}\cdot\vct{E}\) and \(\dd\vct{P}/\dd t=q(\vct{E}+\vct{v}\times\vct{B})\) are obtained by division by \(\gamma\).

\subsection{Lagrangian and energy-momentum remarks}
The electromagnetic particle Lagrangian is defined by
\begin{equation}
\mathcal L:=-\frac{m}{\gamma}-q(V-\vct{v}\cdot\vct{A}),
\label{eq:em-lagrangian}
\end{equation}
and $\dd\vct{P}/\dd t=q(\vct{E}+\vct{v}\times\vct{B})$ is recovered from its Euler--Lagrange equations.
The free action from Eq.\ \eqref{eq:free-lagrangian} is supplemented by the potential term, and the field-strength form of the force law is recovered after variation.

The frame-dependent field energy-density and flux paravector is defined by
\begin{equation}
\pv{T}:=\frac{1}{2}\pv{F}\herm{\pv{F}}.
\label{eq:energy-momentum-paravector}
\end{equation}
Its scalar and vector parts are
\begin{equation}
\pv{T}
=\frac{1}{2}(\vct{E}^2+\vct{B}^2)
+(\vct{E}\times\vct{B})\cdot\sigv.
\label{eq:energy-momentum-components}
\end{equation}
The energy density and Poynting vector are packaged by Eq.\ \eqref{eq:energy-momentum-paravector} into one frame-dependent paravector. The precise caveat and the corresponding covariant stress-energy map are recorded in the supporting subsection.

\subsection{Supporting Electromagnetic Derivations}
The field decomposition, potential reconstruction, source formulas, gauge transformations, force-law expansions, and variational identities underlying Eqs.\ \eqref{eq:gradient-notation}--\eqref{eq:energy-momentum-paravector} are written out in the supporting derivations.
\subsubsection*{Electric and Magnetic Fields}
The next step is to spell out the electromagnetic field content that was previously presented in compact STA form. These identities make the relation between the vector fields and the single field paravector fully explicit. See Eqs.\ \eqref{eq:dalembert}, \eqref{eq:F-Phi}, and \eqref{eq:maxwell}. Here \(V,\rho,S,q\in\R\), \(\mathbf A,\mathbf E,\mathbf B,\mathbf J\in\R^3\), \(\Phi=V+\mathbf A\cdot\sigv\in\R\oplus\R^3\), and \(F=(\mathbf E+i\mathbf B)\cdot\sigv\in\{0\}\oplus\C^3\).
\noindent\textit{Introduce the following vector notation.}\par
\begin{align}
\delr{} &= \{\partial /\partial r_{1}, \partial /\partial r_{2}, \partial /\partial r_{3}\},\notag \\
\partial _{i} &= \partial / \partial r_{i},\notag \\
\delr{}^{2} &= \partial _{1}^{2} + \partial _{2}^{2} + \partial _{3}^{2},\notag \\
\partial &= \partial _{t} + \sigv \cdot \delr{},\notag \\
&= \partial _{t} + \sigma _{1}\partial _{1}+ \sigma _{2}\partial _{2} + \sigma _{3}\partial _{3}.
\end{align}
\noindent\textit{Introduce the following vectors.}\par
\begin{align}
\mathbf E &= \{\mathbf E_{1}, \mathbf E_{2}, \mathbf E_{3}\}\qquad \text{electric field},\notag \\
\mathbf B &= \{\mathbf B_{1}, \mathbf B_{2}, \mathbf B_{3}\}\qquad \text{magnetic field},\notag \\
\mathcal{F} &= \mathbf E + i \mathbf B\qquad \text{electromagnetic field},\notag \\
F &= (\mathbf E + i \mathbf B) \cdot \sigv , i = \sigma _{1}\sigma _{2}\sigma _{3},\notag \\
&= \mathcal{F} \cdot \sigv.
\end{align}
\begin{align}
F^{*} &= (-\mathbf E + i \mathbf B) \cdot \sigv,\notag \\
&= -\mathcal{F}^{*} \cdot \sigv.
\end{align}

\noindent\textit{Introduce the source-density paravector.}\par
\begin{align}
D &= \rho + \mathbf J \cdot \sigv, \\
D^{*} &= \rho - \mathbf J \cdot \sigv.
\end{align}
Let \(\rho\) denote the scalar charge density.\par
Let \(\mathbf J\) denote the current-density three-vector.\par
\begin{align}
\partial F &= D^{*}\qquad \text{Maxwell's equations},\notag \\
&= \rho - \mathbf J \cdot \sigv.
\end{align}
\begin{equation}
\partial ^{*} F^{*} = D.
\end{equation}
\noindent\textit{Expanding the left-hand side gives.}\par
\begin{align}
\partial F &= (\partial _{t} + \sigv \cdot \delr{}) ( (\mathbf E + i \mathbf B) \cdot \sigv ),\notag \\
&= ( \delr{} \cdot \mathbf E ),\notag \\
&\quad + ( \delr{} \cdot \mathbf B ) i,\notag \\
&\quad + (\partial _{t} \mathbf E - \delr{} \times \mathbf B) \cdot \sigv,\notag \\
&\quad + (\partial _{t} \mathbf B + \delr{} \times \mathbf E) \cdot (i \sigv ).
\end{align}
\noindent\textit{Rearranging the vector equations gives.}\par
\begin{align}
\delr{} \times \mathbf B &= \mathbf J + \partial _{t} \mathbf E, \\
\delr{} \times \mathbf E &= -\partial _{t} \mathbf B.
\end{align}

\subsubsection*{Lorentz Transformation of the Field}
Once the electromagnetic field has been assembled into the single paravector \(F\), the boost law from Eq.\ \eqref{eq:field-transform} can be expanded directly in components. For a boost
\(L=\exp(\tfrac12\theta\,\mathbf u\cdot\sigv)\) with \(\mathbf u^{2}=1\), \(\beta=\tanh(\theta)\), and \(\gamma=\cosh(\theta)\), the parallel and perpendicular field parts relative to \(\mathbf u\) obey
\begin{align}
F' &= L^{*} F L, \\
\eqtext{par}( \mathbf E ' ) &= \eqtext{par}( \mathbf E ), \\
\eqtext{par}( \mathbf B ' ) &= \eqtext{par}( \mathbf B ), \\
\eqtext{prp}( \mathbf E ' ) &= \gamma ( \eqtext{prp}( \mathbf E ) + \beta \mathbf u \times \mathbf B ), \\
\eqtext{prp}( \mathbf B ') &= \gamma ( \eqtext{prp}( \mathbf B ) - \beta \mathbf u \times \mathbf E ).
\end{align}
\noindent\textit{These identities can be compressed into the single relation.}\par
\begin{align}
\mathbf E ' &= \eqtext{par}(\mathbf E) + \gamma ( \eqtext{prp}(\mathbf E) + \beta \mathbf u \times \mathbf B ), \\
\mathbf B ' &= \eqtext{par}(\mathbf B) + \gamma ( \eqtext{prp}(\mathbf B) - \beta \mathbf u \times \mathbf E ), \\
\partial ' &= L^{*} \partial L, \\
\partial ' F' &= L^{*} (\rho - \mathbf J \cdot \sigv ) L.
\end{align}

\subsubsection*{Wave Operator and d’Alembertian}
The wave operator is written out here in full so that its factorization into the d'Alembertian and the component field equations can be read directly. See Eq.\ \eqref{eq:dalembert}.
\begin{align}
\partial &= \partial _{t} + \sigv \cdot \delr{}, \\
\partial ^{*} &= \partial _{t} - \sigv \cdot \delr{}.
\end{align}
\noindent\textit{The d'Alembertian is.}\par
\begin{align}
\square &:= \partial _{t}^{2} - \delr{}^{2},\notag \\
&= \partial \partial ^{*} = \partial ^{*} \partial,\notag \\
&= \langle\!\langle \partial \rangle\!\rangle^{2}.
\end{align}
\noindent\emph{Note.} This expansion uses the sigma-algebra conventions fixed earlier in the manuscript.\par
\begin{equation}
\partial ^{2} = \partial _{t}^{2} + \delr{}^{2} + 2 \partial _{t} \delr{} \cdot \sigv.
\end{equation}
\noindent\textit{This differs from }\(\square\)\textit{.}\par
\noindent\emph{Note.} Sign conventions for \(F\) vary across the literature; the present manuscript keeps the convention fixed by Eqs.\ \eqref{eq:F-Phi} and \eqref{eq:maxwell}.\par
\noindent\emph{Note.} Some authors instead use the metric signature \((- ,+,+,+)\).\par
\noindent\emph{Note.} Some authors denote the wave operator with a different box-symbol convention.\par
\noindent\textit{The corresponding wave equation is.}\par
\begin{align}
\square F &= \partial ^{*} \partial F,\notag \\
&= \partial ^{*} (\rho - \mathbf J \cdot \sigv ).
\end{align}
\begin{align}
\square F^{*} &= \partial (\rho + \mathbf J \cdot \sigv ),\notag \\
( \square \mathbf E + i \square \mathbf B ) \cdot \sigv,\notag \\
&= (\partial _{t} \rho + \delr{} \cdot \mathbf J) + (-\partial _{t} \mathbf J - \delr{} \rho + i (\delr{} \times \mathbf J)) \cdot \sigv.
\end{align}
\begin{equation}
\partial _{t} \rho = - \delr{} \cdot \mathbf J\qquad \text{conservation of charge}.
\end{equation}
If \(\rho = 0\) and \(\mathbf J = 0\), then the following relation holds.\par
\begin{align}
\square \mathbf E &= 0 \Rightarrow \partial _{t}^{2} \mathbf E = \delr{}^{2} \mathbf E, \\
\square \mathbf B &= 0 \Rightarrow \partial _{t}^{2} \mathbf B = \delr{}^{2} \mathbf B.
\end{align}

\subsubsection*{Potential Fields}
The potential-based formulation is expanded here because it is the point at which gauge freedom, source reconstruction, and the emergence of the electric and magnetic fields all become transparent at once. See Eqs.\ \eqref{eq:potential-gradient} and \eqref{eq:E-B-from-potential}. Here \(V,S,\rho,\lambda\in\R\), \(\mathbf A,\mathbf E,\mathbf B,\mathbf J\in\R^3\), \(\Phi\in\R\oplus\R^3\), and \(F\in\{0\}\oplus\C^3\subset\C\oplus\C^3\).
\begin{align}
F^{*} &= (-\mathbf E + i \mathbf B) \cdot \sigv,\notag \\
\Phi &= V + \mathbf A \cdot \sigv,\notag \\
\partial \Phi &= \partial _{t} V + \delr{} \cdot \mathbf A,\notag \\
&\quad + ( \partial _{t} \mathbf A + \delr{} V ) \cdot \sigv,\notag \\
&\quad + i ( \delr{} \times \mathbf A ) \cdot \sigv.
\end{align}
\begin{align}
S &:= \eqtext{scalar}( \partial \Phi ),\notag \\
&= \partial _{t} V + \delr{} \cdot \mathbf A.
\end{align}
\noindent\textit{The target component identities are.}\par
\begin{align}
\mathbf E &= - \partial _{t} \mathbf A - \delr{} V, \\
\mathbf B &= \delr{} \times \mathbf A.
\end{align}
\noindent\textit{Consider the following identity.}\par
\begin{align}
\partial \Phi &= S + F^{*}, \\
F^{*} &= \partial \Phi - S, \\
\eqtext{scalar}(F^{*}) &= \eqtext{scalar}(\partial \Phi - S) = 0, \\
\eqtext{vector}(F^{*}) &= \eqtext{vector}( \partial \Phi ).
\end{align}
\noindent\emph{Note.} always differentiate first.\par
\begin{equation}
(\partial \Phi )^{*\mathrm{R}} \neq \Phi ^{*\mathrm{R}} \partial ^{*\mathrm{R}}.
\end{equation}
\noindent\textit{Equivalently, without introducing \(S\), one may write.}\par
\begin{equation}
F^{*} = (\partial \Phi - (\partial \Phi )^{*\mathrm{R}}) / 2.
\end{equation}

\subsubsection*{Sources Derived from the Potential}
\noindent\textit{Recall the previously established definitions.}\par
\begin{align}
F &= (\mathbf E + i \mathbf B) \cdot \sigv,\notag \\
\Phi &= V + \mathbf A \cdot \sigv,\notag \\
S &= \partial _{t} V + \delr{} \cdot \mathbf A,\notag \\
F^{*} &= \partial \Phi - S \Rightarrow F = \partial ^{*} \Phi ^{*} - S,\notag \\
\partial F &= \partial \partial ^{*} \Phi ^{*} - \partial S,\notag \\
&= \square \Phi ^{*} - \partial S,\notag \\
&= ( \square V - \partial _{t} S) + (- \square \mathbf A - \delr{} S) \cdot \sigv,\notag \\
&= \rho - \mathbf J \cdot \sigv.
\end{align}
\begin{align}
\rho &= \square V - \partial _{t} S, \\
\mathbf J &= \square \mathbf A + \delr{} S.
\end{align}
If \(S = 0,\), then the following relation holds.\par
\noindent\textit{Lorenz gauge condition.}\par
\begin{align}
\partial _{t} V &= -\delr{} \cdot \mathbf A, \\
\rho &= \square V, \\
\mathbf J &= \square \mathbf A.
\end{align}

\subsubsection*{Gauge Transform of the Potential}
\begin{align}
\eqtext{scalar}( \partial \Phi ) &= S, \\
S &= \partial _{t} V + \delr{} \cdot \mathbf A.
\end{align}
Let \(\lambda\) denote scalar field.\par
\begin{align}
\Phi ' &= \Phi + \partial ^{*} \lambda,\notag \\
\partial \Phi ' &= \partial (\Phi + \partial ^{*}\lambda ),\notag \\
&= \partial \Phi + \partial (\partial ^{*}\lambda ),\notag \\
&= \partial \Phi + \square \lambda.
\end{align}
\noindent\emph{Note.} The quantity \(\square\lambda\) is again a scalar field.\par
\begin{align}
\eqtext{scalar}( \partial \Phi ' ) &= S + \square \lambda, \\
\eqtext{vector}( \partial \Phi ' ) &= \eqtext{vector}( \partial \Phi ), \\
F' &= F = (\mathbf E + i \mathbf B) \cdot \sigv.
\end{align}

\subsubsection*{Proper Force and Lorentz Dynamics}
\begin{align}
m U &= m dX/ds,\notag \\
&= m \gamma + m \gamma \mathbf v \cdot \sigv.
\end{align}
\begin{equation}
d/ds = \gamma d/dt.
\end{equation}
\noindent\textit{The time component gives the relativistic work law.}\par
\begin{equation}
d/dt (m \gamma ) = q \mathbf v \cdot \mathbf E.
\end{equation}
\noindent\emph{Note.} Here \(\mathbf E\) denotes the electric field, not the energy operator \(E\).\par
\noindent\emph{Note.} The magnetic field contributes no work term in the scalar component.\par
\noindent\textit{The spatial component gives the relativistic Lorentz force law.}\par
\begin{align}
\mathbf P &:= m \gamma \mathbf v, \\
d\mathbf P/dt &= q \mathbf E + q \mathbf v \times \mathbf B.
\end{align}
\noindent\textit{These terms combine into the single relation.}\par
\begin{equation}
d/dt (m U) = q \mathbf v \cdot \mathbf E + (q \mathbf E + q \mathbf v \times \mathbf B) \cdot \sigv.
\end{equation}
\noindent\textit{Consider the following identity.}\par
\begin{align}
F U / \gamma &= ((\mathbf E + i \mathbf B) \cdot \sigv ) ( 1 + \mathbf v \cdot \sigv ),\notag \\
&= \mathbf E \cdot \sigv + i \mathbf B \cdot \sigv + (\mathbf E \cdot \sigv ) (\mathbf v \cdot \sigv ),\notag \\
&\quad + i (\mathbf B \cdot \sigv )(\mathbf v \cdot \sigv ),\notag \\
&= \mathbf E \cdot \sigv + i \mathbf B \cdot \sigv + \mathbf E \cdot \mathbf v,\notag \\
&\quad + i (\mathbf E \times \mathbf v) \cdot \sigv + i \mathbf B \cdot \mathbf v - (\mathbf B \times \mathbf v) \cdot \sigv.
\end{align}
\begin{align}
\eqtext{real}( F U ) &= \tfrac{1}{2} ( F U + U F^{\mathrm{R}} ),\notag \\
&= \gamma (\mathbf v \cdot \mathbf E + ( \mathbf E + \mathbf v \times \mathbf B ) \cdot \sigv ).
\end{align}
\noindent\emph{Note.} The imaginary part of \(FU\) exchanges the electric and magnetic roles and is therefore not the sector used in the physical force law.\par
\begin{align}
m dU/ds &= q \eqtext{real}( F U ),\notag \\
&= \tfrac{1}{2} q ( F U + U F^{\mathrm{R}} ).
\end{align}
\begin{equation}
m ddX/ds^{2} = q \eqtext{real}( F dX/ds )\qquad \text{4-force}.
\end{equation}

\subsubsection*{Variational Derivation of the Lorentz Force}
\noindent\textit{Recall the previously established definitions.}\par
\begin{align}
\Phi &= V + \mathbf A \cdot \sigv,\notag \\
U &= dX/ds,\notag \\
\gamma &= dt/ds,\notag \\
\text{action} &= -m \int ds - q \int \eqtext{scalar}( \Phi ^{*} U ) ds,\notag \\
&= \int ( -m / \gamma - q \eqtext{scalar}( \Phi ^{*} U / \gamma )) dt,\notag \\
&= \int ( -m / \gamma,\notag \\
&\quad - q \eqtext{scalar}( (V - \mathbf A \cdot \sigv ) (1 + \mathbf v \cdot \sigv )) ) dt,\notag \\
&= \int ( -m / \gamma - q V + q \mathbf v \cdot \mathbf A ) dt.
\end{align}
\begin{equation}
\mathcal{L} = -m / \gamma - q ( V - \mathbf v \cdot \mathbf A )\qquad \text{Lagrangian}.
\end{equation}
\noindent\textit{Verification.}\par
\begin{align}
\delv{} \mathcal{L} &= m \gamma \mathbf v + q \mathbf A,\notag \\
&= \mathbf P + q \mathbf A.
\end{align}
\begin{align}
d/dt \delv{} \mathcal{L} &= d\mathbf P/dt + q \partial _{t} \mathbf A + q (\mathbf v \cdot \delr{}) \mathbf A, \\
\delr{} \mathcal{L} &= - q \delr{} V + q \delr{} (\mathbf v \cdot \mathbf A).
\end{align}
\noindent\textit{Use the following vector-calculus identity.}\par
\begin{equation}
(\mathbf v(t) \cdot \delr{}) \mathbf A(t, \mathbf r) = \delr{} ( \mathbf v \cdot \mathbf A ) - \mathbf v \times ( \delr{} \times \mathbf A ).
\end{equation}
Solving \(d/dt \delv{} \mathcal{L} = \delr{} \mathcal{L}\) for \(d\mathbf P/dt\) gives the following relation.\par
\begin{align}
d\mathbf P/dt &= - q \delr{} V - q \partial _{t} \mathbf A,\notag \\
&\quad + q ( \delr{} (\mathbf v \cdot \mathbf A) - (\mathbf v \cdot \delr{}) \mathbf A ),\notag \\
&= q ( -\delr{}V - \partial _{t}\mathbf A + \mathbf v \times ( \delr{} \times \mathbf A ) ),\notag \\
&= q ( \mathbf E + \mathbf v \times \mathbf B ),\notag \\
&= m \eqtext{vector}( dU/ds ) / \gamma.
\end{align}
\begin{align}
\mathcal{H} &= (\delv{} \mathcal{L}) \cdot \mathbf v - \mathcal{L},\notag \\
&= m \gamma \mathbf v \cdot \mathbf v + q \mathbf v \cdot \mathbf A + m / \gamma.
\end{align}
\noindent\textit{Adding the potential term contributes } \(qV-q\,\mathbf v\cdot\mathbf A\)\textit{.}\par
\begin{align}
&= m \gamma ( \mathbf v^{2} + 1 / \gamma ^{2} ) + q V,\notag \\
&= m \gamma ( \mathbf v^{2} + 1 - \mathbf v^{2} ) + q V,\notag \\
&= m \gamma + q V.
\end{align}
\noindent\textit{Recall the Hamiltonian property.}\par
\begin{align}
dH/dt &= - \partial _{t} \mathcal{L},\notag \\
d\mathcal{H}/dt &= d(m \gamma )/dt + q dV/dt,\notag \\
&= d(m \gamma )/dt + q \partial _{t}V + q \mathbf v \cdot \delr{} V.
\end{align}
\begin{align}
\partial _{t} H &= q \partial _{t} V - q \mathbf v \cdot \partial _{t} \mathbf A,\notag \\
d(m \gamma )/dt &= q \mathbf v \cdot ( -\delr{} V - \partial _{t} \mathbf A ),\notag \\
&= q \mathbf v \cdot \mathbf E.
\end{align}

\subsubsection*{Electromagnetic Energy-Momentum}
\begin{align}
F &= (\mathbf E + i \mathbf B) \cdot \sigv,\notag \\
F^{\mathrm{R}} &= (\mathbf E - i \mathbf B) \cdot \sigv,\notag \\
T &= \tfrac{1}{2} F F^{\mathrm{R}},\notag \\
&= \tfrac{1}{2} (\mathbf E^{2} + \mathbf B^{2}) + (\mathbf E \times \mathbf B) \cdot \sigv,\notag \\
&= T_{s} + T_{\mathbf v} \cdot \sigv.
\end{align}
\begin{equation}
\partial = \partial _{t} + \sigv \cdot \delr{}.
\end{equation}
\noindent\flag{This is a frame-dependent energy-density/Poynting paravector, not by itself the full electromagnetic stress-energy tensor.}\par
\noindent\textit{Recall for any paravectors Z and Z'.}\par
\begin{align}
\partial (ZZ') &= (\partial Z) Z' + Z (\partial Z') + [\sigv , Z] \cdot \delr{} Z',\notag \\
\partial T &= \tfrac{1}{2} \partial (FF^{\mathrm{R}}),\notag \\
&= \tfrac{1}{2}(\partial F) F^{\mathrm{R}} + \tfrac{1}{2}F (\partial F^{\mathrm{R}}) + \tfrac{1}{2}[\sigv , F] \cdot \delr{} F^{\mathrm{R}}.
\end{align}
\begin{align}
D &= \rho + \mathbf J \cdot \sigv\qquad \text{density paravector}, \\
D^{*} &= \rho - \mathbf J \cdot \sigv, \\
\partial &= \partial _{t} + \sigv \cdot \delr{}, \\
\partial F &= D^{*}, \\
F^{2} &= (\mathbf E^{2} - \mathbf B^{2}) + 2 i (\mathbf E \cdot \mathbf B)\qquad \text{scalar}.
\end{align}

\noindent\textit{Recall for any paravectors Z and Z'.}\par
\begin{align}
\partial (ZZ') &= (\partial Z) Z' + Z (\partial Z') + [\sigv , Z] \cdot \delr{} Z',\notag \\
\partial (F^{2}) &= (\partial F) F + F (\partial F) + [\sigv , F] \cdot \delr{} F,\notag \\
&= D^{*} F + F D^{*} + [\sigv , F] \cdot \delr{} F.
\end{align}
\begin{align}
F &= F_{\mathbf v} \cdot \sigv,\notag \\
[\sigv , F] &= 2 i F_{\mathbf v} \times \sigv,\notag \\
[\sigv , F] \cdot \delr{} F &= -2 i F_{\mathbf v} \cdot (\delr{} \times F_{\mathbf v}),\notag \\
&\quad -2 ( (\delr{} \cdot F_{\mathbf v}) F_{\mathbf v} - \tfrac{1}{2} \delr{} F_{\mathbf v}^{2} ) \cdot \sigv.
\end{align}
\begin{equation}
\square F^{2} = \partial ^{*} \partial (F^{2})\qquad \text{scalar}.
\end{equation}


\section{Relativistic Wave Equations and Spin Dynamics}\label{sec:qm}
The standard relativistic wave equations are obtained from the same ordered product calculus. Three logically different carriers are separated: a complex scalar for Klein--Gordon theory, a fixed projected column space for a Dirac field, and the surviving two-entry field in the Pauli limit. By keeping these distinctions explicit, algebraic degrees of freedom are prevented from being mistaken for independent physical wave-function components \cite{klein1926,gordon1926,dirac1928,pauli1927}.
Every spatial integral in this section is evaluated at fixed time. Unless a
region is named, the integration is extended over all real values of the three spatial
components; the measure is written directly as
\(\dd r_1\,\dd r_2\,\dd r_3\).

\subsection{Klein--Gordon field and frequency splitting}
The form to be represented is fixed by the classical mass-shell relation,
while the replacement of classical energy and momentum by differential
operators is supplied as the standard quantization input; it is not obtained from the determinant
identity by algebra alone. For a complex scalar field \(\psi\in\C\), the free
Klein--Gordon equation is
\begin{equation}
(\pv{\partial}\conj{\pv{\partial}}+m^2)\psi=(\Box+m^2)\psi=0.
\label{eq:kg-free}
\end{equation}
On a suitable spatial domain, the positive frequency operator is defined by
\begin{equation}
\Omega:=\sqrt{m^2-\gradv^2}.
\label{eq:frequency-operator}
\end{equation}
Here \(\Omega\) is defined mode by mode: a spatial Fourier mode with wave vector \(\vct{k}\) is multiplied by the positive real number
\(\sqrt{m^2+\vct{k}^2}\). When \(\Omega^{-1}\) exists on the chosen data,
the two frequency functions are defined by
\begin{equation}
\psi(\pm):=\frac12\bigl(\psi\pm i\Omega^{-1}\partial_t\psi\bigr).
\label{eq:frequency-functions}
\end{equation}
They satisfy
\begin{equation}
\begin{gathered}
(\partial_t+i\Omega)\psi(+)=0,\qquad
(\partial_t-i\Omega)\psi(-)=0,\\
\psi=\psi(+)+\psi(-).
\end{gathered}
\label{eq:frequency-split}
\end{equation}
Thus
\[
\psi(\pm,t)=\exp(\mp i\Omega t)\psi(\pm,0).
\]
A massless zero mode is treated separately. The scalar \(|\psi|^2\) is not the relativistic Klein--Gordon probability density; derivatives are contained in the conserved current, which is not positive on the full positive/negative-frequency solution space.
Frequency is labeled in this subsection by the sign in \(\psi(\pm)\); that
input is unrelated to the plus/minus input of the spin-axis projectors in
Eq.\ \eqref{eq:projectors}.

\subsection{Minimal coupling and ordered factorization}
Under the standard minimal-coupling rule, the free operators are replaced by
the following covariant operators. With the real potential
\(\pv{\Phi}=V+\vct{A}\cdot\sigv\), the energy operator is defined by
\begin{equation}
E:=i\partial_t-qV.
\label{eq:energy-operator}
\end{equation}
The momentum operator is defined by
\begin{equation}
\vct{P}:=-i\gradv-q\vct{A}.
\label{eq:momentum-operator}
\end{equation}
The minimally coupled scalar equation is
\begin{equation}
(-E^2+\vct{P}^2+m^2)\psi=0,
\label{eq:kg-coupled}
\end{equation}
or, after applying the Leibniz rule,
\begin{equation}
\Bigl(\Box+iqS-q^2\detf{\pv{\Phi}}+m^2
 +2iq(V\partial_t+\vct{A}\cdot\gradv)\Bigr)\psi=0.
\label{eq:kg-expanded}
\end{equation}
Gauge covariance of the scalar equation is also obtained from the same phase
law used for the Dirac field. Under Eq.\ \eqref{eq:gauge}, the transformed
scalar field is defined by
\begin{equation}
\psi(\lambda):=\exp(-iq\lambda)\psi.
\label{eq:scalar-gauge-field}
\end{equation}
The gauge-dependent energy and momentum operators are defined by
\begin{equation}
E(\lambda):=i\partial_t-qV(\lambda).
\label{eq:gauge-energy-operator}
\end{equation}
\begin{equation}
\vct{P}(\lambda):=-i\gradv-q\vct{A}(\lambda).
\label{eq:gauge-momentum-operator}
\end{equation}
The following relations are obtained by direct application of the product rule:
\begin{equation}
E(\lambda)\psi(\lambda)
=\exp(-iq\lambda)E\psi,\qquad
\vct{P}(\lambda)\psi(\lambda)
=\exp(-iq\lambda)\vct{P}\psi,
\label{eq:kg-gauge}
\end{equation}
so the same form of Eq.\ \eqref{eq:kg-coupled} is retained in every gauge.
For smooth fields,
\begin{align}
\com{E}{P_m}&=iqE_m,\notag\\
\com{P_m}{P_n}&=iq(\partial_mA_n-\partial_nA_m).
\label{eq:commutators}
\end{align}
Consequently, the written order must be retained:
\begin{align}
(E+\sigv\cdot\vct{P})(E-\sigv\cdot\vct{P})
 &=E^2-\vct{P}^2-iq\pv{F},
\label{eq:EP-factorization}\\
(E-\sigv\cdot\vct{P})(E+\sigv\cdot\vct{P})
 &=E^2-\vct{P}^2-iq\conj{\pv{F}}.
\label{eq:EP-factorization-conjugate}
\end{align}
Here
\[
\begin{aligned}
\conj{\pv{F}}&=(-\vct{E}+i\vct{B})\cdot\sigv,\\
-iq\pv{F}&=(-iq\vct{E}+q\vct{B})\cdot\sigv,\\
-iq\conj{\pv{F}}&=(iq\vct{E}+q\vct{B})\cdot\sigv.
\end{aligned}
\]
These are ordered star--\(H\) products of differential operators, not
applications of the scalar-coefficient determinant formula. In a background
field the full products are not scalar factorizations of the Klein--Gordon
operator: the displayed \(\pv{F}\) or \(\conj{\pv{F}}\) term is retained.
The scalar equation is recovered by taking the scalar part or by subtracting
that residual term explicitly. The same residuals are converted into the
established spin--field terms when the Dirac system is squared; they are not additional
scalar-wave components.

\subsection{Klein--Gordon current and pairing}
For two scalar solutions \(\phi\) and \(\psi\) of
Eq.\ \eqref{eq:kg-coupled} in the same real potential, the mixed density
function is defined by
\begin{equation}
\rho(\phi,\psi)
:=\frac{\phi^*E\psi+(E\phi)^*\psi}{2m}.
\label{eq:kg-mixed-current}
\end{equation}
The associated spatial-current function is defined by
\begin{equation}
\vct{j}(\phi,\psi)
:=\frac{\phi^*\vct{P}\psi+(\vct{P}\phi)^*\psi}{2m}.
\label{eq:kg-mixed-spatial-current}
\end{equation}
Star here is ordinary complex conjugation because both fields are scalars.
Gauge invariance of this mixed current and its pairing is established
immediately by Eq.\ \eqref{eq:kg-gauge} when the same phase is acquired by both solutions.
The local identity is obtained from the product rules:
\begin{align}
\partial_t\rho(\phi,\psi)
+\gradv\cdot\vct{j}(\phi,\psi)
&=\frac{i}{2m}\bigl((E^2-\vct{P}^2)\phi\bigr)^*\psi\notag\\
&\quad-\frac{i}{2m}\phi^*(E^2-\vct{P}^2)\psi\notag\\
&=\frac{i}{2m}(m^2\phi^*\psi-\phi^*m^2\psi)=0.
\label{eq:kg-continuity-proof}
\end{align}
The conserved Klein--Gordon current is obtained in the equal-field case; its
electric current is obtained by multiplication by \(q\). The mixed solution
pairing is defined by its spatial integral:
\begin{equation}
\langle\phi\mid\psi\rangle
:=\iiint\rho(\phi,\psi)\,\dd r_1\,\dd r_2\,\dd r_3.
\label{eq:kg-pairing}
\end{equation}
For free positive-frequency solutions, the pairing is
\[
\langle\phi\mid\psi\rangle
=\frac1m\iiint\phi^*\Omega\psi\,\dd r_1\,\dd r_2\,\dd r_3,
\]
which is positive for equal nonzero fields. The local density can nevertheless
be negative at some points in a positive-frequency superposition. Therefore
a positive local position probability is not supplied by
\(\rho(\psi,\psi)\) in the scalar theory; on the full positive- and
negative-frequency solution space, even the integrated pairing is indefinite.

\subsection{Projected carrier for one Dirac field}
In the Pauli realization, a paravector is represented by an arbitrary two-by-two complex matrix. Its two columns are acted upon independently by left multiplication. Therefore two degenerate Dirac fields, rather than one, are described by two unrestricted paravector blocks. A single Dirac field is obtained by selecting one matrix column. The fixed-column function is defined intrinsically by
\begin{equation}
\pv{\Pi}(\eqtext{column}):=\frac12(1+\pv{\sigma}_3).
\label{eq:fixed-column-carrier}
\end{equation}
One Dirac field is selected by the constraint
\begin{equation}
\bPsi_n=\bPsi_n\pv{\Pi}(\eqtext{column}),
\qquad n\in\{1,2\}.
\label{eq:fixed-column-constraint}
\end{equation}
Every operator below is applied from the left, so the restriction is preserved:
\((\pv{Z}\bPsi_n)\pv{\Pi}(\eqtext{column})
=\pv{Z}(\bPsi_n\pv{\Pi}(\eqtext{column}))=\pv{Z}\bPsi_n\).
Selection of one matrix column by this fixed right factor is denoted by the
input \(\eqtext{column}\). No measurement of spin along the third axis is
represented by it. In the matrix realization, each \(\bPsi_n\) is therefore
represented by one ordinary two-entry complex column, and the usual four
complex entries are contained in the block column
\(\bPsi=\{\bPsi_1;\bPsi_2\}\).
The apparent scalar and three-vector coefficients are correspondingly
constrained. If
\(\vct{\psi}_{\vct{v}}:=\{\psi_1,\psi_2,\psi_3\}\) and
\(\pv{\psi}=\psi_s+\vct{\psi}_{\vct{v}}\cdot\sigv\), then
\[
\begin{aligned}
\pv{\psi}=\pv{\psi}\pv{\Pi}(\eqtext{column})
&\quad\eqtext{if and only if}\quad
\psi_3=\psi_s,\qquad \psi_2=-i\psi_1,\\
\mat{M}&=\{2\psi_s,0;2\psi_1,0\}.
\end{aligned}
\]
Thus exactly the two complex amplitudes in the selected column are contained
in a projected block, not four independent paravector amplitudes.

For any projected block \(\pv{\psi}\), the positive density function is
defined by
\begin{equation}
\rho(\pv{\psi})
:=2\,\eqtext{scalar}(\herm{\pv{\psi}}\pv{\psi}).
\label{eq:projected-block-density-function}
\end{equation}
The factor \(2\) converts scalar extraction into the ordinary norm of the
selected two-entry column.

Using the block map \(\Wmat\) from Eq.\ \eqref{eq:W-map}, the free and coupled
Dirac equations are
\begin{align}
(\Wmat(i\pv{\partial})-m\I)\bPsi&=0,
\label{eq:dirac-free}\\
(\Wmat(i\pv{\partial}-q\conj{\pv{\Phi}})-m\I)\bPsi&=0.
\label{eq:dirac-coupled}
\end{align}
The coupled system is equivalently
\begin{equation}
(E-\sigv\cdot\vct{P})\bPsi_2=m\bPsi_1,\qquad
(E+\sigv\cdot\vct{P})\bPsi_1=m\bPsi_2.
\label{eq:dirac-coupled-blocks}
\end{equation}
In the matrix realization, the selected columns are represented directly as
\[
\begin{aligned}
\bPsi_1&\quad\eqtext{is represented by}\quad
\begin{Bmatrix}a\\b\end{Bmatrix},\\
\bPsi_2&\quad\eqtext{is represented by}\quad
\begin{Bmatrix}c\\d\end{Bmatrix}.
\end{aligned}
\]
Equation \eqref{eq:dirac-coupled-blocks} is then exactly
\begin{equation}
\begin{aligned}
(E-P_3)c-(P_1-iP_2)d&=ma,\\
-(P_1+iP_2)c+(E+P_3)d&=mb,\\
(E+P_3)a+(P_1-iP_2)b&=mc,\\
(P_1+iP_2)a+(E-P_3)b&=md.
\end{aligned}
\label{eq:dirac-four-components}
\end{equation}
These are the four conventional first-order Dirac component equations in
the two-block form. The free system is obtained by setting \(q=0\).
For a constant determinant-one transform \(\pv{T}\), the two blocks are
fixed differently by covariance:
\begin{equation}
\bPsi_1'(\pv{X}')=\pv{T}^{\mathsf H}\bPsi_1(\pv{X}),
\label{eq:dirac-transform}
\end{equation}
\begin{equation}
\bPsi_2'(\pv{X}')=\pv{T}^{-1}\bPsi_2(\pv{X}).
\label{eq:dirac-transform-lower}
\end{equation}
As shown by the supporting derivation of
Eqs.\ \eqref{eq:dirac-transform} and \eqref{eq:dirac-transform-lower}, the
first block is therefore acted upon by \(\pv{L}\) under a boost, whereas the
second is acted upon by \(\conj{\pv{L}}\); both are acted upon by
\(\pv{R}^{\mathsf H}\) under a rotation. This is the wave-function law, distinct from
both event congruence and field similarity.
The mass-conjugate block operator is defined by
\begin{equation}
\Dmat':=\Wmat(i\pv{\partial}-q\conj{\pv{\Phi}})+m\I.
\label{eq:dirac-mass-conjugate-operator}
\end{equation}
By multiplication by this operator and use of the two ordered factorizations,
the squared system is obtained:
\begin{equation}
\begin{cases}
(-E^2+\vct{P}^2+m^2+iq\conj{\pv{F}})\bPsi_1=0,\\[2pt]
(-E^2+\vct{P}^2+m^2+iq\pv{F})\bPsi_2=0.
\end{cases}
\label{eq:dirac-square}
\end{equation}
With
\(iq\conj{\pv{F}}=(-iq\vct{E}-q\vct{B})\cdot\sigv\) and
\(iq\pv{F}=(iq\vct{E}-q\vct{B})\cdot\sigv\), the same equations are written as
\begin{equation}
\begin{cases}
\bigl(-E^2+\vct{P}^2+m^2
-q\vct{B}\cdot\sigv-iq\vct{E}\cdot\sigv\bigr)\bPsi_1=0,\\[2pt]
\bigl(-E^2+\vct{P}^2+m^2
-q\vct{B}\cdot\sigv+iq\vct{E}\cdot\sigv\bigr)\bPsi_2=0.
\end{cases}
\label{eq:dirac-square-components}
\end{equation}
Thus the expected Klein--Gordon operator, the same magnetic spin coupling in
both blocks, and the opposite electric spin couplings of the two blocks are
obtained component by component by squaring the first-order system.

Under the gauge law \eqref{eq:gauge},
\begin{equation}
\bPsi(\lambda):=\exp(-iq\lambda)\bPsi.
\label{eq:dirac-gauge}
\end{equation}
The form of Eq.\ \eqref{eq:dirac-coupled} is unchanged.

A positive time-slice norm is supplied by the Hermitian product, but its
vector part is not a probability current. With the fixed-column convention
above, the Dirac density function is defined by
\begin{equation}
\rho(\bPsi):=2\,\eqtext{scalar}\!\left(
 \herm{\bPsi_1}\bPsi_1+\herm{\bPsi_2}\bPsi_2\right).
\label{eq:dirac-density}
\end{equation}
The spatial-current function is defined by
\begin{equation}
\vct{j}(\bPsi):=2\,\eqtext{vector}\!\left(
 \bPsi_2\herm{\bPsi_2}-\bPsi_1\herm{\bPsi_1}\right).
\label{eq:dirac-spatial-current}
\end{equation}
The full current paravector is defined by
\begin{equation}
\pv{J}(\bPsi):=\rho(\bPsi)+\vct{j}(\bPsi)\cdot\sigv.
\label{eq:dirac-current-paravector}
\end{equation}
The two definitions satisfy the continuity equation
\begin{equation}
\partial_t\rho(\bPsi)+\gradv\cdot\vct{j}(\bPsi)=0.
\label{eq:dirac-current}
\end{equation}
The normalization condition is the spatial integral
\[
\iiint\rho(\bPsi)\,\dd r_1\,\dd r_2\,\dd r_3=1,
\]
not a pointwise equation. By
Eqs.\ \eqref{eq:matrix-operation-dictionary} and
\eqref{eq:hermitian-column-pairing}, \(H\) is the ordinary conjugate
transpose and \(\rho(\bPsi)\) is the sum of the two conventional column norm
squares. The familiar Dirac spatial current is supplied by the displayed
opposite-order bilinears, rather than by the vector part of a single Hermitian
norm product.
The real current paravector from Eq.\ \eqref{eq:dirac-current-paravector}
obeys the event congruence, as shown in the supporting derivation.
The component balances are obtained directly from the two time-evolution equations:
\begin{align}
\partial_t(|a|^2+|b|^2)
&=\gradv\cdot\left(
\begin{Bmatrix}a^*&b^*\end{Bmatrix}\sigv
\begin{Bmatrix}a\\b\end{Bmatrix}\right)\notag\\
&\quad+im\bigl(c^*a+d^*b-a^*c-b^*d\bigr),\notag\\
\partial_t(|c|^2+|d|^2)
&=-\gradv\cdot\left(
\begin{Bmatrix}c^*&d^*\end{Bmatrix}\sigv
\begin{Bmatrix}c\\d\end{Bmatrix}\right)\notag\\
&\quad+im\bigl(a^*c+b^*d-c^*a-d^*b\bigr).
\label{eq:dirac-continuity-components}
\end{align}
By addition of the lines, the mass exchange is canceled and the following expressions are obtained exactly:
\begin{align*}
\rho(\bPsi)&=|a|^2+|b|^2+|c|^2+|d|^2,\\
\vct{j}(\bPsi)
&=\begin{Bmatrix}c^*&d^*\end{Bmatrix}\sigv
  \begin{Bmatrix}c\\d\end{Bmatrix}
 -\begin{Bmatrix}a^*&b^*\end{Bmatrix}\sigv
  \begin{Bmatrix}a\\b\end{Bmatrix},
\end{align*}
and hence the continuity equation in Eq.\ \eqref{eq:dirac-current}. Neither
block norm is conserved separately: norm is exchanged between them by the
mass terms.

\subsection{Probabilities and measurement expectations}
For two projected one-block fields, the bracket is defined by
\begin{equation}
\langle\pv{\phi}\mid\pv{\psi}\rangle
:=2\iiint\eqtext{scalar}(\herm{\pv{\phi}}\pv{\psi})\,
\dd r_1\,\dd r_2\,\dd r_3.
\label{eq:adjoint-paravector-braket}
\end{equation}
For two projected two-block fields, the complete Dirac pairing is defined by
\begin{equation}
\begin{aligned}
&\left\langle\bPsi(\eqtext{left})\mid
\bPsi(\eqtext{right})\right\rangle\\
&\quad:=2\iiint\sum_{n=1}^{2}
\eqtext{scalar}\!\left(
 \herm{\bPsi_n(\eqtext{left})}\bPsi_n(\eqtext{right})
\right)\,\dd r_1\,\dd r_2\,\dd r_3.
\end{aligned}
\label{eq:dirac-pairing}
\end{equation}
The two inputs are identified by the words \(\eqtext{left}\) and
\(\eqtext{right}\); no new kind of field is introduced. When both inputs are the same,
the shorter notation \(\langle\bPsi\mid\bPsi\rangle\) is used.
The conversion of intrinsic scalar extraction into the ordinary
selected-column product is supplied by the factor \(2\). For a specified
block operator \(\mat{\mathcal{O}}\), the operator pairing is defined by
\begin{equation}
\begin{aligned}
&\left\langle\bPsi(\eqtext{left})\mid
\mat{\mathcal{O}}\bPsi(\eqtext{right})\right\rangle\\
&:=2\iiint\sum_{n=1}^{2}
\eqtext{scalar}\!\left(
\herm{\bPsi_n(\eqtext{left})}
\bigl(\mat{\mathcal{O}}\bPsi(\eqtext{right})\bigr)_n
\right)\,\dd r_1\,\dd r_2\,\dd r_3.
\end{aligned}
\label{eq:dirac-operator-pairing}
\end{equation}
Its normalized matrix element is
\begin{equation}
\frac{\langle\bPsi\mid\mat{\mathcal{O}}\bPsi\rangle}
{\langle\bPsi\mid\bPsi\rangle}.
\label{eq:dirac-expectation}
\end{equation}
It is a real expectation when \(\mat{\mathcal{O}}\) is Hermitian under this
spatial pairing. For multiplication operators, Hermitian
diagonal entries and
\(\pv{\mathcal{O}}_{21}=\herm{\pv{\mathcal{O}}_{12}}\) are required for the
diagonal and off-diagonal entries, respectively; for derivatives, the stated
integration-by-parts boundary conditions and a common operator domain are
required. If a Hermitian operator is also idempotent, the
probability of the outcome selected by that projector is given by the same
ratio.

For a normalized Dirac field, the probability of finding the system in a
specified spatial region at the chosen time is defined by
\begin{equation}
p(\eqtext{region})
:=\iiint\rho(\bPsi)\,\dd r_1\,\dd r_2\,\dd r_3,
\label{eq:region-probability}
\end{equation}
where the integration is restricted to the specified region.
The density is nonnegative, but it is the time component of the covariant
current, not a Lorentz scalar. The normalization is kept time-independent by
conservation when the boundary flux vanishes.

For a projected Pauli block \(\pv{\psi}\), the factor \(2\) is canceled from
normalized ratios. The probability function supplied by the axis projectors
is therefore defined by
\begin{equation}
p(\pm)
:=\frac{\langle\pv{\psi}(\pm)\mid\pv{\psi}(\pm)\rangle}
{\langle\pv{\psi}\mid\pv{\psi}\rangle}.
\label{eq:pauli-projector-probability}
\end{equation}
By Eqs.\ \eqref{eq:projectors} and \eqref{eq:projected-part-function}, this
can equally be written with
\(\pv{\psi}(\pm)=\pv{\Pi}(\pm)\pv{\psi}\), and
\begin{equation}
p(+)+p(-)=1.
\label{eq:pauli-probability-sum}
\end{equation}
The expectation of \(\vct{u}\cdot\sigv\) is \(p(+)-p(-)\); the physical
spin component along \(\vct{u}\) is one half of this value in the natural
units used here. The corresponding free two-block energy and momentum-axis
probability is defined after those fields have been constructed below.

The Klein--Gordon carrier is instead paired by Eq.\ \eqref{eq:kg-pairing}.
On a chosen positive-frequency solution sector, the normalized matrix element
of an operator \(\mathcal{O}\) is
\[
\frac{\langle\psi\mid\mathcal{O}\psi\rangle}
{\langle\psi\mid\psi\rangle}
\]
only when that sector is preserved and the operator is Hermitian under the
Klein--Gordon pairing. The density \(\rho(\psi,\psi)\) is not thereby
converted into a pointwise position probability.
Likewise, nonscalar channels of a generic paravector sandwich are algebraic
bilinears, not additional measured quantities unless they are identified by
a field equation and a measurement rule.

\subsection{Pauli limit}
A single block by itself is not the nonrelativistic large component. First,
the slowly varying blocks are defined by removal of the rest phase:
\begin{equation}
\bPsi_n(\eqtext{slow}):=\exp(imt)\bPsi_n.
\label{eq:slow-field-definition}
\end{equation}
The large component is then defined by
\begin{equation}
\pv{\phi}(\eqtext{large})
:=\frac{\bPsi_1(\eqtext{slow})+\bPsi_2(\eqtext{slow})}{\sqrt{2}}.
\label{eq:large-small}
\end{equation}
The small component is defined independently by
\begin{equation}
\pv{\phi}(\eqtext{small})
:=\frac{\bPsi_2(\eqtext{slow})-\bPsi_1(\eqtext{slow})}{\sqrt{2}}.
\label{eq:small-component}
\end{equation}
Equation \eqref{eq:dirac-coupled-blocks} is then transformed exactly into
\begin{equation}
\begin{aligned}
E\pv{\phi}(\eqtext{large})
&=(\vct{P}\cdot\sigv)\pv{\phi}(\eqtext{small}),\\
(E+2m)\pv{\phi}(\eqtext{small})
&=(\vct{P}\cdot\sigv)\pv{\phi}(\eqtext{large}).
\end{aligned}
\label{eq:large-small-coupled}
\end{equation}
At energies small compared with \(m\),
\[
\pv{\phi}(\eqtext{small})
=\frac{\vct{P}\cdot\sigv}{2m}\pv{\phi}(\eqtext{large})+O(m^{-2}).
\]
By use of
\[
(\vct{P}\cdot\sigv)^2
=\vct{P}^2-q\,\vct{B}\cdot\sigv,
\]
the following result is obtained:
\begin{equation}
E\pv{\phi}(\eqtext{large})
=\frac{\vct{P}^2-q\,\vct{B}\cdot\sigv}{2m}\,\pv{\phi}(\eqtext{large})
+O(m^{-2}).
\label{eq:pauli-operator}
\end{equation}
Equivalently,
\begin{equation}
\begin{aligned}
i\partial_t\pv{\phi}(\eqtext{large})
&=\left(\frac{(-i\gradv-q\vct{A})^2}{2m}+qV\right)
\pv{\phi}(\eqtext{large})\\
&\quad-\frac{q}{2m}(\vct{B}\cdot\sigv)
\pv{\phi}(\eqtext{large})+O(m^{-2}).
\end{aligned}
\label{eq:pauli}
\end{equation}
The spin--orbit, Darwin, and time-dependent corrections are generated at the next order by electric-field commutators; no term of that order is retained selectively here.
The selected two-entry column representing the large component is defined by
\begin{equation}
\vct{\phi}(\eqtext{large})
:=\eqtext{selected column representing }\pv{\phi}(\eqtext{large}).
\label{eq:large-column-definition}
\end{equation}
The corresponding small-component column is defined by
\begin{equation}
\vct{\phi}(\eqtext{small})
:=\eqtext{selected column representing }\pv{\phi}(\eqtext{small}).
\label{eq:small-column-definition}
\end{equation}
The following expressions are obtained by the exact change of basis:
\begin{align*}
\rho(\bPsi)
&=\mherm{\vct{\phi}(\eqtext{large})}
  \vct{\phi}(\eqtext{large})
+\mherm{\vct{\phi}(\eqtext{small})}
  \vct{\phi}(\eqtext{small}),\\
\vct{j}(\bPsi)
&=\mherm{\vct{\phi}(\eqtext{large})}\sigv
  \vct{\phi}(\eqtext{small})\\
&\quad+\mherm{\vct{\phi}(\eqtext{small})}\sigv
  \vct{\phi}(\eqtext{large}).
\end{align*}
The projected-block density from
Eq.\ \eqref{eq:projected-block-density-function} is represented here by
\begin{equation}
\begin{aligned}
\rho(\pv{\phi}(\eqtext{large}))
&=\mherm{\vct{\phi}(\eqtext{large})}
  \vct{\phi}(\eqtext{large})\\
&=2\,\eqtext{scalar}\bigl(
 \herm{\pv{\phi}(\eqtext{large})}\pv{\phi}(\eqtext{large})\bigr).
\end{aligned}
\label{eq:pauli-density}
\end{equation}
The Pauli spatial-current function is defined by
\begin{align}
\vct{j}(\pv{\phi}(\eqtext{large}))
&:=\frac{\mherm{\vct{\phi}(\eqtext{large})}
  \vct{P}\vct{\phi}(\eqtext{large})}{2m}\notag\\
&\quad+\frac{\mherm{(\vct{P}\vct{\phi}(\eqtext{large}))}
  \vct{\phi}(\eqtext{large})}{2m}\notag\\
&\quad+\frac{1}{2m}\gradv\times
 \bigl(\mherm{\vct{\phi}(\eqtext{large})}\sigv
 \vct{\phi}(\eqtext{large})\bigr).
\label{eq:pauli-current}
\end{align}
Because \(\vct{\phi}(\eqtext{small})=O(m^{-1})\), the original Dirac
density and current are reduced to
\(\rho(\bPsi)=\rho(\pv{\phi}(\eqtext{large}))+O(m^{-2})\) and
\(\vct{j}(\bPsi)=\vct{j}(\pv{\phi}(\eqtext{large}))+O(m^{-2})\).
The first numerator is read component by component in the three entries of
\(\vct{P}\).
The standard convection current plus the divergence-free magnetization
current is thereby obtained. The equation
\(\partial_t\rho(\pv{\phi}(\eqtext{large}))
+\gradv\cdot\vct{j}(\pv{\phi}(\eqtext{large}))=0\)
is obeyed through the retained order.

If \(\vct{B}\ne0\) is uniform and constant, the magnetic-field axis is
defined by
\begin{equation}
\vct{u}:=\frac{\vct{B}}{\|\vct{B}\|}.
\label{eq:magnetic-axis}
\end{equation}
Equation \eqref{eq:projectors} is then specialized as
\(\pv{\Pi}(\pm)=\tfrac12(1\pm\vct{u}\cdot\sigv)\).
The two spin-axis parts are defined by
\begin{equation}
\pv{\phi}(\eqtext{large},\pm)
:=\pv{\Pi}(\pm)\pv{\phi}(\eqtext{large}).
\label{eq:pauli-spin-decomposition}
\end{equation}
They reconstruct the large component:
\begin{equation}
\pv{\phi}(\eqtext{large})
=\pv{\phi}(\eqtext{large},+)+\pv{\phi}(\eqtext{large},-).
\label{eq:pauli-spin-reconstruction}
\end{equation}
Then
\begin{align}
i\partial_t\pv{\phi}(\eqtext{large},\pm)
&=\left(\frac{\vct{P}^2}{2m}+qV\right)
\pv{\phi}(\eqtext{large},\pm)\notag\\
&\quad\mp\frac{q\|\vct{B}\|}{2m}
\pv{\phi}(\eqtext{large},\pm).
\label{eq:pauli-spin-split}
\end{align}
For a varying axis, the two projected fields are coupled by derivatives of
\(\pv{\Pi}(\pm)\) and its commutator with the time-evolution operator on the
right of Eq.\ \eqref{eq:pauli}, so this decoupling is lost.

\subsection{Free energy and momentum-axis split}
For a free momentum mode, the positive energy function is defined by
\begin{equation}
E(\vct{p}):=\sqrt{m^2+\vct{p}^2}.
\label{eq:free-energy-function}
\end{equation}
The block Hamiltonian is defined by
\begin{equation}
\mat{\mathcal H}:=
\begin{pmatrix}
-\sigv\cdot\vct{p}&m\\
m&\sigv\cdot\vct{p}
\end{pmatrix}.
\label{eq:dirac-hamiltonian}
\end{equation}
Equation \eqref{eq:dirac-coupled-blocks} then gives
\begin{equation}
\begin{aligned}
i\partial_t\bPsi&=\mat{\mathcal H}\bPsi,\\
\mat{\mathcal H}\bPsi
&=
\begin{Bmatrix}
-(\sigv\cdot\vct{p})\bPsi_1+m\bPsi_2\\
m\bPsi_1+(\sigv\cdot\vct{p})\bPsi_2
\end{Bmatrix},\\
\mat{\mathcal H}^{\,2}&=E(\vct{p})^2\I.
\end{aligned}
\label{eq:dirac-hamiltonian-action}
\end{equation}
The block matrix \(\mat{\mathcal H}\) is called the Hamiltonian because
\(i\partial_t\bPsi\) is supplied by exactly this operator. It is not a new
field: the two free Dirac equations are given by its two displayed rows.
The calligraphic symbol \(\mat{\mathcal H}\) is distinct from the superscript
\(H\), which is reserved for the Hermitian operation.

The energy sign \(e\) is taken from \(\{+1,-1\}\). The corresponding part of the
field is defined directly, without naming another operator:
\begin{equation}
\bPsi(e)
:=\frac12\left(\bPsi+e\frac{\mat{\mathcal H}\bPsi}{E(\vct{p})}\right).
\label{eq:energy-projected-field}
\end{equation}
Its two blocks are displayed explicitly as
\begin{equation}
\bPsi(e)
=\frac12
\begin{Bmatrix}
\bPsi_1+\dfrac{e}{E(\vct{p})}
 \bigl(-(\sigv\cdot\vct{p})\bPsi_1+m\bPsi_2\bigr)\\[3pt]
\bPsi_2+\dfrac{e}{E(\vct{p})}
 \bigl(m\bPsi_1+(\sigv\cdot\vct{p})\bPsi_2\bigr)
\end{Bmatrix}.
\label{eq:energy-projected-components}
\end{equation}
The eigenvalue and reconstruction identities are
\begin{equation}
\begin{aligned}
\mat{\mathcal H}\bPsi(e)&=eE(\vct{p})\bPsi(e),
\qquad
\bPsi=\bPsi(+)+\bPsi(-).
\end{aligned}
\label{eq:energy-projection-properties}
\end{equation}
Thus the two choices of \(e\) are the positive- and negative-energy parts of
the same free field. The two blocks displayed in \(\bPsi(e)\) are denoted
below by \(\bPsi_1(e)\) and \(\bPsi_2(e)\).

For \(\vct{p}\ne0\), the momentum direction is defined by
\begin{equation}
\vct{u}:=\frac{\vct{p}}{\|\vct{p}\|}.
\label{eq:momentum-axis}
\end{equation}
For \(s\in\{+1,-1\}\), the existing axis projector from
Eq.\ \eqref{eq:projectors} is then
\begin{equation}
\pv{\Pi}(s)
=\frac12(1+s\vct{u}\cdot\sigv)
=\frac12\left(1+s\frac{\vct{p}\cdot\sigv}{\|\vct{p}\|}\right).
\label{eq:momentum-axis-projector}
\end{equation}
The simultaneous energy and momentum-axis part is defined by
\begin{equation}
\bPsi(e,s)
:=
\begin{Bmatrix}
\pv{\Pi}(s)\bPsi_1(e)\\
\pv{\Pi}(s)\bPsi_2(e)
\end{Bmatrix}.
\label{eq:energy-spin-projectors}
\end{equation}
It satisfies
\begin{equation}
\begin{aligned}
(\vct{p}\cdot\sigv)\pv{\Pi}(s)
&=s\|\vct{p}\|\pv{\Pi}(s),\\
\bPsi&=\sum_{e\in\{+1,-1\}}\sum_{s\in\{+1,-1\}}\bPsi(e,s).
\end{aligned}
\label{eq:energy-spin-reconstruction}
\end{equation}
Here \(s\) is the eigenvalue of
\(\vct{u}\cdot\sigv\). In the natural units used here, the physical spin
component along momentum is \(s/2\); its sign is commonly called helicity.
The Hamiltonian is Hermitian, and the two energy parts are orthogonal. The
already established \(\pv{\Pi}(s)\) projectors are Hermitian,
idempotent, and complementary. Because the momentum axis is used, commutation
of their left action with \(\mat{\mathcal H}\) is obtained, so the four
displayed fields are simultaneous energy and momentum-axis parts. At rest any spin axis may
be chosen. Negative-frequency classical solutions are not identified with
antiparticles until the field is quantized.

The joint probability function is now defined without referring to a field
before its construction:
\begin{equation}
p(e,s):=
\frac{\langle\bPsi(e,s)\mid\bPsi(e,s)\rangle}
{\langle\bPsi\mid\bPsi\rangle}.
\label{eq:dirac-projector-probability}
\end{equation}
Completeness gives
\begin{equation}
\sum_{e\in\{+1,-1\}}\sum_{s\in\{+1,-1\}}p(e,s)=1.
\label{eq:dirac-probability-sum}
\end{equation}
For a normalizable free field this split is applied to each momentum mode
before the spatial pairing is evaluated. For a single plane mode, the usual
box or delta-function normalization is required instead.

\subsection{Supporting Wave-Equation Derivations}
The ordered operator expansions, scalar frequency evolution, projected-carrier interpretation, conserved pairings, gauge covariance, squared Dirac system, and controlled Pauli approximation are given in detail in the supporting derivations.
\subsubsection{Relativistic Quantum-Mechanical Setting}
The wave-equation section uses one fixed operator language throughout. With the differential conventions from Section \ref{sec:differential},
\begin{align*}
\partial&=\partial_t+\sigv\cdot\gradv,\\
\partial^*&=\partial_t-\sigv\cdot\gradv,\\
(i\partial)^*&=-i\partial^*.
\end{align*}
Recall also that the d'Alembertian is
\begin{align*}
\Box
&:=\partial_t^2-\gradv^2\notag\\
&=\partial\partial^*=\partial^*\partial\notag\\
&=\langle\!\langle\partial\rangle\!\rangle^2.
\end{align*}
All differential operators act left to right.

\subsubsection{What a Complex-Paravector State Represents}
Throughout the spinful part of the wave-equation section, the state is not a space-time event but an algebra-valued amplitude field:
\begin{equation*}
\psi:\R\times\R^3\to\C\oplus\C^3,
\qquad
\psi(t,\mathbf r)=\psi_s(t,\mathbf r)+\psi_{\mathbf v}(t,\mathbf r)\cdot\sigv.
\end{equation*}
The argument \((t,\mathbf r)\) tells where the state is evaluated in space-time, while the value \(\psi(t,\mathbf r)\) lives in the internal complex-paravector algebra. This distinction matters. The real paravector
\begin{equation*}
X=t+\mathbf r\cdot\sigv
\end{equation*}
describes an event or interval variable, whereas \(\psi\) is a quantum state amplitude defined \emph{over} such events.

At each fixed point \((t,\mathbf r)\), the field carries one complex scalar channel \(\psi_s\) and one complex vector channel \(\psi_{\mathbf v}\). Equivalently, it carries four complex local amplitudes, i.e. eight real local degrees of freedom. This is why the complex-paravector form can play the role usually assigned to a spinor carrier: it stores the same amount of local data, but arranges it as a scalar-plus-vector object adapted to the sigma algebra rather than as a column with separately named components.

The scalar and vector pieces should not be interpreted as separately measurable particles or classical subfields. They are amplitude channels that mix under the wave operators. The scalar part is the sigma-number-zero channel, while the vector part is the sigma-number-one channel that couples directly to \(\sigv\cdot\gradv\), to spin projectors, and to the first-order Dirac factorization. In that sense the vector part of \(\psi\) is not a literal velocity field; it is the part of the quantum amplitude that carries directional and spin-sensitive information.

The quantum-mechanical interpretation is therefore bilinear rather than componentwise. The raw field \(\psi\) is not itself an observable. Observable densities are extracted from sandwiches such as
\begin{equation*}
\rev{\psi}\psi
=\rho_\psi+\mathbf J_\psi\cdot\sigv,
\qquad
\rev{\psi}O\psi.
\end{equation*}
The scalar part gives the ordinary probability or expectation density, while the vector part records the associated flow, transport, or spin-orientation information. Thus the extra vector channels are not discarded because they are meaningless; they are discarded only when one asks for a scalar expectation value rather than for a current-like or polarization-like quantity.

This viewpoint also explains the later division of labor inside the section. The spin-0 Klein--Gordon subsections use \(\psi\in\C\) because a scalar carrier is sufficient there. Once one passes to spin-\(\tfrac12\), axis projectors, or first-order Dirac factorization, the complex-paravector carrier becomes the natural one because the sigma algebra must act nontrivially on the state itself. Later, when the manuscript writes a two-component field
\begin{equation*}
\bPsi=
\begin{Bmatrix}
\bPsi_1\\[2pt]
\bPsi_2
\end{Bmatrix},
\end{equation*}
it is packaging algebra-valued amplitudes into the auxiliary \(W\)-map factorization of the Dirac system; that notation is a structural reorganization of the same state content, not a redefinition of what the state represents.

\subsubsection{Complex-Paravector Bilinears and Bra-Ket Extraction}
To make the bra-ket notation precise in the paravector setting, write
\begin{equation*}
\begin{aligned}
\psi&=\psi_s+\psi_{\mathbf v}\cdot\sigv,
&
O&=O_s+O_{\mathbf v}\cdot\sigv,\\
Z&=Z_s+Z_{\mathbf v}\cdot\sigv,
&
X&=X_s+X_{\mathbf v}\cdot\sigv,
\end{aligned}
\end{equation*}
with \(\psi_s,Z_s\in\C\), \(\psi_{\mathbf v},Z_{\mathbf v}\in\C^3\), \(X_s\in\R\), \(X_{\mathbf v}\in\R^3\), and with operator-valued components acting left to right. The three involutions on the state are
\begin{equation*}
\begin{aligned}
\conj{\psi}&=\psi_s^*-\psi_{\mathbf v}^*\cdot\sigv,\\
\rev{\psi}&=\psi_s^*+\psi_{\mathbf v}^*\cdot\sigv,\\
\crev{\psi}&=\psi_s-\psi_{\mathbf v}\cdot\sigv.
\end{aligned}
\end{equation*}
The basic component products are collected below.
% These tables are algebraic component identities.  Their vector rows are not
% identified with probability currents without an accompanying field equation.
\newenvironment{bilineartable}[2][]{%
  \begin{table*}[p]
  \centering
  \caption{Components of $#2$.}%
  \if\relax\detokenize{#1}\relax\else\label{#1}\fi
  \renewcommand{\arraystretch}{1.10}%
  \setlength{\tabcolsep}{2pt}%
  \small
  \begin{tabular}{@{}r@{:\quad}>{\raggedright\arraybackslash}p{0.86\textwidth}@{}}
  \toprule
}{%
  \bottomrule
  \end{tabular}
  \end{table*}
}
\newcommand{\bcomponentrow}[2]{\eqtext{#1} & \(\displaystyle #2\) \\}

\noindent The direct action of a complex paravector operator on a complex paravector state is recorded in Table \ref{tab:components-O-psi}.
\begin{bilineartable}[tab:components-O-psi]{\pv{\mathcal{O}}\pv{\psi}}
\bcomponentrow{scalar}{\eqtext{scalar}(\pv{\mathcal{O}})\eqtext{scalar}(\pv{\psi})+\eqtext{vector}(\pv{\mathcal{O}})\cdot\eqtext{vector}(\pv{\psi})}
\bcomponentrow{vector}{\eqtext{scalar}(\pv{\mathcal{O}})\eqtext{vector}(\pv{\psi})+\eqtext{vector}(\pv{\mathcal{O}})\eqtext{scalar}(\pv{\psi})+i\bigl(\eqtext{vector}(\pv{\mathcal{O}})\times\eqtext{vector}(\pv{\psi})\bigr)}
\end{bilineartable}

\noindent The three basic ordered two-sided products with a generic complex
paravector
\(\pv{Z}=\eqtext{scalar}(\pv{Z})+\eqtext{vector}(\pv{Z})\cdot\sigv\) are given in Tables
\ref{tab:components-psistar-Z-psi}--\ref{tab:components-psistarR-Z-psi}.
\begin{bilineartable}[tab:components-psistar-Z-psi]{\conj{\pv{\psi}}\pv{Z}\pv{\psi}}
\bcomponentrow{scalar}{\begin{aligned}[t]
&\eqtext{scalar}(\pv{Z})\bigl(|\eqtext{scalar}(\pv{\psi})|^2-\eqtext{vector}(\pv{\psi})^*\cdot\eqtext{vector}(\pv{\psi})\bigr)\\
&\quad+\eqtext{scalar}(\pv{\psi})^*(\eqtext{vector}(\pv{Z})\cdot\eqtext{vector}(\pv{\psi}))
-\eqtext{scalar}(\pv{\psi})(\eqtext{vector}(\pv{Z})\cdot\eqtext{vector}(\pv{\psi})^*)\\
&\quad-i\,\eqtext{vector}(\pv{\psi})^*\cdot(\eqtext{vector}(\pv{Z})\times\eqtext{vector}(\pv{\psi}))
\end{aligned}}
\bcomponentrow{vector}{\begin{aligned}[t]
&\eqtext{scalar}(\pv{Z})\bigl(\eqtext{scalar}(\pv{\psi})^*\eqtext{vector}(\pv{\psi})-\eqtext{scalar}(\pv{\psi})\eqtext{vector}(\pv{\psi})^*
-i\,\eqtext{vector}(\pv{\psi})^*\times\eqtext{vector}(\pv{\psi})\bigr)\\
&\quad+\bigl(|\eqtext{scalar}(\pv{\psi})|^2+\eqtext{vector}(\pv{\psi})^*\cdot\eqtext{vector}(\pv{\psi})\bigr)\eqtext{vector}(\pv{Z})\\
&\quad-(\eqtext{vector}(\pv{Z})\cdot\eqtext{vector}(\pv{\psi})^*)\eqtext{vector}(\pv{\psi})
-(\eqtext{vector}(\pv{Z})\cdot\eqtext{vector}(\pv{\psi}))\eqtext{vector}(\pv{\psi})^*\\
&\quad+i\bigl(\eqtext{scalar}(\pv{\psi})^* \eqtext{vector}(\pv{Z})\times\eqtext{vector}(\pv{\psi})
+\eqtext{scalar}(\pv{\psi}) \eqtext{vector}(\pv{Z})\times\eqtext{vector}(\pv{\psi})^*\bigr)
\end{aligned}}
\end{bilineartable}

\begin{bilineartable}[tab:components-psiR-Z-psi]{\herm{\pv{\psi}}\pv{Z}\pv{\psi}}
\bcomponentrow{scalar}{\begin{aligned}[t]
&\eqtext{scalar}(\pv{Z})\bigl(|\eqtext{scalar}(\pv{\psi})|^2+\eqtext{vector}(\pv{\psi})^*\cdot\eqtext{vector}(\pv{\psi})\bigr)\\
&\quad+\eqtext{scalar}(\pv{\psi})^*(\eqtext{vector}(\pv{Z})\cdot\eqtext{vector}(\pv{\psi}))
+\eqtext{scalar}(\pv{\psi})(\eqtext{vector}(\pv{Z})\cdot\eqtext{vector}(\pv{\psi})^*)\\
&\quad+i\,\eqtext{vector}(\pv{\psi})^*\cdot(\eqtext{vector}(\pv{Z})\times\eqtext{vector}(\pv{\psi}))
\end{aligned}}
\bcomponentrow{vector}{\begin{aligned}[t]
&\eqtext{scalar}(\pv{Z})\bigl(\eqtext{scalar}(\pv{\psi})^*\eqtext{vector}(\pv{\psi})+\eqtext{scalar}(\pv{\psi})\eqtext{vector}(\pv{\psi})^*
+i\,\eqtext{vector}(\pv{\psi})^*\times\eqtext{vector}(\pv{\psi})\bigr)\\
&\quad+\bigl(|\eqtext{scalar}(\pv{\psi})|^2-\eqtext{vector}(\pv{\psi})^*\cdot\eqtext{vector}(\pv{\psi})\bigr)\eqtext{vector}(\pv{Z})\\
&\quad+(\eqtext{vector}(\pv{Z})\cdot\eqtext{vector}(\pv{\psi})^*)\eqtext{vector}(\pv{\psi})
+(\eqtext{vector}(\pv{Z})\cdot\eqtext{vector}(\pv{\psi}))\eqtext{vector}(\pv{\psi})^*\\
&\quad+i\bigl(\eqtext{scalar}(\pv{\psi})^* \eqtext{vector}(\pv{Z})\times\eqtext{vector}(\pv{\psi})
-\eqtext{scalar}(\pv{\psi}) \eqtext{vector}(\pv{Z})\times\eqtext{vector}(\pv{\psi})^*\bigr)
\end{aligned}}
\end{bilineartable}

\begin{bilineartable}[tab:components-psistarR-Z-psi]{\cherm{\pv{\psi}}\pv{Z}\pv{\psi}}
\bcomponentrow{scalar}{\eqtext{scalar}(\pv{Z})\bigl(\eqtext{scalar}(\pv{\psi})^2-\eqtext{vector}(\pv{\psi})\cdot\eqtext{vector}(\pv{\psi})\bigr)}
\bcomponentrow{vector}{\begin{aligned}[t]
&\bigl(\eqtext{scalar}(\pv{\psi})^2+\eqtext{vector}(\pv{\psi})\cdot\eqtext{vector}(\pv{\psi})\bigr)\eqtext{vector}(\pv{Z})\\
&\quad-2(\eqtext{vector}(\pv{Z})\cdot\eqtext{vector}(\pv{\psi}))\eqtext{vector}(\pv{\psi})
+2i\,\eqtext{scalar}(\pv{\psi}) \eqtext{vector}(\pv{Z})\times\eqtext{vector}(\pv{\psi})
\end{aligned}}
\end{bilineartable}

\noindent Tables \ref{tab:components-psistar-X-psi}--\ref{tab:components-psistarR-X-psi} are obtained by specialization to a real paravector \(\pv{X}=\eqtext{scalar}(\pv{X})+\eqtext{vector}(\pv{X})\cdot\sigv\), with \(\eqtext{scalar}(\pv{X})\in\R\) and \(\eqtext{vector}(\pv{X})\in\R^3\).
\begin{bilineartable}[tab:components-psistar-X-psi]{\conj{\pv{\psi}}\pv{X}\pv{\psi}}
\bcomponentrow{scalar}{\begin{aligned}[t]
&\eqtext{scalar}(\pv{X})\bigl(|\eqtext{scalar}(\pv{\psi})|^2-\eqtext{vector}(\pv{\psi})^*\cdot\eqtext{vector}(\pv{\psi})\bigr)\\
&\quad+\eqtext{scalar}(\pv{\psi})^*(\eqtext{vector}(\pv{X})\cdot\eqtext{vector}(\pv{\psi}))
-\eqtext{scalar}(\pv{\psi})(\eqtext{vector}(\pv{X})\cdot\eqtext{vector}(\pv{\psi})^*)\\
&\quad-i\,\eqtext{vector}(\pv{\psi})^*\cdot(\eqtext{vector}(\pv{X})\times\eqtext{vector}(\pv{\psi}))
\end{aligned}}
\bcomponentrow{vector}{\begin{aligned}[t]
&\eqtext{scalar}(\pv{X})\bigl(\eqtext{scalar}(\pv{\psi})^*\eqtext{vector}(\pv{\psi})-\eqtext{scalar}(\pv{\psi})\eqtext{vector}(\pv{\psi})^*
-i\,\eqtext{vector}(\pv{\psi})^*\times\eqtext{vector}(\pv{\psi})\bigr)\\
&\quad+\bigl(|\eqtext{scalar}(\pv{\psi})|^2+\eqtext{vector}(\pv{\psi})^*\cdot\eqtext{vector}(\pv{\psi})\bigr)\eqtext{vector}(\pv{X})\\
&\quad-(\eqtext{vector}(\pv{X})\cdot\eqtext{vector}(\pv{\psi})^*)\eqtext{vector}(\pv{\psi})
-(\eqtext{vector}(\pv{X})\cdot\eqtext{vector}(\pv{\psi}))\eqtext{vector}(\pv{\psi})^*\\
&\quad+i\bigl(\eqtext{scalar}(\pv{\psi})^* \eqtext{vector}(\pv{X})\times\eqtext{vector}(\pv{\psi})
+\eqtext{scalar}(\pv{\psi}) \eqtext{vector}(\pv{X})\times\eqtext{vector}(\pv{\psi})^*\bigr)
\end{aligned}}
\end{bilineartable}

\begin{bilineartable}[tab:components-psiR-X-psi]{\herm{\pv{\psi}}\pv{X}\pv{\psi}}
\bcomponentrow{scalar}{\begin{aligned}[t]
&\eqtext{scalar}(\pv{X})\bigl(|\eqtext{scalar}(\pv{\psi})|^2+\eqtext{vector}(\pv{\psi})^*\cdot\eqtext{vector}(\pv{\psi})\bigr)\\
&\quad+\eqtext{scalar}(\pv{\psi})^*(\eqtext{vector}(\pv{X})\cdot\eqtext{vector}(\pv{\psi}))
+\eqtext{scalar}(\pv{\psi})(\eqtext{vector}(\pv{X})\cdot\eqtext{vector}(\pv{\psi})^*)\\
&\quad+i\,\eqtext{vector}(\pv{\psi})^*\cdot(\eqtext{vector}(\pv{X})\times\eqtext{vector}(\pv{\psi}))
\end{aligned}}
\bcomponentrow{vector}{\begin{aligned}[t]
&\eqtext{scalar}(\pv{X})\bigl(\eqtext{scalar}(\pv{\psi})^*\eqtext{vector}(\pv{\psi})+\eqtext{scalar}(\pv{\psi})\eqtext{vector}(\pv{\psi})^*
+i\,\eqtext{vector}(\pv{\psi})^*\times\eqtext{vector}(\pv{\psi})\bigr)\\
&\quad+\bigl(|\eqtext{scalar}(\pv{\psi})|^2-\eqtext{vector}(\pv{\psi})^*\cdot\eqtext{vector}(\pv{\psi})\bigr)\eqtext{vector}(\pv{X})\\
&\quad+(\eqtext{vector}(\pv{X})\cdot\eqtext{vector}(\pv{\psi})^*)\eqtext{vector}(\pv{\psi})
+(\eqtext{vector}(\pv{X})\cdot\eqtext{vector}(\pv{\psi}))\eqtext{vector}(\pv{\psi})^*\\
&\quad+i\bigl(\eqtext{scalar}(\pv{\psi})^* \eqtext{vector}(\pv{X})\times\eqtext{vector}(\pv{\psi})
-\eqtext{scalar}(\pv{\psi}) \eqtext{vector}(\pv{X})\times\eqtext{vector}(\pv{\psi})^*\bigr)
\end{aligned}}
\end{bilineartable}

\begin{bilineartable}[tab:components-psistarR-X-psi]{\cherm{\pv{\psi}}\pv{X}\pv{\psi}}
\bcomponentrow{scalar}{\eqtext{scalar}(\pv{X})\bigl(\eqtext{scalar}(\pv{\psi})^2-\eqtext{vector}(\pv{\psi})\cdot\eqtext{vector}(\pv{\psi})\bigr)}
\bcomponentrow{vector}{\begin{aligned}[t]
&\bigl(\eqtext{scalar}(\pv{\psi})^2+\eqtext{vector}(\pv{\psi})\cdot\eqtext{vector}(\pv{\psi})\bigr)\eqtext{vector}(\pv{X})\\
&\quad-2(\eqtext{vector}(\pv{X})\cdot\eqtext{vector}(\pv{\psi}))\eqtext{vector}(\pv{\psi})
+2i\,\eqtext{scalar}(\pv{\psi}) \eqtext{vector}(\pv{X})\times\eqtext{vector}(\pv{\psi})
\end{aligned}}
\end{bilineartable}

\noindent For a scalar function \(S\), the same three products are reduced to
Tables \ref{tab:components-psistar-S-psi}--\ref{tab:components-psistarR-S-psi}.
\begin{bilineartable}[tab:components-psistar-S-psi]{\conj{\pv{\psi}}S\pv{\psi}}
\bcomponentrow{scalar}{S\bigl(|\eqtext{scalar}(\pv{\psi})|^2-\eqtext{vector}(\pv{\psi})^*\cdot\eqtext{vector}(\pv{\psi})\bigr)}
\bcomponentrow{vector}{S\bigl(\eqtext{scalar}(\pv{\psi})^*\eqtext{vector}(\pv{\psi})-\eqtext{scalar}(\pv{\psi})\eqtext{vector}(\pv{\psi})^*-i\,\eqtext{vector}(\pv{\psi})^*\times\eqtext{vector}(\pv{\psi})\bigr)}
\end{bilineartable}

\begin{bilineartable}[tab:components-psiR-S-psi]{\herm{\pv{\psi}}S\pv{\psi}}
\bcomponentrow{scalar}{S\bigl(|\eqtext{scalar}(\pv{\psi})|^2+\eqtext{vector}(\pv{\psi})^*\cdot\eqtext{vector}(\pv{\psi})\bigr)}
\bcomponentrow{vector}{S\bigl(\eqtext{scalar}(\pv{\psi})^*\eqtext{vector}(\pv{\psi})+\eqtext{scalar}(\pv{\psi})\eqtext{vector}(\pv{\psi})^*+i\,\eqtext{vector}(\pv{\psi})^*\times\eqtext{vector}(\pv{\psi})\bigr)}
\end{bilineartable}

\begin{bilineartable}[tab:components-psistarR-S-psi]{\cherm{\pv{\psi}}S\pv{\psi}}
\bcomponentrow{scalar}{S\bigl(\eqtext{scalar}(\pv{\psi})^2-\eqtext{vector}(\pv{\psi})\cdot\eqtext{vector}(\pv{\psi})\bigr)}
\bcomponentrow{vector}{0}
\end{bilineartable}


The reverse sandwich with a real scalar is the key probabilistic object. For \(S\in\R\),
\begin{equation*}
\rev{\psi}S\psi
=S\bigl(\rho_\psi+\mathbf J_\psi\cdot\sigv\bigr),
\end{equation*}
where
\begin{equation*}
\rho_\psi:=|\psi_s|^2+\psi_{\mathbf v}^*\cdot\psi_{\mathbf v},
\qquad
\mathbf J_\psi:=\psi_s^*\psi_{\mathbf v}
+\psi_s\psi_{\mathbf v}^*
+i\,\psi_{\mathbf v}^*\times\psi_{\mathbf v}.
\end{equation*}
Thus \(\rev{\psi}\psi\) is naturally a density-current paravector: the scalar part is the positive density, while the vector part carries the associated directional information. More generally, the raw bilinear \(\rev{\psi}O\psi\) is a complex paravector, so before any extraction it contains one complex scalar and one complex vector, that is, four real component channels. The scalar channel gives the ordinary expectation value; the remaining channels can encode currents, fluxes, spin-orientation data, or other geometric bilinears depending on the operator.

At fixed time, the natural normalization is therefore
\begin{equation*}
\langle \psi\mid\psi\rangle
:=\int_{\R^3}\eqtext{scalar}(\rev{\psi}\psi)\,\dd^3r
=\int_{\R^3}\rho_\psi\,\dd^3r
=1.
\end{equation*}
For a complex-paravector operator \(O\), define
\begin{equation*}
\langle \psi\mid O\mid\psi\rangle
:=\int_{\R^3}\eqtext{scalar}(\rev{\psi}O\psi)\,\dd^3r.
\end{equation*}
If the state is not already normalized, the corresponding expectation map is
\begin{equation*}
\mathbb E_\psi[O]
:=
\frac{\int_{\R^3}\eqtext{scalar}(\rev{\psi}O\psi)\,\dd^3r}
{\int_{\R^3}\eqtext{scalar}(\rev{\psi}\psi)\,\dd^3r}.
\end{equation*}
If \(O\psi=\phi_s+\phi_{\mathbf v}\cdot\sigv\), then the extracted scalar channel is simply
\begin{equation*}
\eqtext{scalar}(\rev{\psi}O\psi)
=\psi_s^*\phi_s+\psi_{\mathbf v}^*\cdot\phi_{\mathbf v}.
\end{equation*}
When \(O\) is reverse-self-adjoint under the spatial integral and the boundary terms are benign, this scalar expectation is real.

By contrast,
\begin{equation*}
\crev{\psi}\psi
=\psi_s^2-\psi_{\mathbf v}\cdot\psi_{\mathbf v}
=\langle\!\langle\psi\rangle\!\rangle^2
\end{equation*}
is the STA metric bilinear. It is algebraically useful, but it is not positive-definite and so it should not be used as the probabilistic normalization. In this framework the single bra-ket bracket \(\langle\cdot\mid\cdot\rangle\) is therefore best reserved for the scalar expectation map built from reversion, while \(\langle\!\langle\cdot\rangle\!\rangle^2\) continues to denote the paravector metric.

\subsubsection{NEW CONTENT FOR REVIEW}
Standard physics does not assign spin from component count alone. The physically relevant data are
\begin{equation*}
(\text{transformation law},\ \text{field equation},\ \text{constraints}).
\end{equation*}
In standard Lorentz language,
\begin{align*}
\text{spin }0:&\quad (0,0),\\
\text{spin }\tfrac12:&\quad (\tfrac12,0)\oplus(0,\tfrac12),\\
\text{spin }1:&\quad (1,0)\oplus(0,1)\quad \text{for }F,\\
&\quad (\tfrac12,\tfrac12)\quad \text{for }A^\mu.
\end{align*}
So the algebraic carrier spaces
\begin{equation*}
\C,\qquad
\C\oplus\C^3,\qquad
\varnothing\oplus\C^3
\end{equation*}
do not determine the physical spin by themselves. They become spin sectors only after the operator law is fixed.

\noindent\textbf{Bilinear pairings.}
The full reverse product is the one that directly produces the density-current paravector:
\begin{equation*}
\rev{\psi}\psi
=\rho_\psi+\mathbf J_\psi\cdot\sigv,
\qquad
\rho_\psi
=|\psi_s|^2+\psi_{\mathbf v}^*\cdot\psi_{\mathbf v}.
\end{equation*}
For two fields \(\phi,\psi\in\C\oplus\C^3\), the three scalar pairings are
\begin{align*}
\eqtext{scalar}(\rev{\phi}\psi)
&=\phi_s^*\psi_s+\phi_{\mathbf v}^*\cdot\psi_{\mathbf v}
&& \text{Hermitian-type pairing},\\
\eqtext{scalar}(\crev{\phi}\psi)
&=\phi_s\psi_s-\phi_{\mathbf v}\cdot\psi_{\mathbf v}
&& \text{metric pairing},\\
\eqtext{scalar}(\conj{\phi}\psi)
&=\phi_s^*\psi_s-\phi_{\mathbf v}^*\cdot\psi_{\mathbf v}
&& \text{indefinite companion}.
\end{align*}
For \(\phi=\psi\), the metric pairing reduces to
\begin{equation*}
\eqtext{scalar}(\crev{\psi}\psi)
=\crev{\psi}\psi
=\langle\!\langle\psi\rangle\!\rangle^2.
\end{equation*}
Only the reverse pairing is being used here as the quantum bra-ket pairing. The metric pairing is the intrinsic STA quadratic form used in factorization identities, and the companion pairing is algebraically well-defined but does not yet carry an independent standard quantum interpretation in this manuscript.

Accordingly, for the spatially integrated bracket
\begin{equation}
\langle \phi \mid \psi \rangle
:=\int_{\R^3}\eqtext{scalar}(\rev{\phi}\psi)\,\dd^3r,
\qquad
\langle \phi \mid O\psi \rangle
=\langle O^\dagger\phi \mid \psi \rangle,
\label{eq:adjoint-paravector-braket}
\end{equation}
the operator adjoint \(O^\dagger\) is the reverse-adjoint determined by this pairing, up to the usual boundary terms when derivatives are present. For left multiplication by a fixed paravector \(Z\), the adjoint is left multiplication by \(\rev{Z}\), since \(\rev{(\rev{Z}\phi)}=\rev{\phi}Z\). The companion vector channel is also worth naming:
\begin{equation*}
\mathbf J_{O,\psi}
:=\eqtext{vector}(\rev{\psi}O\psi).
\end{equation*}
The case \(O=1\) gives the density-current pair \((\rho_\psi,\mathbf J_\psi)\), while more general \(O\) produce current-, flux-, or polarization-type vector bilinears alongside the ordinary scalar expectation value.

\noindent\textbf{Generalized operator products.}
The following identities are algebraic. They describe operator composition on a paravector-valued state, but by themselves they do not assign a physical spin.

For operator-valued components, the commutator no longer collapses to the pure cross-product law of Eq.\ \eqref{eq:paravector-commutator}. Write
\begin{equation*}
O=O_s+O_{\mathbf v}\cdot\sigv,
\qquad
O'=O_s'+O_{\mathbf v}'\cdot\sigv,
\end{equation*}
with component operators commuting with the fixed basis elements and acting left to right on the state. Then, with dot and cross products taken componentwise and in the written operator order,
\begin{align}
[O,O']\psi
&=\Bigl([O_s,O_s']
+O_{\mathbf v}\cdot O_{\mathbf v}'
-O_{\mathbf v}'\cdot O_{\mathbf v}\Bigr)\psi\notag\\
&\quad+\Bigl([O_s,O_{\mathbf v}']
+[O_{\mathbf v},O_s']\notag\\
&\qquad\qquad
+i\bigl(O_{\mathbf v}\times O_{\mathbf v}'
-O_{\mathbf v}'\times O_{\mathbf v}\bigr)\Bigr)\cdot\sigv\,\psi,
\label{eq:operator-commutator-action}
\end{align}
where the scalar-vector commutators are understood componentwise. Unlike the purely algebraic paravector commutator, Eq.\ \eqref{eq:operator-commutator-action} generally has both scalar and vector parts. Only when all component operators mutually commute does it reduce to the simpler law
\begin{equation*}
[O,O']\psi
=2i\bigl(O_{\mathbf v}\times O_{\mathbf v}'\bigr)\cdot\sigv\,\psi.
\end{equation*}

The matching anticommutator is
\begin{align}
\operatorname{acom}(O,O')\psi
&:=(OO'+O'O)\psi\notag\\
&=\Bigl(\operatorname{acom}(O_s,O_s')
+O_{\mathbf v}\cdot O_{\mathbf v}'
+O_{\mathbf v}'\cdot O_{\mathbf v}\Bigr)\psi\notag\\
&\quad+\Bigl(\operatorname{acom}(O_s,O_{\mathbf v}')\notag\\
&\qquad\qquad
+\operatorname{acom}(O_{\mathbf v},O_s')\notag\\
&\qquad\qquad
+i\bigl(O_{\mathbf v}\times O_{\mathbf v}'
+O_{\mathbf v}'\times O_{\mathbf v}\bigr)\Bigr)\cdot\sigv\,\psi,
\label{eq:operator-anticommutator-action}
\end{align}
so that
\begin{align*}
OO'
&=\tfrac12\bigl(\operatorname{acom}(O,O')+[O,O']\bigr).
\end{align*}
This is the natural companion to the operator commutator formula and is the form from which covariance and uncertainty identities would be built if that development is expanded later.

\noindent\textbf{Carrier closure.}
The same algebra shows which carrier spaces are actually closed under which operators:
\begin{align}
\partial:\C&\to\C\oplus\C^3,\notag\\
\partial:\varnothing\oplus\C^3&\to\C\oplus\C^3,\notag\\
\partial:\C\oplus\C^3&\to\C\oplus\C^3.
\label{eq:carrier-closure}
\end{align}
Indeed,
\begin{align*}
\partial\psi_s
&=(\partial_t\psi_s)+(\gradv\psi_s)\cdot\sigv,\\
\partial(\psi_{\mathbf v}\cdot\sigv)
&=(\gradv\cdot\psi_{\mathbf v})
+\bigl(\partial_t\psi_{\mathbf v}+i\,\gradv\times\psi_{\mathbf v}\bigr)\cdot\sigv.
\end{align*}
So the full complex-paravector carrier is the minimal space closed under first-order paravector operators such as \(\partial\) and \(E\pm\sigv\cdot\mathbf P\). The scalar subspace closes only after passing to scalar second-order combinations such as \(\Box+m^2\). A pure-vector specialization is algebraically meaningful, but it is not preserved by generic first-order paravector action and therefore requires Maxwell- or Proca-type constraints if it is to describe a spin-1 field. A real-paravector specialization \(\psi\in\R\oplus\R^3\) is also possible algebraically, but it is not a general quantum carrier because it is not closed under multiplication by the phase \(e^{i\alpha}\).

With those algebraic preliminaries fixed, the organizing principle is this: one first chooses a carrier field for \(\psi\), then restricts to the operators that preserve that carrier or define its field equation, and then lists the bilinear pairings that extract the physically relevant observables from that carrier. The bilinears need not take values in the same carrier field; in general they land in scalar, vector, or paravector observable channels.

With that understood, the physically consistent sector-by-sector summary is the following.

\noindent\textbf{1. Spin \(0\): complex scalar Klein--Gordon sector}
\begin{align*}
\text{carrier:}\quad
&\psi:\R^{1,3}\to\C,\\
\text{representation:}\quad
&(0,0),\\
\text{field equation:}\quad
&(\Box+m^2)\psi=0,\\
\text{operator class:}\quad
&O_s:\C\to\C,\\
\text{examples:}\quad
&\Box+m^2,\ \partial_t\pm i\Omega,\\
&-E^2+\mathbf P^2+m^2,\ldots,\\
\text{spin projector:}\quad
&\Pi(\pm)\psi\notin\C.
\end{align*}
The scalar pairings reduce to
\begin{align*}
\eqtext{scalar}(\rev{\phi}\psi)
&=\phi^*\psi,\\
\eqtext{scalar}(\crev{\phi}\psi)
&=\phi\psi,\\
\eqtext{scalar}(\conj{\phi}\psi)
&=\phi^*\psi.
\end{align*}
So the reverse and sigma-conjugate pairings coincide on the scalar sector, while the metric pairing is different. The ordinary complex sesquilinear product \(\phi^*\psi\) is the known scalar-field pairing, but the standard relativistic conserved Klein--Gordon current is derivative-built rather than simply \(\psi^*\psi\):
\begin{equation*}
j^\mu_{\mathrm{KG}}
:=\frac{i}{2m}\bigl(\psi^*\partial^\mu\psi-(\partial^\mu\psi^*)\psi\bigr).
\end{equation*}
Thus \(\phi^*\psi\) is the natural Hermitian scalar pairing, while \(\phi\psi\) is an algebraic metric pairing and not a probability density.

\noindent\textbf{2. Spin \(\tfrac12\): Dirac/Pauli sector}
\begin{align*}
\text{algebraic block:}\quad
&\psi:\R^{1,3}\to\C\oplus\C^3,\\
\text{physical field:}\quad
&\bPsi=
\begin{Bmatrix}
\bPsi_1\\[2pt]
\bPsi_2
\end{Bmatrix},
\qquad
\bPsi_1,\bPsi_2\in\C\oplus\C^3,\\
\text{representation:}\quad
&(\tfrac12,0)\oplus(0,\tfrac12),\\
\text{field equations:}\quad
&(\Wmat(i\partial)-m\I)\bPsi=0,\\
&(\Wmat(i\partial-q\Phi^*)-m\I)\bPsi=0,\\
\text{operator class:}\quad
&O=O_s+O_{\mathbf v}\cdot\sigv,\\
&\text{acting on each paravector block},\\
\text{examples:}\quad
&\partial,\ \partial^*,\\
&E\pm\sigv\cdot\mathbf P,\ \Pi(\pm),\\
\text{spin projector:}\quad
&\Pi(\pm):\C\oplus\C^3\to\C\oplus\C^3,\\
&\bPsi(\pm):=\Pi(\pm)\bPsi,\\
&(\mathbf u\cdot\sigv)\bPsi(\pm)=\pm\bPsi(\pm).
\end{align*}
The basic one-block pairings are
\begin{align*}
\eqtext{scalar}(\rev{\phi}\psi)
&=\phi_s^*\psi_s+\phi_{\mathbf v}^*\cdot\psi_{\mathbf v},\\
\mathbf J_{\phi,\psi}
&:=\eqtext{vector}(\rev{\phi}\psi),\\
\eqtext{scalar}(\crev{\phi}\psi)
&=\phi_s\psi_s-\phi_{\mathbf v}\cdot\psi_{\mathbf v},\\
\eqtext{scalar}(\conj{\phi}\psi)
&=\phi_s^*\psi_s-\phi_{\mathbf v}^*\cdot\psi_{\mathbf v}.
\end{align*}
For the two-block Dirac carrier, the natural extension is
\begin{equation*}
\langle \{\Xi_1;\Xi_2\} \mid \{\Psi_1;\Psi_2\} \rangle
:=\int_{\R^3}
\sum_{n=1}^{2}\eqtext{scalar}(\rev{\Xi_n}\Psi_n)\,\dd^3r.
\end{equation*}
Here the reverse pairing is the Hermitian-type pairing used for expectation values and for current or polarization extraction. The metric pairing is the one used in factorization identities such as \(Z\crev{Z}\) and \(\Wmat(Z)^2\). The companion pairing is indefinite and no unique standard physical observable has been attached to it here.

The key physics point is that the one-block complex paravector is the minimal algebraic block preserved by first-order paravector action, but its scalar-plus-vector component count does not by itself make it an irreducible spin-\(\tfrac12\) field in the standard Lorentz sense. In this manuscript the physical spin-\(\tfrac12\) assignment belongs to the later Dirac system built from those paravector blocks.

\noindent\textbf{3. Spin \(1\): vector or field-strength sector}
\begin{align*}
\text{carrier:}\quad
&F:\R^{1,3}\to\varnothing\oplus\C^3,\\
\text{main example:}\quad
&F=(\mathbf E+i\mathbf B)\cdot\sigv,\\
\text{representation:}\quad
&(1,0)\oplus(0,1)\quad\text{for }F,\\
&(\tfrac12,\tfrac12)\quad\text{for a vector potential }A^\mu,\\
\text{field equation:}\quad
&\partial F=D^*,\\
\text{operator class:}\quad
&\partial,\ \partial^*,\ \Box,\ \text{and Maxwell/Proca constraints},\\
\text{closure facts:}\quad
&\partial,\partial^*:\varnothing\oplus\C^3\to\C\oplus\C^3,\\
&\Box:\varnothing\oplus\C^3\to\varnothing\oplus\C^3,\\
\text{spin projector:}\quad
&\Pi(\pm)F\notin\varnothing\oplus\C^3.
\end{align*}
So the two-projector split \(\Pi(\pm)\) is not the standard spin-\(1\) decomposition. The standard polarization counts are
\begin{align*}
\text{massive spin }1:&\quad m_s=-1,0,1,\\
\text{massless spin }1:&\quad \lambda=\pm1.
\end{align*}
The bilinear quantities already present in the manuscript are
\begin{align*}
\tfrac12F^2
&=\tfrac12(\mathbf E^2-\mathbf B^2)+i\,\mathbf E\cdot\mathbf B,\\
\tfrac12F F^{\mathrm R}
&=\tfrac12(\mathbf E^2+\mathbf B^2)+(\mathbf E\times\mathbf B)\cdot\sigv.
\end{align*}
These have direct standard physical interpretations: \(\tfrac12F^2\) packages the Lorentz invariants, while \(\tfrac12F F^{\mathrm R}\) packages the energy density and the Poynting vector. A simple bra-ket built from \(\rev{F}F\) is not the standard Maxwell-field pairing used here; the physically relevant quantities are the field bilinears together with the field equations and their gauge or subsidiary constraints.

The consistent summary is therefore this:
\begin{align*}
\text{spin}
&\neq
\text{component count alone},\\
\text{spin}
&=
(\text{carrier},\ \text{operator law},\ \text{physical bilinears},\ \text{constraints}).
\end{align*}
Within that summary, \((\cdot)^{\mathrm R}\) is the pairing used for quantum expectations, \((\cdot)^{*\mathrm R}\) is the metric pairing used for algebraic factorization, and the pure sigma-conjugate pairing remains an auxiliary companion unless and until a further physical identification is supplied.

\subsubsection{Spin-0 Wave Mechanics}
The scalar relativistic wave equation can be read directly from the mass-shell identity developed earlier in the kinematic section. Here \(\psi\in\C\) is a complex scalar field, \(m\in\R\), and the free-particle operators are defined using the differential conventions from Section \ref{sec:differential}. Start from
\begin{equation*}
\langle\!\langle mU\rangle\!\rangle^2=m^2
\qquad\Rightarrow\qquad
-E^2+\mathbf P^2+m^2=0.
\end{equation*}
For a free particle,
\begin{equation*}
E=i\partial_t,
\qquad
\mathbf P=-i\gradv,
\end{equation*}
and equivalently
\begin{equation*}
E+\mathbf P\cdot\sigv=i\partial^*.
\end{equation*}
Acting on a scalar field \(\psi\), the operator chain is
\begin{align*}
(-E^2+\mathbf P^2+m^2)\psi
&=\bigl(-(i\partial_t)^2-\gradv^2+m^2\bigr)\psi\notag\\
&=\bigl(\partial_t^2-\gradv^2+m^2\bigr)\psi\notag\\
&=(\partial\partial^*+m^2)\psi\notag\\
&=(\Box+m^2)\psi.
\end{align*}
Thus the free spin-0 relativistic wave equation is
\begin{equation*}
(\Box+m^2)\psi=0
\qquad\text{Klein--Gordon equation}.
\end{equation*}
This is the scalar wave equation recorded in Eq.\ \eqref{eq:kg-free}. The dispersion relation \(\omega^2=\mathbf k^2+m^2\) and the later positive/negative-frequency split are consequences of this same operator identity.

\subsubsection{Positive and Negative Frequency Sectors}
The frequency decomposition is one of the central interpretive steps in the scalar theory. Here \(\psi,\psi(\pm),c(\pm)\in\C\), while \(m\in\R\) and \(\gradv\) is the spatial derivative operator fixed earlier in Section \ref{sec:differential}. The frequency operator is
\begin{align*}
\Omega
&=\sqrt{m^2-\gradv^2}\notag\\
&=\langle\!\langle m+\sigv\cdot\gradv\rangle\!\rangle.
\end{align*}
Accordingly,
\begin{align*}
\Box+m^2
&=\partial_t^2+\Omega^2\notag\\
&=\partial_t^2+(m^2-\gradv^2)\notag\\
&=(\partial_t-i\Omega)(\partial_t+i\Omega)\notag\\
&=(\partial_t+i\Omega)(\partial_t-i\Omega).
\end{align*}
The positive- and negative-frequency sectors are then defined by
\begin{align*}
(\partial_t\pm i\Omega)\psi(\pm)&=0,\\
\psi(\pm)&=c(\pm)\exp(\mp i\Omega t),\\
\psi&=\psi(+)+\psi(-).
\end{align*}
This decomposition is not merely notation: it separates the two energy branches of the scalar theory in the same way that the Fourier transform separates positive and negative frequencies \cite{greiner2000}.

\subsubsection{Interpreting Root Operators with Fourier}
Square-root operators are easiest to interpret one Fourier mode at a time. In this subsection \(\psi\in\C\), the coefficients \(c_n\in\C\), the frequencies \(\omega_n\in\R\), and the wave vectors \(\mathbf k_n\in\R^3\). Write
\begin{align*}
\psi&=\sum_{n} c_{n}f_{n},\\
f_{n}&=\exp\!\bigl(i(\omega_{n}t+\mathbf k_{n}\cdot\mathbf r)\bigr).
\end{align*}
Then the d'Alembertian and its square root act modewise:
\begin{align*}
\Box\psi
&=\sum_{n}(\mathbf k_{n}^{2}-\omega_{n}^{2})c_{n}f_{n},\\
\pm\sqrt{\Box}\,\psi
&=\sum_{n}\pm\sqrt{\mathbf k_{n}^{2}-\omega_{n}^{2}}\,c_{n}f_{n},\\
(\Box+m^{2})\psi
&=\sum_{n}(\mathbf k_{n}^{2}+m^{2}-\omega_{n}^{2})c_{n}f_{n}.
\end{align*}
For on-shell Klein--Gordon modes,
\begin{equation*}
\mathbf k_{n}^{2}+m^{2}=\omega_{n}^{2},
\qquad
\omega_{n}=\pm\sqrt{\mathbf k_{n}^{2}+m^{2}},
\end{equation*}
so the operator factorization can be read as the modewise identity
\begin{equation*}
\Box+m^{2}=(m\mp i\sqrt{\Box})(m\pm i\sqrt{\Box}).
\end{equation*}
Likewise, if \(\Omega=\sqrt{m^{2}-\gradv^{2}}\) as in Eq.\ \eqref{eq:frequency-split}, then
\begin{align*}
\Omega^{2}\psi
&=\sum_{n}(m^{2}+\mathbf k_{n}^{2})c_{n}f_{n},\\
\Omega\psi
&=\sum_{n}\sqrt{m^{2}+\mathbf k_{n}^{2}}\,c_{n}f_{n}.
\end{align*}
This is the concrete meaning of the root operators used above: they do not represent a formal symbolic manipulation, but multiplication of each Fourier mode by the corresponding dispersion function.

\subsubsection{Klein--Gordon with Potential Fields}
To keep the minimally coupled scalar theory in one self-contained place, start from the charged-particle Lagrangian
\begin{equation*}
\mathcal L=-\frac{m}{\gamma}-q(V-\mathbf v\cdot\mathbf A).
\end{equation*}
Here \(E\) denotes the energy operator, not the electric field \(\mathbf E\). The Hamiltonian and canonical momentum are then
\begin{equation*}
i\partial_t=\mathcal H=E+qV,
\qquad
-i\gradv=\delv\mathcal L=\mathbf P+q\mathbf A,
\end{equation*}
so the covariant energy and momentum operators are
\begin{equation*}
E=i\partial_t-qV,
\qquad
\mathbf P=-i\gradv-q\mathbf A.
\end{equation*}

Acting on a smooth test field \(\psi\), the Leibniz rules needed below are
\begin{align*}
\partial_t(V\psi)&=(\partial_tV)\psi+V\partial_t\psi,\\
\partial_m(A_n\psi)&=(\partial_mA_n)\psi+A_n\partial_m\psi.
\end{align*}
From these one obtains
\begin{align*}
E^2\psi
&=E(i\partial_t\psi-qV\psi)\notag\\
&=\left(-\partial_t^2-iq\,\partial_tV-2iqV\partial_t+q^2V^2\right)\psi,\\
\mathbf P^2\psi
&=\mathbf P\cdot(-i\gradv\psi-q\mathbf A\psi)\notag\\
&=\left(-\gradv^2+iq\,\gradv\cdot\mathbf A+2iq\,\mathbf A\cdot\gradv+q^2\mathbf A^2\right)\psi.
\end{align*}

Recall also
\begin{align*}
\Box&=\partial_t^2-\gradv^2,\\
\Phi&=V+\mathbf A\cdot\sigv,\\
S&=\partial_tV+\gradv\cdot\mathbf A,\\
\langle\!\langle\partial\rangle\!\rangle^2&=\Box,\\
\langle\!\langle\Phi\rangle\!\rangle^2&=V^2-\mathbf A^2.
\end{align*}
Substituting the operator expansions gives the minimally coupled Klein--Gordon equation in explicit form:
\begin{align*}
0&=(-E^2+\mathbf P^2+m^2)\psi\notag\\
&=\Bigl(\Box+iq(\partial_tV+\gradv\cdot\mathbf A)+2iq(V\partial_t+\mathbf A\cdot\gradv)\notag\\
&\qquad\qquad-q^2(V^2-\mathbf A^2)+m^2\Bigr)\psi\notag\\
&=\Bigl(\Box+iqS-q^2\langle\!\langle\Phi\rangle\!\rangle^2+m^2+2iq(V\partial_t+\mathbf A\cdot\gradv)\Bigr)\psi.
\end{align*}

The electromagnetic fields and commutators enter through
\begin{align*}
\mathbf E&=-\partial_t\mathbf A-\gradv V,\\
\mathbf B&=\gradv\times\mathbf A,\\
F&=(\mathbf E+i\mathbf B)\cdot\sigv,
\end{align*}
and
\begin{align*}
[E,P_m]\psi
&=(i\partial_t-qV)(-i\partial_m\psi-qA_m\psi)\notag\\
&\quad-(-i\partial_m-qA_m)(i\partial_t\psi-qV\psi)\notag\\
&=-iq\bigl((\partial_tA_m)\psi+A_m\partial_t\psi\bigr)\notag\\
&\quad-iq\bigl((\partial_mV)\psi+V\partial_m\psi\bigr)+iqV\partial_m\psi+iqA_m\partial_t\psi\notag\\
&=-iq(\partial_tA_m+\partial_mV)\psi\notag\\
&=iqE_m\psi,\\
[P_m,P_n]\psi
&=(-i\partial_m-qA_m)\bigl(-i\partial_n\psi-qA_n\psi\bigr)\notag\\
&\quad-(-i\partial_n-qA_n)\bigl(-i\partial_m\psi-qA_m\psi\bigr)\notag\\
&=iq(\partial_mA_n-\partial_nA_m)\psi.
\end{align*}
Hence
\begin{equation*}
[E,\mathbf P]=iq\mathbf E,
\qquad
\mathbf P\times\mathbf P=iq\mathbf B.
\end{equation*}
Substituting those commutators into the ordered products yields
\begin{align*}
(E+\sigv\cdot\mathbf P)(E-\sigv\cdot\mathbf P)
&=E^2-\mathbf P^2\notag\\
&\qquad+\sigv\cdot\bigl(-[E,\mathbf P]-i\,\mathbf P\times\mathbf P\bigr)\notag\\
&=E^2-\mathbf P^2-iqF,\\
(E-\sigv\cdot\mathbf P)(E+\sigv\cdot\mathbf P)
&=E^2-\mathbf P^2\notag\\
&\qquad+\sigv\cdot\bigl([E,\mathbf P]-i\,\mathbf P\times\mathbf P\bigr)\notag\\
&=E^2-\mathbf P^2-iqF^*.
\end{align*}
Since
\begin{equation*}
E+\sigv\cdot\mathbf P=i\partial^*-q\Phi,
\qquad
E-\sigv\cdot\mathbf P=i\partial-q\Phi^*
=(E+\sigv\cdot\mathbf P)^{*\mathrm R},
\end{equation*}
their scalar paravector metric is
\begin{equation*}
\langle\!\langle E\pm\sigv\cdot\mathbf P\rangle\!\rangle^2
=E^2-\mathbf P^2+\sigv\cdot(\pm [E,\mathbf P]-i\,\mathbf P\times\mathbf P).
\end{equation*}
This is not an ordinary scalar operator: operator ordering still matters, so the two first-order products remain distinct.
Consequently,
\begin{align*}
-(E+\sigv\cdot\mathbf P)(E+\sigv\cdot\mathbf P)^{*\mathrm R}
&=(\partial^*+iq\Phi)(\partial+iq\Phi^*),\\
-(E+\sigv\cdot\mathbf P)^{*\mathrm R}(E+\sigv\cdot\mathbf P)
&=(\partial+iq\Phi^*)(\partial^*+iq\Phi).
\end{align*}
The same factorizations can be rewritten directly in derivative form:
\begin{align*}
(E+\sigv\cdot\mathbf P)(E-\sigv\cdot\mathbf P)\psi
&=-(\partial^*+iq\Phi)(\partial+iq\Phi^*)\psi\notag\\
&=-\Box\psi-iq\,\partial^*(\Phi^*\psi)\notag\\
&\qquad-iq\,\Phi\,\partial\psi\notag\\
&\qquad+q^2\langle\!\langle\Phi\rangle\!\rangle^2\psi\notag\\
&=-\Box\psi-iqS\psi+q^2\langle\!\langle\Phi\rangle\!\rangle^2\psi\notag\\
&\qquad-2iq(V\partial_t+\mathbf A\cdot\gradv)\psi-iqF\psi,\\
(E-\sigv\cdot\mathbf P)(E+\sigv\cdot\mathbf P)\psi
&=-(\partial+iq\Phi^*)(\partial^*+iq\Phi)\psi\notag\\
&=-\Box\psi-iq\,\partial(\Phi\psi)\notag\\
&\qquad-iq\,\Phi^*\,\partial^*\psi\notag\\
&\qquad+q^2\langle\!\langle\Phi\rangle\!\rangle^2\psi\notag\\
&=-\Box\psi-iqS\psi+q^2\langle\!\langle\Phi\rangle\!\rangle^2\psi\notag\\
&\qquad-2iq(V\partial_t+\mathbf A\cdot\gradv)\psi-iqF^*\psi.
\end{align*}
Equivalently,
\begin{align*}
\partial^*(\Phi^*\psi)+\Phi\,\partial\psi
&=S\psi+2(V\partial_t+\mathbf A\cdot\gradv)\psi+F\psi,\\
\partial(\Phi\psi)+\Phi^*\,\partial^*\psi
&=S\psi+2(V\partial_t+\mathbf A\cdot\gradv)\psi+F^*\psi,
\end{align*}
and their difference isolates the electric contribution:
\begin{align*}
\bigl(\partial^*(\Phi^*\psi)+\Phi\,\partial\psi\bigr)
&-\bigl(\partial(\Phi\psi)+\Phi^*\,\partial^*\psi\bigr)\notag\\
&=\bigl([\Phi,\partial]+[\partial,\Phi]^*\bigr)\psi\notag\\
&=(F-F^*)\psi\notag\\
&=2(\mathbf E\cdot\sigv)\psi.
\end{align*}

In summary,
\begin{align*}
E&:=i\partial_t-qV, &
\mathbf P&:=-i\gradv-q\mathbf A,\\
\partial&:=\partial_t+\sigv\cdot\gradv, &
\partial^*&:=\partial_t-\sigv\cdot\gradv,\\
\Phi&:=V+\mathbf A\cdot\sigv, &
F&:=(\mathbf E+i\mathbf B)\cdot\sigv,\\
S&:=\partial_tV+\gradv\cdot\mathbf A, &
\mathbf E&:=-\partial_t\mathbf A-\gradv V,\\
\mathbf B&:=\gradv\times\mathbf A,
\end{align*}
and the scalar Klein--Gordon operator can be recovered from either ordered first-order product:
\begin{align*}
0&=(-E^2+\mathbf P^2+m^2)\psi\notag\\
&=\Bigl(\Box+iqS-q^2\langle\!\langle\Phi\rangle\!\rangle^2+m^2+2iq(V\partial_t+\mathbf A\cdot\gradv)\Bigr)\psi\notag\\
&=(\partial^*+iq\Phi)(\partial+iq\Phi^*)\psi-iqF\psi+m^2\psi\notag\\
&=(\partial+iq\Phi^*)(\partial^*+iq\Phi)\psi-iqF^*\psi+m^2\psi.
\end{align*}
The ordered first-order products themselves obey
\begin{align*}
(E\pm\sigv\cdot\mathbf P)(E\mp\sigv\cdot\mathbf P)
&=E^2-(\sigv\cdot\mathbf P)^2\mp[E,\sigv\cdot\mathbf P]\notag\\
&=E^2-\mathbf P^2-iq(\pm\mathbf E+i\mathbf B)\cdot\sigv.
\end{align*}
Taking the scalar part therefore recovers
\begin{equation*}
0=\eqtext{scalar}\!\bigl((\partial+iq\Phi^*)(\partial^*+iq\Phi)\bigr)\psi+m^2\psi.
\end{equation*}
Only the scalar part reproduces the Klein--Gordon operator itself; the full ordered products retain the non-scalar field terms \(-iqF\) and \(-iqF^*\) that distinguish the two factorizations.

\subsubsection{Spin-\texorpdfstring{$\frac12$}{1/2} Wave Mechanics}
Before the electromagnetic field is turned on, the same auxiliary \(2\times2\) map from Eq.\ \eqref{eq:W-map} gives a self-contained first-order version of the relativistic mass shell. Consider the two-component field
\begin{equation*}
\bPsi=
\begin{Bmatrix}
\bPsi_1\\[2pt]
\bPsi_2
\end{Bmatrix}.
\end{equation*}
For any complex paravector or paravector-valued operator \(Z\), define
\begin{align*}
\Wmat(Z)&:=
\begin{bmatrix}
0 & Z \\
\crev{Z} & 0
\end{bmatrix},\\
\Wmat(Z)^2&=\operatorname{diag}(Z\crev{Z},\crev{Z}Z).
\end{align*}
If \(a,b\in\R\), then \(\Wmat((a+ib)Z)=(a+ib)\Wmat(Z)\). In particular,
\begin{align*}
\Wmat(i\partial)
&=i
\begin{bmatrix}
0 & \partial \\
\partial^* & 0
\end{bmatrix}
=i\Wmat(\partial),\\
\Wmat(i\partial)^2
&=-\operatorname{diag}(\partial\partial^*,\partial^*\partial)\notag\\
&=-\Box\I.
\end{align*}

Define the free Dirac operator matrix and its mass conjugate by
\begin{equation*}
\Dmat:=\Wmat(i\partial)-m\I,
\qquad
\Dmat':=\Wmat(i\partial)+m\I.
\end{equation*}
Then
\begin{align*}
\Dmat'\Dmat
&=(\Wmat(i\partial)+m\I)(\Wmat(i\partial)-m\I)\notag\\
&=\Wmat(i\partial)^2-m^2\I\notag\\
&=-(\Box+m^2)\I.
\end{align*}
Thus the free Dirac equation
\begin{equation*}
\Dmat\bPsi=0
\end{equation*}
forces each component to satisfy the Klein--Gordon operator. Formally, its characteristic polynomial is
\begin{align*}
0
&=\det(\Dmat-s\I)\notag\\
&=\det
\begin{bmatrix}
-(s+m) & i\partial \\
i\partial^* & -(s+m)
\end{bmatrix}\notag\\
&=(s+m)^2+\Box,
\end{align*}
so the modewise eigenvalue branches are \(s=-m\pm i\sqrt{\Box}\), with the square-root operator interpreted in the Fourier sense from the previous subsection.

Writing out \(\Dmat\bPsi=0\) gives the two coupled equations
\begin{align*}
i\partial\bPsi_2-m\bPsi_1&=0,\\
i\partial^*\bPsi_1-m\bPsi_2&=0.
\end{align*}
Solving the first equation for the upper component and substituting into the second yields
\begin{align*}
\bPsi_1&=\frac{i\partial\bPsi_2}{m},\\
0&=i\partial^*\bPsi_1-m\bPsi_2\notag\\
&=-\frac{\partial^*\partial\bPsi_2}{m}-m\bPsi_2\notag\\
&=-\frac{(\Box+m^2)\bPsi_2}{m}.
\end{align*}
Hence the lower component is a Klein--Gordon solution, and the upper component is obtained from it by one first-order derivative.

Conversely, let \(\psi'\) be any scalar Klein--Gordon solution and let \(a,b\in\C\) be constants. Define
\begin{equation*}
\bPsi_1=\frac{(ma+ib\,\partial)\psi'}{m},
\qquad
\bPsi_2=\frac{(ia\,\partial^*+mb)\psi'}{m}.
\end{equation*}
Because \(a\) and \(b\) are constant and \((\Box+m^2)\psi'=0\),
\begin{align*}
i\partial\bPsi_2-m\bPsi_1
&=\frac{(-\partial a\,\partial^*+im\,\partial b-m^2a-imb\,\partial)\psi'}{m}\notag\\
&=-\frac{a}{m}(\Box+m^2)\psi'=0,\\
i\partial^*\bPsi_1-m\bPsi_2
&=\frac{(mi\,\partial^*a-\partial^*b\,\partial-mi\,a\,\partial^*-m^2b)\psi'}{m}\notag\\
&=-\frac{b}{m}(\Box+m^2)\psi'=0.
\end{align*}
The family generated from a scalar Klein--Gordon solution therefore solves the free two-component system for any constant \(a\) and \(b\). The simplest choice is \(a=1\) and \(b=0\), for which
\begin{equation*}
\bPsi=
\begin{Bmatrix}
\psi'\\[2pt]
\dfrac{i}{m}\,\partial^*\psi'
\end{Bmatrix}.
\end{equation*}
This is the free spin-\(\tfrac12\) wave-mechanics factorization in its most explicit form: the Dirac operator linearizes the Klein--Gordon mass shell, and scalar solutions generate two-component spinor solutions by a first-order lift.

\subsubsection{Dirac with Potential Field}
Consider again the two-component field
\begin{equation*}
\bPsi=
\begin{Bmatrix}
\bPsi_1\\[2pt]
\bPsi_2
\end{Bmatrix}.
\end{equation*}
For any complex paravector or paravector-valued operator \(Z\), define
\begin{align*}
\Wmat(Z)&:=
\begin{bmatrix}
0 & Z \\
\crev{Z} & 0
\end{bmatrix},\\
\Wmat(Z)^2&=\operatorname{diag}(Z\crev{Z},\crev{Z}Z),\\
\Wmat(Z)^*&=\Wmat(Z^*),\\
\Wmat(Z)^{\mathrm R}&=\Wmat(Z^{\mathrm R}).
\end{align*}
The Dirac operator with potential is
\begin{align*}
\Dmat&:=\Wmat(i\partial-q\Phi^*)-m\I\notag\\
&=i\Wmat(\partial)-q\Wmat(\Phi^*)-m\I.
\end{align*}
Because
\begin{align*}
\Wmat(\partial)^2&=\Box\I,\\
\Wmat(\partial)\Wmat(\Phi^*)
&=
\begin{bmatrix}
0 & \partial\\
\partial^* & 0
\end{bmatrix}
\begin{bmatrix}
0 & \Phi^*\\
\Phi & 0
\end{bmatrix}
=\operatorname{diag}(\partial\Phi,\partial^*\Phi^*),\\
\Wmat(\Phi^*)\Wmat(\partial)
&=
\begin{bmatrix}
0 & \Phi^*\\
\Phi & 0
\end{bmatrix}
\begin{bmatrix}
0 & \partial\\
\partial^* & 0
\end{bmatrix}
=\operatorname{diag}(\Phi^*\partial^*,\Phi\partial),\\
\Wmat(\Phi^*)^2
&=
\begin{bmatrix}
0 & \Phi^*\\
\Phi & 0
\end{bmatrix}^{2}
=\operatorname{diag}(\Phi^*\Phi,\Phi\Phi^*)\notag\\
&=\langle\!\langle\Phi\rangle\!\rangle^2\I,
\end{align*}
one obtains
\begin{align*}
\Wmat(i\partial-q\Phi^*)^2
&=-\Wmat(\partial)^2\notag\\
&\qquad-iq\bigl(\Wmat(\partial)\Wmat(\Phi^*)+\Wmat(\Phi^*)\Wmat(\partial)\bigr)\notag\\
&\qquad+q^2\Wmat(\Phi^*)^2\notag\\
&=-\Box\I-iq\,\operatorname{diag}(\partial\Phi+\Phi^*\partial^*,\partial^*\Phi^*+\Phi\partial)\notag\\
&\qquad+q^2\langle\!\langle\Phi\rangle\!\rangle^2\I\notag\\
&=-\operatorname{diag}\!\Big(
(\partial+iq\Phi^*)(\partial^*+iq\Phi),\notag\\
&\qquad\qquad\qquad\quad
(\partial^*+iq\Phi)(\partial+iq\Phi^*)
\Big).
\end{align*}
The Dirac equation
\begin{equation*}
\Dmat\bPsi=0
\end{equation*}
therefore becomes the two coupled equations
\begin{align*}
-m\bPsi_1+(i\partial-q\Phi^*)\bPsi_2&=0,\\
(i\partial^*-q\Phi)\bPsi_1-m\bPsi_2&=0.
\end{align*}
Solving the first equation for the upper component and substituting into the second gives
\begin{align*}
\bPsi_1&=\frac{(i\partial-q\Phi^*)\bPsi_2}{m},\\
0&=(i\partial^*-q\Phi)\bPsi_1-m\bPsi_2\notag\\
&=-\frac{\bigl((\partial^*+iq\Phi)(\partial+iq\Phi^*)+m^2\bigr)\bPsi_2}{m}.
\end{align*}
Thus each component satisfies a Klein--Gordon-type operator with the appropriate ordered field coupling.

The mass-conjugate operator is
\begin{equation*}
\Dmat':=\Wmat(i\partial-q\Phi^*)+m\I.
\end{equation*}
Multiplying by \(\Dmat'\) yields
\begin{align*}
0=\Dmat'\Dmat\bPsi
&=\bigl(\Wmat(i\partial-q\Phi^*)^2-m^2\I\bigr)\bPsi\notag\\
&=-\operatorname{diag}\!\Big(
(\partial+iq\Phi^*)(\partial^*+iq\Phi)+m^2,\notag\\
&\qquad\qquad\qquad\quad
(\partial^*+iq\Phi)(\partial+iq\Phi^*)+m^2
\Big)\bPsi.
\end{align*}
Equivalently,
\begin{align*}
0&=\bigl((\partial+iq\Phi^*)(\partial^*+iq\Phi)+m^2\bigr)\bPsi_1\notag\\
&=\bigl(-E^2+\mathbf P^2+m^2+iqF^*\bigr)\bPsi_1,\\
0&=\bigl((\partial^*+iq\Phi)(\partial+iq\Phi^*)+m^2\bigr)\bPsi_2\notag\\
&=\bigl(-E^2+\mathbf P^2+m^2+iqF\bigr)\bPsi_2.
\end{align*}
Once the square is taken, the coupled Dirac system separates into two Klein--Gordon-type equations whose field terms differ only by the conjugation pattern fixed by the operator ordering.

\subsubsection{Gauge Transform of the Dirac System}
Let \(\lambda\in\R\) be a real scalar gauge field and define
\begin{equation*}
\Phi'=\Phi+\partial^*\lambda.
\end{equation*}
Because \(\lambda\) is real,
\begin{equation*}
\partial\lambda=(\Phi'-\Phi)^*.
\end{equation*}
Transform the two-component field by
\begin{equation*}
\bPsi=e^{iq\lambda}\bPsi'.
\end{equation*}
Then the derivative of the phase factor is
\begin{align*}
\partial\bigl(e^{iq\lambda}\bPsi'\bigr)
&=e^{iq\lambda}\bigl(iq\,\partial\lambda+\partial\bigr)\bPsi'.
\end{align*}
Therefore,
\begin{align*}
0&=(\Wmat(i\partial-q\Phi^*)-m\I)\bPsi\notag\\
&=(\Wmat(i\partial-q\Phi^*)-m\I)e^{iq\lambda}\bPsi'\notag\\
&=e^{iq\lambda}(\Wmat(i\partial-q\,\partial\lambda-q\Phi^*)-m\I)\bPsi'\notag\\
&=e^{iq\lambda}(\Wmat(i\partial-q(\Phi'-\Phi)^*-q\Phi^*)-m\I)\bPsi'\notag\\
&=e^{iq\lambda}(\Wmat(i\partial-q\Phi'^*)-m\I)\bPsi'.
\end{align*}
Hence
\begin{equation*}
(\Wmat(i\partial-q\Phi'^*)-m\I)\bPsi'=0.
\end{equation*}
This is the gauge-covariance statement: the phase change of the Dirac field is exactly compensated by the shift of the potential, so the coupled Dirac equation keeps the same form in the primed gauge.

\subsubsection{Dirac Equation Components}
For the component form of the coupled Dirac equation, write
\begin{align*}
\partial&:=\partial_t+\sigv\cdot\gradv,\\
\partial^*&:=\partial_t-\sigv\cdot\gradv,\\
\Phi&:=V+\mathbf A\cdot\sigv\in\R\oplus\R^3,\\
\Phi^*&:=V-\mathbf A\cdot\sigv,\\
\psi&:=\psi_s+\psi_{\mathbf v}\cdot\sigv\in\C\oplus\C^3.
\end{align*}
For the field term, use
\begin{align*}
F&:=\mathbf F\cdot\sigv,\\
\eqtext{scalar}(F)&=0,\\
\mathbf F&:=\mathbf E+i\mathbf B\in\C^3,\\
F^*&=-\mathbf F^*\cdot\sigv.
\end{align*}
The supporting tables then give the scalar and vector parts of \(\partial\psi\), \(\partial^*\psi\), \(\Phi\psi\), \(\Phi^*\psi\), \(iqF\psi\), and \(iqF^*\psi\) explicitly.
\noindent The direct component action of \(\pv{\partial}\) and \(\conj{\pv{\partial}}\) on a complex paravector test field is recorded in Tables \ref{tab:components-partial-psi} and \ref{tab:components-partialstar-psi}.
\begin{componenttable}[tab:components-partial-psi]{\pv{\partial} \pv{\psi}_{\eqtext{block}}}
\componentrow{scalar}{\partial_t\eqfunc{scalar}(\pv{\psi}_{\eqtext{block}})+\gradv\cdot\eqfunc{vector}(\pv{\psi}_{\eqtext{block}})}{}
\componentrow{vector}{\partial_t\eqfunc{vector}(\pv{\psi}_{\eqtext{block}})+\gradv\eqfunc{scalar}(\pv{\psi}_{\eqtext{block}})+i(\gradv\times\eqfunc{vector}(\pv{\psi}_{\eqtext{block}}))}{}
\end{componenttable}

\begin{componenttable}[tab:components-partialstar-psi]{\conj{\pv{\partial}}\pv{\psi}_{\eqtext{block}}}
\componentrow{scalar}{\partial_t\eqfunc{scalar}(\pv{\psi}_{\eqtext{block}})-\gradv\cdot\eqfunc{vector}(\pv{\psi}_{\eqtext{block}})}{}
\componentrow{vector}{\partial_t\eqfunc{vector}(\pv{\psi}_{\eqtext{block}})-\gradv\eqfunc{scalar}(\pv{\psi}_{\eqtext{block}})-i(\gradv\times\eqfunc{vector}(\pv{\psi}_{\eqtext{block}}))}{}
\end{componenttable}

\noindent The component actions of the potential \(\pv{\Phi}\) and its conjugate on the same field are then shown in Tables \ref{tab:components-phi-psi} and \ref{tab:components-phistar-psi}.
\begin{componenttable}[tab:components-phi-psi]{\pv{\Phi}\pv{\psi}_{\eqtext{block}}}
\componentrow{scalar}{V\eqfunc{scalar}(\pv{\psi}_{\eqtext{block}})+\vct{A}\cdot\eqfunc{vector}(\pv{\psi}_{\eqtext{block}})}{}
\componentrow{vector}{V\eqfunc{vector}(\pv{\psi}_{\eqtext{block}})+\eqfunc{scalar}(\pv{\psi}_{\eqtext{block}})\,\vct{A}+i(\vct{A}\times\eqfunc{vector}(\pv{\psi}_{\eqtext{block}}))}{}
\end{componenttable}

\begin{componenttable}[tab:components-phistar-psi]{\conj{\pv{\Phi}}\pv{\psi}_{\eqtext{block}}}
\componentrow{scalar}{V\eqfunc{scalar}(\pv{\psi}_{\eqtext{block}})-\vct{A}\cdot\eqfunc{vector}(\pv{\psi}_{\eqtext{block}})}{}
\componentrow{vector}{V\eqfunc{vector}(\pv{\psi}_{\eqtext{block}})-\eqfunc{scalar}(\pv{\psi}_{\eqtext{block}})\,\vct{A}-i(\vct{A}\times\eqfunc{vector}(\pv{\psi}_{\eqtext{block}}))}{}
\end{componenttable}

\noindent The Dirac-block operator tables are completed by Tables
\ref{tab:components-iqf-psi} and \ref{tab:components-iqfstar-psi}, in which
the field-coupling terms \(iq\pv{F}\pv{\psi}_{\eqtext{block}}\) and
\(iq\conj{\pv{F}}\pv{\psi}_{\eqtext{block}}\) are recorded. In the second table,
\(\conj{\pv{F}}=-\vct{F}^{*}\cdot\sigv
=-(\vct{E}-i\vct{B})\cdot\sigv\), so no additional conjugation operation is
hidden in the notation.
\begin{componenttable}[tab:components-iqf-psi]{iq\pv{F}\pv{\psi}_{\eqtext{block}}}
\componentrow{scalar}{iq\bigl(\vct{F}\cdot\eqfunc{vector}(\pv{\psi}_{\eqtext{block}})\bigr)}{}
\componentrow{vector}{iq\,\eqfunc{scalar}(\pv{\psi}_{\eqtext{block}})\,\vct{F}
-q\bigl(\vct{F}\times\eqfunc{vector}(\pv{\psi}_{\eqtext{block}})\bigr)}{}
\end{componenttable}

\begin{componenttable}[tab:components-iqfstar-psi]{iq\conj{\pv{F}}\pv{\psi}_{\eqtext{block}}}
\componentrow{scalar}{-iq\bigl(\vct{F}^{*}\cdot\eqfunc{vector}(\pv{\psi}_{\eqtext{block}})\bigr)}{}
\componentrow{vector}{-iq\,\eqfunc{scalar}(\pv{\psi}_{\eqtext{block}})\,\vct{F}^{*}
+q\bigl(\vct{F}^{*}\times\eqfunc{vector}(\pv{\psi}_{\eqtext{block}})\bigr)}{}
\end{componenttable}


\subsubsection{Low-Energy Spin-0 Limit}
For the minimally coupled scalar theory, use the covariant operators
\begin{equation*}
E=i\partial_t-qV,
\qquad
\mathbf P=-i\gradv-q\mathbf A.
\end{equation*}
Remove the fast rest-mass oscillation by writing
\begin{equation*}
\psi'=e^{-imt}\psi.
\end{equation*}
Then
\begin{align*}
E^2\psi'&=e^{-imt}(E+m)^2\psi,\\
\mathbf P^2\psi'&=e^{-imt}\mathbf P^2\psi.
\end{align*}
Substituting into the minimally coupled Klein--Gordon equation gives
\begin{align*}
0&=(-E^2+\mathbf P^2+m^2)\psi'\notag\\
&=e^{-imt}\bigl(-(E+m)^2+\mathbf P^2+m^2\bigr)\psi,
\end{align*}
hence
\begin{equation*}
0=(-E^2-2mE+\mathbf P^2)\psi.
\end{equation*}
Formally iterating this relation gives
\begin{align*}
E&=\frac{\mathbf P^2-E^2}{2m},\\
&=\frac{\mathbf P^2}{2m}-\frac{\mathbf P^4}{8m^3}+\cdots.
\end{align*}
At low energy,
\begin{equation*}
E\psi\approx \frac{\mathbf P^2}{2m}\psi.
\end{equation*}
Solving for \(i\partial_t\psi\) yields the Schr\"odinger equation
\begin{equation*}
i\partial_t\psi
=\left(\frac{(i\gradv+q\mathbf A)^2}{2m}+qV\right)\psi.
\end{equation*}
If \(\mathbf A=0\), this reduces to
\begin{equation*}
i\partial_t\psi
=\left(-\frac{\gradv^2}{2m}+qV\right)\psi.
\end{equation*}
The scalar low-energy limit is therefore already encoded in the relativistic Klein--Gordon theory.

\subsubsection{Low-Energy Spin-\texorpdfstring{$\frac12$}{1/2} and Pauli Limit}
Start from the two-component field
\begin{equation*}
\bPsi'=
\begin{Bmatrix}
\bPsi_1'\\[2pt]
\bPsi_2'
\end{Bmatrix}
\end{equation*}
and the potential-coupled Dirac equation
\begin{equation*}
\bigl(\Wmat(i\partial-q\Phi^*)-m\I\bigr)\bPsi'=0.
\end{equation*}
Its two component equations are
\begin{align*}
(i\partial-q\Phi^*)\bPsi_2'-m\bPsi_1'&=0,\\
(i\partial^*-q\Phi)\bPsi_1'-m\bPsi_2'&=0.
\end{align*}
Write
\begin{align*}
\partial&:=\partial_t+\sigv\cdot\gradv,
&
\partial^*&:=\partial_t-\sigv\cdot\gradv,\\
\Phi&:=V+\mathbf A\cdot\sigv,
&
\Phi^*&:=V-\mathbf A\cdot\sigv.
\end{align*}
Then the two equations become
\begin{align*}
(i\partial_t-qV)\bPsi_2'
+\bigl((i\gradv+q\mathbf A)\cdot\sigv\bigr)\bPsi_2'
-m\bPsi_1'&=0,\\
(i\partial_t-qV)\bPsi_1'
-\bigl((i\gradv+q\mathbf A)\cdot\sigv\bigr)\bPsi_1'
-m\bPsi_2'&=0.
\end{align*}
Introduce the covariant operators
\begin{equation*}
E:=i\partial_t-qV,
\qquad
\mathbf P:=-i\gradv-q\mathbf A.
\end{equation*}
The system is then
\begin{align*}
(E-\mathbf P\cdot\sigv)\bPsi_2'-m\bPsi_1'&=0,\\
(E+\mathbf P\cdot\sigv)\bPsi_1'-m\bPsi_2'&=0.
\end{align*}
Peel off the fast rest-mass oscillation with
\begin{equation*}
\bPsi_n'=e^{-imt}\bPsi_n.
\end{equation*}
Since
\begin{equation*}
E\bPsi_n'=e^{-imt}(E+m)\bPsi_n,
\end{equation*}
the reduced system becomes
\begin{align*}
(E+m-\mathbf P\cdot\sigv)\bPsi_2-m\bPsi_1&=0,\\
(E+m+\mathbf P\cdot\sigv)\bPsi_1-m\bPsi_2&=0.
\end{align*}
Solve the first equation for the upper component:
\begin{equation*}
\bPsi_1=\frac{(E+m-\mathbf P\cdot\sigv)\bPsi_2}{m}.
\end{equation*}
Substituting this into the second equation and multiplying by \(m\) gives
\begin{equation*}
0=\bigl((E+m+\mathbf P\cdot\sigv)(E+m-\mathbf P\cdot\sigv)-m^2\bigr)\bPsi_2.
\end{equation*}
Expand the operator product as
\begin{align*}
&(E+m+\mathbf P\cdot\sigv)\notag\\
&\qquad\times(E+m-\mathbf P\cdot\sigv)-m^2\notag\\
&=(E+m)^2-m^2-(\mathbf P\cdot\sigv)^2\notag\\
&\qquad+[\mathbf P\cdot\sigv,E].
\end{align*}
Now
\begin{align*}
-(\mathbf P\cdot\sigv)^2
&=-\mathbf P^2-i(\mathbf P\times\mathbf P)\cdot\sigv,\\
[P_m,P_n]
&=[-i\partial_m-qA_m,-i\partial_n-qA_n]\notag\\
&=iq(\partial_mA_n-\partial_nA_m),\\
\mathbf P\times\mathbf P
&:=([P_2,P_3],-[P_1,P_3],[P_1,P_2])\notag\\
&=iq\,\gradv\times\mathbf A\notag\\
&=iq\,\mathbf B,
\end{align*}
so
\begin{equation*}
-i(\mathbf P\times\mathbf P)\cdot\sigv=q\,\mathbf B\cdot\sigv.
\end{equation*}
Also,
\begin{align*}
[\mathbf P\cdot\sigv,E+m]
&=[\mathbf P,E]\cdot\sigv\notag\\
&=iq(\gradv V+\partial_t\mathbf A)\cdot\sigv\notag\\
&=-iq\,\mathbf E\cdot\sigv,
\end{align*}
where
\begin{equation*}
\mathbf E:=-\partial_t\mathbf A-\gradv V,
\qquad
\mathbf B:=\gradv\times\mathbf A.
\end{equation*}
Defining
\begin{equation*}
F:=(\mathbf E+i\mathbf B)\cdot\sigv,
\end{equation*}
one has
\begin{equation*}
-iqF=q\,\mathbf B\cdot\sigv-iq\,\mathbf E\cdot\sigv.
\end{equation*}
Therefore the exact lower-component equation is
\begin{align*}
0
&=\bigl((E+m)^2-m^2-\mathbf P^2+q\,\mathbf B\cdot\sigv+[\mathbf P\cdot\sigv,E]\bigr)\bPsi_2\notag\\
&=\bigl(E^2+2mE-\mathbf P^2+q\,\mathbf B\cdot\sigv-iq\,\mathbf E\cdot\sigv\bigr)\bPsi_2\notag\\
&=\bigl(E^2+2mE-\mathbf P^2-iqF\bigr)\bPsi_2.
\end{align*}
This can be iterated recursively as
\begin{align*}
E
&=\frac{\mathbf P^2-E^2-q\,\mathbf B\cdot\sigv-[\mathbf P\cdot\sigv,E]}{2m}\notag\\
&=\frac{\mathbf P^2-q\,\mathbf B\cdot\sigv}{2m}+O(1/m^2).
\end{align*}
At low energy, the dominant terms satisfy
\begin{equation*}
0\approx(2mE-\mathbf P^2+q\,\mathbf B\cdot\sigv)\bPsi_2.
\end{equation*}
Solving for \(i\partial_t\bPsi_2\) yields the Pauli equation
\begin{equation*}
i\partial_t\bPsi_2
=\left(\frac{(-i\gradv-q\mathbf A)^2}{2m}+qV\right)\bPsi_2
-\frac{q}{2m}\,(\mathbf B\cdot\sigv)\bPsi_2.
\end{equation*}
The nonrelativistic spin coupling is therefore not an extra postulate: it is the leading field correction that survives when the massive Dirac system is expanded around its rest-energy oscillation \cite{pauli1927,sakurai2017}.

\subsubsection{Spin Up and Spin Down}
Let the magnetic-field direction define the spin axis,
\begin{equation*}
\mathbf u:=\frac{\mathbf B}{\|\mathbf B\|},
\end{equation*}
and use the projection operators
\begin{align*}
\Pi(\pm)&=\tfrac12(1\pm \mathbf u\cdot\sigv),\\
\Pi(\pm)^2&=\Pi(\pm),\\
\Pi(+)\Pi(-)&=0,\\
\Pi(+)+\Pi(-)&=1,\\
(\mathbf u\cdot\sigv)\Pi(\pm)&=\pm\Pi(\pm),\\
\Pi(\pm)^{\mathrm R}&=\Pi(\pm),\\
\Pi(\pm)^*&=\Pi(\mp).
\end{align*}
Accordingly,
\begin{equation*}
\mathbf u\cdot\sigv=(+1)\Pi(+)+(-1)\Pi(-).
\end{equation*}
Decompose the Pauli state by
\begin{align*}
\bPsi(\pm)&:=\Pi(\pm)\bPsi,\\
(\mathbf u\cdot\sigv)\bPsi(\pm)&=\pm\bPsi(\pm),\\
\bPsi&=\bPsi(+)+\bPsi(-),
\end{align*}
so \(\bPsi(+)\) and \(\bPsi(-)\) are the spin-up and spin-down parts along the field axis.

Write the Pauli equation as
\begin{align*}
\mathbf P&:=-i\gradv-q\mathbf A,\\
i\partial_t\bPsi
&=\left(\frac{\mathbf P^2}{2m}+qV\right)\bPsi
-\frac{q}{2m}\,\|\mathbf B\|(\mathbf u\cdot\sigv)\bPsi.
\end{align*}
Equivalently,
\begin{equation*}
i\partial_t\bPsi=H\bPsi,
\end{equation*}
with
\begin{equation*}
H:=\frac{\mathbf P^2}{2m}+qV-\frac{q}{2m}\,\|\mathbf B\|\,\mathbf u\cdot\sigv.
\end{equation*}
Applying \(\Pi(\pm)\) to the time evolution gives
\begin{align*}
i\partial_t\bPsi(\pm)
&=i(\partial_t\Pi(\pm))\bPsi+\Pi(\pm)(i\partial_t\bPsi)\notag\\
&=i(\partial_t\Pi(\pm))\bPsi+\Pi(\pm)(H\bPsi).
\end{align*}
Add zero in the form \(H\Pi(\pm)\bPsi-H\Pi(\pm)\bPsi\):
\begin{align*}
\Pi(\pm)(H\bPsi)
&=\Pi(\pm)(H\bPsi)+H\Pi(\pm)\bPsi-H\Pi(\pm)\bPsi\notag\\
&=H\bPsi(\pm)+[\Pi(\pm),H]\bPsi.
\end{align*}
Therefore
\begin{align*}
i\partial_t\bPsi(\pm)
&=H\bPsi(\pm)+\bigl(i\partial_t\Pi(\pm)+[\Pi(\pm),H]\bigr)\bPsi\notag\\
&=:H\bPsi(\pm)+S(\pm).
\end{align*}
In explicit form,
\begin{align*}
i\partial_t\bPsi(\pm)
&=\left(\frac{\mathbf P^2}{2m}+qV\right)\bPsi(\pm)\notag\\
&\qquad-\frac{q}{2m}\,\|\mathbf B\|(\mathbf u\cdot\sigv)\bPsi(\pm)+S(\pm),
\end{align*}
with coupling term
\begin{equation*}
S(\pm):=\left(i\partial_t\Pi(\pm)+\frac{[\mathbf P^2,\Pi(\pm)]}{2m}\right)\bPsi.
\end{equation*}
If \(\mathbf B\) is uniform and static, then \(\Pi(\pm)\) is constant and
\begin{equation*}
S(+)=S(-)=0,
\end{equation*}
so the two spin sectors decouple. In that case
\begin{equation*}
(-\mathbf B\cdot\sigv)\bPsi(\pm)=\mp \|\mathbf B\|\,\bPsi(\pm),
\end{equation*}
and the Zeeman effect appears as
\begin{equation*}
i\partial_t\bPsi(\pm)
=\left(\frac{\mathbf P^2\mp q\|\mathbf B\|}{2m}+qV\right)\bPsi(\pm).
\end{equation*}
The field direction therefore selects a preferred axis, and the two projector parts evolve with opposite Zeeman shifts.


\section{Conclusions}\label{sec:conclusions}\label{sec:synthesis}
The algebra used here is the real Pauli algebra, written as one complex scalar plus one complex three-vector. Its economy is supplied by a structural fact rather than by a new physical postulate: the complex unit is supplied by the central pseudoscalar, dot and cross products are combined by the sigma product, and the determinant is obtained by multiplication by the adjugate. On real event paravectors, that determinant is the Minkowski interval.

Within this framework, spatial rotations, Lorentz boosts, four-velocity and four-momentum, Maxwell's equations, the Lorentz force, scalar relativistic waves, the Dirac factorization, and the leading Pauli equation are derived in one notation. The distinctions between object types are essential. Events are transformed by the determinant-preserving congruence derived from the general paravector law, the electromagnetic field is transformed by similarity, and one Dirac field is carried by a fixed projected column space. Once those distinctions are enforced, agreement between the compact notation and the standard component theories is obtained.

Several complementary structures are also exposed by the reformulation. The four Maxwell equations are obtained from the real and imaginary parts of \(\pv{\partial}\pv{F}\); the conjugate spin couplings \(\pv{F}\) and \(\conj{\pv{F}}\) are obtained from the two ordered momentum products; and elementary functions, Pauli spin sectors, and free helicity sectors are organized by the same axis projectors. These are connections among established equations, not evidence for an additional interaction.

In particular, no independent measured components are added by the
paravector wave carriers. Each Dirac or Pauli block is reduced by the fixed-column
condition to one ordinary two-entry state, its opposite-order bilinear is the familiar
rank-one density matrix, and the two squared field terms are the familiar
two Dirac spin couplings. Other nonscalar sandwich channels remain algebraic
until a physical meaning is assigned by a separately derived observable.

\subsection{Scope of the Established Results}
The results in this paper are established under the assumptions of flat Minkowski space, smooth fields, and domains on which the displayed differential operators and square roots are defined. The electromagnetic discussion is classical, and the wave-equation discussion is first-quantized. An antiparticle interpretation of a negative-frequency solution is obtained only after second quantization. Likewise, \(\pv{F}\herm{\pv{F}}/2\) is the energy-density and flux paravector seen in one frame; the full electromagnetic stress-energy object is the linear function \(\pv{T}(\pv{a})=\pv{F}\pv{a}\herm{\pv{F}}/2\) on real paravectors \(\pv{a}\).

The displayed components are changed by a nonlinear coordinate map or algebraic change of basis in flat space, but curvature is not thereby created. Additional geometric structure---a position-dependent frame, a compatible connection, curvature, and covariant field equations---would be required for a genuine extension to general relativity. Complex space-time, dark-sector, and cosmological claims are outside what has been derived here.

These and other uncompleted directions are recorded in the following chapter without being treated as results.

\section{Future Work}\label{sec:future-work}
The developed results constitute a reformulation of established physics in
space-time without gravitational curvature. Extensions not yet constructed
are separated from points requiring
greater mathematical precision. A proposed extension is marked
\textbf{TODO}. A gap requiring resolution before physical claims can be
supported is marked \textbf{FIXME}. No premise for the established results is
taken from this section.

\subsection{TODO: Local Frames and Gravitation}
One fixed sigma basis is used throughout space and time in the present calculations.
A gravitational extension would instead need a position-dependent local
basis and a rule for comparing algebra elements defined at neighboring
positions.  That comparison rule would have to distinguish an ordinary
change of description from a physical gravitational field.  In particular,
a basis written everywhere as a single transform of a fixed basis is locally
only a change of description. A failure to return an algebra element
unchanged after comparison around a sufficiently small closed path is called
curvature \cite{doran2003}; it is not produced by that basis
change alone.

A TODO question is therefore whether one comparison rule can be defined and
used consistently in every equation. Such a rule is conventionally called a
connection \cite{doran2003}. Curvature would then have to be calculated by
the closed-path comparison just defined, rather than by a change of basis
alone. If two successive comparisons depend on their
order, they are said not to commute. The interval, freely falling motion,
Maxwell equation, Dirac system, and energy--momentum source rule would all
have to be rebuilt with that same comparison rule. Agreement with the
constant-basis formulas would be a required check.

It remains to be determined whether translational motion can be combined
with rotation-rate and boost-rate cross-position terms and consistently
interpreted as a velocity. Multiplication by mass and construction of a
transformed interval would follow only if that interpretation were proved.
Those physical identifications are not justified by the moving-basis generator derived in
Sec.\ \ref{sec:transform-algebra} alone. In particular, surrounding each
basis element by a transform and its inverse is not the same operation as the
two-sided event transformation used for a Lorentz boost, and a complex
velocity-like expression is generally produced by the proposed boost-rate
term. It must therefore be determined whether a real event law and an
observer's measured time and distance can be defined, whether an ordinary
derivative can be corrected for the changing local basis so that the
corrected derivative transforms with the field it acts on, and whether real
measured velocity, energy, momentum, and interval can then be derived. The
corrected derivative is conventionally called a covariant derivative
\cite{doran2003}. Recovery of the inertial formulas when the
basis-change rate is constant or zero would be required.

Space-time curvature is also not yet supplied by a plane-curve analogy.
A numerator obtained by differentiating a coordinate quotient is only one
piece of a curve formula. The remaining TODO questions are whether the curve
parameter has a fixed physical meaning, whether the vector pointing along
the curve, conventionally called its tangent, has a fixed length, how
directions are compared at neighboring positions, and whether the proposed
result works for motion in more than one direction. Curvature
cannot be inferred from that numerator alone.

\subsection{TODO: Additional Field-Value Transformations}
One phase freedom and one potential paravector are used in the electromagnetic
coupling. Here an internal field value means a label carried by the field
that is changed without changing the space-time event at which the field is
evaluated. Other established field theories use several coupled internal
values and transformations for which applying the first and then the second
can differ from applying them in the opposite order \cite{doran2003}. The
TODO question is whether such transformations can be added while keeping the
space-time product and the internal action separate.

The allowed internal field values, their transformation rule, the
corresponding potentials, and the quantities preserved by the resulting
equations would have to be defined before any field equation is written.
The products applied to space-time paravectors and those applied only to the
internal field values would also have to be shown explicitly. It must be
determined whether the existing algebra can express both actions without
ambiguity. Merely replacing the electromagnetic potential by a larger array
would not answer that question.

\subsection{TODO: Many-Particle and Particle-Number-Changing Fields}
Classical fields and single-particle wave functions are described by the wave
equations in this manuscript. The term ``first quantized'' refers to this use
of an ordinary wave function while the number of particles remains fixed. In
particular, a negative-frequency solution is not by itself an antiparticle
\cite{greiner2000}. A future theory would require field operators whose
action can change the number of particles, operations that add or remove one
particle, a transformation that reverses charge, and a no-particle reference
state called the vacuum. The TODO question is whether all of these objects
can be represented without obscuring the present algebra, and whether the
particle--antiparticle interpretation can then be derived rather than
assumed \cite{peskin1995}.

Further TODO questions concern how the state is changed when two identical
particles are exchanged, how a state is represented when it cannot be
written as one independent state for each particle, and how the remaining
state is described after some particles are left unmeasured. The condition
that the state cannot be separated is conventionally called entanglement,
and the last partial
description is conventionally called a reduced state. The rule for extracting
quantities measured near one position would have to be specified in the same
development \cite{peskin1995}.
The fixed projected column space used for one Dirac field is a useful
starting point, but these many-particle constructions are not supplied by it
alone.

The indefinite Klein--Gordon pairing and the absence of a positive local
one-particle position density are also not repaired by a change of algebraic
packaging. In a theory with particle-number-changing field operators, the
origin of a positive squared state length, an operation that measures charge,
and the distinction between particles and antiparticles must be shown. The
one-particle formulas must then be recovered when particle-number changes can
be neglected.

\subsection{TODO: Electromagnetic Response of Materials}
Empty-space fields and prescribed sources are treated in the electromagnetic section.
A material requires an additional rule that converts the electromagnetic
field into the response of that material \cite{jackson1998}. Such a rule is
conventionally called a constitutive map. Its outputs can include electric
dipole moment per unit volume, called polarization; magnetic dipole moment
per unit volume, called magnetization; and electric current. An electric
dipole moment records the amount and direction of separated positive and
negative charge. A magnetic dipole moment records the strength and direction
of a small magnet or current loop. Direction
dependence is called anisotropy, frequency dependence is called dispersion,
current produced by an electric field is described by conductivity, and
conversion of field energy into heat or other unavailable forms is called
dissipation.

The first TODO question is whether simple media whose response respects
addition and multiplication by an ordinary number can be written clearly.
Such a response rule is called linear. Direction-dependent and
time-dependent responses can then be tested. The distinction between a field
paravector and an operation acting on that field must be preserved, and the
energy balance must be derived rather than assumed. It must then be
determined which material laws fit into one paravector product and which
require a separately defined map.

\subsection{FIXME: Operator Foundations}
Differential, inverse, and square-root operators are used in several
factorizations in the wave-mechanics sections. Their algebraic
forms are correct on the stated test fields, for which all displayed
derivatives exist, but the allowed set of fields must be specified
completely. That allowed set is the operator's domain. Conditions imposed at
the edge of the spatial region or at infinite distance are its boundary
conditions \cite{greiner2000,sakurai2017}. The conditions under which
integration by parts---moving a derivative from one factor in an integral to
the other factor---adds no boundary term, an inverse exists, and an energy
operator has real measured values are determined by those choices.

A consistent definition of the square-root frequency operator is also
required for every plane-wave term. The nonnegative square root must be
selected explicitly, and the spatially constant term with
\(\vct{k}=\vct{0}\) must be treated separately when the inverse is used.
Enough derivatives and integrals must exist for both written operation orders
in every use of the commutator function. These conditions should eventually
be stated once in a common operator-domain section and then cited wherever
they are used in a factorization or low-energy expansion.

Likewise, the determinant of a block differential operator minus a proposed
measured value is not defined merely because the paravector determinant is
defined. An ordinary matrix determinant whose zeros select possible values
may be used only after the equation has been reduced to a finite matrix, for
example one plane wave at a time, a finite set of basis functions, or a
finite numerical grid. Merely restricting the spatial region does not make
the operator finite-dimensional. At the operator level, only the ordered
factorizations and explicitly proved
polynomial identities are presently justified. It remains to be determined
under what stated conditions an infinite-dimensional determinant can be
defined, if one is needed at all.

For scalar relativistic waves, the positive-energy solution space for each
permitted background that does not change with time must be constructed, and
the measurement operators that preserve it must be determined. General
time-dependent backgrounds can mix the positive- and negative-frequency
parts, and a position probability cannot be obtained merely by declaring the
local Klein--Gordon density to be a position-probability density. For Dirac
and Pauli wave-function
spaces, the shared operator domains must be specified. It must then be proved
that moving an operator from one side of the bracket to the other gives the
displayed \(H\) operation, that the energy-selecting operators preserve their
selected states, and that the calculated measurement averages are real.

\subsection{FIXME: Bounds on Measurement Spread}
A spin-measurement uncertainty bound cannot be inferred directly from the
axis commutator function and a determinant inequality. For a measured
quantity, its centered operator is obtained by subtracting its average from
the operator. Its variance is the average square of that difference, and its
uncertainty is the nonnegative square root of the variance
\cite{robertson1929,sakurai2017}. The normalized projected-wave bracket in
Eq.\ \eqref{eq:adjoint-paravector-braket} must be used to define those
averages. A shared set of fields on which both operator orders and their
\(H\) operations exist must also be required. The standard lower bound built
from the commutator function, and the stronger bound that also retains the
anticommutator function, should then be proved in the manuscript's notation.
Spin-axis measurement operators should be checked as a worked example rather
than treated as the proof.

\subsection{TODO: Higher-Order Pauli Reduction}
The leading kinetic and magnetic spin terms are proved by the present Dirac
reduction. The next TODO question is whether the next-order terms can be
derived with the same sign and charge conventions. These include the coupling
between spin and motion in an electric field, called the spin--orbit term,
and the local correction produced by spatial variation of the electric
field, often called the contact or Darwin term
\cite{greiner2000,sakurai2017}. Corrections required for time-dependent
potentials must also be derived. The density and current must be expanded to
the same order as the energy operator; retaining higher-order energy terms
while leaving the measurement rule at leading order would be inconsistent.
These terms are established consequences of the Dirac equation, but their
derivation in this notation remains future work.

\subsection{TODO: Classify Companion Components}
Paired real and imaginary components and paired scalar and vector components
are repeatedly produced by the algebra. Some of these are familiar physical
laws: the magnetic-divergence and induction equations, for example, are
encoded by the imaginary parts of the Maxwell equation. Other components, such as the
imaginary part of the force product, are valid algebraic companions without
a separate physical interpretation established here.

Two cases are already identified with established physics: a projected
block's opposite-order product is the standard density matrix for one fully
specified two-entry state \cite{sakurai2017}. Its rank is one, meaning that
every column is a multiple of the one selected state column. The remaining
\(\pv{F}\) and \(\conj{\pv{F}}\) terms in the two ordered products are the
standard two Dirac spin couplings \cite{greiner2000}. For the remaining
nonscalar channels of unrestricted or differently ordered products, a
separate derivation as a current, spin-axis quantity, or measurement result
is still required before physical meaning is assigned. For a fixed-column
state and a left-acting operator that preserves that column, the product
\(\herm{\pv{\psi}}\pv{\mathcal{O}}\pv{\psi}\) is already fixed by its one
complex row--column matrix element and contains no additional independent
channels.

The following exact identity may also be studied:
\begin{equation*}
\Box(\pv{F}^{\,2})
=\conj{\pv{\partial}}\pv{\partial}(\pv{F}^{\,2}).
\end{equation*}
Because \(\pv{F}^{\,2}\) is a complex scalar, this equality is an exact use
of the defined wave operator, but it is not an independent field equation.
Any useful source expansion must be derived from Maxwell's equation with the
ordered product rule of Eq.\ \eqref{eq:paravector-leibniz}. A TODO question
is whether the resulting scalar identity describes a balance unchanged by
the stated transformations or merely repeats information already contained
in Maxwell's equation. It must not be presented as an independent law of
field evolution without such a derivation.

Each companion formula should be placed by a systematic classification into
one of three categories: an alternate form of an established law, redundant
information implied by another displayed equation, or an algebraic channel
whose physical meaning is presently unassigned. The operations that return
the original input when applied twice, the order functions,
\(\pv{\Pi}(\pm)\) selection functions, and field products should all be
included in that classification. No unassigned
component should be advertised as a new
interaction merely because it is made available by the algebra.
Maps that preserve the defined addition and multiplication, together with
the selection, factorization, exponential, and logarithmic constructions,
should be separated into those proved directly from the algebra and those
that depend on the chosen matrix realization.

\subsection{TODO: Complex Space-Time and Imaginary Coordinates}
A complex paravector \(\pv{Z}\in\Pspace(\C)\) is already permitted by the
generic algebra: \(\pv{Z}=\pv{X}+i\pv{Y}\), with
\(\pv{X},\pv{Y}\in\Pspace(\R)\). Treating such an element as a physical event, rather than as an
algebraic field value, would be a new hypothesis. The following self-contained
consistency identity is obtained from the existing determinant and adjugate
definitions:
\begin{equation*}
\detf{\pv{Z}}
=\detf{\pv{X}}-\detf{\pv{Y}}
+i\Bigl(
  \pv{X}\adjf{\pv{Y}}
 +\pv{Y}\adjf{\pv{X}}
\Bigr).
\end{equation*}
The quantity in parentheses is a scalar, but the determinant of
\(\pv{Z}\) is generally complex. No interval interpretation is fixed by this
identity. In a proposed complex-event theory, it would have to be stated which,
if any, real quantity is measured as an interval and why that choice is
preserved by the allowed transformations.

The scalar part of \(\pv{Y}\) could be described as an imaginary-time
coordinate, and its vector part could be described as imaginary space.
Those descriptions alone carry no physical content. A metric is a stated
rule that assigns an interval to a coordinate difference
\cite{minkowski1908}. The TODO question is whether a real interval rule can
be selected for a complex paravector and preserved by the allowed
transformations. Taking the real part, absolute magnitude, or squared
absolute magnitude of a complex determinant would give different interval
rules, so none can be selected by terminology alone. The ordinary real
interval would have to be recovered when the imaginary part vanishes.

Only after that choice is justified could the time measured along an
observer's path, the possible time order of events connected by a signal,
velocity, momentum, or force be defined for a complex event. For real
space-time, the zero-interval directions form the light cone and separate
event differences that can be joined by a signal no faster than light from
those that cannot \cite{einstein1905,minkowski1908}. A TODO question is
whether that same measured boundary is recovered to experimental accuracy.
Complex mass, unobserved paths, or new force channels should not be introduced
before these more basic questions are resolved.

Faster-than-light propagation is not implied by complex coordinates. A
propagation speed can be assigned only after field equations, the surfaces
along which disturbances propagate, observers, and measured time and distance
have been defined. If
the measured projection remains the real paravector \(\pv{X}\), the ordinary
light cone may remain unchanged. If \(\pv{Y}\) changes how disturbances
propagate, the new boundary between permitted and forbidden signal directions
and the possibility of controllable faster-than-light signaling would have
to be derived explicitly.

An algebraic relationship between the real and imaginary parts is shown by
the mixed terms in \(\detf{\pv{Z}}\), but a weak physical interaction is not
established. Such an interaction would require an equation of evolution and
a source law. Only then could constants to be measured experimentally be
identified. TODO questions could then be asked about a real--imaginary mixing
strength, a mass or length scale that reduces the imaginary part, or a
relative propagation or decay scale. None of these constants is predicted by
the algebra alone; each would have to be connected to a definite measured
quantity and constrained by experiment.

Allowing complex coefficients in the algebra cannot by itself be tested
experimentally because it is a mathematical extension. A proposed physical
interpretation could be rejected if measured intervals or energies became
unavoidably complex, probabilities failed to remain real and nonnegative,
time evolution failed to preserve normalization, source-free solutions grew
without bound, or predicted departures from the observed light cone and
ordinary Lorentz behavior exceeded experimental limits. The final TODO
question is whether any nontrivial complex-event evolution law can pass these
tests or whether every consistently measured result reduces to the
real-paravector theory already developed.

\subsection{FIXME: Proposed Connections to Large-Scale Observations}
No derivation is established here for energy assigned to otherwise empty
space, for the observed speeding-up of large-scale expansion, or for the
additional gravitational effects commonly attributed to unseen matter
\cite{weinberg2008}. These topics are collectively described here as
dark-sector questions. A contribution to gravity, a tendency to gather with
matter, or a measured energy is not implied by the existence of imaginary
algebra components.

A consistent local gravitational equation and a conserved physical source
must first be supplied by any such model. Ordinary local gravity must then be
recovered. The TODO question is whether quantitative predictions can be made
for the changing distance between far-separated galaxies, the bending of
light by gravity, the formation of large matter concentrations, and the
motion of matter \cite{weinberg2008}. A nearly constant contribution spread
over large distances and a contribution that gathers into dense regions are
physically different requirements and must not be inferred from the same
algebraic term without separate derivations.

\subsection{TODO: Limits of the Notation-Light Form}
The compact notation is most effective when the relevant scalar and vector
data are carried by one algebra element. In some extensions, maps with several
independent inputs, information that is valid only on overlapping local
regions, or conditions on the whole spatial region that cannot be recovered
from one local formula
are instead required. That mathematical structure is not removed by avoidance of indices;
only its recorded form is changed.

The conditions under which a paravector is sufficient, a map that preserves
addition and multiplication by ordinary numbers is required, or
several-input or whole-region data must be introduced should be stated
explicitly. Any new notation should be justified by a construction that
cannot be expressed clearly with the existing functions and operators.
The raw-text character of the framework is thereby preserved while
notational economy is prevented from being mistaken for a reduction in
physical content.


\clearpage
\onecolumn
\appendix
\section{Verification and Formula Status}\label{sec:verification}
Exact algebraic identities are separated in Table \ref{tab:verification-status} from physical statements for which a stated representation or approximation is required. The accompanying symbolic report is generated by \texttt{tools/verify\_sta\_results.py}.
\begin{longtable}{@{}L{0.22\linewidth}L{0.24\linewidth}L{0.14\linewidth}L{0.30\linewidth}@{}}
\caption{Status of the principal constructions.}\label{tab:verification-status}\\
\toprule
Statement & Standard correspondence & Status & Qualification \\
\midrule
\endfirsthead
\toprule
Statement & Standard correspondence & Status & Qualification \\
\midrule
\endhead
Eqs.\ \eqref{eq:operator-rules} and \eqref{eq:paravector-product}
& Real Pauli algebra and three-dimensional real Clifford algebra \cite{hestenes1966,doran2003}
& \verifiedmark
& Associative scalar--vector form with eight real degrees of freedom.\\
Eqs.\ \eqref{eq:transform-intertwining}--\eqref{eq:general-event-congruence}
& Similarity transport and the standard double-sided Lorentz action
& \verifiedmark
& Generic algebra elements are transported by similarity; real events are transported by congruence. The adjugate equals the inverse only after determinant-one normalization.\\
Eqs.\ \eqref{eq:euclidean-form} and \eqref{eq:minkowski-form}
& Positive coefficient norm and ordinary matrix determinant
& \verifiedmark
& The first is nonnegative and real; the second is generally complex and is the exact inverse divisor.\\
Eqs.\ \eqref{eq:rotor} and \eqref{eq:boost}
& Spatial rotations and Lorentz boosts \cite{lorentz1904,einstein1905}
& \verifiedmark
& For spatial rotors, the event congruence is reduced to similarity; for event boosts, the congruence form is retained.\\
Eq.\ \eqref{eq:interval}
& Minkowski interval \cite{minkowski1908}
& \verifiedmark
& Valid on real event paravectors with signature \((+,-,-,-)\).\\
Eqs.\ \eqref{eq:maxwell} and \eqref{eq:field-transform}
& Maxwell electrodynamics \cite{maxwell1865,jackson1998}
& \verifiedmark
& Signs are fixed by \(\pv{F}=(\vct{E}+i\vct{B})\cdot\sigv\) and the passive-boost convention.\\
Eq.\ \eqref{eq:lorentz-force}
& Relativistic Lorentz force \cite{jackson1998}
& \verifiedmark
& The coordinate-time power and force equations are obtained by division by \(\gamma\).\\
Eq.\ \eqref{eq:energy-momentum-paravector}
& Electromagnetic energy density and flux
& \verifiedmark
& \(\pv{T}\) is the chosen-observer paravector; the full stress-energy object is the linear map on real paravectors.\\
Eqs.\ \eqref{eq:kg-free} and \eqref{eq:kg-coupled}
& Klein--Gordon equation \cite{klein1926,gordon1926}
& \verifiedmark
& The frequency split is defined on the stated spatial operator domain.\\
Eqs.\ \eqref{eq:kg-mixed-current}--\eqref{eq:kg-pairing}
& Conserved Klein--Gordon solution pairing
& \verifiedmark
& Its integral is positive on the free positive-frequency sector; its local density is not a general position probability.\\
Eqs.\ \eqref{eq:dirac-coupled} and \eqref{eq:fixed-column-carrier}
& Standard Dirac equation in two coupled two-entry blocks \cite{dirac1928,greiner2000}
& \verifiedmark
& One Dirac field is represented by one fixed projected matrix column; two copies are represented by unrestricted algebra blocks.\\
Eqs.\ \eqref{eq:dirac-four-components}--\eqref{eq:dirac-expectation}
& Four Dirac components, conserved current, and Born expectation rule
& \verifiedmark
& The positive norm is conserved only for the complete two-block field; normalized outcome probabilities are supplied by Hermitian projectors.\\
Eq.\ \eqref{eq:dirac-square}
& Squared minimally coupled Dirac system
& \verifiedmark
& The terms \(+iq\conj{\pv{F}}\) and \(+iq\pv{F}\) are carried by the upper and lower blocks, respectively.\\
Eqs.\ \eqref{eq:large-small-coupled}--\eqref{eq:pauli-current}
& Leading Pauli reduction \cite{pauli1927,sakurai2017}
& \verifiedmark
& The expansion is in inverse mass after the large/small energy-basis split; convection and magnetization terms are contained in the inherited current.\\
Eq.\ \eqref{eq:energy-spin-projectors}
& Free energy and helicity projectors
& \verifiedmark
& Commutation of joint sectors is obtained for helicity (or any spin axis at rest), but not for an arbitrary fixed axis.\\
\bottomrule
\end{longtable}

% Verification report generated by tools/verify_sta_results.py via build_latex.sh.
\input{generated/python_verification.tex}

\phantomsection
\addcontentsline{toc}{section}{References}
\begin{thebibliography}{99}
\bibitem{hamilton1844}
W. R. Hamilton, ``On quaternions; or on a new system of imaginaries in algebra,'' \emph{Philosophical Magazine} 25 (1844), 489--495. DOI: \url{https://doi.org/10.1080/14786444408644923}.

\bibitem{maxwell1865}
J. C. Maxwell, ``A dynamical theory of the electromagnetic field,'' \emph{Philosophical Transactions of the Royal Society of London} 155 (1865), 459--512. DOI: \url{https://doi.org/10.1098/rstl.1865.0008}.

\bibitem{clifford1878}
W. K. Clifford, ``Applications of Grassmann's extensive algebra,'' \emph{American Journal of Mathematics} 1 (1878), 350--358. DOI: \url{https://doi.org/10.2307/2369379}.

\bibitem{lorentz1904}
H. A. Lorentz, ``Electromagnetic phenomena in a system moving with any velocity less than that of light,'' \emph{Proceedings of the Royal Netherlands Academy of Arts and Sciences} 6 (1904), 809--831.

\bibitem{einstein1905}
A. Einstein, ``Zur Elektrodynamik bewegter K\"orper,'' \emph{Annalen der Physik} 322 (1905), 891--921. DOI: \url{https://doi.org/10.1002/andp.19053221004}.

\bibitem{minkowski1908}
H. Minkowski, ``Raum und Zeit,'' lecture delivered September 21, 1908; published in \emph{Physikalische Zeitschrift} 10 (1909), 104--111.

\bibitem{klein1926}
O. Klein, ``Quantentheorie und f\"unfdimensionale Relativit\"atstheorie,'' \emph{Zeitschrift f\"ur Physik} 37 (1926), 895--906. DOI: \url{https://doi.org/10.1007/BF01397481}.

\bibitem{gordon1926}
W. Gordon, ``Der Comptoneffekt nach der Schr\"odingerschen Theorie,'' \emph{Zeitschrift f\"ur Physik} 40 (1926), 117--133. DOI: \url{https://doi.org/10.1007/BF01390840}.

\bibitem{pauli1927}
W. Pauli, ``Zur Quantenmechanik des magnetischen Elektrons,'' \emph{Zeitschrift f\"ur Physik} 43 (1927), 601--623. DOI: \url{https://doi.org/10.1007/BF01397326}.

\bibitem{dirac1928}
P. A. M. Dirac, ``The quantum theory of the electron,'' \emph{Proceedings of the Royal Society A} 117 (1928), 610--624. DOI: \url{https://doi.org/10.1098/rspa.1928.0023}.

\bibitem{robertson1929}
H. P. Robertson, ``The uncertainty principle,'' \emph{Physical Review} 34 (1929), 163--164. DOI: \url{https://doi.org/10.1103/PhysRev.34.163}.

\bibitem{hestenes1966}
D. Hestenes, \emph{Space-Time Algebra}, Gordon and Breach, 1966. See also the 2nd edition, Birkh\"auser, 2015, DOI: \url{https://doi.org/10.1007/978-3-319-18413-5}.

\bibitem{baylis1999}
W. E. Baylis, \emph{Electrodynamics: A Modern Geometric Approach}, Birkh\"auser, 1999.

\bibitem{doran2003}
C. Doran and A. Lasenby, \emph{Geometric Algebra for Physicists}, Cambridge University Press, 2003.

\bibitem{jackson1998}
J. D. Jackson, \emph{Classical Electrodynamics}, 3rd ed., Wiley, 1998.

\bibitem{greiner2000}
W. Greiner, \emph{Relativistic Quantum Mechanics: Wave Equations}, 3rd ed., Springer, 2000.

\bibitem{sakurai2017}
J. J. Sakurai and J. Napolitano, \emph{Modern Quantum Mechanics}, 2nd ed., Cambridge University Press, 2017.
\end{thebibliography}
\end{document}
